\documentclass[10pt]{amsart}
\usepackage{amsmath,amssymb,amsthm,mathtools}
\newtheorem{theorem}{Theorem}[section]
\newtheorem{lemma}[theorem]{Lemma}
\newtheorem{proposition}[theorem]{Proposition}
\newtheorem{corollary}[theorem]{Corollary}
\theoremstyle{definition}
\newtheorem{definition}[theorem]{Definition}
\newtheorem{problem}[theorem]{Problem}
\theoremstyle{remark}
\newtheorem{remark}[theorem]{Remark}
\newcommand{\NA}{\operatorname{NA}}
\newcommand{\supp}{\operatorname{supp}}
\newcommand{\sgn}{\operatorname{sign}}
\newcommand{\spn}{\operatorname{span}}
\newcommand{\dist}{\operatorname{dist}}
\newcommand{\R}{\mathbb R}
\newcommand{\N}{\mathbb N}
\newcommand{\Cset}{\mathcal C}
\newcommand{\Mset}{\mathcal M}
\newcommand{\Cert}{\operatorname{Cert}}
\newcommand{\Li}{\operatorname*{Li}}
\newcommand{\cl}{\operatorname{cl}}
\newcommand{\gap}{\operatorname{gap}}
\newcommand{\Def}{\operatorname{Def}}
\newcommand{\Rec}{\mathcal R}
\newcommand{\zh}{\hat z}
\newcommand{\ip}[2]{\langle #1,#2\rangle}
\newcommand{\lin}{\mathfrak l}
\newcommand{\Exc}{\operatorname{Exc}}
\setlength{\emergencystretch}{3em}
\title[Recovery of mates on Mart\'in norms III]{Recovery of mates on Mart\'in norms III:\\finite base support and truncation}
\author{}
\date{Research continuation, October 2026}
\begin{document}
\begin{abstract}
This is the third part of a study of the density of norm-attaining
operators from Mart\'in's renormings of $c_0$ into $\ell_2^2$.  Part~I
reduced density, for a designed operator and $N$ blocks, to an
approximation statement (Lemma~Z) about first rows, and Part~II split
Lemma~Z into a statement about first rows with finite base support and an
approximation statement (LSC-trunc) for infinite base support; at finite
base support it left open the rows of the residual (C$^*$).  Here we settle
the finite-support half for a new designed operator.  We show that the lever hypothesis of Part~II is a property
of the first row and not of the design, and replace it by a choice of
\emph{generic} weights: if the weights are algebraically independent over
the field generated by the other data of the design, then the threshold
value is never a local-minimum value of the floor statuses on an exact zero
set, which gives status coherence with design constants.  We construct
violation-tolerant $d$-consistent engineered approximants, whose per-piece
bounds have radius proportional to the scale with window-independent
constants (for a fixed lower bound of the near-threshold gaps), and average
at the norm-attaining approximant itself.  We then fix one operator
$T_{\rm fin}$, independent of $N$, in two versions: with option~(a) it has
norm $c_1<1$ and satisfies all other conclusions of Mart\'in's Lemma~B, so
that it defines Mart\'in-type norms; with option~(b) it has norm one and
satisfies the conclusion of Mart\'in's Lemma~B.  We prove Master
Theorem~VI: with option~(a), for the norms $p_N$ built with $T_{\rm fin}$,
every first row with finite base support is recoverable, for every
$N\ge1$ (with option~(b): for $N=1$, and for $N\ge2$ at the rows with
$\varrho_1(f)\ne1$).  Consequently, for $T_{\rm fin}$ with option~(a) and
every $N$ (with option~(b): for $N=1$), density of
$\NA((c_0,p_N),\ell_2^2)$ is equivalent to (LSC-trunc)$_N$.  For infinite
base support we prove a lowering transfer for contact truncations, an
engineered bound and a windowed recovery theorem at arbitrary base
support, and, for a variant $T^\infty_{\rm fin}$ of $T_{\rm fin}$
(option~(a), over the doubly exponential base of Part~II, with finitely
many additional design quantities per stage), cushion domination at
deep-raised rows and uniform solvability of the $d$-consistency system,
and obtain a transport theorem for $T^\infty_{\rm fin}$: (LSC-trunc)$_N$ holds at every triple
$(f,g,\rho)$ that satisfies a joint push condition and an A-regularity
condition at infinitely many main stages.  Triples not covered satisfy, at
almost all main stages, one of two explicit conditions, (W$^*$)
or (A$^*$/U$^*$).  The density of $\NA((c_0,p_N),\ell_2^2)$, for any $N$,
and of $\NA((c_0,p),\ell_2^2)$ remains open.
\end{abstract}
\maketitle
\tableofcontents

\section{Introduction}\label{sec:III-intro}

\subsection{The problem}
Mart\'in \cite{Martin} constructed renormings $p$ of $c_0$ whose
norm-attaining functionals form an algebraically trivial set.  The
construction starts from the canonical base norm $q$ of $c_0$ (a smoothing
of the supremum norm in the sense of Debs, Godefroy and Saint-Raymond
\cite{DGS}) and an injective operator $T$ of norm one given by his
Lemma~B, which rests on \cite[Proposition~2.8]{KLMW20}; we call such an
operator \emph{admissible} [I,~Definition~1.3] ($q^*(Te_{n,m})\le1$ with
equality for some $(n,m)$, injectivity, $Y\cap c_{00}=\{0\}$, and density
of the normalized vectors in every block).  For a set $I$ of blocks,
$p_N$ denotes the norm built with $I=\{1,\dots,N\}$ and $p$ the norm built
with $I=\N$.  The question studied in Part~I \cite{PartI}, Part~II
\cite{PartII} and here is whether $\NA((c_0,p_N),\ell_2^2)$ and
$\NA((c_0,p),\ell_2^2)$ are norm dense in the corresponding spaces of
operators.  It is a test case for Lindenstrauss' property~B
\cite{Lindenstrauss} of the two-dimensional Hilbert space (compare
\cite{JW,KLMW}): a negative answer for one $N$ and one admissible operator
would give a two-dimensional range space without property~B
[I,~Section~9].

An operator into $\ell_2^2$ is a pair $(f,g)$ of functionals.  For a
\emph{first row} $f\in S_{p^*}$ the \emph{mate fibre} is
$\Cset(f)=\{g:(f,g)\text{ is contractive}\}$, and the \emph{recoverable set}
is
\[
 \Rec(p):=\{f\in S_{p^*}:\ (f,g)\in\cl\NA((c_0,p),\ell_2^2)\text{ for every }
 g\in\Cset(f)\}
\]
[I,~Definition~1.15].  Density holds if and only if $\Rec(p)=S_{p^*}$
[I,~Proposition~1.13], and density for $p_N$ for infinitely many $N$
implies density for $p$ [I,~Lemma~1.5].

\subsection{What Parts~I and~II proved}
Part~I proved the transport of finite certificates along sequences of
first rows and an averaging theorem over scales [I,~Theorems~3.8
and~4.1]; the resonance obstruction, an admissible operator and a first
row at which the intrinsically recoverable mates lie on one line while
$\Cset(f)$ is infinite dimensional [I,~Theorem~6.6]; engineered recovery at
first rows with finite base support [I,~Theorem~7.18]; and, for
signature-ladder operators, recovery under signature room
[I,~Theorem~8.18] and the reduction of density to an approximation
statement [I,~Theorem~8.20, Problem~8.21]:

\begin{itemize}
\item[] \emph{Lemma~Z.}  For every $f\in S_{p_N^*}$, $g\in\Cset(f)$,
$\rho<1$ and $\varepsilon>0$ there is a first row $f'$ with finite base
support and signature room such that $p_N^*(f'-f)<\varepsilon$ and
$\dist(\rho g,\Cset(f'))<\varepsilon$.
\end{itemize}

Part~II fixed one admissible operator $T_{\rm final}$, independent of $N$,
with a doubly exponentially decaying diagonal base, bounded-gap signature
sets and windows split into pigeonhole sub-windows at which every rate
object of a first row is either tiny or robust ([II,~Definition~3.5,
Theorem~3.12, Theorem~3.19]).  For $T_{\rm final}$ it proved that a first
row with finite base support which is not recoverable is a (C$^*$) row, at
which a near-exact coherent shift resonance occurs at every clean
sub-window of every large level [II,~Theorem~5.15, Corollary~5.17], and,
for a modified design with absorber targets, recovery whenever at
infinitely many main stages some activity class has no active shift or
admits a $\varkappa$-neutral lever (KN$_{w,a}$) [II,~Theorem~5.68].  It
split Lemma~Z into its two halves [II,~Proposition~6.48]: density of
$\NA((c_0,p_N),\ell_2^2)$ is equivalent to recovery at first rows with
finite base support together with
\begin{itemize}
\item[] \textup{(LSC-trunc)$_N$}: for every $f\in S_{p_N^*}$, $g\in\Cset(f)$,
$\rho\in(0,1)$ and $\varepsilon>0$ there is $f'\in S_{p_N^*}$ with finite
base support such that $p_N^*(f'-f)<\varepsilon$ and
$\dist(\rho g,\Cset(f'))<\varepsilon$.
\end{itemize}
Lemma~Z, and density for every $N$ and for $p$, were left open in Part~II;
the open part consisted of the (C$^*$) rows at finite base support and of
(LSC-trunc) for the rows with infinite base support not covered by
[II,~Section~6].

\subsection{Main results: finite base support}
\emph{Status coherence by generic weights (Section~\ref{sec:III-generic}).}
Inside one block the exact shifted coarse system and the relative positions
of all switching carriers depend on the block only through the ratios
$u_l(\zh)/\varkappa$ (Lemma~\ref{lem:III-Rkt}), so peak pushes alone cannot
achieve exactness and threshold protection independently; Part~II supplied
the second lever by the hypothesis (KN$_{w,a}$).  We show that (KN) is a
property of the first row, not of the design (Lemma~\ref{lem:III-floor},
Remark~\ref{rem:III-KNf}), and that a second lever is nevertheless always
available.  Writing the inward rows of the coarse cone at the \emph{floor}
of the relative positions removes the weights of the near-threshold
switching carriers from the exact zero set of the tiny minors and from the
floor functions (Lemma~\ref{lem:III-weightfree}).  If all weights of the
design are algebraically independent over the countable field generated by
its other data, then the threshold value $1$ is never a local-minimum value
of a floor status on an exact zero set (Lemmas~\ref{lem:III-NLM}
and~\ref{lem:III-GEN}); the proof uses the Tarski--Seidenberg principle
\cite{Tarski,Seidenberg}, the semialgebraic Sard theorem \cite{BCR} and
Puiseux expansions \cite{vdDries}.  Compactness and the {\L}ojasiewicz
inequality turn this into status coherence with design constants
(Lemma~\ref{lem:III-QC}), and Proposition~\ref{prop:III-KNtr} is the
resulting exactification with threshold protection, without (KN).

\emph{The violation-tolerant $d$-consistent engineered approximant
(Section~\ref{sec:III-UE}).}  To average over a window of scales at a row
near the target row, one needs per-piece bounds with radius proportional to
the scale and window-independent constants.  The obstruction recorded in
[II,~Remark~2.31(c)], the junction mismatch between the
$\theta$-decomposition and the $\pm$-decompositions of a piece
(Proposition~\ref{prop:III-J}), is removed exactly: at the approximant the
normalized values of the finitely many switching carriers are re-tuned to
their values at the base row by small levers ($d$-consistency,
Lemma~\ref{lem:III-DC}, proved with the Sherman--Morrison formula and the
Poincar\'e--Miranda theorem \cite{Miranda}).  This gives the uniform
engineered bound (Theorem~\ref{thm:III-UE}), windowed recovery at the
norm-attaining approximant itself (Theorem~\ref{thm:III-Eeng}), the
threshold comparison (Proposition~\ref{prop:III-S2c}), and recovery of
every activity class under (KN) (Theorem~\ref{thm:III-Cmix}).

\emph{The operator $T_{\rm fin}$ and Master Theorem~VI
(Section~\ref{sec:III-MTVI}).}  $T_{\rm fin}$ is the design of
[II,~Definition~5.48] on the ladder of [II,~Definition~3.5] with the
changes listed in Remark~\ref{rem:III-Tfincomp}, chiefly: the base is
$\mu_s=2^{-s^2-1}$; there are no fresh-dominated stages; the weights are
generic; the rate scheme and the design factor contain the objects of
Sections~\ref{sec:III-generic} and~\ref{sec:III-UE}; and the window factor
$Q(w)$ does not contain Part~II's $\bar D(w)^3$ and has exponent
$\omega(l)+l^2+20$
(Definitions~\ref{def:III-Tfindata}--\ref{def:III-Tfinstage}).  The
recursion is well founded and does not refer to $N$
(Lemma~\ref{lem:III-Wfin}), and $T_{\rm fin}$ satisfies (T-b)--(T-d) of
[I,~Definition~1.3]; with option (b) it is admissible and satisfies the
conclusion of Mart\'in's Lemma~B; with option (a) its norm is $c_1<1$ (the
only departure from Lemma~B), so the norms it defines are Mart\'in-type
norms (Theorem~\ref{thm:III-Afin}).  The results of Parts~I and~II used here are
transferred to $T_{\rm fin}$ through an explicit list of interface
properties (Lemmas~\ref{lem:III-partI} and~\ref{lem:III-interface},
Proposition~\ref{prop:III-transfer}).  The recovery follows a single route
for every kind of activity class (route~V, Subsection~\ref{sec:III-routeV}):
the companion is the exactification of Proposition~\ref{prop:III-KNtr}, the
approximant is the $d$-consistent engineered approximant of
Section~\ref{sec:III-UE}, at which the block scalars are frozen up to the
weight of the fine carriers (Lemma~\ref{lem:III-PM}), and the averaging is
Theorem~\ref{thm:III-Eeng}.  The main result of this part is the following.

\begin{itemize}
\item[] \emph{Master Theorem~VI} (Theorem~\ref{thm:III-MTVI}).  Let $p_N$ be
the norm on $c_0$ built with $T_{\rm fin}$ and the blocks $\{1,\dots,N\}$.
With option \textup{(a)}, for every $N\ge1$ every $f\in S_{p_N^*}$ with
finite base support lies in $\Rec(p_N)$.  With option \textup{(b)}, the same
holds for $N=1$, and for $N\ge2$ at every $f$ with finite base support and
$\varrho_1(f)\ne1$, where $\varrho_1(f)$ is the relative position of the
first stage.
\end{itemize}

The operator is fixed before $N$ and before the row; this is legitimate
because the coincidence sets of the generic-weight argument contain no
datum of a row (Lemma~\ref{lem:III-GENfin}).  Combined with the splitting
of [II,~Proposition~6.48] this gives:

\begin{itemize}
\item[] \emph{Density criterion} (Theorem~\ref{thm:III-VI2},
Corollary~\ref{cor:III-VI3}).  Let $T_{\rm fin}$ be taken with option
\textup{(a)}.  For every $N\ge1$, $\NA((c_0,p_N),\ell_2^2)$ is dense if and
only if \textup{(LSC-trunc)$_N$} holds.  If \textup{(LSC-trunc)$_N$} holds
for infinitely many $N$, then $\NA((c_0,p),\ell_2^2)$ is dense for the
Mart\'in-type norm $p$ built with $T_{\rm fin}$.
\end{itemize}

The same equivalence holds with option (b) for $N=1$.  Thus, for this
operator, the density problem is reduced exactly to the approximation of
arbitrary first rows by first rows with finite base support while keeping
the mate approximately in the fibre.

\subsection{Main results: infinite base support}
Section~\ref{sec:III-inf} treats (LSC-trunc).  Truncating the base support
to finitely many coordinates costs nothing at finitely many scales fixed in
advance (the lowering transfer, Lemma~\ref{lem:III-LT} and
Corollary~\ref{cor:III-LTwin}); the factor $2$ in the lowering transfer is
sharp, and the cushion loss can be linear in the scale $r$ rather than
quadratic (Remark~\ref{rem:III-LTsharp}), so truncation is free only at
finitely many scales fixed in advance and must be built into the
approximant rather than performed afterwards.  Inside the violation-tolerant
engineered approximant the truncated support becomes a set of contacts whose
two-sided use is a tolerated violation; this gives an engineered bound and a
windowed recovery theorem with raise transfer at arbitrary base support
(Theorems~\ref{thm:III-UEinf} and~\ref{thm:III-Eenginf}), so that
(LSC-trunc) at $(f,g,\rho)$ follows from window families with side cushion
domination and exact $d$-consistency at base rows of arbitrary support.
For the design $T^\infty_{\rm fin}$ of Definition~\ref{def:III-designinf}
($T_{\rm fin}$ with option (a), over the doubly exponential diagonal base of
Part~II, with finitely many additional quantities per stage), the route-V
data at deep-raised rows are cushion dominated with an absolute constant
$C_+=5$ on the whole support of the companion
(Proposition~\ref{prop:III-cushion}); after a cheap preparation of the base
row the $d$-consistency system is uniformly solvable at every row
(Theorem~\ref{thm:III-SysDC}); and Step~A of the companion construction is
solvable whenever a cone condition (A-regularity,
Definition~\ref{def:III-Areg}) holds, while the obstructions to it have a
rigid sign structure (Lemma~\ref{lem:III-cone}, Theorem~\ref{thm:III-Aobs},
Corollary~\ref{cor:III-Asuff}).  The result is:

\begin{itemize}
\item[] \emph{Transport theorem} (Theorem~\ref{thm:III-transport}).  Let
the design be $T^\infty_{\rm fin}$, $N\ge1$, $f\in S_{p_N^*}$ with finite
or infinite base support, $g\in\Cset(f)$ and $\rho\in(0,1)$.  If for
infinitely many main stages $L$ there are a clean sub-window $w$ of $L$ and
an activity class $a$ with at least $n(w)/D_{\rm cl}(w)$ scales such that
the push condition (H1$^*$) holds and Step~A is A-regular (H2) at
$(L,w,a)$, then $(f,\rho g)$ is a limit of norm-attaining
pairs whose first rows have finite base support; in particular
\textup{(LSC-trunc)$_N$} holds at $(f,g,\rho)$.
\end{itemize}

Explicit sufficient conditions are given in
Corollary~\ref{cor:III-transport}; for instance, every first row whose base
support meets no signature set beyond $s_{\max}(L)$, for infinitely many
main stages $L$, lies in $\Rec(p_N)$ for the norms $p_N$ built with
$T^\infty_{\rm fin}$.

\subsection{What remains open}
\emph{The density of $\NA((c_0,p_N),\ell_2^2)$, for every $N\ge1$, and of
$\NA((c_0,p),\ell_2^2)$ remains open, for $T_{\rm fin}$, for
$T^\infty_{\rm fin}$ and for every admissible operator, in particular for
the operators of Mart\'in's Lemma~B.}  For
$T_{\rm fin}$ with option (a) the problem is equivalent to
(LSC-trunc)$_N$ (Theorem~\ref{thm:III-VI2}), and for $T^\infty_{\rm fin}$ the
triples $(f,g,\rho)$ at which (LSC-trunc)$_N$ is not established are
described by Proposition~\ref{prop:III-residual}: if the hypothesis of the
transport theorem fails, then at every clean sub-window of all but finitely
many main stages and for every admissible class one of the following holds.
\begin{itemize}
\item[(W$^*$)] The push condition (H1$^*$) fails: the active weak peaks of
the blocks of kind (T) cannot be pushed jointly below the threshold by
admissible moves of the allowed size.
\item[(A$^*$/U$^*$)] (H1$^*$) holds, and for every pre-placement there is a
near-certificate of Step~A with a robustly negative common mode $\gamma$;
it is of type (A$^*$) if $\gamma\le-\gamma_0$, and of type (U$^*$) if
$-\gamma_0<\gamma\le-c_{\rm A}\nu_1/(8H(L))$, in which case the
coefficients of the carriers whose near coordinates are all thick carry
$\ell_2$-mass at least $c_{\rm A}/2$.  Here $c_{\rm A}$ is the Step-A
constant of Theorem~\ref{thm:III-Aobs} and $\gamma_0=c_{\rm A}/(4\|h\|_2)$
is the constant of Definition~\ref{def:III-residual}; the type may depend
on the pre-placement.
\end{itemize}
In the case (A$^*$/U$^*$) the thick part of the base support at design
depth agrees, up to an error $\varepsilon\,\mathrm{Des}(L)/|\varkappa_\Psi|$
($\varepsilon$ the size of the certificate), with a fixed multiple
$\Psi_y/\varkappa_\Psi$ of the carrier combination of the certificate,
with a robustly negative factor $\varkappa_\Psi$
(Proposition~\ref{prop:III-residual}(i)).  These are necessary conditions for
not being covered, not characterizations; it is not known whether rows
satisfying them at all large main stages exist, nor whether such rows
carry a mate that is not recovered (Remark~\ref{rem:III-residual}).  At
finite base support one case remains open for option (b): $N\ge2$ and
$\varrho_1(f)=1$ (Remark~\ref{rem:III-MTVIopen}(b)).  Master Theorem~VI is
proved for the designed operator only, not for an arbitrary admissible
operator.  No counterexample to density is known, and none of the
verifications produced one.

\subsection{Verification status}
The results of Sections~\ref{sec:III-generic} and~\ref{sec:III-UE} on which
Master Theorem~VI rests were checked by one referee each and then
re-derived in a second, independent verification organized in four
clusters (the route-V companion; the generic-weight chain; the engineered
approximant and the averaging; the transfer to $T_{\rm fin}$ and the window
bookkeeping); the combination in Section~\ref{sec:III-MTVI} was checked by
two referees and re-derived in the same verification.  The corrections
introduced in the second verification (the common pull window, the
transfer interface of Lemma~\ref{lem:III-interface} and
Proposition~\ref{prop:III-transfer}, the $s$-free coefficients, the gap
criterion of Lemma~\ref{lem:III-dcons}(a), the bounds of
Lemmas~\ref{lem:III-DC} and~\ref{lem:III-UE1}(g)) were each made and
checked once.  Proposition~\ref{prop:III-J}, Lemma~\ref{lem:III-noB},
Corollary~\ref{cor:III-noB} and Lemma~\ref{lem:III-nocube} were checked by
one referee.  The proof of Master Theorem~VI given here uses
Lemma~\ref{lem:III-nocube}; a version that does not use it is described in
Remark~\ref{rem:III-cert}.  The present text was then read by three further
referees, whose corrections are incorporated; these corrections have not
been through a second independent verification.  The results of
Section~\ref{sec:III-inf} were each checked by one adversarial referee,
and the corrections found there are incorporated; they have not yet been
through a second independent verification.  They are proved for the
design $T^\infty_{\rm fin}$; they use the route-V companion and the window
count of Section~\ref{sec:III-MTVI} for that design, that is, the transfer
of Section~\ref{sec:III-MTVI} from the base $\mu$ to the base $\mu^\star$,
which rests on the base independence of the window theorems of Part~II
and has not been re-verified line by line (Remark~\ref{rem:III-infdep}).  Statements of earlier drafts that were found to be false or not proved as
stated are not used: the master theorem without a lever in the design of
an earlier draft (see [II,~Remark~5.70]), the two-sidedness of a bank placed
with an offset at the row where the $d$-consistency system is solved
(Lemma~\ref{lem:III-offset}), the sufficiency of a push condition for each
weak peak separately, characterizations of the non-covered rows by
carrier types (Remark~\ref{rem:III-residual}(e)), and levers for tuned
zero-value absorbers (an earlier Lemma~LV(d); see
Remarks~\ref{rem:III-routedecision} and~\ref{rem:III-abs}); the exact route with
zero-value absorbers of [II,~Theorem~5.68] does not transfer to
$T_{\rm fin}$ and is not used (Remark~\ref{rem:III-routeX}).  Finite numerical
models were used as sanity checks of several identities; they are evidence
only and are reported as such.

\subsection{Organization and conventions}
Section~\ref{sec:III-generic} proves status coherence by generic weights,
Section~\ref{sec:III-UE} constructs the violation-tolerant $d$-consistent
engineered approximant, Section~\ref{sec:III-MTVI} defines $T_{\rm fin}$ and
proves Master Theorem~VI and the density criterion,
Section~\ref{sec:III-inf} treats infinite base support, and
Section~\ref{sec:III-status} summarizes the status and states the open
problems.

The notation of Parts~I and~II is used without comment, in particular the
forced data $(\xi,q_0,a,w,F,z,\zh,e,\nu)$ of a first row
[I,~Proposition~1.9], two-piece data [I,~Definition~7.7], the
signature-ladder notation of [I,~Section~8], and the windows, rate objects
and clean sub-windows of [II,~Section~3].  References of the form
[I,~Theorem~8.20] and [II,~Proposition~6.48] are to Part~I \cite{PartI} and
Part~II \cite{PartII}.  Constants written $C_f$ depend on the first row
and on $N$ (and on the mate where said) but not on the level, the sub-window or the
scale; ``design constants'' are powers of the design factor
$\mathrm{Des}(L)$.  Every lemma, proposition, theorem and corollary is
given with a proof; material that is only sketched or heuristic appears
only in remarks that say so explicitly.
\section{Generic weights and status coherence}\label{sec:III-generic}

This section removes the lever hypothesis (KN$_{w,a}$) of
[II,~Definition~5.55] from the exactification with threshold protection of
[II,~Proposition~5.56].  The difficulty, isolated in [II,~Proposition~5.46]
and [II,~Remark~5.47], is that inside one block the exact shifted coarse
system and the relative positions of all switching carriers depend on the
block only through the ratios $u_l(\zh)/\varkappa$; peak pushes act on all of
them in the same way, so exactness and threshold protection cannot be
achieved independently by peak pushes (Subsection~\ref{subsec:III-twolevers}).
The hypothesis (KN) supplied a second, $\varkappa$-neutral lever.  We show
that (KN) is a property of the first row and not of the design
(Lemma~\ref{lem:III-floor}), and that a second lever is nevertheless always
available: after the inward rows of the coarse cone are written at the
\emph{floor} of the relative positions (Definition~\ref{def:III-cone}), the
exact zero set of the tiny minors and the floor functions contain no weight
of the near-threshold switching carriers (Lemma~\ref{lem:III-weightfree}).
These weights then enter only as denominators of the floor statuses, and if
all weights of the design are algebraically independent over the countable
field of its other data, the threshold value $1$ is never a local-minimum
value of the floor statuses on an exact zero set
(Lemmas~\ref{lem:III-NLM} and~\ref{lem:III-GEN}).  Compactness, the
{\L}ojasiewicz inequality and Puiseux expansions turn this into status
coherence with design constants (Lemma~\ref{lem:III-QC}), and
Proposition~\ref{prop:III-KNtr} is the resulting exactification.

The section incorporates the corrections found in the two independent
verifications of this material: the floor convention for active weak peaks
(Lemma~\ref{lem:III-K}(b) and Definition~\ref{def:III-pattern}), the
weight-free compact region (Definition~\ref{def:III-region}), the exact hold
of every non-active coarse strict non-peak of a (T)-block
(Proposition~\ref{prop:III-KNtr}(b) and Remark~\ref{rem:III-KN1}), the
common pull window of the realization (Lemma~\ref{lem:III-window}), the
Lipschitz constant of the floor statuses in the design thresholds
(Definition~\ref{def:III-QCreq}), and the $s$-free aggregated coefficients of
the rate objects (Lemma~\ref{lem:III-weightfree}(c)).

\subsection{Setting and notation}\label{subsec:III-gen-setting}

\subsubsection*{Blocks}
We use the block notation of [I,~Lemma~1.7] and [II,~Subsection~2.2].  For a
first row $f'$ and a block $m$ let $\zeta=\hat\zeta_m=R_m^{**}\zh'$,
$|\zeta|=|\zeta|_m$, let $\bar\vartheta=\bar\vartheta(\zeta)$ be the zero of
the threshold equation [II,~Lemma~2.6], $M=\bar\vartheta/(\bar\vartheta+
|\zeta|)$, $C=|\zeta|/(\bar\vartheta+|\zeta|)$ (so $M+C=1$ and $\bar\vartheta=
|\zeta|M/C$), $\mathsf v_k=|\zeta(k)|/\Phi(k)^2$ and $\varrho_k=\mathsf v_k/
\bar\vartheta$ the relative position of the coordinate $k$; $P=\{k:\varrho_k\ge1\}$
is the set of peaks and $Q$ its complement, the strict non-peaks.  We index
carriers by their stages $l$, with block $m(l)$, coordinate $k(l)$, weight
$\Phi_l:=\Phi_{m(l)}(k(l))=2^{-m(l)-k(l)}c_l$, $\lambda_l:=m(l)\Phi_l$ and
normalized target vector $u_l$; then $|\zeta(k(l))|=\lambda_l|u_l(\zh')|$ and
$\mathsf v_{k(l)}\Phi_l=m(l)|u_l(\zh')|$, and we write $\varrho_l:=
\varrho_{k(l)}$.  For a set $S$ of coordinates, $\Phi_S^2:=\sum_{k\in S}
\Phi(k)^2$.  Put
\[
 a:=\frac{|\zeta|}{\bar\vartheta}=\frac CM .
\]

\subsubsection*{The block scalar for an arbitrary finite switching set}
For a finite set $\Omega$ of coordinates of block $m$ (not necessarily
contained in $Q$) put, as in [II,~Subsection~2.2],
\[
 \varkappa(\zeta;\Omega):=\frac{|\zeta|\big(|\zeta|+\bar\vartheta\big)-
 \sum_{k\in\Omega}\mathsf v_k^2\Phi(k)^2}{\bar\vartheta},
\]
and, for the carriers $l$ with $k(l)\in\Omega$, the \emph{signed ratios}
$r_l:=u_l(\zh')/\varkappa(\zeta;\Omega)$ (when $\varkappa\ne0$), together with
\[
 R^2:=m^2\sum_{k(l)\in\Omega}r_l^2,\qquad \mathsf k:=
 \frac{\varkappa(\zeta;\Omega)}{\bar\vartheta},\qquad s_Q:=\frac1{\bar\vartheta^2}
 \sum_{k\in Q\setminus\Omega}\mathsf v_k^2\Phi(k)^2 .
\]
The ratios of [II,~Corollary~2.11] are $\mathsf r_l=m|r_l|$.  For
$s>0$ and $x\in[0,1)$ let
\[
 \mathbf K(s,x):=\frac{s+\sqrt{s^2x+s(1-x)}}{1-x}.
\]

\subsubsection*{Windows}
In Subsections~\ref{subsec:III-kinds}--\ref{subsec:III-KNtr} the following
is fixed: $N\ge1$, the design $T_{\rm fin}$ of
Section~\ref{sec:III-design} (besides the properties of $T_{\rm final}$
and $D^{U1'}$ used through Part~II, namely (T-a), (W1), (W5), the
signature sets with bounded gaps, the window properties of
[II,~Corollary~3.20] and the tools of [II,~Subsections~4.2--4.3
and~5.6], the new requirements are listed in
Definition~\ref{def:III-QCreq}), a row $f\in S_{p_N^*}$, a main stage
$L\ge l_f$, a sub-window $w$ of $L$ that is clean for $f$
([II,~Definition~3.18]), and an activity class $a$ of $w$
([II,~Lemma~5.53]).  We write $b:=b(w)$, $u:=u(w)$, $T_{\rm lo}:=T_{\rm lo}(w)$,
$\mathrm{Des}:=\mathrm{Des}(L)$, $D:=D(L)$, and
\[
 \eta':=\eta'(w):=\frac{T_{\rm lo}(w)^4}{L\,\mathrm{Des}(L)} .
\]
Carriers $l\le L$ are \emph{coarse}, carriers $l>L$ are \emph{fine}.  By
cleanness [II,~Corollary~3.20(b)] every coarse carrier of a block $m\le N$
has $\varrho_l(f)\in[0,b]\cup[u,1-u]\cup[1-b,1+b]\cup[1+u,\infty)$.  A coarse
carrier is \emph{nearly neutral} if $\varrho_l(f)\le b$, \emph{near-threshold}
if $|\varrho_l(f)-1|\le b$, a \emph{weak peak} if $\varrho_l(f)\in[1,1+b]$ and
a \emph{robust peak} if $\varrho_l(f)\ge1+u$.  For a block $m$ let
$P_{\rm rob}(m)$ be its robust peaks and $\sigma_{\rm rob}(m):=
\Phi_{P_{\rm rob}(m)}^2$.  Constants $C$, $c$ are numerical, $C_f$, $c_f$
depend only on $f$ (and $N$), and $K$ denotes quantities of the form
$C_f(\mathrm{Des}/u)^C$ as in [II,~Subsection~5.6].

\subsection{Why two levers are needed}\label{subsec:III-twolevers}

\begin{lemma}[Status rigidity in ratio space]\label{lem:III-Rkt}
Let $\zeta$ be the block vector of block $m$ of a first row, let
$\Omega\subseteq Q$ be finite, $\varkappa:=\varkappa(\zeta;\Omega)$, and let
$r_l$, $R$, $a$, $\mathsf k$, $s_Q$ be as in
Subsection~\textup{\ref{subsec:III-gen-setting}}.  Then:
\begin{enumerate}
\item[(a)] $a^2=\Phi_P^2+\mathsf k^2R^2+s_Q$, $\mathsf k=a+\Phi_P^2+s_Q$, and
$\varrho_l=m|r_l|\mathsf k/\Phi_l$ for every carrier $l$ with $k(l)\in\Omega$;
moreover $R<1$ and $a>\mathsf kR$.
\item[(b)] For fixed $(\Phi_P^2,s_Q)$ the pair $(a,\mathsf k)$ is a
real-analytic function of $R^2$, and $d\mathsf k/d(R^2)=\mathsf k^2/
(2(a-\mathsf kR^2))>0$.
\item[(c)] After division by $\varkappa$, the coupling row \textup{(X3)} of
[II,~Definition~5.52] reads $\Delta'_m=\sum_{k(l)\in\Omega}r_l\gamma_l$; hence
the exact shifted coarse system depends on block $m$ only through the
ratios $(r_l)$, and by \textup{(a)}, \textup{(b)} so do the relative
positions of all carriers of $\Omega$, for fixed $(\Phi_P^2,s_Q)$.
\item[(d)] A push of a peak, with the values on $\Omega$ fixed, multiplies
all ratios of the block by one common factor; if the ratios are restored
afterwards, every $\varrho_l$, $k(l)\in\Omega$, returns to its old value up
to the change of $(\Phi_P^2,s_Q)$.
\end{enumerate}
\end{lemma}

\begin{proof}
(a) is [II,~Corollary~2.11] with $\mathsf r_l=m|r_l|$.  (b) With $x:=R^2$ and
$c_0:=\Phi_P^2+s_Q>0$ fixed, $a$ is the root of $G(a):=a^2-(a+c_0)^2x-c_0$ on
the half-line $\{a>(a+c_0)\sqrt x\}$, where $G'(a)=2(a-\mathsf kx)\ge
2(a-\mathsf k\sqrt x)>0$; the implicit function theorem gives analyticity,
and differentiating $a^2=c_0+\mathsf k^2x$ with $da=d\mathsf k$ gives
$2a\,d\mathsf k=\mathsf k^2dx+2\mathsf kx\,d\mathsf k$.  (c) is
[II,~Proposition~5.46(i)], and (d) is [II,~Proposition~5.46(ii)].
\end{proof}

\begin{remark}[Two levers]\label{rem:III-twolevers}
Lemma~\ref{lem:III-Rkt} explains the gap of the one-lever exactification
recorded in [II,~Remark~5.47]: if the ratios of the active switching
carriers are moved to an exact point of the tiny minors and $\varkappa$ is set
by one peak push per block, all relative positions of the switching carriers
are determined by the exact point and by $(\Phi_P^2,s_Q)$, and peak pushes
cannot change them independently.  Threshold protection therefore needs a
second, independent lever.  In [II,~Proposition~5.56] this lever is an
inactive switching carrier of robust relative position, whose inward push
raises $\bar\vartheta$ at first order without moving $\varkappa$
([II,~Proposition~5.46(iii)]); its existence is the hypothesis
(KN$_{w,a}$) of [II,~Definition~5.55].  Lemma~\ref{lem:III-floor} below
shows that (KN) is a property of the row $f$, which no move made at a
companion can create.  In this section the second lever is the position of
the exact point on the exact zero set, decoupled from the statuses by
generic weights.
\end{remark}

\subsection{The block identities}\label{subsec:III-blockid}

\begin{lemma}[The $\alpha$-mass identity]\label{lem:III-alpha}
\textup{(a)} For every $\zeta\in\ell_1\setminus\{0\}$,
\[
 \sum_{k\in P}(\varrho_k-1)\Phi(k)^2=\frac{|\zeta|}{\bar\vartheta}=a=\frac CM .
\]
In particular every block has a peak $k$ with $\varrho_k>1$.

\textup{(b)} In the window setting of
Subsection~\textup{\ref{subsec:III-gen-setting}} let $m\le N$ and assume
\begin{equation}\label{eq:III-alphalf}
 \frac b{12}+\frac{2C_U\,c_{L+1}}{3\bar\vartheta_m(f)}<a_m(f),\qquad
 C_U:=\sup_l|u_l(\zh)| ,
\end{equation}
where $\zh$, $\bar\vartheta_m(f)$, $a_m(f)$ refer to $f$.  Then $P_{\rm rob}(m)
\ne\emptyset$, and $\sigma_{\rm rob}(m)\ge\Phi_{\min}(L)^2$ with $\Phi_{\min}(L):=
\min_{l\le L}\Phi_l$.  Condition \eqref{eq:III-alphalf} holds for every block
$m\le N$ as soon as $L$ exceeds an $f$-constant, which we include in $l_f$.
\end{lemma}

\begin{proof}
(a) By [II,~Lemma~2.6(c)], $\sum_{k\in P}\Phi(k)^2(\mathsf v_k-\bar\vartheta)=
|\zeta|$; divide by $\bar\vartheta$ and use $\varrho_k=\mathsf v_k/\bar\vartheta$ and
$|\zeta|/\bar\vartheta=C/M$ ([II,~Lemma~2.6(b)]).  Since $|\zeta|>0$, some term is
positive.  (Equivalently: pairing the representation $\zeta/|\zeta|=\alpha+
D^2w/C$ of [I,~Lemma~1.7] with $w$ gives $1=M\|\alpha\|_1+C$, so $\|\alpha\|_1
=1$, and $\varrho_k-1=|\alpha(k)|C/(M\Phi(k)^2)$ on $P$.)

(b) Split $P$ into the robust coarse peaks, the weak coarse peaks and the
fine peaks.  A weak peak contributes $(\varrho_k-1)\Phi(k)^2\le b\Phi(k)^2$, and
$\|\Phi_m\|_2^2\le\sum_k4^{-m-k}\le\frac1{12}$, so the weak peaks contribute
at most $b/12$.  A fine peak $l$ contributes at most $\varrho_l\Phi_l^2=
|\zeta(k(l))|/\bar\vartheta_m=\lambda_l|u_l(\zh)|/\bar\vartheta_m\le C_U\lambda_l/
\bar\vartheta_m$; as $\lambda_l=m(l)2^{-m(l)-k(l)}c_l\le c_l/2$ and $c_{l+1}\le
c_l/4$, the fine carriers contribute at most $\frac23C_Uc_{L+1}/
\bar\vartheta_m$.  By (a) and \eqref{eq:III-alphalf} the robust peaks carry a
positive mass, so $P_{\rm rob}(m)\ne\emptyset$ and $\sigma_{\rm rob}(m)\ge
\Phi_{\min}(L)^2$.  $C_U$ is finite because the vectors $u_l$ are bounded in
$\ell_1$ and $\zh$ is bounded; $b=b(w)\le\eta'(w)\to0$ and $c_{L+1}\to0$ as
$L\to\infty$, while $a_m(f)>0$ and $\bar\vartheta_m(f)>0$ do not depend on $L$.
\end{proof}

\begin{lemma}[Closed form of $\varkappa/\bar\vartheta$]\label{lem:III-K}
Let $\zeta$ be the block vector of block $m$ of a first row, let $\Omega$ be a
finite set of coordinates of block $m$, and put
\[
 \sigma':=\Phi_{P\setminus\Omega}^2+\sum_{k\in\Omega\cap P}\big(1-\varrho_k^2\big)
 \Phi(k)^2+s_Q .
\]
\begin{enumerate}
\item[(a)] $a^2=\sum_k\Phi(k)^2\min(\varrho_k,1)^2\le\|\Phi_m\|_2^2\le\frac1{12}$
and $\sigma'\le\|\Phi_m\|_2^2$; if $\Omega\subseteq Q$, then $\sigma'=\sigma:=
\Phi_P^2+s_Q$.
\item[(b)] $\mathsf k=a+\sigma'$ and $a^2=\sigma'+\mathsf k^2R^2$; if
$\varkappa(\zeta;\Omega)\ne0$, then $\varrho_l=m|r_l|\mathsf k/\Phi_l$ for every
carrier $l$ with $k(l)\in\Omega$.
\item[(c)] If $\sigma'>0$, then $R<1$, $a>\mathsf kR$, $\mathsf k=\mathbf
K(\sigma',R^2)$ and
\[
 1-R^2=\frac{\sigma'(1+2a+\sigma')}{(a+\sigma')^2}\ge\sigma' .
\]
\item[(d)] $\mathbf K$ is real-analytic on a neighbourhood of $(0,\infty)\times
[0,1)$; with $\mathbf A:=\mathbf K-s$ one has $\mathbf A>0$, $\mathbf A^2=s+
\mathbf K^2x$, $\partial_s\mathbf K>0$ and $\partial_x\mathbf K=\mathbf K^2/
(2(\mathbf A-\mathbf Kx))>0$ there.
\item[(e)] For $0<s_0\le1$ and $(s,x)$ in $\mathcal K(s_0):=\{s_0\le s\le1,\
0\le x\le1-s_0/4\}$:
\[
 2s\le\mathbf K\le\frac{10}{s_0},\qquad 0<\partial_s\mathbf K\le\frac{16}{s_0^2},
 \qquad0<\partial_x\mathbf K\le\frac{1000}{s_0^4} .
\]
\end{enumerate}
\end{lemma}

\begin{proof}
(a) Divide the second identity of [II,~Lemma~2.6(c)] by $\bar\vartheta^2$:
$a^2=\Phi_P^2+\sum_{k\in Q}\varrho_k^2\Phi(k)^2$, which is the first formula.
By (T-a), $\Phi_m(k)\le2^{-m-k}$, so $\|\Phi_m\|_2^2\le4^{-m}/3$.  Each of the
three terms of $\sigma'$ is at most the corresponding part of
$\sum_k\Phi(k)^2$ (note $1-\varrho_k^2\le1$ and $\varrho_k<1$ on $Q$).

(b) By [II,~Lemma~2.6(c)], $|\zeta|^2=\bar\vartheta^2\Phi_P^2+\sum_{k\in Q}
\mathsf v_k^2\Phi(k)^2$.  Inserting this into the definition of $\varkappa$ and
splitting $\Phi_P^2=\Phi_{P\setminus\Omega}^2+\Phi_{\Omega\cap P}^2$,
$\sum_Q=\sum_{Q\setminus\Omega}+\sum_{\Omega\cap Q}$,
\[
 \bar\vartheta\varkappa=|\zeta|\bar\vartheta+\bar\vartheta^2\Phi_P^2+\sum_{k\in Q}
 \mathsf v_k^2\Phi(k)^2-\sum_{k\in\Omega}\mathsf v_k^2\Phi(k)^2
 =|\zeta|\bar\vartheta+\bar\vartheta^2\sigma' ;
\]
divided by $\bar\vartheta^2$ this is $\mathsf k=a+\sigma'$.  Next,
$\sum_{k\in\Omega}\mathsf v_k^2\Phi(k)^2=\sum_{k(l)\in\Omega}m^2u_l(\zh)^2=
\varkappa^2R^2$, i.e.\ $\sum_{k\in\Omega}\varrho_k^2\Phi(k)^2=\mathsf k^2R^2$, and
the same splitting in the formula of (a) gives $a^2=\sigma'+\sum_{k\in
\Omega}\varrho_k^2\Phi(k)^2=\sigma'+\mathsf k^2R^2$.  Finally $\varrho_l=
\mathsf v_{k(l)}/\bar\vartheta=m|u_l(\zh)|/(\Phi_l\bar\vartheta)=m|r_l|\mathsf k/
\Phi_l$.

(c) $a>0$ because $\zeta\ne0$.  If $\sigma'>0$, then $\mathsf k>a>0$ and
$a^2>\mathsf k^2R^2$, so $\mathsf kR<a<\mathsf k$ and $R<1$.  Substituting
$a=\mathsf k-\sigma'$ into $a^2=\sigma'+\mathsf k^2R^2$ gives $(1-R^2)\mathsf k^2-
2\sigma'\mathsf k+\sigma'^2-\sigma'=0$, with roots $\mathsf k_\pm:=[\sigma'\pm
(\sigma'^2R^2+\sigma'(1-R^2))^{1/2}]/(1-R^2)$.  The inequality $\mathsf k_-\le
\sigma'$ is equivalent to $\sigma'R^2\le(\sigma'^2R^2+\sigma'(1-R^2))^{1/2}$, i.e.\
to $0\le\sigma'(1-R^2)(1+\sigma'R^2)$, which holds; so $\mathsf k_-$ would give
$a\le0$, and $\mathsf k=\mathsf k_+=\mathbf K(\sigma',R^2)$.  From $R^2=(a^2-
\sigma')/(a+\sigma')^2$ we get the formula for $1-R^2$; by (a), $a+\sigma'\le
12^{-1/2}+\frac1{12}<1$, hence $(a+\sigma')^2\le a+\sigma'\le1+2a+\sigma'$, which
gives $1-R^2\ge\sigma'$.

(d) For $s>0$ and $x\in[0,1)$ the radicand $s^2x+s(1-x)$ is positive, so
$\mathbf K$ is real-analytic near such points, and $\mathbf A=[sx+(s^2x+
s(1-x))^{1/2}]/(1-x)>0$.  The identity $\mathbf A^2=s+\mathbf K^2x$ is the
quadratic equation of (c).  The $s$-derivative of the numerator of $\mathbf
K$ is $1+(2sx+1-x)/(2(s^2x+s(1-x))^{1/2})>0$.  Differentiating $\mathbf A^2=s+
\mathbf K^2x$ in $x$ (note $\partial_x\mathbf A=\partial_x\mathbf K$) gives
$2\mathbf A\,\partial_x\mathbf K=\mathbf K^2+2\mathbf Kx\,\partial_x\mathbf K$;
and $\mathbf A^2>\mathbf K^2x$ gives $\mathbf A>\mathbf K\sqrt x\ge\mathbf Kx$.

(e) On $\mathcal K(s_0)$ the radicand lies in $[s^2,\,s^2+s]\subseteq[s^2,2]$
(for $s\le1$, $s(1-x)\ge s^2(1-x)$), so the numerator lies in $[2s,\frac52]$
and $1-x\in[s_0/4,1]$; this gives the bounds on $\mathbf K$.  Since the
radicand is $\ge s(1-x)\ge s_0^2/4$ and $2sx+1-x\le3$, $\partial_s\mathbf K\le
(1+3/s_0)\cdot4/s_0\le16/s_0^2$.  Finally $\mathbf A-\mathbf Kx\ge\mathbf A-
\mathbf K\sqrt x=s/(\mathbf A+\mathbf K\sqrt x)\ge s/(2\mathbf K)$ (as $\mathbf
A<\mathbf K$), so $\partial_x\mathbf K\le\mathbf K^3/s\le1000/s_0^4$.
\end{proof}

\begin{remark}[Convention for active weak peaks]\label{rem:III-K1}
Lemma~\ref{lem:III-K}(b) holds for every finite $\Omega$; this is used with
the \emph{pattern's} switching set (Definition~\ref{def:III-pattern}), which
contains the active weak peaks of $f$.  For a weak peak, $1-\varrho_k^2\in
[-3b,0]$, so $\sigma'\ge\sigma_{\rm rob}(m)-3b\|\Phi_m\|_2^2\ge\sigma_{\rm rob}(m)-
b/4$ whenever $P\setminus\Omega\supseteq P_{\rm rob}(m)$ and $\Omega\cap P$ consists
of weak peaks.  If instead the active
weak peaks were treated as peaks (``peak convention''), all ratios of the
block would change by a relative amount of order $\Phi_l^2/\mathsf k$, which is
not small compared with the margins below, and their weights would enter
the peak traces of the cone (Lemma~\ref{lem:III-weightfree}); the formula-1
convention with the pattern's $\Omega$ avoids both.
\end{remark}

\begin{lemma}[Floor lemma]\label{lem:III-floor}
Let $\zeta$ be the block vector of block $m$ of a first row, $\Omega\subseteq Q$
finite with $\varkappa(\zeta;\Omega)\ne0$, $\Omega_a\subseteq\Omega$, and
$P_0\subseteq P$ with $\sigma_0:=\Phi_{P_0}^2>0$.  Put $R_a^2:=m^2\sum_{k(l)\in
\Omega_a}r_l^2$.  Then $R_a<1$, and for every carrier $l$ with $k(l)\in\Omega_a$
\[
 \varrho_l\ \ge\ \varrho^0_l:=\frac{m|r_l|\,\mathbf K(\sigma_0,R_a^2)}{\Phi_l}.
\]
If $r_l\ne0$, equality holds if and only if $P=P_0$, $s_Q=0$ and $r_{l'}=0$
for every $l'$ with $k(l')\in\Omega\setminus\Omega_a$.
\end{lemma}

\begin{proof}
By Lemma~\ref{lem:III-K}(a)--(c) with $\Omega\subseteq Q$: $\sigma=\Phi_P^2+s_Q\ge
\sigma_0>0$, $R<1$ and $\varrho_l=m|r_l|\mathbf K(\sigma,R^2)/\Phi_l$.  Moreover
$R^2=R_a^2+m^2\sum_{k(l')\in\Omega\setminus\Omega_a}r_{l'}^2\ge R_a^2$.  $\mathbf K$
is strictly increasing in both variables (Lemma~\ref{lem:III-K}(d)), and
$\sigma=\sigma_0$ iff $P=P_0$ and $s_Q=0$.
\end{proof}

\begin{remark}[(KN) is a property of the row]\label{rem:III-KNf}
(a) Let $f'$ be any first row (for instance a companion of $f$) at which
block $m$ has $P(f')\supseteq P_{\rm rob}(m)$, $\Omega\subseteq Q(f')$, and at
which the ratios of the active carriers $\Omega_a$ equal a given vector $r$.
By Lemma~\ref{lem:III-floor} no active relative position at $f'$ is below its
floor $\varrho^0_l(r)=m|r_l|\mathbf K(\sigma_{\rm rob}(m),R_a(r)^2)/\Phi_l$.  An
inactive carrier, a fine carrier or an additional peak only adds to
$\sigma$ or to $R^2$; a lever that is created at a companion (a carrier with
value $0$ at $f$ that is given a value and then pushed inward) can therefore
lower the active statuses at fixed active ratios only back towards the
floor.  At fixed active ratios the largest possible decrease of $\mathsf k$ is
$I_m(f'):=\mathbf K(\sigma(f'),R(f')^2)-\mathbf K(\sigma_{\rm rob}(m),R_a^2)$,
which at $f'=f$ is computed from $f$.  If every inactive carrier of
$\Omega_m$ is nearly neutral and no weak peak of block $m$ lies outside
$\Omega$, then $I_m(f)\le C\sigma_{\rm rob}(m)^{-5}(b+c_{L+1}^2)$ (proof of
Lemma~\ref{lem:III-Tblock} below); a lever of robust size exists at $f$ only
if some inactive carrier of $\Omega_m$ has $\varrho(f)\ge u$, which is the
condition of [II,~Definition~5.55] for block $m$.  In this sense (KN) is a
property of $f$.

(b) \emph{Heuristic.}  We do not see how a design could supply such levers
for every row: fine carriers have capacity at most of order $\bar\vartheta
c_{L+1}^2$, far below the displacements of the {\L}ojasiewicz step;
levers tuned at a companion give back what they add, by (a); and families
of near-parallel ``sibling'' carriers leave the coarsest member of every
family without a sibling.  This is not used anywhere; the rest of the
section shows that no lever of this kind is needed.
\end{remark}

\subsection{Patterns, kinds of blocks and the floor-form cone}\label{subsec:III-kinds}

\begin{definition}[Switching set, kinds of blocks, pattern]\label{def:III-pattern}
In the window setting of Subsection~\textup{\ref{subsec:III-gen-setting}}:
\begin{enumerate}
\item[(a)] For a block $m$ let $\Omega_m$ be the set of coordinates $k(l)$ of
the coarse carriers $l$ of block $m$ with $\varrho_l(f)\le1+b$, i.e.\ of all
coarse carriers except the robust peaks, and $\Omega:=\bigcup_m\Omega_m$.  The
ratios of the pattern are always computed with $\varkappa(\cdot\,;\Omega_m)$
(formula of Subsection~\ref{subsec:III-gen-setting}), also when $\Omega_m$
contains weak peaks.  The class $a$ declares each component $\gamma_l$
($k(l)\in\Omega$) and each shift $\Delta'_m$ active or inactive
([II,~Lemma~5.53]); $\Omega_a(m)$ denotes the active carriers of $\Omega_m$,
and $\mathcal N(m)\subseteq\Omega_a(m)$ the active near-threshold ones.
\item[(b)] A block $m$ is of kind
 \textup{(U)} if $\Delta'_m$ is inactive (this includes the blocks without
 active components);
 \textup{(KN)} if $\Delta'_m$ is active and some inactive carrier of $\Omega_m$
 has $\varrho(f)\ge u$;
 \textup{(T)} if $\Delta'_m$ is active and every inactive carrier of $\Omega_m$ is
 nearly neutral.
By cleanness every carrier of $\Omega$ has $\varrho(f)\in[0,b]\cup[u,1+b]$, so
every block has exactly one kind.  Since $\Omega_m$ contains the weak peaks
of block $m$, an inactive weak peak makes a block with active shift of kind
(KN); in particular every weak peak of a block of kind (T) is active.  Put $\mathcal N_T:=\bigcup\{\mathcal N(m):
m\text{ of kind (T)}\}$.
\item[(c)] The \emph{pattern} $\kappa$ of $(f,w,a)$ consists of the finite
combinatorial data: the sets $\Omega_m$, $\Omega_a(m)$, $\mathcal N(m)$,
$P_{\rm rob}(m)$, the signs $\mathrm{vs}_l$ of the coarse carriers, the classes
R/G and the contact/free types of the coordinates of $E_c(w)$, the type of
each aggregated signature row (sign row, equality row or absent), the kinds
of the blocks, and the labels \emph{tiny}/\emph{robust} of the rate objects of
Definition~\ref{def:III-cone}.  Its \emph{level} is $L$.  There are finitely
many patterns of each level, and they do not depend on $N$ (blocks $m>N$
are absent and contribute no row).
\end{enumerate}
\end{definition}

The kind (U) needs no lever (Lemma~\ref{lem:III-Uscale}), the kind (KN) is
handled by the lever of [II,~Proposition~5.56], and only the kind (T) is new.

\begin{lemma}[Floor estimate in (T)-blocks]\label{lem:III-Tblock}
Let $m$ be a block of kind \textup{(T)}, let $r(f)$ be its ratio vector
\textup{(}computed with $\varkappa(\hat\zeta_m(f);\Omega_m)$\textup{)}, and
assume $b\le\sigma_{\rm rob}(m)/8$.  For every carrier $l$ of $\Omega_a(m)$ put
$\varrho^0_l(r):=m|r_l|\mathbf K(\sigma_{\rm rob}(m),R_a(m)(r)^2)/\Phi_l$, where
$R_a(m)(r)^2:=m^2\sum_{l\in\Omega_a(m)}r_l^2$.  Then:
\begin{enumerate}
\item[(a)] $|\varrho_l(f)-\varrho^0_l(r(f))|\le C\sigma_{\rm rob}(m)^{-6}\big(b+
c_{L+1}^2\big)$;
\item[(b)] $1-R_a(m)(r(f))^2\ge\sigma_{\rm rob}(m)-b/4$.
\end{enumerate}
\end{lemma}

\begin{proof}
Write $\sigma_0:=\sigma_{\rm rob}(m)$.  At $f$, $P\setminus\Omega_m$ consists of
$P_{\rm rob}(m)$ and the fine peaks, $\Omega_m\cap P$ consists of weak peaks,
and $Q\setminus\Omega_m$ consists of fine carriers.  Hence, in
Lemma~\ref{lem:III-K},
\[
 \sigma'-\sigma_0=\Phi^2_{\text{fine peaks}}+s_Q+\sum_{k\in\Omega_m\cap P}
 (1-\varrho_k^2)\Phi(k)^2\in\Big[-\frac b4,\ \sum_{l>L}\Phi_l^2\Big],
\]
and $\sum_{l>L}\Phi_l^2\le\sum_{l>L}c_l^2/16\le c_{L+1}^2/8$ (using
$1-\varrho^2\ge-3b$ on weak peaks and $\|\Phi_m\|_2^2\le\frac1{12}$).  So
$\sigma'\ge\frac78\sigma_0>0$, and by Lemma~\ref{lem:III-K}(c) $1-R^2\ge\sigma'
\ge\sigma_0-b/4$; as $R_a\le R$ this gives (b).  Also $\mathsf k=a+\sigma'\ge
\sigma_0/2$.  The inactive carriers of $\Omega_m$ are nearly neutral, so
$R^2-R_a^2=\sum_{\rm inactive}\varrho^2\Phi^2/\mathsf k^2\le b^2/(3\sigma_0^2)$.
Both points $(\sigma',R^2)$ and $(\sigma_0,R_a^2)$ lie in the convex set
$\mathcal K(\sigma_0/2)$ of Lemma~\ref{lem:III-K}(e) (since $1-R_a^2\ge1-R^2
\ge\sigma_0/2$ and $\sigma_0,\sigma'\le1$), so
\[
 |\mathbf K(\sigma',R^2)-\mathbf K(\sigma_0,R_a^2)|\le\frac{64}{\sigma_0^2}
 \Big(\frac b4+c_{L+1}^2\Big)+\frac{16000}{\sigma_0^4}\cdot\frac{b^2}{3\sigma_0^2}
 \le C\sigma_0^{-5}(b+c_{L+1}^2),
\]
using $b\le\sigma_0/8$.  Finally $\varrho_l(f)=(m|r_l|/\Phi_l)\mathbf K(\sigma',
R^2)$ by Lemma~\ref{lem:III-K}(b),(c), and $m|r_l|/\Phi_l=\varrho_l(f)/
\mathbf K(\sigma',R^2)\le(1+b)/(2\sigma')\le2/\sigma_0$ (Lemma~\ref{lem:III-K}(e):
$\mathbf K\ge2s$).  Multiplying gives (a).
\end{proof}

The hypothesis $b\le\sigma_{\rm rob}(m)/8$ holds for $L\ge l_f$: by
Lemma~\ref{lem:III-alpha}(b), $\sigma_{\rm rob}(m)\ge\Phi_{\min}(L)^2\ge D(L)^{-2}$,
while $b\le\eta'\le T_{\rm lo}^4$ and $\mathrm{Des}(L)^2T_{\rm lo}\le1$
([II,~Corollary~3.20(c)]).

\begin{lemma}[Effective columns]\label{lem:III-EC}
Let $f'$ be a first row at which the coarse peaks of every block $m$ are
exactly $P_{\rm rob}(m)$ and every carrier of $\Omega$ is a strict non-peak, and
let $\mathcal L$ be the coarse certificate of [II,~Definition~5.52] at $f'$,
\[
 \mathcal L(\Delta',\gamma)=\Big(\sum_{k(l)\in\Omega}\gamma_lu_l-\sum_m\Delta'_m
 \tau_m\Big)\Big|_{E_c(w)\setminus F'},\qquad
 \tau_m:=\sum_{p\in P_{\rm rob}(m)}\mathrm{vs}_p\lambda_pu_p .
\]
If $(\Delta',\gamma)$ satisfies \textup{(X3)} and \textup{(X5)} of
[II,~Definition~5.52], then
\[
 \mathcal L(\Delta',\gamma)=\sum_{l\in\Omega_a}\gamma_l\big(u_l-r_l\tau_{m(l)}
 \big)\Big|_{E_c(w)\setminus F'},\qquad\Omega_a:=\bigcup_m\Omega_a(m),
\]
with the ratios $r_l$ of $f'$.  The column of $\gamma_l$ depends only on
$u_l$, on its own ratio $r_l$ and on $\tau_{m(l)}$, which contains the weights
of the robust peaks of block $m(l)$ only; the weights of the carriers of
$\Omega$ do not occur.
\end{lemma}

\begin{proof}
By [II,~Lemma~2.22(a)] and the definition of $\mathcal L$ in [II,~Definition~5.52],
the coefficient of a coarse peak $p$ is $-\Delta'_m\mathrm{vs}_p\lambda_p$, which
produces $-\Delta'_m\tau_m$.  By (X5), $\gamma_l=0$ for inactive carriers, and
$\Delta'_m=0$ for inactive shifts; by (X3) divided by $\varkappa_m$
(Lemma~\ref{lem:III-Rkt}(c)), $\Delta'_m=\sum_{l\in\Omega_a(m)}r_l\gamma_l$ for
active shifts.  For an inactive shift the term $\Delta'_m\tau_m$ vanishes,
and (X3) with $\Delta'_m=0$ gives $\sum_{l\in\Omega_a(m)}r_l\gamma_l=0$, so
$\sum_l\gamma_lr_l\tau_m=0$ as well.
Substitute.
\end{proof}

\begin{definition}[The weight-free region]\label{def:III-region}
For a pattern $\kappa$ let
\[
 \mathcal K_\kappa:=\Big\{r\in\R^{\Omega_a}:\ R_a(m)(r)^2\le1-\tfrac14
 \sigma_{\rm rob}(m)\ \text{for every block $m$ with}\ \Omega_a(m)\ne\emptyset
 \Big\},
\]
where $R_a(m)(r)^2:=m^2\sum_{l\in\Omega_a(m)}r_l^2$.  $\mathcal K_\kappa$ is
compact (on it $|r_l|\le1/m(l)$) and semialgebraic, and it is defined by
the robust peak weights only.  For a carrier $l\in\Omega_a$ let $\varrho^0_l(r):=
m|r_l|\mathbf K(\sigma_{\rm rob}(m),R_a(m)(r)^2)/\Phi_l$ $(m=m(l))$ be its floor
status, and let $\Lambda_\kappa$ be the largest Lipschitz constant on
$\mathcal K_\kappa$ of the functions $\varrho^0_l$, $l\in\Omega_a$; by
Lemma~\ref{lem:III-K}(e), $\Lambda_\kappa\le C\max_lm(l)\sigma_{\rm rob}(m(l))^{-4}
\Phi_l^{-1}<\infty$.
\end{definition}

\begin{lemma}[Interior points]\label{lem:III-interior}
Assume $b\le\sigma_{\rm rob}(m)/4$ for every block $m$.  Then $1-R_a(m)(r(f))^2
\ge\sigma_{\rm rob}(m)-b/4$ for every block $m$, where $r(f)$ is the ratio vector
of the pattern at $f$.  Here and in Lemmas~\ref{lem:III-QC}, \ref{lem:III-loj}
and Proposition~\ref{prop:III-KNtr}, $|\cdot|$ denotes the Euclidean norm (on
$\R^{\Omega_a}$, resp.\ $\R^n$).  If $\rho_0>0$ and $3L\rho_0\le\sigma_{\rm rob}(m)/4$ for
every $m$, then every $r$ with $|r-r(f)|\le\rho_0$ satisfies $R_a(m)(r)^2\le
1-\sigma_{\rm rob}(m)/2$ for every $m$; in particular $r\in\mathcal K_\kappa$.
\end{lemma}

\begin{proof}
For every block, $P\setminus\Omega_m\supseteq P_{\rm rob}(m)$ and $\Omega_m\cap P$
consists of weak peaks, so $\sigma'\ge\sigma_{\rm rob}(m)-b/4>0$
(Remark~\ref{rem:III-K1}); Lemma~\ref{lem:III-K}(c) gives $1-R_a^2\ge1-R^2\ge
\sigma'$.  The function $r\mapsto R_a(m)(r)=m\|r^{(m)}\|_2$ is $m$-Lipschitz and
$m\le L$, so for $|r-r(f)|\le\rho_0$ (and $L\rho_0\le1$), $R_a(m)(r)^2\le
R_a(m)(r(f))^2+3L\rho_0\le1-\sigma_{\rm rob}(m)+b/4+\sigma_{\rm rob}(m)/4$.
\end{proof}

\begin{definition}[The floor-form cone and its rate objects]\label{def:III-cone}
Let $\kappa$ be a pattern of level $L$ with class $a$.  For $r\in\mathcal
K_\kappa$, the polyhedral cone $\Gamma^\flat(\kappa,a;r)$ in the variables
$(\Delta',\gamma)\in\R^I\times\R^\Omega$ is defined by the rows (X1), (X3),
(X5) of [II,~Definition~5.52], with the coarse peaks $P_{\rm rob}(m)$, the
switching set $\Omega$ of Definition~\ref{def:III-pattern}, the ratio
variables $r$ in (X3), and the aggregated rows of (X1) on $S_l\cap(L,\infty)$
taken with the $s$-free coefficient $\alpha_l:=\|v_l1_{S_l\cap(L,\infty)}\|_1$;
and with (X4) of [II,~Definition~5.52] replaced, for every carrier $l\in
\mathcal N(m)$, by
\begin{enumerate}
\item[(X4$^0$)] $\mathrm{vs}_l\gamma_l+\Delta'_m\,m^2|r_l|\,\mathbf K\big(
\sigma_{\rm rob}(m),R_a(m)(r)^2\big)\ge0$ \ if $m$ is of kind (T);
\item[(X4)] $\mathrm{vs}_l\gamma_l/\lambda_l+\Delta'_m\ge0$ \ if $m$ is of kind
(KN) (as in [II,~Definition~5.52]);
\item[(X4$^U$)] $\mathrm{vs}_l\gamma_l\ge0$ \ if $m$ is of kind (U).
\end{enumerate}
The \emph{matrix of $\Gamma^\flat(\kappa,a;r)$} is formed in the active
variables only: the columns of the components declared inactive by (X5)
(the inactive $\gamma_l$ and the inactive shifts $\Delta'_m$) are deleted
and the (X5) rows are omitted.  The cone, as a set in the active
coordinates, and its Hoffman constant are those of $\Gamma^\flat$.  The
\emph{rate objects of $\kappa$} are the minors of this matrix and the
minors of the matrices $g_J$ formed by the
$\Delta'$-parts of the cofactor vectors of its $(n-1)$-row subsystems
\textup{(}the cofactor polynomials of [II,~Definition~5.52]\textup{)}, as functions
of $r\in\mathcal K_\kappa$.  They are continuous semialgebraic functions on
$\mathcal K_\kappa$ ($\mathbf K$ is semialgebraic: its graph is $\{(s,x,k):
(1-x)k^2-2sk+s^2-s=0,\ k>s\}$).  The pattern labels each of them
\emph{tiny} or \emph{robust} (its modulus at $r(f)$ is $\le b$, resp.\
$\ge u$; cleanness, [II,~Definition~3.18, Theorem~3.19], concerns moduli).
The \emph{exact zero set} of $\kappa$ is
\[
 Z_\kappa:=\{r\in\mathcal K_\kappa:\ \text{every tiny rate object of $\kappa$
 vanishes at}\ r\}.
\]
\end{definition}

In (X4$^U$) the shift $\Delta'_m$ is inactive, i.e.\ $0$ by (X5), so (X4$^U$) is the row (X4)
of [II] multiplied by $\lambda_l>0$.  The only modification of the cone of
[II] is therefore (X4$^0$) in the blocks of kind (T).

\begin{lemma}[The floor-form inward row]\label{lem:III-floorrow}
\textup{(a) (Meaning.)}  Let $f'$ be a first row, $m$ a block, $l$ a carrier
of block $m$ with $k(l)\in Q(f')$ and $u_l(\zh')\ne0$, $\mathrm{vs}_l:=\sgn
w'_m(k(l))$, and let $\Delta_m\in\R$, $\Delta\omega(l)\in\R$, $\gamma_l:=
\lambda_l(\Delta\omega(l)-\Delta_mw'_m(k(l)))$, $\Delta'_m:=\Delta_mM'_m$.  Then
\[
 \lambda_l\mathrm{vs}_l\Delta\omega(l)=\mathrm{vs}_l\gamma_l+\Delta'_m\lambda_l
 \varrho_l(f'),\qquad\lambda_l\varrho_l(f')=m^2|r_l|\,\mathsf k(f').
\]
So \textup{(X4$^0$)} says ``$\mathrm{vs}_l\Delta\omega(l)\ge0$, with $\mathsf k$
replaced by its floor $\mathbf K(\sigma_{\rm rob},R_a^2)$''.

\textup{(b) (The data satisfy (X4$^0$).)}  Let $m$ be of kind \textup{(T)}, $l\in
\mathcal N(m)$, and assume $b\le\sigma_{\rm rob}(m)/8$ and $\Lambda_\kappa\le
\mathrm{Des}(L)$.  For every scale $t\in\mathcal W(w)$ and every $r\in
\mathcal K_\kappa$ whose block-$m$ coordinates satisfy
$|r^{(m)}-r(f)^{(m)}|\le2\eta'$, the coarse decomposition data
$(\Delta'_{\rm dec}(t),\gamma_{\rm dec}(t))$ of [II,~Subsection~5.6] satisfy
\[
 \mathrm{vs}_l\gamma_{{\rm dec},l}+\Delta'_{{\rm dec},m}\,m^2|r_l|\,\mathbf K
 \big(\sigma_{\rm rob}(m),R_a(m)(r)^2\big)\ \ge\ -Kt .
\]

\textup{(c) (Inward direction at a companion.)}  Let $f^\#$ be a first row at
which block $m$ \textup{(}of kind \textup{(T))} has the coarse peaks
$P_{\rm rob}(m)$, every carrier of $\Omega_m$ is a strict non-peak, every
carrier of $\Omega_m\setminus\Omega_a(m)$ has value $0$, and the active ratios
\textup{(}computed with $\varkappa(\hat\zeta^\#_m;\Omega_m)$\textup{)} form a vector
$r^\#\in\mathcal K_\kappa$.  Then for $l\in\mathcal N(m)$ with $r^\#_l\ne0$,
\[
 0\le\varrho_l(f^\#)-\varrho^0_l(r^\#)\le\sigma_{\rm rob}(m)^{-3}c_{L+1}^2,
\]
and every $(\Delta',\gamma)$ satisfying \textup{(X4$^0$)} at $r^\#$ satisfies
\[
 \lambda_l\mathrm{vs}_l\Delta\omega^\#(l)\ge-|\Delta'_m|\lambda_l\sigma_{\rm rob}
 (m)^{-3}c_{L+1}^2,
\]
where $\Delta\omega^\#(l):=\gamma_l/\lambda_l+\Delta_mw^\#_m(k(l))$ and
$\Delta_m:=\Delta'_m/M^\#_m$.
\end{lemma}

\begin{proof}
(a) By [II,~Lemma~2.6(b)], $w'_m(k)=\sgn(\zeta(k))M'\min(1,\varrho_k)$, so
$w'_m(k(l))=\mathrm{vs}_lM'\varrho_l$ off $P$.  Hence $\lambda_l\Delta\omega(l)=
\gamma_l+\Delta_m\lambda_lw'_m(k(l))=\gamma_l+\Delta'_m\lambda_l\mathrm{vs}_l
\varrho_l$; multiply by $\mathrm{vs}_l$.  By Lemma~\ref{lem:III-K}(b),
$\lambda_l\varrho_l=m\Phi_l\cdot m|r_l|\mathsf k/\Phi_l$.

(b) The proof of [II,~Proposition~5.54(a)] (via [I,~Lemma~8.8] at strict
non-peaks of $f$ and [I,~(16)] at peaks of $f$) gives $\mathrm{vs}_l\gamma_{{\rm
dec},l}/\lambda_l+\Delta'_{{\rm dec},m}\ge-(4\gap(f)/t+C_fb)$ with $\gap(f)\le
bM\le b$.  Multiply by $\lambda_l$ and write
\[
 \mathrm{vs}_l\gamma_{\rm dec}+\Delta'_{\rm dec}\lambda_l\varrho^0_l(r)=\big(
 \mathrm{vs}_l\gamma_{\rm dec}+\Delta'_{\rm dec}\lambda_l\big)+\Delta'_{\rm dec}
 \lambda_l\big(\varrho^0_l(r)-1\big).
\]
Here $|\Delta'_{\rm dec}|\le C_f$ ([II,~Subsection~5.6]) and
\[
 |\varrho^0_l(r)-1|\le\Lambda_\kappa|r-r(f)|+|\varrho^0_l(r(f))-\varrho_l(f)|+
 |\varrho_l(f)-1|\le2\mathrm{Des}\,\eta'+C\sigma_{\rm rob}^{-6}(b+c_{L+1}^2)+b
\]
by Lemma~\ref{lem:III-Tblock}(a); here only the block-$m$ coordinates of
$r$ enter, because $\varrho^0_l$ depends on the ratios of block $m$ alone,
and $\Lambda_\kappa$ bounds its Lipschitz constant in these coordinates.
Since $t\ge T_{\rm lo}$, $b\le\eta'\le
T_{\rm lo}^4$, $\mathrm{Des}\,\eta'\le T_{\rm lo}^4$, $c_{L+1}\le b$ and
$\sigma_{\rm rob}^{-1}\le D(L)^2$, all terms are $\le C_f\lambda_lD(L)^{12}t\le Kt$.

(c) By Lemma~\ref{lem:III-floor} with $P_0:=P_{\rm rob}(m)$, $\varrho_l(f^\#)\ge
\varrho^0_l(r^\#)$.  As the carriers of $\Omega_m\setminus\Omega_a(m)$ have value
$0$, $R^2=R_a(m)(r^\#)^2$, and $\sigma-\sigma_{\rm rob}(m)=\Phi^2_{\text{fine
peaks}}+s_Q\le\sum_{l'>L}\Phi_{l'}^2\le c_{L+1}^2/8$.  The points $(\sigma,
R^2)$ and $(\sigma_{\rm rob},R^2)$ lie in $\mathcal K(\sigma_{\rm rob})$ because
$r^\#\in\mathcal K_\kappa$, so by Lemma~\ref{lem:III-K}(e) the two values of
$\mathbf K$ differ by at most $2c_{L+1}^2/\sigma_{\rm rob}^2$; and $m|r_l|/\Phi_l=
\varrho^0_l/\mathbf K(\sigma_{\rm rob},R^2)\le1/(2\sigma_{\rm rob})$, since
$\varrho^0_l\le\varrho_l(f^\#)<1$ and $\mathbf K\ge2s$.  This gives the bound.
For the last claim apply (a) at $f^\#$: $\lambda_l\mathrm{vs}_l\Delta\omega^\#(l)
=\big(\mathrm{vs}_l\gamma_l+\Delta'_m\lambda_l\varrho^0_l(r^\#)\big)+\Delta'_m
\lambda_l(\varrho_l(f^\#)-\varrho^0_l(r^\#))$, and the first bracket is $\ge0$
by (X4$^0$), since $\lambda_l\varrho^0_l(r^\#)=m^2|r^\#_l|\mathbf K(\sigma_{\rm rob},
R_a^2)$.
\end{proof}

\begin{remark}[Use of Lemma~\ref{lem:III-floorrow}]\label{rem:III-kind3}
Part (b) replaces, in (T)-blocks, the data estimate for the row (X4) in the
proof of [II,~Proposition~5.54(a)]; for $\Delta'_{\rm dec}<0$ only the
inequality $\varrho^0\le\varrho$ of Lemma~\ref{lem:III-floor} is needed, while for
$\Delta'_{\rm dec}>0$ the smallness of the inactive mass of a (T)-block
(Lemma~\ref{lem:III-Tblock}) is used.  Part (c) shows that at the companion
of Proposition~\ref{prop:III-KNtr} every point of the cone, and every
increment, is inward at the protected carriers up to $C_f\lambda_l
\sigma_{\rm rob}^{-3}c_{L+1}^2$; with the shift trick
this makes these coordinates of kind [3] (the kinds are defined in
[II,~Subsection~2.4]; the shift trick itself is in the proof of
[II,~Proposition~4.24(iii)]).  For the $+$ side one needs
$\gap(f)\le\frac32\gap^\#$, which holds because $\gap(f)\le Mb\le M\mu/8$ and
$\gap^\#\ge M\mu/4$ (Proposition~\ref{prop:III-KNtr} and
Lemma~\ref{lem:III-drift}); the inequality ``$Mb\le\gap^\#/4$'' used in an
earlier version is not needed and is not claimed.  The assembly with the
true shift is carried out in Section~\ref{sec:III-routeV}.
\end{remark}

\begin{lemma}[Rescaling in (U)-blocks]\label{lem:III-Uscale}
Let $m$ be a block of kind \textup{(U)} with $\Omega_a(m)\ne\emptyset$, and
write $r=(r^{(m)},r^{\rm rest})$ with $r^{(m)}:=(r_l)_{l\in\Omega_a(m)}$.
\begin{enumerate}
\item[(a)] Every minor of the matrix of $\Gamma^\flat(\kappa,a;r)$ is
homogeneous of degree $0$ or $1$ in $r^{(m)}$, and every cofactor polynomial
is homogeneous in $r^{(m)}$ of a degree $e\le n$, where $n$ is the number of
variables of the cone.  Hence for $s\in(0,1]$ the substitution $r^{(m)}\mapsto
sr^{(m)}$ preserves the vanishing of every rate object and multiplies every
rate object by $s^e\ge s^n$; it maps $Z_\kappa$ into itself and does not
change the functions $\varrho^0_d$, $d\in\mathcal N_T$.
\item[(b)] If at a first row the carriers of $\Omega_m$ are strict non-peaks
and the active ratios of block $m$ are $sr^{(m)}$, then every active status
of block $m$ is at most $s$ times the status obtained, with the same
$\sigma$ and the same inactive ratios, at the ratios $r^{(m)}$.
\end{enumerate}
\end{lemma}

\begin{proof}
(a) The shift $\Delta'_m$ is inactive, so its column is not part of the
matrix (Definition~\ref{def:III-cone}).  Then $r^{(m)}$ occurs only in the
row (X3)$_m$, $\sum_{l\in\Omega_a(m)}r_l\gamma_l=0$, whose entries are linear
in $r^{(m)}$: the rows (X4$^U$) contain no ratio, the rows (X3) and
(X4$^0$) of the other blocks contain only their own ratios, and the (X1)
rows contain no ratio at all (their entries are $u_l(j)$, $-\tau_m(j)$,
$\alpha_l$ and $-\mathrm{vs}_l\lambda_l\alpha_l$); the ratios enter only
through (X3) and (X4$^0$).  A minor is a sum
of products with one entry from each of its rows, hence homogeneous of
degree $1$ if it uses the row (X3)$_m$ and of degree $0$ otherwise.  The
cofactor vectors of an $(n-1)$-row subsystem $J$ have as components
$(n-1)$-minors, so each column of $g_J$ is homogeneous of a degree
$e_J\in\{0,1\}$, and each minor of $g_J$ is homogeneous of degree $\sum e_J\le
n$.  The functions $\varrho^0_d$ ($d\in\mathcal N_T$) involve only ratios of
blocks of kind (T).  (b) By Lemma~\ref{lem:III-K}(b),(c), $\varrho_l=m|r_l|\mathbf
K(\sigma,R^2)/\Phi_l$; replacing $r^{(m)}$ by $sr^{(m)}$ multiplies $|r_l|$ by $s$
and does not increase $R^2$, and $\mathbf K$ is increasing in $R^2$.
\end{proof}

\begin{lemma}[Weight-freeness]\label{lem:III-weightfree}
Let $\kappa$ be a pattern of level $L$ with class $a$, and for $d\in\mathcal
N_T$ put
\[
 \psi_d(r):=m(d)\,|r_d|\,\mathbf K\big(\sigma_{\rm rob}(m(d)),R_a(m(d))(r)^2\big)
 \qquad(r\in\mathcal K_\kappa),
\]
so that $\varrho^0_d=\psi_d/\Phi_d$.
\begin{enumerate}
\item[(a)] Every entry of the matrix of $\Gamma^\flat(\kappa,a;r)$, every rate
object of $\kappa$, the region $\mathcal K_\kappa$, the set $Z_\kappa$ and the
functions $\psi_d$ are defined, by explicit formulas, from: the values
$u_l(j)$ of target vectors and the coefficients $\alpha_l$; signs and the
combinatorial data of $\kappa$; the weights of the robust peaks of the
blocks with an active component; the weights of the carriers of the blocks of kind (KN); and
the variables $r$.  In particular no weight $\Phi_d$ with $d\in\mathcal N_T$
occurs in any of them.
\item[(b)] The only row of the cone of [II,~Definition~5.52] that contains
the weight of a near-threshold switching carrier is the row (X4), through
$1/\lambda_l$; in $\Gamma^\flat$ it is kept only in blocks of kind (KN).  An
active weak peak $d$ of $f$ belongs to $\Omega$, so its weight occurs neither
in $\tau_{m(d)}$ nor in its aggregated signature row.
\item[(c)] For the clean sub-window $w$, let $\Gamma^\#$ be the cone with
the aggregated coefficients
\[
 \|v_l1_{S_l\cap(L,s_{\rm far}(w)]}\|_1
\]
actually used at $w$ ([II,~Definition~5.52], [II,~Lemma~5.50]).  Then
$\Gamma^\#$ arises
from $\Gamma^\flat$ by multiplying each aggregated row by a factor $t_l\in
[1-\varepsilon_w,1]$, where
\[
 \varepsilon_w:=\max_l2^{1+s_1(l)-s_{\rm far}(w)}
\]
and $s_1(l)$ is the least element of $S_l$ beyond $L$.  Every rate object of
$\Gamma^\#$ equals the corresponding one of $\Gamma^\flat$ times a product of at
most $n^2$ factors $t_l$; in particular the zero sets coincide, and robust
values are multiplied by a factor $\ge(1-\varepsilon_w)^{n^2}\ge\frac12$.
\end{enumerate}
\end{lemma}

\begin{proof}
(a) Row by row.  (X1): the column of $\gamma_l$ has entries $u_l(j)$; the
column of $\Delta'_m$ has entries $-\tau_m(j)$, with $\tau_m$ as in
Lemma~\ref{lem:III-EC} (robust peak weights); an aggregated row has the single
entry $\alpha_l$ (multiplying $\gamma_l$ if $k(l)\in\Omega$) or $-\mathrm{vs}_l
\lambda_l\alpha_l$ (multiplying $\Delta'_m$ if $l\in P_{\rm rob}(m)$).  (X3) in
ratio variables: entries $1$ and $-r_l$ (the (X5) rows and the columns of
the inactive variables are not part of the matrix,
Definition~\ref{def:III-cone}).  (X4$^0$): entries
$\mathrm{vs}_l$ and $m^2|r_l|\mathbf K(\sigma_{\rm rob}(m),R_a(m)(r)^2)$, with
$\sigma_{\rm rob}(m)$ a sum of squares of robust peak weights.  (X4) in
(KN)-blocks: entries $\mathrm{vs}_l/\lambda_l$ and $1$, with $l$ a carrier of a
block of kind (KN).  (X4$^U$): entries $\mathrm{vs}_l$.  Minors and cofactor
polynomials are polynomials in these entries; $\mathcal K_\kappa$, $Z_\kappa$ and
$\psi_d$ are defined from them and from $\sigma_{\rm rob}$.  A carrier $d\in
\mathcal N_T$ is not a robust peak, and it is not a carrier of a block of
kind (KN), since each carrier belongs to exactly one block.  (b) is read off
from the list in (a) and Definition~\ref{def:III-pattern}(a).  (c) For $l$
with $S_l\cap(L,\infty)\ne\emptyset$, $\alpha_l\ge v_l(s_1(l))=\delta_l2^{-s_1(l)}/
n_l$ and $\alpha_l-\|v_l1_{S_l\cap(L,s_{\rm far}]}\|_1\le\sum_{s>s_{\rm far}}
\delta_l2^{-s}/n_l\le2\delta_l2^{-s_{\rm far}}/n_l$, so the actual coefficient
is $t_l\alpha_l$ with $t_l$ as stated.  Multiplying a row by $t_l>0$
multiplies every minor containing it by $t_l$, every cofactor vector of a
subsystem containing it by $t_l$, and hence every minor of $g_J$ by a
product of at most $n$ such column factors, each a product of at most $n$
factors $t_l$.  Finally $s_{\rm far}(w)$ is the least $s$ with $2^{-s}\le
c_{L+1}^2$ ([II,~Lemma~5.50]); the weight bound (W5) of the design
(Section~\ref{sec:III-design}) gives $c_{L+1}\le2^{-\min S_{L+1}}=2^{-3\cdot
2^{L+1}}$, so $s_{\rm far}(w)\ge12\cdot2^L$, while $s_1(l)\le L+3\cdot2^L$ for
$l\le L$.  So $\varepsilon_w\le2^{1+L-9\cdot2^L}$, and $n\le2L$
([II,~Lemma~5.53]) gives $(1-\varepsilon_w)^{n^2}\ge1-4L^2\varepsilon_w\ge\frac12$.
\end{proof}

\begin{remark}[Why the $s$-free coefficients]\label{rem:III-sfar}
The rate objects of level $L$ must be computed before the sub-windows of
level $L$, while $s_{\rm far}(w)$ depends on $c_{L+1}$, which is chosen after
them.  Lemma~\ref{lem:III-weightfree}(c) removes this forward reference: the
rate objects are those of $\Gamma^\flat$, and the cone actually used at $w$ has
the same zero sets, robust values within a factor $2$, and Hoffman and
{\L}ojasiewicz constants within a factor $2$.  The same one-line change
applies to the rate objects of [II,~Definition~5.52].
\end{remark}

\subsection{Generic weights}\label{subsec:III-NLMGEN}

Fix a pattern $\kappa$ of level $L$ with $\mathcal N_T\ne\emptyset$, and let
$\psi_d$ ($d\in\mathcal N_T$) be as in Lemma~\ref{lem:III-weightfree}.  For
$\Phi=(\Phi_d)_{d\in\mathcal N_T}\in(0,\infty)^{\mathcal N_T}$ put
\[
 F_\Phi(r):=\max_{d\in\mathcal N_T}\frac{\psi_d(r)}{\Phi_d}\qquad(r\in
 Z_\kappa).
\]
At the true weights, $F_\Phi(r)=\max_{d\in\mathcal N_T}\varrho^0_d(r)$ is the
largest floor status of the protected carriers.

\begin{definition}[Coincidences]\label{def:III-coinc}
$\Phi\in(0,\infty)^{\mathcal N_T}$ is a \emph{coincidence} for $\kappa$ if some
$r_0\in Z_\kappa$ with $F_\Phi(r_0)=1$ is a (not necessarily strict) local
minimum of $F_\Phi$ on $Z_\kappa$.  $B_\kappa\subseteq(0,\infty)^{\mathcal N_T}$
denotes the set of coincidences.
\end{definition}

By Lemma~\ref{lem:III-weightfree}(a), $Z_\kappa$ and the $\psi_d$ do not depend
on $\Phi$; this is what makes the following two lemmas possible.

\begin{lemma}[The threshold is not a local-minimum value]\label{lem:III-NLM}
$B_\kappa$ is a semialgebraic set of dimension $<|\mathcal N_T|$.
\end{lemma}

\begin{proof}
$Z_\kappa$ is a compact semialgebraic set and the $\psi_d$ are continuous
semialgebraic functions on it, all independent of $\Phi$.  The set
\[
 B_\kappa=\big\{\Phi:\ \exists r_0\in Z_\kappa\ \big[F_\Phi(r_0)=1\ \wedge\
 \exists\varepsilon>0\ \forall r\in Z_\kappa\ (|r-r_0|<\varepsilon\Rightarrow
 F_\Phi(r)\ge1)\big]\big\}
\]
is first-order definable, hence semialgebraic (Tarski--Seidenberg \cite{Tarski,Seidenberg}).

Take a finite partition of $Z_\kappa$ into connected Nash submanifolds $S$
compatible with the sets $\{r_d>0\}$, $\{r_d=0\}$, $\{r_d<0\}$, $d\in\mathcal
N_T$ \cite[Proposition~9.1.8]{BCR}.  On a stratum $S$ each $|r_d|$ equals
$r_d$, $-r_d$ or $0$, and $R_a^2<1$ on $\mathcal K_\kappa$, so each $\psi_d|_S$ is
a Nash function (Lemma~\ref{lem:III-K}(d)).  For a nonempty $D\subseteq\mathcal
N_T$ and a stratum $S\subseteq\{r_d\ne0:d\in D\}$ let $\Psi_{D,S}:=(\psi_d)_{d\in
D}|_S:S\to\R^D$ and let $CV_{D,S}$ be its set of critical values; by the
semialgebraic Sard theorem \cite[Theorem~9.6.2]{BCR}, $\dim CV_{D,S}<|D|$
(this includes the case $\dim S<|D|$, where every point is critical).

Let $\Phi\in B_\kappa$, with a witness $r_0$ lying in the stratum $S$, and
$D:=\{d:\psi_d(r_0)=\Phi_d\}$.  $D$ is nonempty because $F_\Phi(r_0)=1$, and
$r_{0,d}\ne0$ for $d\in D$ because $\psi_d(r_0)=\Phi_d>0$.  A local minimum of
$F_\Phi$ on $Z_\kappa$ is a local minimum on $S$.  Suppose $r_0$ were a regular
point of $\Psi_{D,S}$.  Then there is $v\in T_{r_0}S$ with $d\psi_d(r_0)v=-1$
for all $d\in D$, and a Nash arc $\gamma$ in $S$ with $\gamma(0)=r_0$, $\gamma'(0)=v$.
For small $s>0$, $\psi_d(\gamma(s))=\Phi_d-s+o(s)<\Phi_d$ for $d\in D$, and
$\psi_d(\gamma(s))<\Phi_d$ for $d\notin D$ by continuity; so $F_\Phi(\gamma(s))<1$,
contradicting the local minimality.  Hence $(\Phi_d)_{d\in D}\in CV_{D,S}$, and
\[
 B_\kappa\subseteq\bigcup_{D,S}\big\{\Phi:\ (\Phi_d)_{d\in D}\in CV_{D,S}\big\},
\]
a finite union of semialgebraic sets of dimension $<|D|+|\mathcal N_T\setminus
D|=|\mathcal N_T|$.
\end{proof}

\begin{remark}\label{rem:III-NLMcases}
The proof covers isolated points of $Z_\kappa$ (zero-dimensional strata, at
which every point is a local minimum: for an isolated point $r^*$ the
corresponding $\Phi$ satisfy $\max_d\psi_d(r^*)/\Phi_d=1$, a set of codimension
one), kinks of the maximum
(through the active set $D$), minima on the boundary of $\mathcal K_\kappa$
(boundary strata are strata of $Z_\kappa$), non-strict minima, and the
non-smoothness of $|r_d|$ at $r_d=0$ (witnesses have $r_d\ne0$ for $d\in D$,
and for $d\notin D$ only continuity is used).  If the region were defined
with bounds involving the weights $\Phi_d$, $d\in\mathcal N_T$, the set
$Z_\kappa$ would depend on $\Phi$ and the argument would fail; this is why
$\mathcal K_\kappa$ is defined by the robust peak weights only.
\end{remark}

\begin{lemma}[Algebraically independent weights avoid all coincidences]\label{lem:III-GEN}
Let $k_0\subseteq\R$ be a subfield containing all values $u_l(j)$ of the target
vectors and all coefficients $\alpha_l$, let $\kappa$ be a pattern with
$\mathcal N_T\ne\emptyset$, and let $k:=k_0(\Phi_l:\ l\le L,\ l\notin\mathcal
N_T)$.  If the weights $(\Phi_d)_{d\in\mathcal N_T}$ are algebraically
independent over $k$, then $(\Phi_d)_{d\in\mathcal N_T}\notin B_\kappa$.
\end{lemma}

\begin{proof}
Put $n:=|\mathcal N_T|$.  By Lemma~\ref{lem:III-weightfree}(a), $\mathcal K_\kappa$,
$Z_\kappa$ and the $\psi_d$ are defined by first-order formulas whose
parameters are values $u_l(j)$, coefficients $\alpha_l$, rational numbers
(signs; the polynomial defining the graph of $\mathbf K$) and weights of
carriers outside $\mathcal N_T$; so $B_\kappa$ is defined by a first-order
formula $\varphi$ with parameters in $k$.  Let $k^{\rm rc}\subseteq\R$ be the
field of real numbers algebraic over $k$; it is real closed.  Let $B_0\subseteq
(k^{\rm rc})^n$ be the set defined by $\varphi$ in $k^{\rm rc}$.  By the
Tarski--Seidenberg transfer, $B_\kappa$ is the extension of $B_0$ to $\R$, and
$\dim B_0=\dim B_\kappa<n$ by Lemma~\ref{lem:III-NLM} and the invariance of
dimension under extension of real closed fields \cite[Chapter~5]{BCR}.  The
Zariski closure of $B_0$ has the same dimension \cite[Section~2.8]{BCR}, so
some nonzero $P\in k^{\rm rc}[X_1,\dots,X_n]$ vanishes on $B_0$; by transfer
of the sentence $\forall x\,(\varphi(x)\Rightarrow P(x)=0)$, $P$ vanishes on
$B_\kappa$.  The coefficients of $P$ lie in a finite extension $k'\subseteq
k^{\rm rc}$ of $k$.  Let $\sigma_1=\mathrm{id},\sigma_2,\dots,\sigma_q$ be the
$k$-embeddings of $k'$ into an algebraic closure $\bar k\supseteq k'$, and
$\mathcal NP:=\prod_i\sigma_i(P)$.  Its coefficients are fixed by every
automorphism of $\bar k$ over $k$, hence lie in $k$ (characteristic $0$);
$\mathcal NP\ne0$ as a product of nonzero polynomials; and $\mathcal
NP(x)=P(x)\prod_{i\ge2}\sigma_i(P)(x)=0$ for every $x\in B_\kappa$.  If
$(\Phi_d)_{d\in\mathcal N_T}$ were in $B_\kappa$, then $\mathcal NP((\Phi_d)_d)=0$
would be a nontrivial algebraic relation over $k$.
\end{proof}

\begin{remark}[Order of choices]\label{rem:III-GENorder}
If a family $(c_l)_{l\in\Lambda}$ is algebraically independent over $k_0$, then
for disjoint finite sets $A,B\subseteq\Lambda$ the family $(c_a)_{a\in A}$ is
algebraically independent over $k_0(c_b:b\in B)$, and the same holds for
$\Phi_l=2^{-m(l)-k(l)}c_l$.  So Lemma~\ref{lem:III-GEN} applies to every
pattern simultaneously, whatever the order in which the weights were chosen;
what matters is that $k_0$ contains no weight.  In particular, no datum of
the row $f$ enters $B_\kappa$: $f$ only selects the pattern $\kappa$ and the
point $r(f)$.  This is why weights fixed in the design, before $f$, avoid
the coincidence sets of every row.
\end{remark}

\subsection{Quantitative coherence}\label{subsec:III-QC}

\begin{lemma}[Quantitative coherence]\label{lem:III-QC}
Let $\kappa$ be a pattern with $\mathcal N_T\ne\emptyset$ and $Z_\kappa\ne
\emptyset$, and let $\Phi=(\Phi_d)_{d\in\mathcal N_T}\notin B_\kappa$.  Write $F:=
F_\Phi$, $Y_1:=\{r\in Z_\kappa:F(r)\le1\}$ and, for $h>0$ and $r\in Z_\kappa$
\textup{(}$|\cdot|$ the Euclidean norm\textup{)},
\[
 G_h(r):=\inf\{F(r''):\ r''\in Z_\kappa,\ |r''-r|<h\} .
\]
\begin{enumerate}
\item[(a)] $\mu_\kappa(h):=1-\max_{r\in Y_1}G_h(r)>0$ for every $h>0$
\textup{(}$\mu_\kappa:\equiv1$ if $Y_1=\emptyset$\textup{)}.  For every $r\in Y_1$
there is $r''\in Z_\kappa$ with $|r''-r|<h$ and $F(r'')\le1-\frac34\mu_\kappa(h)$.
\item[(b)] If $Y_1=\emptyset$, then $F\ge1+m_0$ on $Z_\kappa$ for some $m_0>0$.
If $Y_1\ne\emptyset$, there are $C_Y>0$ and $\alpha_Y\in(0,1]$ with
$\dist(r,Y_1)\le C_Y(F(r)-1)_+^{\alpha_Y}$ for all $r\in Z_\kappa$.
\item[(c)] $h\mapsto\mu_\kappa(h)$ is semialgebraic, positive, nondecreasing and
$\le1$, and there are $c_\kappa,\beta_\kappa,h_\kappa>0$ with $\mu_\kappa(h)\ge
c_\kappa h^{\beta_\kappa}$ for $0<h\le h_\kappa$.
\end{enumerate}
\end{lemma}

\begin{proof}
(a) $G_h$ is upper semicontinuous on $Z_\kappa$: if $r_n\to r$ and $r''\in
Z_\kappa$ has $|r''-r|<h$, then $|r''-r_n|<h$ for large $n$, so $\limsup_n
G_h(r_n)\le F(r'')$, and taking the infimum over $r''$ gives $\limsup_nG_h(r_n)
\le G_h(r)$.  $Y_1$ is compact.  For $r\in Y_1$: if $F(r)<1$, then $G_h(r)\le
F(r)<1$; if $F(r)=1$, then $r$ is not a local minimum of $F$ on $Z_\kappa$
(as $\Phi\notin B_\kappa$), so there are points $r''\in Z_\kappa$ arbitrarily close
to $r$ with $F(r'')<1$, and $G_h(r)<1$.  An upper semicontinuous function
attains its maximum on a compact set, so $\max_{Y_1}G_h<1$.  The second
statement follows from $G_h(r)\le1-\mu_\kappa(h)<1-\frac34\mu_\kappa(h)$ and the
definition of the infimum (which need not be attained on the open ball).

(b) If $Y_1=\emptyset$, $F-1$ is positive and continuous on the compact set
$Z_\kappa$.  Otherwise $f_1:=\dist(\cdot,Y_1)$ and $f_2:=(F-1)_+$ are continuous
semialgebraic functions on the compact semialgebraic set $Z_\kappa$ with
$f_2^{-1}(0)=Y_1=f_1^{-1}(0)$; by the {\L}ojasiewicz inequality
\cite[Corollary~2.6.7]{BCR} there are an integer $N_Y\ge1$ and $C>0$ with
$f_1^{N_Y}\le Cf_2$ on $Z_\kappa$.  Put $\alpha_Y:=1/N_Y$, $C_Y:=C^{1/N_Y}$.

(c) $\mu_\kappa$ is defined by a first-order formula (with the parameters of
$\kappa$ and $\Phi$), hence semialgebraic; it is positive by (a),
nondecreasing because $G_h$ is nonincreasing in $h$, and $\le1$ because
$F\ge0$.  A semialgebraic function of one variable that is positive on
some $(0,h_1)$ has, near $0^+$, a convergent Puiseux expansion $\mu_\kappa(h)=
ch^{p/q}(1+o(1))$ with $c>0$ and $p/q\ge0$ rational (the real closure of
$\R(t)$ is the field of algebraic Puiseux series; \cite[Chapter~7]{vdDries}).
Take $c_\kappa:=c/2$, $\beta_\kappa:=\max(p/q,\,1)$ and $h_\kappa\in(0,1]$ small.
\end{proof}

\begin{remark}[Definite constants]\label{rem:III-QCchoice}
The design needs definite values of the constants of Lemma~\ref{lem:III-QC};
for instance $N_Y$ is the least admissible integer, $C_Y$ the least admissible
constant for it, $p/q$ the Puiseux exponent, $c_\kappa$ half the leading
coefficient and $h_\kappa:=\min\{1,\sup\{h':\mu_\kappa\ge c_\kappa h^{\beta_\kappa}
\text{ on }(0,h']\}\}$; the same device fixes the constants $C_L$, $N_L$ of
Lemma~\ref{lem:III-loj}.  Only their existence is used.  The normalization of
(b) is $\dist\le C_Y(F-1)_+^{\alpha_Y}$; it enters the threshold (D2) of
Definition~\ref{def:III-QCreq} in exactly this form.  The exponent
$\beta_\kappa$ is not explicit, and $\mu_\kappa(h)$ may be much smaller than
$h^p$ for any exponent $p$ fixed in advance (for instance $p=4$); this is the reason for the exact holds in
Proposition~\ref{prop:III-KNtr} (Remark~\ref{rem:III-KN1}).
\end{remark}

\begin{lemma}[{\L}ojasiewicz displacement on a compact semialgebraic set]\label{lem:III-loj}
Let $\mathcal K\subseteq\R^n$ be compact and semialgebraic and let $\pi_1,\dots,
\pi_s$ be continuous semialgebraic functions on $\mathcal K$.  There are
$C_L\ge1$ and an integer $N_L\ge1$ such that for every $S\subseteq\{1,\dots,
s\}$, every $\beta\in[0,1]$ and every $v\in\mathcal K$ with $\max_{i\in S}
|\pi_i(v)|\le\beta$, at least one of the following holds: \textup{(i)} $\beta\ge
1/C_L$; \textup{(ii)} there is $v'\in\mathcal K$ with $\pi_i(v')=0$ for all $i\in
S$ and $|v'-v|\le C_L\beta^{2/N_L}$ \textup{(}Euclidean norm\textup{)}.
\end{lemma}

\begin{proof}
The proof of [II,~Lemma~5.3] applies verbatim with $[-2,2]^n$ replaced by
$\mathcal K$ and $Z_S:=\Pi_S^{-1}(0)\cap\mathcal K$, $\Pi_S:=\sum_{i\in S}\pi_i^2$:
the {\L}ojasiewicz inequality \cite[Corollary~2.6.7]{BCR} holds for continuous
semialgebraic functions on a compact semialgebraic set, and a nearest point
of $Z_S$ lies in $\mathcal K$.
\end{proof}

\subsection{Exactification with threshold protection}\label{subsec:III-KNtr}

\begin{definition}[Requirements on the design at level $L$]\label{def:III-QCreq}
Let $\mathfrak P(L)$ be the set of patterns of level $L$; under option (b)
($c_1=1$) only the patterns with $1\notin\mathcal N_T(\kappa)$ are included.
For $\kappa\in\mathfrak P(L)$ let $\mathrm{Lip}^*_\kappa$ be the largest of: the
Lipschitz constants on $\mathcal K_\kappa$ of the rate objects of $\kappa$; the
constant $\Lambda_\kappa$ of Definition~\ref{def:III-region} (which dominates
the Lipschitz constant of $F_\Phi$); and the Lipschitz constants of the
functions $r\mapsto m|r_l|\mathbf K(\sigma,R_a(m)(r)^2+\rho)/\Phi_l$, $l\in\Omega_a$,
on $\{r:R_a(m)(r)^2+\rho\le1-\sigma_{\rm rob}(m)/8\}$, uniformly in $\sigma\in
[\sigma_{\rm rob}(m)/2,1]$ and $\rho\ge0$ (these are the actual statuses at fixed
$\sigma$ and fixed inactive ratios).  Let $C_L(\kappa)$, $N_L(\kappa)$ be the
constants of Lemma~\ref{lem:III-loj} for the rate objects of $\kappa$ on
$\mathcal K_\kappa$, and, when $\mathcal N_T(\kappa)\ne\emptyset$ and
$Z_\kappa\ne\emptyset$, let $m_0,C_Y,\alpha_Y,\mu_\kappa,c_\kappa,\beta_\kappa,h_\kappa$
be the constants of Lemma~\ref{lem:III-QC} at the actual weights (with the
conventions $m_0:=\infty$ if $Y_1\ne\emptyset$, $C_Y:=0$ if $Y_1=\emptyset$, and
the choices of Remark~\ref{rem:III-QCchoice}).  The design satisfies the
\emph{QC requirements at level $L$} if:
\begin{enumerate}
\item[(QC1)] \textup{(genericity)} there is a countable subfield $k_0\subseteq\R$,
containing every value $u_l(j)$ and every coefficient $\alpha_l$, over which
the family $(c_l)_{l\ge2}$ (option (b)), resp.\ $(c_l)_{l\ge1}$ (option (a)),
is algebraically independent;
\item[(QC2)] \textup{(rate objects)} the rate objects of every $\kappa\in
\mathfrak P(L)$ belong to the rate scheme of level $L$, with values computed
at the ratio vector of the row;
\item[(QC3)] \textup{(constants)} $\mathrm{Des}(L)\ge D(L)\ge\Phi_{\min}(L)^{-1}$,
$\mathrm{Des}(L)\ge C_*(L)$, and for every $\kappa\in\mathfrak P(L)$:
$\mathrm{Des}(L)\ge\max\{\mathrm{Lip}^*_\kappa,C_L(\kappa),N_L(\kappa),C_Y,1/m_0,
1/c_\kappa,\beta_\kappa,1/h_\kappa\}$; here $C_*(L):=|T(L)|+2$ is the bound,
in units of values, of [II,~Definition~3.6] (see also (A1) of
[II,~Theorem~5.7] and Step~0 of the proof of [II,~Proposition~5.56]): the
closing moves (C1), (C2) of the assembly ([II,~Subsection~4.3]) change
every value $\mathrm{val}_{l''}$ by at most $C_*(L)b$.  We put
$C_{**}(L):=\mathrm{Des}(L)^3$; for $L\ge l_f$ the closing moves change the
active ratio vector (in the Euclidean norm) and every status of a coarse
carrier by at most $C_{**}(L)b$ (Step~0 of the proof of
Proposition~\ref{prop:III-KNtr});
\item[(QC4)] \textup{(thresholds)} for every sub-window $w$ of $L$, with
$\mathrm{Lip}^*:=\max_\kappa\mathrm{Lip}^*_\kappa$, $C_L:=\max_\kappa C_L(\kappa)$,
$N_L:=\max_\kappa N_L(\kappa)$, $\beta(b):=(1+\mathrm{Lip}^*C_{**}(L))b$ and
\[
 \delta_0(b):=\beta(b)+C_0D(L)^{12}b+\mathrm{Lip}^*C_L\beta(b)^{2/N_L}
\]
($C_0$ twice the numerical constant of Lemma~\ref{lem:III-Tblock}(a)), the
number $b=b(w)$ satisfies $b\le\eta'(w)$, $\beta(b)<1/C_L$,
$\beta(b)\le\eta'(w)/3$ and
\begin{enumerate}
\item[(D1)] $C_L\beta(b)^{2/N_L}\le\eta'(w)/3$;
\item[(D2)] for every $\kappa\in\mathfrak P(L)$ with $\mathcal N_T\ne\emptyset$:
$\delta_0(b)\le m_0/2$, $C_Y\delta_0(b)^{\alpha_Y}\le\eta'(w)/3$, and $b\le
\mu_\kappa(\eta'(w)/3)/8$;
\end{enumerate}
\item[(QC5)] \textup{(weights after $L$)} $c_{L+1}\le b(w)^2$ for every
sub-window $w$ of $L$.
\end{enumerate}
\end{definition}

\begin{remark}[Consistency of the requirements]\label{rem:III-QCwell}
(i) The constants in (QC3) are finite: $\Lambda_\kappa$ and the Lipschitz
constants of the status functions in $\mathrm{Lip}^*_\kappa$ are $\le CL\,D(L)^9$
by Lemma~\ref{lem:III-K}(e) and Lemma~\ref{lem:III-alpha}(b); the
Lipschitz constants of the rate objects (minors of matrices of size at most
$2L$ whose entries are Lipschitz on the compact set $\mathcal K_\kappa$) are
finite; the {\L}ojasiewicz constants exist by
Lemma~\ref{lem:III-loj}; the QC constants by Lemma~\ref{lem:III-QC}, which
applies because, by (QC1) and Lemma~\ref{lem:III-GEN}
(Remark~\ref{rem:III-GENorder}), the actual weights are not a coincidence
for any $\kappa\in\mathfrak P(L)$.  A pattern of level $L$ involves only
carriers $\le L$, so only $c_1,\dots,c_L$ enter; these are chosen before the
level-$L$ data.  (ii) In (QC4) the left sides are continuous, increasing in
$b$ and tend to $0$ as $b\to0$, while the right sides are positive numbers
computed before $b(w)$; so the admissible $b$ form an interval $(0,b_{\rm
QC}(w)]$ with $b_{\rm QC}(w)>0$.  The clause $b\le\mu_\kappa(\eta'/3)/8$ is
positive exactly because no pattern is a coincidence.  The number $b(w)$
enters later quantities only through upper bounds, so the non-explicit
exponents $1/\alpha_Y$, $N_L/2$ need not be dominated by $\mathrm{Des}$.
(iii) $\eta'/3\le1/\mathrm{Des}(L)\le h_\kappa$ by (QC3), so Lemma~\ref{lem:III-QC}(c)
gives $\mu_\kappa(\eta'/3)\ge c_\kappa(\eta'/3)^{\beta_\kappa}\ge\mathrm{Des}^{-1}
(\eta'/3)^{\mathrm{Des}}$.  (iv) That $T_{\rm fin}$ satisfies (QC1)--(QC5) at
every main stage is part of its construction in Section~\ref{sec:III-design}
((QC1) is Lemma~GEN$_{\rm fin}$ there, (QC5) is the weight bound (W2)).
\end{remark}

\begin{lemma}[Common pull window]\label{lem:III-window}
In the situation of [II,~Lemma~4.15], let $\eta_*$ be a number with
$\|x\|_\infty\le\eta_*\le c_T$, and choose the pull coordinate $j_{l''}$ as the
largest element of $S_{l''}$ with $v_{l''}(j_{l''})\ge\eta_*$.  Then the
conclusions of [II,~Lemma~4.15] hold with $\varepsilon_x$ replaced by $\eta_*$;
in particular $v_{l''}(j_{l''})\in[\eta_*,2^{G(l)}\eta_*)$ for every $l''\in L_0$,
the banks $\mathfrak n_{l''}$ are positive, and the pull masses are $\mathfrak
p_{l''}=24\lambda_{l''}v_{l''}(j_{l''})\ge24\lambda_{l''}\eta_*$.  For $\eta_*:=
\max(\|x\|_\infty,c_{L+1})$ the pulls also satisfy $\mathfrak p_{l''}\ge24
\lambda_{l''}c_{L+1}$.
\end{lemma}

\begin{proof}
The proof of [II,~Lemma~4.15] uses the level $\varepsilon_x$ only through the
inequality $|x_{l''}|\le\varepsilon_x$ and through the window bounds of the
pull coordinates; both hold with $\eta_*$ in place of $\varepsilon_x$.  For the
positivity of the banks we recall the mechanism.  After the pulls, the
bank equation of carrier $l''$ reads $\mathfrak n_{l''}=(\Delta_{l''}\mathcal N+
\beta_{l''}(\mathcal N-\nu_p))/((\mu^\star_{j'})^2v_{l''}(j'))$, where $\Delta_{l''}=
x_{l''}+2v_{l''}(j_{l''})+O(C_T\eta_*^2)$ is the correction still to be made,
$\beta_{l''}=\epsilon_{l''}\ip{U^*u_{l''}}{e^p}$ has either sign, $\nu_p$ is
the Hilbert norm after the pulls and $\mathcal N\ge\nu_p$ the norm after all
banks ([II,~Lemma~4.14]), with $\mathcal N-\nu_p\le C_T\eta_*^2$.  Since
$v_{l''}(j_{l''})\ge\eta_*\ge|x_{l''}|$, $\Delta_{l''}\ge\eta_*/2$ and hence
$\mathfrak n_{l''}\ge c_f\eta_*>0$ for $\eta_*\le c_T$.
\end{proof}

\begin{remark}[Why a common window]\label{rem:III-window}
If the pull of a carrier is placed at a level adapted to its own first-order
correction (for instance at the larger of that correction and $c_{L+1}$), a
carrier whose correction is tiny (such as a nearly neutral carrier that is
to be held at value $0$) gets $\Delta_{l''}$ of the order of $c_{L+1}$, while
$\beta_{l''}(\mathcal N-\nu_p)$ is of the order of $\mathrm{Des}\,\eta'^2$, created
by the other carriers' banks; if $\beta_{l''}<0$, the bank equation then
demands a negative mass, which is not admissible.  The common window
$[\eta_*,2^{G(l)}\eta_*)$ of Lemma~\ref{lem:III-window} removes this.  The lower
bound $\mathfrak p\ge24\lambda c_{L+1}$ is the one needed for the side condition
(U3) at the supports created by the companion (Section~\ref{sec:III-routeV}).
\end{remark}

\begin{lemma}[Statuses at exact ratios]\label{lem:III-status}
Let $m$ be a block of kind \textup{(T)} and $f^\#$ a first row such that:
\textup{(i)} the coarse carriers of block $m$ outside $\Omega_m$ are the
robust peaks $P_{\rm rob}(m)$, and they are peaks of $f^\#$; \textup{(ii)}
every carrier of $\Omega_m\setminus\Omega_a(m)$ has value $0$ at $f^\#$;
\textup{(iii)} $\varkappa^\#:=\varkappa(\hat\zeta^\#_m;\Omega_m)>0$ and the active
ratios $r^\#$ of block $m$ at $f^\#$ lie in $\mathcal K_\kappa$; \textup{(iv)}
$\varrho^0_l(r^\#)+\sigma_{\rm rob}(m)^{-3}c_{L+1}^2<1$ for every $l\in\Omega_a(m)$.
Then every carrier of $\Omega_m$ is a strict non-peak of $f^\#$, and
$\varrho^0_l(r^\#)\le\varrho_l(f^\#)\le\varrho^0_l(r^\#)+\sigma_{\rm rob}(m)^{-3}
c_{L+1}^2$ for every $l\in\Omega_a(m)$.
\end{lemma}

\begin{proof}
Apply Lemma~\ref{lem:III-K} at $f^\#$ with $\Omega:=\Omega_m$.  By (i), $P\setminus
\Omega_m$ consists of $P_{\rm rob}(m)$ and of fine peaks, and $Q\setminus\Omega_m$
of fine carriers, so $\sigma'\le\sigma_0+c_{L+1}^2/8$ with $\sigma_0:=\sigma_{\rm
rob}(m)$ (the terms $1-\varrho_k^2$, $k\in\Omega_m\cap P$, are $\le0$).  By (ii),
$R^2=R_a(m)(r^\#)^2\le1-\sigma_0/4$.  If $\sigma'\le0$, then $\mathsf k=a+\sigma'\le a$
and $\mathsf k>0$ by (iii), so $a^2=\sigma'+\mathsf k^2R^2\le a^2R^2$, which
forces $R\ge1$; this is impossible.  So $\sigma'>0$, and $\mathsf k=\mathbf K(\sigma',
R^2)\le\mathbf K(\sigma_0+c_{L+1}^2/8,R^2)$.  By Lemma~\ref{lem:III-K}(e) on
$\mathcal K(\sigma_0)$, this exceeds $\mathbf K(\sigma_0,R^2)$ by at most $2c_{L+1}^2/
\sigma_0^2$, and $m|r_l|/\Phi_l=\varrho^0_l(r^\#)/\mathbf K(\sigma_0,R^2)\le1/(2\sigma_0)$
by (iv).  Hence $\varrho_l(f^\#)=m|r^\#_l|\mathsf k/\Phi_l\le\varrho^0_l(r^\#)+
\sigma_0^{-3}c_{L+1}^2<1$ for $l\in\Omega_a(m)$, and $\varrho_l(f^\#)=0$ for the
other carriers of $\Omega_m$.  So $\Omega_m\subseteq Q(f^\#)$, and
Lemma~\ref{lem:III-floor} with $P_0:=P_{\rm rob}(m)$ gives the lower bound.
\end{proof}

\begin{proposition}[Exactification with threshold protection without (KN)]\label{prop:III-KNtr}
Let the design satisfy the QC requirements of
Definition~\textup{\ref{def:III-QCreq}} at the main stage $L$.  Let $N\ge1$, let
$f\in S_{p_N^*}$ have finite base support, let $L\ge l_f$, let $w$ be a
sub-window of $L$ that is clean for $f$, let $a$ be an activity class of $w$,
and let $\kappa$ be the pattern of $(f,w,a)$; under option \textup{(b)} assume
$1\notin\mathcal N_T(\kappa)$.  Put $\mu:=\mu_\kappa(\eta'/3)$ if $\mathcal N_T\ne
\emptyset$.  Then there is a companion $f^{(1a)}$ of $f$ with the following
properties.
\begin{enumerate}
\item[(a)] \textup{(Exactness.)}  For every block $m$, $\varkappa(\hat\zeta^{(1a)}_m;
\Omega_m)=\varkappa^*_m>0$, where $\varkappa^*_m$ is the value after the closing
moves; the active ratios $r''$ of $f^{(1a)}$ lie in $Z_\kappa$, every robust rate
object of $\kappa$ has modulus $\ge u/4$ at $r''$, $|r''-r(f)|\le5L\mathrm{Des}(L)\eta'$,
and in blocks of kind \textup{(T)} and \textup{(KN)} the block coordinates
satisfy $|r''^{(m)}-r(f)^{(m)}|\le2\eta'$.
\item[(b)] \textup{(Statuses.)}  In every block of kind \textup{(T)}: every
carrier of $\mathcal N(m)$ is a strict non-peak with $\varrho\le1-\mu/2$, every
carrier of $\Omega_m\setminus\Omega_a(m)$ has value exactly $0$, and the coarse
peaks are exactly $P_{\rm rob}(m)$.  In every block of kind \textup{(KN)}: the
conclusions of [II,~Proposition~5.56(b)] hold with $\eta'$ in place of
$\eta(w)$.  In every block of kind \textup{(U)}: every active carrier of
$\Omega_m$ has $\varrho\le1-\eta'$, and every inactive near-threshold carrier is
a strict non-peak with $\varrho\le1-\eta'$, hence with gap $\ge\eta'M$.  Every other coarse carrier keeps
its robust status; relative positions of coarse carriers move by at most
$C\mathrm{Des}(L)\eta'$ except where stated.
\item[(c)] \textup{(Cost.)}  $p^*(f^{(1a)}-f)\le C_f\mathrm{Des}(L)^3\eta'\log(1/\eta')
=o(T_{\rm lo}(w)^2)$.
\end{enumerate}
Moreover $\mu\ge\mathrm{Des}(L)^{-1}(\eta'/3)^{\mathrm{Des}(L)}$ and $8b\le\mu$.
\end{proposition}

\begin{proof}
Write $\sigma_0(m):=\sigma_{\rm rob}(m)$.  For $L\ge l_f$ we have
$b\le\sigma_0(m)/8$ and $3L\cdot2\eta'\le\sigma_0(m)/4$ for every $m$
(Lemma~\ref{lem:III-alpha}(b), (QC3), $\eta'\le T_{\rm lo}^4$ and
$\mathrm{Des}^2T_{\rm lo}\le1$), so Lemmas~\ref{lem:III-Tblock}
and~\ref{lem:III-interior} apply with $\rho_0:=2\eta'$.

\emph{Step 0 (closing moves).}  Apply only the closing moves (C1), (C2) of the
assembly ([II,~Subsection~4.3]); there is no donor raise in active blocks,
and the active weak peaks are \emph{not} pushed at this stage.  Let $f_1$ be
the resulting row and $r_1$ its active ratio vector.  The closing moves
change every value by at most $C_*(L)b\le\mathrm{Des}(L)b$ (QC3).  A ratio is
$r_l=u_l(\zh)/\varkappa_m$, and $\varkappa_m$, $\bar\vartheta_m$ are Lipschitz
functions of the values with $f$-constants, bounded below by $1/C_f$
([II,~Lemma~2.7(a), Lemma~2.10]); a status carries in addition the factor
$1/\Phi_l\le D(L)$.  Hence
\[
 \Delta r_l\le\frac{\Delta u_l}{\varkappa_m}+\frac{|u_l|\,|\Delta\varkappa_m|}
 {\varkappa_m^2}\le C_f\mathrm{Des}\,b,\qquad\Delta\varrho_l\le C_fD(L)
 \mathrm{Des}\,b ,
\]
and since $|\Omega_a|\le2L$, $C_f\le\mathrm{Des}(L)$ and $D(L)\le\mathrm{Des}(L)$
for $L\ge l_f$, we get $|r_1-r(f)|\le C_{**}b$ (Euclidean norm) and statuses
move by at most $C_{**}b$, $C_{**}=\mathrm{Des}^3$.  By (QC4),
$C_{**}b\le\beta(b)\le\eta'/3$.  At $r(f)$ every tiny rate object has
modulus $\le b$, so at $r_1$ its modulus is $\le\beta(b)=(1+\mathrm{Lip}^*C_{**})b$.  If
$\mathcal N_T\ne\emptyset$, then for $d\in\mathcal N_T$, by
Lemma~\ref{lem:III-Tblock}(a) and $\varrho_d(f)\le1+b$,
\[
 \begin{aligned}
 F_\Phi(r_1)&\le\max_d\varrho^0_d(r(f))+\Lambda_\kappa C_{**}b\\
 &\le1+b+\tfrac12C_0D^{12}(b+c_{L+1}^2)+\mathrm{Lip}^*C_{**}b\le1+\beta(b)+
 C_0D^{12}b ,
 \end{aligned}
\]
using $\sigma_{\rm rob}^{-1}\le D^2$, $c_{L+1}^2\le b$ and $\Lambda_\kappa\le
\mathrm{Lip}^*$.

\emph{Step 1 ({\L}ojasiewicz point).}  $r_1\in\mathcal K_\kappa$ by
Lemma~\ref{lem:III-interior} with $\rho_0=2\eta'$, since $|r_1-r(f)|\le
C_{**}b\le\eta'/3$.  Apply Lemma~\ref{lem:III-loj} on $\mathcal
K_\kappa$ to the rate objects of $\kappa$, with $S$ the tiny ones and $\beta:=
\beta(b)$.  Alternative (i) is excluded by $\beta(b)<1/C_L$ (QC4).  So there is
$r'\in Z_\kappa$ with $|r'-r_1|\le C_L\beta^{2/N_L}\le\eta'/3$ (D1), and
$F_\Phi(r')\le F_\Phi(r_1)+\mathrm{Lip}^*C_L\beta^{2/N_L}\le1+\delta_0(b)$.

\emph{Step 2 (a point of $Y_1$).}  If $\mathcal N_T\ne\emptyset$: if $Y_1$ were
empty, Lemma~\ref{lem:III-QC}(b) would give $F_\Phi(r')\ge1+m_0>1+\delta_0(b)$
by (D2).  So $Y_1\ne\emptyset$, and there is $r'_1\in Y_1$ with $|r'_1-r'|=
\dist(r',Y_1)\le C_Y\delta_0(b)^{\alpha_Y}\le\eta'/3$ (D2).  If $\mathcal N_T=
\emptyset$, put $r'_1:=r'$.

\emph{Step 3 (Puiseux point).}  If $\mathcal N_T\ne\emptyset$, the actual weights
are not a coincidence for $\kappa$ (Lemma~\ref{lem:III-GEN} with (QC1); under
option (b) because $1\notin\mathcal N_T$), so Lemma~\ref{lem:III-QC}(a) with
$h:=\eta'/3$ gives $r''\in Z_\kappa$ with $|r''-r'_1|<\eta'/3$ and $F_\Phi(r'')\le
1-\frac34\mu$.  If $\mathcal N_T=\emptyset$, put $r'':=r'_1$.  Then $|r''-r_1|<
\eta'$ and, as $|r_1-r(f)|\le C_{**}b\le\eta'/3$, $|r''-r(f)|\le2\eta'$; so
$r''\in\mathcal K_\kappa$ with $R_a(m)(r'')^2\le1-\sigma_0(m)/2$
(Lemma~\ref{lem:III-interior}).  Every robust rate object has modulus $\ge
u-2\mathrm{Lip}^*\eta'\ge u/2$ at $r''$, every robust status
moves by at most $2\mathrm{Des}\,\eta'\ll u$, and the signs of the ratios of
the carriers of $\mathcal N_T$ are those of $r(f)$, because $|r_d(f)|=
\varrho_d(f)\Phi_d/(m\mathsf k)\ge\Phi_{\min}(L)/(2L)\gg\eta'$ ($\mathsf k\le\frac12$
by Lemma~\ref{lem:III-K}(a)).  By (D2), $8b\le\mu$; by
Remark~\ref{rem:III-QCwell}(iii), $\mu\ge\mathrm{Des}^{-1}(\eta'/3)^{\mathrm{Des}}$.

\emph{Step 4 (realization in one joint fixed point).}  The targets are:
\begin{itemize}
\item active carriers $l\in\Omega_a(m)$: value $u_l(\zh^{(1a)})=r''_l\varkappa^*_m$
if $m$ is of kind (T) or (KN), and $s_mr''_l\varkappa^*_m$ with $s_m:=1-4\mathrm{Des}
\,\eta'$ if $m$ is of kind (U);
\item \emph{holds}: every carrier of $\Omega_m\setminus\Omega_a(m)$, $m$ of kind (T),
gets value exactly $0$ (by (T) these are exactly the nearly neutral coarse
strict non-peaks of the block other than the active ones);
\item $\varkappa(\hat\zeta_m;\Omega_m)=\varkappa^*_m$ for every block $m$.
\end{itemize}
One-shot moves come first: the pulls of all tuned carriers, in the common
window of Lemma~\ref{lem:III-window} with $\eta_*:=\max(\|x\|_\infty,c_{L+1})$,
where $x$ is the vector of corrections.  The corrections include the
(U)-block rescaling $(1-s_m)r''_l\varkappa^*_m=4\mathrm{Des}\,\eta'\,r''_l
\varkappa^*_m$ and the displacements $|r''-r_1|\varkappa^*_m\le\eta'\varkappa^*_m$
with $\varkappa^*_m\le C_f$, so $\eta_*\le C_f\mathrm{Des}\,\eta'$.  In
blocks of kind (KN): if the lever $d_0$ (an inactive carrier of $\Omega_m$
with $\varrho_{d_0}(f)\ge u$, Definition~\ref{def:III-pattern}(b)) is a weak
peak of $f$ ($\varrho_{d_0}(f)\in[1,1+b]$), it is first pushed inward to
$\varrho_{d_0}=1-b$ by a private move; the value change is $\le2b\bar
\vartheta_m\Phi_{d_0}/m$ and the cost is $O(b)$, and $\bar\vartheta_m$ and
$\varkappa_m$ are continuous across the crossing of the threshold
([II,~Lemma~2.7(a), Lemma~2.10(a)]).  Thereafter $d_0$ is a strict non-peak
with $\varrho_{d_0}\ge u$, as required in [II,~Proposition~5.56].  Then the
inward lever push of the proof of [II,~Proposition~5.56] (Step~2), of size
$s_0:=C'\mathrm{Des}\,\eta'\bar\vartheta_m\Phi_{P_m}^2/(\varrho_{d_0}M_m)$, is
made, and the inactive near-threshold carriers of the block are pushed
inward individually to relative position $\le1-\eta'$.  In blocks of kind
(U) the inactive near-threshold carriers are pushed inward individually to
relative position $\le1-\eta'$.  Since such a carrier has $\varrho(f_1)\le1+(1+C_{**})b$ and
$(1+C_{**})b\le\eta'/3$ by (QC4), each push has size $\le2\eta'$ in relative
position.  The remaining unknowns are the banks
of [II,~Lemma~4.15] for the active carriers and for the held carriers of
class G, the sizes of two-sided z-moves on free coordinates of
$S^{\rm nat}_h\setminus(F\cup T(L))$ for the held carriers $h$ of class R
(their two signed rooms are $\ge u\|v_h1_{S^{\rm nat}_h}\|_1\gg b$, and the
moves keep every moved coordinate free, so the $E_c$-types of the pattern
do not change), and the size of one peak push per block.  The equations
are the targets above.  The derivative of each tuned value with respect
to its own bank is $\ge\mathrm{Des}^{-1/6}/2$, with respect to its own z-move
it is $v_h(j)$, and with respect to the other unknowns it is
$O(C\mathrm{Des}\,\eta_*)$ (second order, through the Hilbert normalization,
[II,~Lemma~4.14]); the derivative of $\varkappa_m$ with respect to its peak
push lies in $[\frac34,1]$ ([II,~Lemma~2.10(b)(i)], $X\le CD^2c_{L+1}^2$), and
with respect to the values of $\Omega_m$ it is $O(X)$
([II,~Lemma~2.10(b)(iii)]); $\varkappa_m$ is continuous
([II,~Lemma~2.10(a)]).  So the Jacobian is diagonally dominant, and the
explicit fixed point of [II,~Lemma~4.15], extended by one scalar equation
per block solved by the intermediate value theorem, converges, as in
Step~2 of the proof of [II,~Proposition~5.56].  Later re-solves change the
corrections only by $O(\mathrm{Des}^2\eta_*^2)\ll\eta_*$, which the banks absorb;
the pulls are not moved again.  At the resulting row $f^{(1a)}$ the active
ratios equal $r''$ in blocks of kind (T) and (KN), and $(s_mr''^{(m)},
r''^{\rm rest})$ in blocks of kind (U); we call this vector $r''$ again.  By
Lemma~\ref{lem:III-Uscale}(a) it still lies in $Z_\kappa$, robust rate objects
are multiplied by at least $s_m^n\ge\frac12$ (as $4n\mathrm{Des}\,\eta'\le\frac12$),
hence have modulus $\ge u/4$, and $F_\Phi$ is unchanged.  In blocks of
kind (T) and (KN) the block coordinates of $r''$ are those of the Puiseux
point, so $|r''^{(m)}-r(f)^{(m)}|\le2\eta'$ there.  This proves (a).

\emph{Statuses.}  Blocks of kind (T): hypotheses (i)--(iii) of
Lemma~\ref{lem:III-status} hold by construction (the robust peaks move by
$\ll u$; $\varkappa^*_m>0$).  For $d\in\mathcal N(m)$, $\varrho^0_d(r'')\le F_\Phi(r'')\le
1-\frac34\mu$; for the other active carriers $l$ of the block, $\varrho_l(f)\le1-u$,
so $\varrho^0_l(r'')\le\varrho_l(f)+C_0D^{12}(b+c_{L+1}^2)+2\mathrm{Lip}^*\eta'<1-u/2$.
Since $\sigma_0^{-3}c_{L+1}^2\le D^6b^4\le b/8\le\mu/64$ (QC5), hypothesis (iv)
holds, and Lemma~\ref{lem:III-status} gives $\varrho_d(f^{(1a)})\le1-\frac34\mu+
\frac{\mu}{64}\le1-\mu/2$ for $d\in\mathcal N(m)$, all carriers of $\Omega_m$ strict
non-peaks, and coarse peaks exactly $P_{\rm rob}(m)$.  Blocks of kind (KN):
the proof of [II,~Proposition~5.56](b) applies verbatim, the total
displacement of the ratios being $\le2\eta'$ (including the Puiseux step)
and the lever push dominating it by the choice of $C'$.  Blocks of kind
(U): by (QC3) the actual statuses are $\mathrm{Des}$-Lipschitz in the ratios
at fixed $\sigma$ and fixed inactive ratios, so at the unscaled ratios the
active statuses are $\le1+C_{**}b+2\mathrm{Des}\,\eta'$; by
Lemma~\ref{lem:III-Uscale}(b) the scaled ones are $\le s_m(1+3\mathrm{Des}\,\eta')
\le1-\mathrm{Des}\,\eta'\le1-\eta'$.  The pushed inactive near-threshold carriers
end at relative position $\le1-\eta'$ up to second-order effects (see
below), so they are strict non-peaks with gap $M(1-\varrho)\ge\eta'M$
([II,~Lemma~2.6(b)]); the margin $\eta'$ is kept because the push target
is $1-\eta'-\eta'^{3/2}$, say, and the second-order effects are
$O(\mathrm{Des}^2\eta'^2)\le\eta'^{3/2}$.  Second-order cross
effects are $O(\mathrm{Des}^2\eta'^2)$, small compared with the margins $\eta'$
of (U)-blocks and $c_f\mathrm{Des}\,\eta'M$ of (KN)-blocks; in (T)-blocks they do
not enter at all, because there the statuses depend only on the exact
active ratios, the exact held values and the peak set
(Lemma~\ref{lem:III-status}).  This proves (b).

\emph{Cost.}  By [II,~Lemma~4.15(c)] and Lemma~\ref{lem:III-window}, each
pull has mass $\le6\cdot2^{G(L)}\eta_*$ and each bank mass $\le C_T\eta_*$,
with $C_T\le C_f\mathrm{Des}$ and $2^{G(L)}\le\mathrm{Des}^{1/6}$; the
z-moves of held carriers, the lever push and the inward pushes have size
$\le C_f\mathrm{Des}\,\eta'$, and the weak-peak lever pushes have size
$O(b)\le O(\eta')$.  So every move has size $\le C_f\mathrm{Des}^2\eta'$ in
units of masses and of values (as $\eta_*\le C_f\mathrm{Des}\,\eta'$).  There
are at most $|L_0|\le2L$ tuned carriers (one pull and one bank each), at
most $2L$ held or pushed carriers and one peak push per block, hence at most
$6L$ moves, and $L\le\mathrm{Des}^{1/2}$.  Therefore $\|a^{(1a)}-a\|_1\le
C_f\mathrm{Des}^{5/2}\eta'$ and $\Delta_m(\zh^{(1a)}-\zh)\le C_f
\mathrm{Des}^{5/2}\eta'$ for every block $m$, and the counting bound of
[II,~Subsection~2.3] gives $c(\zh^{(1a)}-\zh)\le C_f\mathrm{Des}^{5/2}\eta'
\log(e/\eta')$.  By [II,~Lemma~2.12(iv)], $p^*(f^{(1a)}-f)\le
C_f\mathrm{Des}^3\eta'\log(1/\eta')$.  Since $\eta'=T_{\rm lo}^4/(L\mathrm{Des})$,
$\log(1/\eta')\le5\log(1/T_{\rm lo})$ and $\mathrm{Des}^2T_{\rm lo}\le1$, this is
$\le C_fT_{\rm lo}^3\log(1/T_{\rm lo})=o(T_{\rm lo}^2)$.
\end{proof}

\begin{remark}[Why the holds must be exact]\label{rem:III-KN1}
The margin $\mu=\mu_\kappa(\eta'/3)$ is only known to be $\ge c_\kappa(\eta'/3)^
{\beta_\kappa}$ with a non-explicit exponent $\beta_\kappa$.  Every quantity that
separates the true status of a protected carrier from its floor
$\psi_d(r'')/\Phi_d$ must therefore be dominated by $\mu$ and not merely by
a power of $\eta'$.  The active ratios are exact (fixed point), $\sigma$
depends only on the peak set up to fine terms $\le c_{L+1}^2/8\le b^4$, and
the remaining term is the contribution of the non-active carriers of
$\Omega_m$ to $R^2$.  A non-active value of size $v$ raises $R^2$ by an amount
of order $v^2$, and the moves made for the other blocks change such values
at second order, by $O(\mathrm{Des}^2\eta'^2)$ ([II,~Lemma~4.14]); the resulting
status error, of order $\mathrm{Des}^4\eta'^4$ times design factors, is not
small compared with $c_\kappa(\eta'/3)^{\beta_\kappa}$ when $\beta_\kappa>4$.
Hence every non-active coarse strict non-peak of a (T)-block is held at
value exactly $0$ inside the joint fixed point, not merely made small.  In
(KN)- and (U)-blocks the margins are of order $\eta'$ and second-order
effects are harmless.
\end{remark}

\begin{lemma}[Persistence of the margin]\label{lem:III-drift}
In the situation of Proposition~\textup{\ref{prop:III-KNtr}}, let $f^\natural$ be
obtained from $f^{(1a)}$ by moves that change the relative position of every
carrier of $\mathcal N_T$ by at most $C_f\mathrm{Des}(L)c_{L+1}$.  If $L\ge l_f$,
then $\varrho_d(f^\natural)\le1-\mu/4$ for every $d\in\mathcal N_T$.
\end{lemma}

\begin{proof}
By (QC5), $c_{L+1}\le b^2$, and $b\le\eta'\le T_{\rm lo}^4/\mathrm{Des}$ gives
$C_f\mathrm{Des}\,b\le C_fT_{\rm lo}^4\le\frac12$ for $L\ge l_f$.  Hence $C_f\mathrm{Des}\,
c_{L+1}\le b/2\le\mu/16<\mu/4$, using $8b\le\mu$.
\end{proof}

\begin{remark}[Scope]\label{rem:III-KNscope}
(i) Lemma~\ref{lem:III-drift} is the form in which the later steps of the
companion (fine structure, release of far parts, and the passage from the
reference shift to the true shift) use Proposition~\ref{prop:III-KNtr}; there
these moves change statuses, shifts and $\varkappa$ by at most $C_f\mathrm{Des}(L)
c_{L+1}$ (Section~\ref{sec:III-routeV}).  The lemma does not use any upper
bound on $\mu$.
(ii) Nothing about $f$ is assumed beyond the clean sub-window and finite base
support; the class is fixed before the exactification, which is legitimate
because classes are read off the decomposition data ([II,~Lemma~5.53]).
(iii) Under option (b), Proposition~\ref{prop:III-KNtr} excludes the patterns
in which stage $1$ (whose weight $c_1=1$ is rational) is an active
near-threshold carrier of a block of kind (T).  If this happens at clean
sub-windows of infinitely many main stages, then $\varrho_1(f)=1$.  If the
weight-one carrier lies in a block $m(1)>N$, no pattern of $p_N$ contains
it; in the remaining case the exception is open.  A slice version of
Lemma~\ref{lem:III-NLM} would apply if $2^{-m(1)-1}$ were a regular value of
$\psi_1$ on every stratum; we do not know whether it can be a critical value
identically in the transcendental parameters, and no such claim is made.
(iv) The finite base support of $f$ is used only through the realization of
Step~4 (far coordinates outside $F$ for pulls, banks and z-moves); the
transfer of this section to infinite base support is not claimed here.
\end{remark}
\section{The violation-tolerant $d$-consistent engineered approximant}\label{sec:III-UE}

The engineered approximants of [I,~Theorem~7.18] are built at a single row and
at late stages of a construction sequence.  To average over a window of scales
at a row \emph{near} the target row $f$ one needs per-piece bounds whose radius
is proportional to the scale of the piece, with constants that do not depend on
the window, and one must tolerate small violations of the sign conditions of
two-piece data.  [II,~Lemma~2.30] tolerates violations at a single stage;
[II,~Remark~2.31(c)] records the obstruction to a uniform use: the
$\theta$-decomposition and the $\pm$-decompositions of a piece are linked by a
first-order term (the \emph{junction mismatch}) that is not uniformly small
along windows.  In this section we remove it exactly: at the approximant $f'$
the normalized values of the finitely many switching carriers are re-tuned to
their values at the base row by small \emph{levers}
(\emph{$d$-consistency}).  With this, the bounds of [I,~Theorem~7.18] and
[II,~Lemma~2.30] hold piece by piece with radius $c_{\rm flat}t$, and the
averaging can be carried out at the norm-attaining approximant itself.

\subsection*{Standing conventions}
Throughout this section $I$ is finite and the base is \emph{diagonal}:
$Uk_j=\mu^\star_je_j$, $U^*e^*_j=\mu^\star_jk_j$ with $\mu^\star_j>0$ and
$(k_j)$ orthonormal, as for $T_{\rm final}$ [II,~Definition~3.1].  For
$A\in\ell_1\setminus\{0\}$ and $z\in B_{\ell_\infty}$ with $z_j=\sgn A_j$ on
$\supp A$ we write $f[A,z]$ for the first row with prescribed data
[II,~Subsection~4.2]; its forced data are $e=U^*A/\|U^*A\|_H$, $\zh=z+Ue$ and
$w_m=J_m(R_m^{**}\zh)$.  For a row $f_0$ with forced data
$(a_0,z_0,e_0,\nu_0,\zh_0,\dots)$ and $F_0:=\supp a_0$ finite we put
$\zeta_m:=R_m^{**}\zh_0$ and $A_m:=|\zeta_m|_m=\sigma_m/q_0$.  For a strict
non-peak $k$ of block $m$ the \emph{normalized value} is
\[
 \mathrm{nv}_k:=\zeta_m(k)/A_m=\lambda_{k,m}u_{k,m}(\zh_0)/A_m .
\]
By [I,~Lemma~1.7], $\mathrm{nv}_k=\Phi_m(k)^2w_m(k)/C_m$, hence for every
$\omega\in c_{00}$ vanishing on $P_m$
\begin{equation}\label{eq:III-nv}
 d_m(\omega)=\frac{\ip{D_mw_m}{D_m\omega}}{C_m}=\sum_k\omega(k)\,\mathrm{nv}_k .
\end{equation}
We also use the clamp variable $v_k:=m|u_{k,m}(\zh_0)|/A_m=|\mathrm{nv}_k|/
\Phi_m(k)$ of the proof of [I,~Lemma~7.14] ($\lambda_{k,m}=m\Phi_m(k)$).

\begin{definition}[Pieces and violations]\label{def:III-pieces}
Let $f_0\in S_{p^*}$ have $F_0$ finite and contact set $K_0$.  A \emph{piece}
at $f_0$ at scale $t\in(0,1]$ is a functional $g$ with $g(\xi_0)=0$ together
with two pairs $(b^\pm,\omega^\pm)$ that represent $g$ at $f_0$ in the sense of
[I,~Definition~7.7] \textup{(}$\omega^\pm_m\in c_{00}$,
$\supp\omega^\pm_m\subseteq Q_m$\textup{)}; they need not be side-admissible.
Put $\Delta\omega:=\omega^--\omega^+$, $\Delta_m:=d_m(\Delta\omega_m)$,
$v:=b^+-b^-=\sum_mR_m^*(\Delta\omega_m-\Delta_mw_m)$,
$b^\theta:=\frac12(b^++b^-)$, $\omega^\theta:=\frac12(\omega^++\omega^-)$.  The
\emph{violation} at $j\notin F_0$ is $\mathrm{viol}(j):=(z_0(j)b^+_j)_-+
(z_0(j)b^-_j)_+$ for $j\in K_0$ and $\mathrm{viol}(j):=|b^+_j|+|b^-_j|$ for
$j\notin F_0\cup K_0$, and the \emph{violation mass} is
$\epsilon:=\sum_{j\notin F_0}\mathrm{viol}(j)$, as in [II,~Lemma~2.30].
\end{definition}

Let $f\in S_{p^*}$ with $F$ finite be the \emph{target row}, and fix
$\rho\in(0,1)$, $\eta_0\in(0,2]$ with $\rho^2(1+\eta_0)\le(1+\rho^2)/2$,
$A_0,A_2\ge1$ and $\gamma_B\in(0,1]$.  Put
\[
 \delta:=\tfrac12\big(1-\rho^2(1+\eta_0/2)\big)\ \ge\ (1-\rho^2)/8 ,
\]
so $\rho^2\kappa\le1-2\delta$ whenever $\kappa\le1+\eta_0/2$.  An
\emph{$f$-constant} is a number that depends only on $(f,I,T,\rho,\eta_0,
A_0,A_2,\gamma_B)$ and on choices made at $f$ (transfer data); a
\emph{window quantity} may in addition depend on a window and its data.

\subsection{Masses off the support, $d$-consistency, uniform neighbourhoods}
\label{subsec:III-prelim}

The next lemma complements [II,~Lemma~4.14].

\begin{lemma}[Masses and values]\label{lem:III-masses}
Let $a\in\ell_1\setminus\{0\}$ with $F:=\supp a$, $\nu:=\|U^*a\|_H$,
$e:=U^*a/\nu$, $z\in B_{\ell_\infty}$ with $z=\sgn a$ on $F$, and $\zh:=z+Ue$.
Let $J$ be finite, $c\in\R^J$, $A:=a+\sum_{s\in J}c_se^*_s$, $\nu_A:=
\|U^*A\|_H$, $e_A:=U^*A/\nu_A$, $X:=U^*(A-a)=\sum_{s\in J}c_s\mu^\star_sk_s$.
\begin{enumerate}
\item[(i)] If $J\cap F=\emptyset$, then $X\perp U^*a$, $\nu_A^2=\nu^2+\|X\|^2$,
$0\le\nu_A-\nu\le\|X\|^2/(2\nu)$, $\|e_A-e\|\le\|X\|/\nu$, and for every
$u\in\ell_1$
\[
 \ip{U^*u}{e_A-e}=\frac1{\nu_A}\sum_{s\in J}c_s(\mu^\star_s)^2u(s)+
 \Big(\frac\nu{\nu_A}-1\Big)\ip{U^*u}{e} .
\]
\item[(ii)] In general $\|e_A-e\|\le2\|X\|/\nu$.  If $s\in J\cap F$ and all
other mass coordinates lie off $F$, the first-order change of
$\ip{U^*u}{e_A}$ caused by $c_s$ is $c_s(\mu^\star_s)^2\big(u(s)-a_s\ip{U^*u}{e}/
\nu\big)/\nu$.
\item[(iii)] If $z^A\in B_{\ell_\infty}$ with $z^A=\sgn A$ on $\supp A$ and
$\hat x_A:=z^A+Ue_A$, then for every $u\in\ell_1$
\[
 |u(\hat x_A)-u(\zh)|\le\|U^*u\|\,\|e_A-e\|+\sum_s|u(s)|\,|z^A_s-z_s| ;
\]
for a carrier $\|U^*u_{k,m}\|\le q^*(u_{k,m})=1$.
\item[(iv)] For every $b\in\ell_1$, $\sum_{s\notin F}(\mu^\star_s)^2b(s)^2\le
\nu h(b)$, where $h(b)=\|P^\perp U^*b\|^2/\nu$ \textup{(}[I,~Definition~3.1]\textup{)}.
Consequently, if pieces $i=1,\dots,n$ at $f[a,z]$ are given and the masses are
$c_s:=4\rho s_1z_s\max_i|b^\theta_i(s)|$ on a set $J$ disjoint from $F$, then
\[
 \|X\|^2\le16\rho^2s_1^2\,\nu\sum_{i=1}^nh(b^\theta_i),\qquad
 \|X\|\le4\rho s_1\max_s\mu^\star_s\sum_{i=1}^n\|b^\theta_i\|_1 .
\]
\end{enumerate}
\end{lemma}

\begin{proof}
(i) $U^*a=\sum_{s\in F}a_s\mu^\star_sk_s$ and $X\in\spn\{k_s:s\in J\}$ with
$J\cap F=\emptyset$, so $X\perp U^*a$; this gives $\nu_A^2=\nu^2+\|X\|^2$ and
$\nu_A-\nu=\|X\|^2/(\nu_A+\nu)$.  The identity is [II,~(5)] rewritten with
$e_A=(\nu e+X)/\nu_A$.  As $\ip{e_A}{e}=\nu/\nu_A$, $\|e_A-e\|^2=2(\nu_A-\nu)
/\nu_A\le\|X\|^2/\nu^2$.  (ii) $\|x/\|x\|-y/\|y\|\|\le2\|x-y\|/\|y\|$; the
first-order formula is the derivative of $(\nu e+c_s\mu^\star_sk_s)/\|\nu e+c_s
\mu^\star_sk_s\|$ at $c_s=0$, using $\ip{k_s}{e}=a_s\mu^\star_s/\nu$.  (iii)
$\hat x_A-\zh=(z^A-z)+U(e_A-e)$ and $u(Uh)=\ip{U^*u}{h}$; $\|U^*u\|\le q^*(u)$.
(iv) Write $U^*b=U^*(b1_{F^c})+U^*(b1_F)$.  The first summand lies in
$\spn\{k_s:s\notin F\}$, which is orthogonal to $e$ and to
$\spn\{k_s:s\in F\}$; hence $P^\perp U^*b=U^*(b1_{F^c})+P^\perp U^*(b1_F)$ is
an orthogonal sum and $\|P^\perp U^*b\|^2\ge\|U^*(b1_{F^c})\|^2=\sum_{s\notin F}
(\mu^\star_s)^2b(s)^2$.  For the masses, $\|X\|^2=\sum_s c_s^2(\mu^\star_s)^2\le
16\rho^2s_1^2\sum_i\sum_{s\notin F}(\mu^\star_s)^2b^\theta_i(s)^2$, and
$\|X\|\le\sum_s|c_s|\mu^\star_s$ with $\max_i|\cdot|\le\sum_i|\cdot|$.
\end{proof}

By (iv) and convexity of $h$, $h(b^\theta)\le(h(b^+)+h(b^-))/2\le
\Gamma_w/q_0$; so for data with $\Gamma_w\le2$ the mass vector is
$O(s_1(n\nu/q_0)^{1/2})$, uniformly in the scales of the pieces.

\begin{definition}[$d$-consistency]\label{def:III-dcons}
Let $f_0,f'\in S_{p^*}$ \textup{(}data of $f'$ primed\textup{)} and let
$\Omega=\bigcup_m\Omega_m$ be a finite set of strict non-peaks of $f_0$,
$\Omega_m\subseteq Q_m$.  $f'$ is \emph{$d$-consistent with $f_0$ on $\Omega$}
if every $k\in\Omega_m$ is a strict non-peak of $f'$ and
$\mathrm{nv}'_k=\mathrm{nv}_k$, i.e.\ $u_{k,m}(\zh')=(A'_m/A_m)u_{k,m}(\zh_0)$.
\end{definition}

\begin{lemma}[Identities of $d$-consistency]\label{lem:III-dcons}
Let $f'$ be $d$-consistent with $f_0$ on $\Omega$.
\begin{enumerate}
\item[(a)] For every $m$ and every $\omega$ supported in $\Omega_m$:
$d'_m(\omega)=d_m(\omega)$ and $H'_m(\omega)=(C_m/C'_m)H_m(\omega)$.  For
$k\in\Omega_m$: $w'_m(k)=(C'_m/C_m)w_m(k)$ and
\begin{equation}\label{eq:III-gapid}
 \gap'_m(k)-\gap_m(k)=\Big(1-\frac{C'_m}{C_m}\Big)\big(1-\gap_m(k)\big).
\end{equation}
In particular $\gap'_m(k)\ge\gap_m(k)/2$ as soon as
$|C'_m-C_m|\le C_m\gap_m(k)/(2(1-\gap_m(k)))$.
\item[(b)] For every piece at $f_0$ with $\omega^\pm$ supported in $\Omega$ the
junction terms $\kappa^\pm_m:=\rho\,(d'_m-d_m)(\omega^\pm_m-\omega^\theta_m)$ of
the proof of [I,~Theorem~7.18] vanish.
\item[(c)] \textup{(The block scalars are frozen up to the carriers outside
$\Omega$.)}  Fix $m$, let $k_\natural\in P_m$ and $r>0$ with
$[C_m-r,C_m+r]\subseteq(0,1)$, $|C'_m-C_m|\le r$, and
$v_{k_\natural},v'_{k_\natural}\ge\Phi_m(k_\natural)(1-c)/c$ for all $c\in
[C_m-r,C_m+r]$.  Put $c_*:=2\Phi_m(k_\natural)^2(1-C_m-r)/(C_m+r)^3$ and
$P^{\rm rob}_m:=\{k\in P_m:v'_k\ge\Phi_m(k)(1-C_m)/C_m\}$.  Then
\[
 |C'_m-C_m|\le\frac{2(1-C_m)}{C_mc_*}\sum_{k\notin\Omega_m\cup P^{\rm rob}_m}
 \Phi_m(k)\,|v'_k-v_k| .
\]
\end{enumerate}
\end{lemma}

\begin{proof}
(a) By \eqref{eq:III-nv} at both rows, $d'_m(\omega)=\sum\omega(k)\mathrm{nv}'_k=
d_m(\omega)$; $H'_m(\omega)=(\|D_m\omega\|^2-d'_m(\omega)^2)/C'_m=(C_m/C'_m)
H_m(\omega)$.  $\Phi^2w'/C'=\mathrm{nv}'=\mathrm{nv}=\Phi^2w/C$ gives
$w'(k)=(C'/C)w(k)$.  With $M=1-C$ and $|w(k)|=M-\gap(k)$:
$\gap'(k)=M'-(C'/C)|w(k)|=(1-C')-(C'/C)(1-C-\gap(k))$, which is
\eqref{eq:III-gapid}; then $\gap'\ge\gap/2$ iff
$(C'/C-1)(1-\gap)\le\gap/2$, which holds under the stated bound.
(b) $\omega^\pm-\omega^\theta=\mp\frac12\Delta\omega$ is supported in $\Omega$;
apply (a).
(c) By the clamp formula of the proof of [I,~Lemma~7.14], $C_m$ and $C'_m$ are
roots of $F(c;v)=1$ and $F(c;v')=1$, $F(c;v):=\sum_k\min(\Phi_m(k)(1-c)/c,
v_k)^2$, each nonincreasing in $c$.  On the interval between $C_m$ and $C'_m$
the $k_\natural$-term of $F(\cdot;v')$ equals $\Phi_m(k_\natural)^2((1-c)/c)^2$,
whose derivative is $\le-c_*$; so $c_*|C'_m-C_m|\le|F(C_m;v')-F(C'_m;v')|=
|F(C_m;v')-F(C_m;v)|$.  Put $x_k:=\Phi_m(k)(1-C_m)/C_m$.  The $k$-th term of
the last difference vanishes for $k\in\Omega_m$ ($v'_k=v_k$ by
$d$-consistency) and for $k\in P^{\rm rob}_m$ (both $v_k\ge x_k$, since $k$ is
a peak of $f_0$, and $v'_k\ge x_k$); otherwise it is at most $2x_k|v'_k-v_k|$.
\end{proof}

\begin{remark}\label{rem:III-freeze}
If $\Omega_m$ contains every strict non-peak of $f_0$ in block $m$ except a set
$Q^{\rm fine}_m$ of carriers of small total weight, and the peaks of $f_0$ stay
peaks with margin above their value changes, Lemma~\ref{lem:III-dcons}(c)
bounds $|C'_m-C_m|$ by the moves of $Q^{\rm fine}_m$ alone; by (a) every
carrier of $\Omega_m$ then keeps its gap and its relative position up to the same
error.  So $d$-consistency on all coarse strict non-peaks also protects the
statuses of the switching carriers.  (Numerically, with $\Omega_m=Q_m$ in a
finite model, $|C'_m-C_m|\le1.1\cdot10^{-16}$; this is evidence only.)
\end{remark}

\begin{lemma}[Uniform neighbourhoods]\label{lem:III-ball}
Let $f\in S_{p^*}$ have $F$ finite.  In every block fix $\gamma_m\in(M_m/2,M_m)$
with $|w_m(k)|\ne\gamma_m$ for all $k$, put $\gamma_0:=\min_m(M_m-\gamma_m)$,
fix a peak $k^\natural_m$ with $\alpha_m(k^\natural_m)\ne0$ and margin
$\mu^\natural_m>0$, and, for $\eta_1\in(0,\frac14]$, transfer data
$(\gamma_m,k^\varsigma_{*,m},\Lambda^\varsigma_m)$ of $f$ by [I,~Lemma~7.4].
There is $r_f(\eta_1)>0$ such that every $f''\in S_{p^*}$ with
$p^*(f''-f)\le r_f(\eta_1)$ \textup{(}data double-primed\textup{)} satisfies:
\begin{enumerate}
\item[(a)] $|C''_m-C_m|\le C_m/8$, $|M''_m-M_m|\le\gamma_0/4$,
$|\sigma''_m-\sigma_m|\le\sigma_m/8$, $|q''_0-q_0|\le q_0/8$,
$|\nu''-\nu|\le\nu/8$, $F\subseteq\supp a''$ with $|a''_j|\ge a_{\min}/2$ on
$F$ \textup{(}$a_{\min}:=\min_F|a_j|$\textup{)}, and $k^\natural_m$ is a peak of
$f''$ with the same sign and margin $\ge\mu^\natural_m/2$;
\item[(b)] the transfer data persist: $k^\varsigma_{*,m}$ is a peak of $f''$
with $w''_m(k^\varsigma_{*,m})=M''_m$ and positive margin, $P''_m\subseteq
L''_m$, and the quantities of [I,~Proposition~7.3] built at $f''$ satisfy
$e''^\varsigma_{y,m}\ge\frac12$, $\iota''^\varsigma_m\le2\eta_1$,
$\|D_my''^\varsigma_m\|_2\le K_y$ and $|Y''^\varsigma_m|/q''_0\le K_Y$, with
$K_y,K_Y$ independent of $\eta_1$.
\end{enumerate}
\end{lemma}

\begin{proof}
Otherwise there are $f_n\to f$ violating (a) or (b).  By
[I,~Proposition~1.11], $a_n\to a$ in $\ell_1$, the scalar data converge,
$w_{n,m}\to w_m$ coordinatewise and $\xi_n\to\xi$ weak$^*$; the strict threshold
inequality of [I,~Lemma~1.7] at $k^\natural_m$ and its margin are continuous,
so $k^\natural_m$ is eventually a peak of $f_n$ with margin $\ge
\mu^\natural_m/2$; and [I,~Lemma~7.5] gives the persistence and convergence of
the transfer quantities.  Since $e^\varsigma_y\ge1-\eta_1\ge\frac34$
([I,~Lemma~7.4(d)]), $e''\ge\frac12$ eventually; $\|D_my^\varsigma\|_2\le
\|\Phi_m\|_2+\Lambda_0/m$ and $|Y^\varsigma|/q_0\le\sigma_me^\varsigma_y/q_0+
\Lambda_0$ do not depend on $\eta_1$.  This contradicts the choice of $f_n$.
\end{proof}

All $f$-constants below are computed at $f$ and are valid at every row of the
ball $B_f:=\{f'':p^*(f''-f)\le r_f(\eta_1)\}$; rows of windows and their
approximants lie in $B_f$ eventually.  The margin clause in (a) is what the
slope argument of Lemma~\ref{lem:III-dcons}(c) needs on the whole interval
between $C_m$ and $C''_m$, and the bound on $|M''_m-M_m|$ is what
[I,~Lemma~7.2] needs at $f''$ (it is not implied by $|C''_m-C_m|\le C_m/8$ when
$\gamma_0<C_m/4$).

\subsection{The junction mismatch for valid data}\label{subsec:III-J}

In the proof of [I,~Theorem~7.18] the $\pm$-decompositions contain the terms
$\kappa^\pm_m=\rho\,(d'_m-d_m)(\omega^\pm_m-\omega^\theta_m)=\mp\frac\rho2
(d'_m-d_m)(\Delta\omega_m)$, which enter the first-order block terms
$\lin'_m$ as $\tau\kappa^\pm_m$.  The following bound corrects the order
$s_1/t^2$ stated in [II,~Remark~2.31(c)], which was computed for a datum
scaled by $1/t$ whose coefficient $\Gamma_w$ grows like $t^{-2}$.

\begin{proposition}[Junction mismatch]\label{prop:III-J}
Let $f_0\in S_{p^*}$ have $F_0$ finite, and let $(b^\pm,\omega^\pm)$ be a piece
at $f_0$ with $\Gamma_w(b^\pm,\omega^\pm)\le\Gamma_{\max}$ \textup{(}computed at
$f_0$\textup{)} and $\supp\omega^\pm_m\subseteq\Omega_m\subseteq Q_m$, $\Omega_m$
finite.  Let $f'=f[A',z']$ be a row such that every $k\in\Omega_m$ is a strict
non-peak of $f'$, put $\Pi_e:=\|e'-e_0\|$ and
$t_k:=\|u_k1_{\{z'\ne z_0\}}\|_1$ \textup{(}for the engineered approximant of
[I,~Definition~7.12], $t_k=t_{k,m}(N'')$ off the window masses\textup{)}.  Then
for every block $m$:
\begin{enumerate}
\item[(a)] $\|D_m\Delta\omega_m\|_2^2=\Delta_m^2+C_mH_m(\Delta\omega_m)\le
\dfrac{4\Gamma_{\max}}{\sigma_m}\Big(C_m+\dfrac{C_m^3}{M_m^2\Phi_{P_m}^2}\Big)$,
where $\Phi_{P_m}^2:=\sum_{k\in P_m}\Phi_m(k)^2$;
\item[(b)] with $\delta u_k:=u_k(\zh')-u_k(\zh_0)$,
\[
 |\kappa^\pm_m|\le\frac{\rho m}2\,\|D_m\Delta\omega_m\|_2\,|\Omega_m|^{1/2}
 \max_{k\in\Omega_m}\frac{|\delta u_k|}{A'_m}+\frac\rho2|\Delta_m|\,
 \frac{|A'_m-A_m|}{A'_m};
\]
\item[(c)] consequently
\[
 |\kappa^\pm_m|\le C\,|\Omega_m|^{1/2}\big(\Pi_e+
\max_{k\in\Omega_m}t_k\big)+C\big(\Pi_e+\textstyle\sum_k\lambda_{k,m}t_k\big),
\]
where $C$ depends only on $\rho$, on $\Gamma_{\max}$ and on the block
constants $C_m$, $M_m$, $\sigma_m$, $\Phi_{P_m}$ and $A'_m$.
\end{enumerate}
\end{proposition}

\begin{proof}
(a) Decompose $D\Delta\omega=(\Delta/C)Dw+y$ with $y\perp Dw$; then
$\|y\|^2=\|D\Delta\omega\|^2-\Delta^2=CH(\Delta\omega)$.  Since $\Delta\omega$
vanishes on $P$ and $|w|=M$ there, $y=-(\Delta/C)Dw$ on $P$, so $\|y\|\ge
|\Delta|M\Phi_P/C$ and $\Delta^2\le C^3H(\Delta\omega)/(M^2\Phi_P^2)$.  Finally
$H(\Delta\omega)\le2H(\omega^+)+2H(\omega^-)\le4\Gamma_{\max}/\sigma$.
(b) By \eqref{eq:III-nv} at $f_0$ and at $f'$ (the carriers of $\Omega_m$ are
strict non-peaks of both rows), $(d'-d)(\Delta\omega)=\sum_{k\in\Omega}
\Delta\omega(k)(\mathrm{nv}'_k-\mathrm{nv}_k)$ and $\mathrm{nv}'_k-\mathrm{nv}_k=
\lambda_k\delta u_k/A'+\lambda_ku_k(\zh_0)(1/A'-1/A)$.  The second parts sum to
$d(\Delta\omega)(A-A')/A'=\Delta(A-A')/A'$.  For the first, $\lambda_k=m\Phi_k$
and Cauchy--Schwarz give $|\sum\Delta\omega(k)\lambda_k\delta u_k|\le
m\|D\Delta\omega\|_2(\sum_{\Omega}\delta u_k^2)^{1/2}$.
(c) Lemma~\ref{lem:III-masses}(iii): $|\delta u_k|\le\Pi_e+t_k$; and
$|A'-A|\le\sum_k\lambda_k|\delta u_k|$ because $|\cdot|_m\le\|\cdot\|_1$
[I,~Remark~1.1(c)] and $\sum_k\lambda_k\le1$.  Insert (a).
\end{proof}

\begin{remark}[Why an exact cancellation is used; heuristic]\label{rem:III-J}
By Lemma~\ref{lem:III-masses}(iv), $\Pi_e=O(s_1(n\nu/q_0)^{1/2})$ for valid
data, uniformly in the scales; so the mismatch is uniform in the scale $t$ of a
piece.  It is \emph{not} uniform along windows: $|\Omega_m|$ and the mass
vector (which carries the masses of all $n$ pieces of a window) grow with the
window, while the rebalancing of [I,~Proposition~7.3] absorbs a first-order
level shift only with a fixed inefficiency $\iota\sim\eta_1$, which cannot be
made window-dependent because [I,~Lemma~7.2] needs $|\varepsilon|\Lambda\le
\gamma/2$ with $\Lambda=\Lambda(\eta_1)$ unbounded.  The ratio of the resulting
cost to the budget $\delta\tau^2$ at $|\tau|=s_1$ is of order
$\eta_1|\Omega|^{1/2}\|X\|/(\delta s_1)$.  That this order is attained in
general is a heuristic statement, supported by a finite-model computation in
which the untuned mismatch is proportional to $\|X\|$ (and Proposition~
\ref{prop:III-J}(b) is attained when $|\Omega_m|=1$); it is not used below.
This is why the mismatch is removed exactly, by $d$-consistency
(Lemma~\ref{lem:III-dcons}(b)).
\end{remark}

\subsection{Levers and exact $d$-consistency}\label{subsec:III-DC}

Let $f_0$ be a row with $F_0$ finite and $\Omega=\bigcup_m\Omega_m$ a finite set
of strict non-peaks of $f_0$.  A \emph{pre-row} is a row $f_1=f[A_1,z_1]$
obtained from $f_0$ by adding finitely many masses (of sign $z_0(j)$ at
contacts $j\notin F_0$, or of the sign of $a_0(j)$ at $j\in F_0$) and by
changing $z_0$ on a set $Z_1$; put $\nu_1:=\|U^*A_1\|_H$,
\[
 \delta_1:=\max_{k\in\Omega}|u_k(\hat x_1)-u_k(\zh_0)|,\qquad
 \Delta_1:=\max_m\sum_k\lambda_{k,m}|u_{k,m}(\hat x_1)-u_{k,m}(\zh_0)| .
\]

\begin{definition}[Levers]\label{def:III-levers}
Let $l\in\Omega$ be a carrier of block $m(l)$ and $\epsilon_l\in\{\pm1\}$.
Lever coordinates $j$ of $l$ satisfy $v_l(j):=u_l(j)>0$.  A \emph{lever} of $l$
is one of the following one-parameter families $P\mapsto(A,z)$, $P\in\R$,
acting only at the lever coordinates $\mathrm{Lev}_l$:
\begin{enumerate}
\item[(Z)] \emph{zero-data lever} at $j\notin\supp A_1$: $z(j):=\mathrm{clamp}_{
[-1,1]}(z_1(j)+\epsilon_lP/v_l(j))$, and the overflow
$|z_1(j)+\epsilon_lP/v_l(j)|-1$, if positive, is placed as a bank at $j$ of mass
$\beta_j\,v_l(j)\cdot(\text{overflow})$ with the sign of the contact reached,
$\beta_j:=\nu_1/((\mu^\star_j)^2v_l(j))$;
\item[(TU)] \emph{tuning lever}: a pull coordinate $j_p\notin\supp A_1$ with
$z_1(j_p)=\epsilon_l$ and a bank coordinate $j'\ne j_p$ with $z_1(j')=
\epsilon_l$ and either $j'\notin\supp A_1$ or $\sgn A_1(j')=\epsilon_l$; the
pull is \emph{fixed}: $z(j_p):=-\epsilon_l$ and a mass $-\epsilon_l\mathfrak p_l$
at $j_p$, $\mathfrak p_l:=24\lambda_lv_l(j_p)$; for $P\ge-2v_l(j_p)$ a bank of
mass $\beta_{j'}(P+2v_l(j_p))$ with sign $\epsilon_l$ at $j'$.  The pull
coordinate is taken from a set $\mathrm{Pull}_l$ of admissible pull
coordinates of $l$ (coordinates $j\notin\supp A_1$ with $z_1(j)=\epsilon_l$
and $v_l(j)>0$), prescribed together with the lever.
\end{enumerate}
A free coordinate with room and a contact with $z_1(j)=\epsilon_l$ (inward
moves for $P\le0$, a bank for $P>0$) are special cases of (Z).  For a family of
levers $(\mathrm{Lev}_l)_{l\in\Omega}$ write $f(P):=f[A(P),z(P)]$ for
$P\in\R^\Omega$, with data $\hat x(P)$, $A_m(P):=|R_m\hat x(P)|_m$.  The
\emph{separation condition} \textup{(SEP)} is: (i) the sets $\mathrm{Lev}_l$ are
pairwise disjoint and disjoint from $\supp A_1$, except TU bank coordinates
with $\sgn A_1(j')=\epsilon_l$; (ii) $\mathrm{Lev}_l\cap\supp u_k=\emptyset$ for
all $k\in\Omega\setminus\{l\}$.  Zero-value absorber pairs tuned by a mass at a
support coordinate are \emph{not} levers in this sense (Remark~\ref{rem:III-abs}).
\end{definition}

Lever constants: $A_{\min}:=\min_mA_m(f_0)$; $G$ with $v_l(j_p)\in[r,2^Gr]$ for
the pulls in use; $\beta_{\max}:=\max\beta_j$ over the bank coordinates;
\[
 c_{\rm tiny}:=\max_{l\in\Omega}\sum_{k\notin\Omega}\lambda_k\sum_{s\in
 \mathrm{Lev}^0_l}\frac{|u_k(s)|}{v_l(s)},\qquad
 W_{\rm pull}:=\sum_{k:\ \supp u_k\ni j_p\text{ for a pull }j_p,\ k\ne l(j_p)}
 \lambda_k ,
\]
where $\mathrm{Lev}^0_l\subseteq\mathrm{Lev}_l$ consists of the (Z) coordinates
and the TU bank coordinates of $l$ (the TU pull coordinates are excluded:
the pull is a fixed move, accounted for by $W_{\rm pull}$).  For $r_{\min}>0$
let $\overline W_{\rm pull}(r_{\min})$ be the total weight $\sum\lambda_k$ of
the carriers $k$ for which there are $l\in\Omega\setminus\{k\}$ and
$j\in\mathrm{Pull}_l$ with $v_l(j)\ge r_{\min}$ and $u_k(j)\ne0$.  It does not
depend on the choice of the pulls, it is nonincreasing in $r_{\min}$, and
$W_{\rm pull}\le\overline W_{\rm pull}(r_{\min})$ whenever every pull in use
has $v_l(j_p)\ge r_{\min}$;
$C_{\rm lev}$ is the largest change of a carrier value per unit of $|P|_\infty$
caused by the lever moves at its own coordinates; $Y=(Y_{k,l})$ is the
first-order effect on $\epsilon_ku_k(\hat x)$ of the banks at support
coordinates $j'\in\supp A_1$ through the normalization
(Lemma~\ref{lem:III-masses}(ii)): $|Y_{k,l}|\le|A_1(j')|/(\nu_1v_l(j'))$.
Put $K_0:=1+2/A_{\min}$ and $K_J:=2K_0$.

\begin{lemma}[Exact $d$-consistency by levers]\label{lem:III-DC}
Let $f_0$, $\Omega$, a pre-row $f_1$ and levers with \textup{(SEP)} be given, and
put, for $l\in\Omega$ of block $m$,
\[
 F_l(P):=\epsilon_l\Big[u_l(\hat x(P))-\frac{A_m(P)}{A_m(f_0)}u_l(\zh_0)\Big].
\]
There are $r_0>0$ and $C_2<\infty$, depending only on the block data of $f_0$,
$\Omega$, the lever data \textup{(}$v_l$ and $\mu^\star$ on $\mathrm{Lev}_l$,
$\beta_{\max}$, $G$\textup{)} and the status margins of the carriers of $\Omega$
and of the peaks of $f_0$, such that the following holds.
\begin{enumerate}
\item[(i)] $|F(0)|_\infty\le\delta_1+\dfrac{\Delta_1+2W_{\rm pull}}{A_{\min}}+
C_2(2^{G+1}r)^2+\dfrac{2^{G+2}r|\Omega|c_{\rm tiny}}{A_{\min}}$.
\item[(ii)] Suppose
\begin{gather*}
 K_J|F(0)|_\infty\le\frac r2,\quad 3K_JC_2r\le\frac18,\quad
 \frac{K_J|\Omega|c_{\rm tiny}}{A_{\min}}\le\frac18,\quad K_0\|Y\|_\infty\le\frac12,\\
 \delta_1\le r,\quad\Delta_1\le r,\quad r\le r_0,
\end{gather*}
and $v_l(j_p)\in[r,2^Gr]$ for the TU levers.  Then $F$ has a zero $P^*$ with
$|P^*|_\infty\le r$.  At $f(P^*)$ every $l\in\Omega$ is a strict non-peak with
$\mathrm{nv}_l(f(P^*))=\mathrm{nv}_l(f_0)$, i.e.\ $f(P^*)$ is $d$-consistent with
$f_0$ on $\Omega$; every peak of $f_0$ whose margin, in value units, exceeds
$C_{\rm lev}(r+\delta_1)$ plus the change of the threshold, is a peak of
$f(P^*)$ with the same sign; the mass vector of the lever masses satisfies
$\|X_{\rm lev}\|\le C_2r$; a carrier $k$ meeting no lever coordinate moves by at
most $|u_k(\hat x_1)-u_k(\zh_0)|+\|e(P^*)-e_1\|$, and a carrier meeting lever
coordinates by at most that plus $C_{\rm lev}r+2|u_k(j_p)|$ \textup{(}the last
term for the pulls it meets\textup{)}.
\item[(iii)] For $l\in\Omega_m$, $\gap_m(l)(f(P^*))\ge\gap_m(l)(f_0)/2$ as soon
as $|C_m(f(P^*))-C_m(f_0)|\le C_m\gap_m(l)/(2(1-\gap_m(l)))$
\textup{(}Lemma~\ref{lem:III-dcons}(a)\textup{)}.
\end{enumerate}
\end{lemma}

\begin{proof}
\emph{(1) Admissibility and continuity.}  For $|P|_\infty\le r_0$ the pair
$(A(P),z(P))$ is admissible: banks have the sign of $z$, clamped $z$-moves keep
$|z|\le1$, a TU bank needs $P\ge-2v_l(j_p)$, which holds since $|P|\le r\le
v_l(j_p)$.  Finitely many coordinates move continuously, so $e(P)=U^*A(P)/
\|U^*A(P)\|_H$, $\hat x(P)=z(P)+Ue(P)$ and $A_m(P)$ are continuous in $P$.

\emph{(2) First-order structure.}  By construction and
Lemma~\ref{lem:III-masses}(i), (ii) with base $A_1$:
$\epsilon_l(u_l(\hat x(P))-u_l(\hat x(0)))=P_l+E_l(P)$.  A $z$-move at $j\in
\mathrm{Lev}_l$ changes $\epsilon_lu_l$ by exactly $P$; a bank of mass
$\beta_jx$ changes it by $x\,\nu_1/\nu_A$ up to $O(\|X_{\rm lev}\|^2/\nu_1^2)$;
by (SEP)(ii) the moves of the other levers change $u_l$ only through the
normalization, at second order, except for the banks at support coordinates,
whose first-order effect is $(YP)_l$.  The bank slope is $\nu_1/\nu_A=1-O(r^2)$,
so $|E_l(P)-(YP)_l|\le C_2'(r+|P|)|P|$ after the fixed moves; this is the form
used in (3).  Next, $A_m$ is a norm with norming functional $w_m(P)$ at
$R_m\hat x(P)$; with $\mathrm d\hat x:=\hat x(P)-\hat x(0)$,
\[
 0\le A_m(P)-A_m(0)-\ip{w_m(0)}{R_m\mathrm d\hat x}\le
 \ip{w_m(P)-w_m(0)}{R_m\mathrm d\hat x}.
\]
For a carrier $k$ whose status does not change, the clamp formula gives
$|w_m(P)(k)-w_m(0)(k)|\le K_{\rm cl}(|u_k(\mathrm d\hat x)|+|C_m(P)-C_m(0)|)$, so
its contribution to the right side is of second order; a carrier that may
change status contributes at most $2\lambda_k|u_k(\mathrm d\hat x)|$, which is
of second order unless $k$ meets a lever coordinate.  Further
$\ip{w_m(0)}{R_m\mathrm d\hat x}=\sum_{k\in\Omega_m}\lambda_kw_m(0)(k)
\epsilon_k(P_k+E_k)+\varrho_m$, where $|\varrho_m|\le|\Omega|c_{\rm tiny}|P|_\infty+
O(|P|^2)$ collects the carriers outside $\Omega$ that meet lever coordinates.
Finally $w_m(0)(k)=w_m(f_0)(k)+O(\delta_1+2^{G+1}r)$ and $u_l(\hat x(0))=
u_l(\zh_0)+O(\delta_1+(2^{G+1}r)^2)$.

\emph{(3) Linearization.}  With $p_l:=\epsilon_lu_l(\zh_0)/A_m(f_0)$ and
$q_k:=\epsilon_k\lambda_kw_m(f_0)(k)$ ($k,l\in\Omega_m$) we obtain
$F(P)=F(0)+(J+Y)P+R(P)$ with $J=I-pq^T$ blockwise and, for $|P|_\infty\le r_0$,
\[
 |R(P)|_\infty\le C_2(|P|_\infty+\delta_1+r)|P|_\infty+\frac{|\Omega|c_{\rm tiny}}
 {A_{\min}}|P|_\infty .
\]
By \eqref{eq:III-nv}, $q^Tp=\sum_{k\in\Omega_m}w_m(k)\mathrm{nv}_k=\sum_{\Omega_m}
\Phi_m(k)^2w_m(k)^2/C_m\le\|D_mw_m\|_2^2/C_m=C_m\le\|\Phi_m\|_2<\frac12$.  By the
Sherman--Morrison formula $(I-pq^T)^{-1}=I+pq^T/(1-q^Tp)$, and since
$|u_l(\zh_0)|\le q^*(u_l)q^{**}(\zh_0)=1$ and $\|q\|_1\le\sum_k\lambda_k\le1$,
$\|J^{-1}\|_\infty\le1+2\|p\|_\infty\|q\|_1\le K_0$.  As $K_0\|Y\|_\infty\le\frac12$,
$\|(J+Y)^{-1}\|_\infty\le2K_0=K_J$.

\emph{(4) Poincar\'e--Miranda.}  $G(P):=(J+Y)^{-1}F(P)=(J+Y)^{-1}F(0)+P+
(J+Y)^{-1}R(P)$ is continuous on $[-r,r]^\Omega$.  On the face $P_l=r$, using
$\delta_1\le r$,
\[
 G_l(P)\ge r-K_J|F(0)|_\infty-K_J\big(3C_2r^2+|\Omega|c_{\rm tiny}r/A_{\min}\big)
 \ge r-\tfrac r2-\tfrac r8-\tfrac r8>0,
\]
and symmetrically $G_l<0$ on $P_l=-r$.  By the Poincar\'e--Miranda theorem \cite{Miranda} $G$
has a zero $P^*$ in the cube, and $F(P^*)=0$.

\emph{(5) Statuses, sizes, and (i).}  Along the cube every value moves by at most
$\delta_1+C_{\rm lev}r$ (plus the fixed pull term for carriers meeting a pull),
and $C_m$, $\bar\vartheta_m$ move by $O((r+\delta_1+\Delta_1)/A_{\min})=
O(3r/A_{\min})$ through the clamp equation (here the hypotheses $\delta_1\le r$
and $\Delta_1\le r$ of (ii) are used; a small $|F(0)|$ alone does not bound
$\Delta_1$, because of cancellations in $A_m(0)-A_m(f_0)$); $r_0$ is chosen below the status margins of the carriers of
$\Omega$ \textup{(}in value units: $l\in\Omega_m$ stays a strict non-peak while
$|u_l|<\bar\vartheta_m\Phi_m(l)/m$, a room of order $\Phi_m(l)\gap_m(l)A_m/
(mC_m)$\textup{)} and of the peaks considered, divided by $C_{\rm lev}$ and the
Lipschitz constant of the threshold.  So no status changes along the cube, which
justifies (2).  At $P^*$, $\epsilon_lu_l(\hat x)=\epsilon_l(A_m/A_m(f_0))u_l(\zh_0)$,
i.e.\ $\mathrm{nv}_l(P^*)=\mathrm{nv}_l(f_0)$.  The lever masses are at most
$\beta_{\max}(r+2v_l(j_p))\le\beta_{\max}2^{G+2}r$ and the pull masses at most
$24\lambda_l2^Gr$, so $\|X_{\rm lev}\|\le C_2r$ after enlarging $C_2$; the value
bounds follow from Lemma~\ref{lem:III-masses}(iii).  For (i): at $P=0$ the TU
levers perform the fixed pull and a bank of mass $2\beta v_l(j_p)$; for $l$ the
flip ($-2v_l(j_p)$) and the bank ($+2v_l(j_p)\nu_1/\nu_A$) cancel up to
$C_2(2^{G+1}r)^2$, other carriers of $\Omega$ move only at second order
(SEP), and $|A_m(0)-A_m(f_0)|\le\Delta_1+2W_{\rm pull}+2^{G+2}r|\Omega|c_{\rm tiny}
+C_2(2^{G+1}r)^2$; divide by $A_{\min}$.  (iii) is
Lemma~\ref{lem:III-dcons}(a).
\end{proof}

\begin{remark}[On the constants]\label{rem:III-DC}
(a) The TU pull is a \emph{fixed} move: a carrier $k\ne l$ meeting $j_p$ moves by
$2|u_k(j_p)|$, not by $O(r)$; this is why $W_{\rm pull}$ appears in (i) and
must be small in applications (there such carriers live at stages beyond
$j_p$).  The fixed pull and bank also contribute the second-order term
$C_2(2^{G+1}r)^2$ to $F(0)$; in a finite model this term dominates $|F(0)|$
(evidence only).  (b) The gap criterion of (iii) is sharp in the sense that the
weaker criterion $|C_m(f(P^*))-C_m(f_0)|\le\gap_m(l)/4$ does not imply
$\gap(f(P^*))\ge\gap(f_0)/2$ when $C_m<1/2$; a finite-model example with
$C_m=0.243$, $\gap=0.01$ turns the carrier into a peak.  (c) The rank-one
structure is exact: a lever acting on one value changes the block norm, hence all
normalized values of the block, and $1-q^Tp\ge1-C_m>\frac12$; no second lever per
block is needed for $d$-consistency.
\end{remark}

\subsection{Window families and the $d$-consistent engineered approximant}
\label{subsec:III-window}

\begin{definition}[Window family]\label{def:III-window}
Let $f$ be the target row.  A \emph{base row} is $f_0\in S_{p^*}$ with
$p^*(f_0-f)\le r_f(\eta_1)/2$ (Lemma~\ref{lem:III-ball}; $\eta_1$ is fixed in
Theorem~\ref{thm:III-UE}), $F_0$ finite, $F\subseteq F_0$; put $E:=F_0\setminus F$.
A \emph{window family} at $f_0$ consists of pieces $(g_i;b^\pm_i,\omega^\pm_i)$,
$i=1,\dots,n$, at scales $t_i\in[T,t_1]$, all supported in one finite set
$\Omega=\bigcup_m\Omega_m$ of strict non-peaks of $f_0$, such that:
\begin{enumerate}
\item[(U1)] $t_i\|b^\pm_i\|_1\le A_0$ and $\Gamma^{(0)}_w(b^\pm_i,\omega^\pm_i)\le
1+\eta_0/2$ (computed at $f_0$);
\item[(U2)] $\supp\omega^\pm_{i,m}\subseteq\Omega_m$, and every $k\in\supp
\omega^\pm_{i,m}$ is of kind \textup{[1$'$]}: $\gap^0(k)>0$ and $|\omega^\pm_i(k)|
\le2\gap^0(k)/t_i$; or \textup{[2]}: $\gap^0(k)\ge\gamma_B$ and
$|\omega^\pm_i(k)|\le A_2/t_i$; or \textup{[3]}: $|\omega^\pm_i(k)|\le A_2/t_i$,
$\mathrm{vs}_k\omega^+_i(k)\le1.5\gap^0(k)/t_i$ and $\mathrm{vs}_k\omega^-_i(k)\ge
-1.5\gap^0(k)/t_i$, $\mathrm{vs}_k:=\sgn w^0_m(k)$;
\item[(U3)] for $j\in E$: $b^\pm_i(j)=\tilde b^\pm_i(j)+r^\pm_i(j)$ with
$r^+_i(j)=-r^-_i(j)$, where $\tilde b^\pm_i(j)$ is contact-like
($z_0(j)\tilde b^+_i(j)\ge0\ge z_0(j)\tilde b^-_i(j)$) or satisfies
$t_i|\tilde b^\pm_i(j)|\le K_3|a_0(j)|$, $K_3$ an $f$-constant;
\item[(U4)] $b^\theta_i(j)=0$ at every free coordinate $j\notin F_0\cup K_0$ with
$\mathrm{viol}_i(j)>0$, and the \emph{total violation}
$\epsilon_i:=\sum_{j\notin F_0}\mathrm{viol}_i(j)+2\sum_{j\in E}|r^+_i(j)|$ is
at most $\epsilon$;
\item[(U5)] for $I_-:=\{m:\Delta_{i,m}<0\text{ for some }i\}$:
$\mathrm{Scr}^{(0)}_m(x)\le C_Sx^2$ for $0<x\le x_0$ and $m\in I_-$
\textup{(}[I,~Definition~7.16] at $f_0$\textup{)};
\item[(U6)] every $l\in\Omega$ has a lever (Definition~\ref{def:III-levers}) at
$f_0$, the levers satisfy \textup{(SEP)}, there is a window quantity
$r_{\rm pull}\ge0$ such that for every $r\in[r_{\rm pull},r_0]$ every TU lever
has a pull coordinate $j_p\in\mathrm{Pull}_l$ with $v_l(j_p)\in[r,2^Gr]$, and at the lever coordinates the data
satisfy: at a (Z) coordinate $j$, $b^\theta_i(j)=0$; at a TU bank coordinate
$j'$, the data are contact-like of sign $\epsilon_l$ up to an antisymmetric
residue; at a TU pull coordinate $j_p$, $t_i|b^\pm_i(j_p)|\le(2A_2+C_\Delta)
\lambda_lv_l(j_p)$ up to an antisymmetric residue; all these residues and all
$|b^\pm_i(j)|$ at (Z) coordinates are included in $\epsilon_i$;
\item[(U7)] $|\Delta_{i,m}|\le C_\Delta$, $C_\Delta$ an $f$-constant.
\end{enumerate}
\emph{Coarse carriers} are those of $\Omega$ and the peaks of $f_0$ with margin
$\ge x_0$; all other carriers are \emph{fine}.  Window quantities: $n$, $T$,
$\gap_{\min}:=\min_{k\in\Omega}\gap^0(k)$, $x_0$, $C_S$, $|\Omega|$, the lever
constants $\Lambda_{\rm lev}:=\max(K_J,C_2,C_{\rm lev},\beta_{\max},2^G,1/r_0,
1/A_{\min},1/C_{\min})$, and the fine weights
\[
 c_{\rm fine}:=\sum_{k\text{ fine}}\lambda_k,\quad
 W_{\rm lev}:=\sum_{\substack{k\notin\Omega\\ k\text{ meets a lever coordinate}}}
 \lambda_k,
\]
and $c_{\rm tiny}$, $W_{\rm pull}$ are as in Subsection~\ref{subsec:III-DC}.
\end{definition}

Kind [1$'$] contains kind [1] of [I,~Lemma~8.42] (constant $2$, without
the condition $\gap\ge t^2$); it does not contain kind [1] of
[II,~Subsection~2.4], whose constant is $3$, and the constant $2$ is used
in Step~3 of the proof of Theorem~\ref{thm:III-UE}.  Kind [1$'$] covers
carriers pushed inward with $\Delta\omega=0$.

\begin{definition}[The $d$-consistent engineered approximant]\label{def:III-approx}
Given a window family, choose in this order:
\begin{enumerate}
\item[(i)] $N_w\ge\max F_0$ with
$\mathrm{tail}_W:=\max_i\big(\|b^\pm_i1_{(N_w,\infty)}\|_1+\|v_i1_{(N_w,\infty)}\|_1\big)$
at most $\min(T^3,c_{\rm late}\delta T/A_0)$;
\item[(ii)] $s_1\in(0,1]$; the \emph{mass set} $W$ consists of the contacts
$j\in[1,N_w]\setminus F_0$ of $f_0$ and of the $j\in E$ with
$\max_i|b^\theta_i(j)|>0$; the masses are $\mathsf m_j:=4\rho s_1\max_i
|b^\theta_i(j)|$ with sign $z_0(j)$;
\item[(iii)] the pre-row $f_1:=f[A_1,z_1]$, $A_1:=a_0+\sum_{j\in W}\mathsf m_j
z_0(j)e^*_j$, $z_1:=z_01_{[1,N'']}$ ($N''$ from (iv)); $\Pi_X:=\|e_1-e_0\|$
(it does not depend on $N''$); with $r_{\min}:=32K_JT^4s_1$ and
$\overline W_{\rm pull}:=\overline W_{\rm pull}(r_{\min})$,
\[
 r:=16K_J\Big(1+\frac1{A_{\min}}\Big)\big(\Pi_X+2T^4s_1\big)+\frac{32K_J
 \overline W_{\rm pull}}{A_{\min}},
\]
and then pull coordinates $j_p\in\mathrm{Pull}_l$ with $v_l(j_p)\in[r,2^Gr]$
\textup{(}they exist by (U6) when $r\ge r_{\rm pull}$\textup{)}.  Since $r\ge
r_{\min}$, these pulls have $v_l(j_p)\ge r_{\min}$, so the actual
$W_{\rm pull}$ satisfies $W_{\rm pull}\le\overline W_{\rm pull}$; the
definition of $r$ is not circular;
\item[(iv)] $N''>N_w$, larger than every lever coordinate, with
\begin{gather*}
 \mathfrak t(N'')\le T^4s_1^2,\qquad\max_{k\in\Omega}t_k(N'')\le T^4s_1,\qquad
 \sum_m\sum_k\Phi_m(k)1[t_{k,m}>s_1T^4]\le s_1^2T^4,\\
 \rho V_{>N''}\le\delta s_1/64,\qquad\|g_i1_{(N'',\infty)}\|_1\le T^4s_1\ \
 \text{for every piece};
\end{gather*}
\item[(v)] $P^*$ by Lemma~\ref{lem:III-DC} at $f_1$, and $f':=f[A(P^*),z(P^*)]$,
$a':=A(P^*)/q^*(A(P^*))$, $\hat x'$, $c'=1/p(\hat x')$, $w'_m$, $d'_m$,
$\sigma'_m$, $\nu'$ its data;
\item[(vi)] $g''_i:=b^\theta_i1_{[1,N_w]}+\sum_mR_m^*(\omega^\theta_{i,m}-
d'_m(\omega^\theta_{i,m})w'_m)$, $c_i:=g''_i(\hat x')$, $g'_i:=\rho(g''_i-c_ia')$.
\end{enumerate}
The \emph{perturbation size} is $\Pi:=\Pi_X+r+\mathfrak t(N'')^{1/2}$.
\end{definition}

$f'$ is norm attaining ($A(P^*)$ and $z(P^*)$ are finitely supported,
[I,~Proposition~1.10(c)]), $g'_i(\hat x')=0$, and, once Lemma~\ref{lem:III-DC}
applies, $f'$ is $d$-consistent with $f_0$ on $\Omega$.  A TU pull coordinate
inside the window that carries a mass gets the total mass $-\epsilon_l
(\mathfrak p_l+\mathsf m_j)$; this only enlarges $|a'_j|$.

\begin{lemma}[Estimates at the approximant]\label{lem:III-UE1}
There is a window constant
\[
 K_w\le C_f\Big(\frac{nA_0|\Omega|\Lambda_{\rm lev}C_S}{T\gap_{\min}x_0}\Big)^{C_{\rm abs}},
\]
with $C_f$ an $f$-constant and $C_{\rm abs}$ an absolute constant, such that the
following holds.  If $K_J2^{G+3}|\Omega|c_{\rm tiny}/A_{\min}\le\frac18$,
$K_0\|Y\|_\infty\le\frac12$ \textup{(}$Y$ computed at $f_1$\textup{)},
$r_{\rm pull}\le32K_JT^4s_1$ and $\Pi\le1/K_w$, then Lemma~\ref{lem:III-DC}\textup{(ii)} applies at $f_1$ with the
radius $r$ of Definition~\ref{def:III-approx}(iii), and:
\begin{enumerate}
\item[(a)] \textup{(sizes)} $\Pi_X\le8\rho s_1\max_j\mu^\star_j\,nA_0/(T\nu_0)$;
$q^*(A(P^*)-a_0)\le K_w(s_1+\overline W_{\rm pull})$; and
$p^*(f'-f_0)\le K_w\big(\Pi\log(e/\Pi)+W_{\rm pull}+W_{\rm lev}\big)$;
\item[(b)] \textup{(values)} a carrier $k$ meeting no lever coordinate satisfies
$|u_k(\hat x')-u_k(\zh_0)|\le2(\Pi_X+C_2r)+t_k(N'')$; carriers of $\Omega$
satisfy $\mathrm{nv}'_k=\mathrm{nv}^0_k$; a carrier meeting lever coordinates
moves by at most $2(\Pi_X+C_2r)+t_k(N'')+C_{\rm lev}r+2\sum_{j_p}|u_k(j_p)|$;
\item[(c)] \textup{(block data)} $|C'_m-C^0_m|+|M'_m-M^0_m|+|A'_m-A^0_m|+|c'-q^0_0|
+|\nu'-\nu_0|+\|e'-e_0\|\le K_w\Pi$; every peak of $f_0$ with margin $>K_w\Pi$
is a peak of $f'$ with the same sign; every $k\in\Omega$ has $\gap'(k)\ge
\gap^0(k)/2$; and $d'_m(\omega)=d_m(\omega)$, $H'_m(\omega)=(C^0_m/C'_m)
H_m(\omega)$ for $\omega$ supported in $\Omega_m$;
\item[(d)] \textup{(Bregman terms)} $0\le\mathrm{Bx}_m:=\ip{w'_m-w^0_m}{R_m\hat x'}
\le K_w\Pi^2+4c_{\rm fine}(\Pi_X+C_2r)+2\mathfrak t(N'')+2C_{\rm lev}W_{\rm lev}r
+4W_{\rm pull}$;
\item[(e)] \textup{(scrambling, $m\in I_-$)} with $A_m$ the scrambled set of
[I,~Lemma~7.15] for $(w^0_m,w'_m)$,
\[
 \mathfrak S_m:=\|D_m(w'_m-w^0_m)\|_2^2+\sum_{k\in A_m}\Big(\Phi_m(k)^2+
 \lambda_k\big(|w'_m(k)-w^0_m(k)|+(C'_m-C^0_m)_+\big)\Big)
\]
satisfies $\mathfrak S_m\le K_w\Pi^2\log(e/\Pi)+4m\big(C_S(K_w\Pi)^2+s_1^2T^4\big)
+4(W_{\rm lev}+W_{\rm pull})$;
\item[(f)] \textup{(base form)} for every piece and $\diamond\in\{+,-,\theta\}$,
with $\beta^\theta:=b^\theta1_{[1,N_w]}$, $\beta^\pm:=b^\pm1_{[1,N_w]}\pm
\frac12v1_{(N_w,\infty)}$ and $h'(\beta):=\|P'^\perp U^*\beta\|^2/\nu'$:
$|c'h'(\beta^\diamond)-q^0_0h^0(b^\diamond)|\le C_f\big(A_0^2K_w\Pi/T^2+A_0\,\mathrm{tail}_W/T\big)$ with $C_f$ an $f$-constant;
\item[(g)] \textup{(targets)} $(g''_i-g_i)(\hat x')=-(b^\theta_i1_{(N_w,\infty)})
(\hat x')-\sum_md^\theta_{i,m}\mathrm{Bx}_m$, hence
\[
 |c_i|\le(1+\|U\|)\mathrm{tail}_W+\|g_i1_{(N'',\infty)}\|_1+K_w(\Pi+W_{\rm pull})
 /t_i,
\]
and $p^*(g'_i-\rho g_i)\le K_w(\Pi+c_{\rm fine}+W_{\rm lev}+W_{\rm pull}+
\mathrm{tail}_W)/t_i$.
\end{enumerate}
\end{lemma}

\begin{proof}
\emph{Lemma~\ref{lem:III-DC} applies.}  By Lemma~\ref{lem:III-masses}(iii),
$\delta_1\le\Pi_X+\max_{\Omega}t_k(N'')\le\Pi_X+T^4s_1$ and $\Delta_1\le\Pi_X+
\mathfrak t(N'')\le\Pi_X+T^4s_1$.  By Lemma~\ref{lem:III-DC}(i), the choice
of $r$ and $W_{\rm pull}\le\overline W_{\rm pull}$, the four terms of
$K_J|F(0)|_\infty$ are bounded by
$K_J(\delta_1+\Delta_1/A_{\min})\le K_J(1+1/A_{\min})(\Pi_X+T^4s_1)\le r/16$,
$2K_JW_{\rm pull}/A_{\min}\le r/16$, $K_JC_22^{2G+2}r\cdot r\le r/8$ and
$K_J2^{G+2}r|\Omega|c_{\rm tiny}/A_{\min}\le r/16$, so
$K_J|F(0)|_\infty\le5r/16\le r/2$, as soon as $K_JC_22^{2G+2}r\le\frac18$;
this, $3K_JC_2r\le\frac18$ and $r\le r_0$ follow from $r\le\Pi\le1/K_w$.
Also $\delta_1\le r$ and $\Delta_1\le r$, since $r\ge16K_J(\Pi_X+T^4s_1)$.
The pulls exist because $r\ge r_{\min}\ge r_{\rm pull}$ by
\textup{(PULL)} of Theorem~\ref{thm:III-UE}.

(a) By Lemma~\ref{lem:III-masses}(ii), (iv) (masses on $E\subseteq F_0$ cost the
factor $2$), with $t_i\ge T$ and $\|b^\theta_i\|_1\le A_0/t_i$.  By
Lemma~\ref{lem:III-DC}, $|P^*|\le r$, bank masses are $\le\beta_{\max}2^{G+2}r$,
pull masses $\le24\lambda_l2^Gr$, $z$-moves $\le r/v_l$; and $r\le K_w(s_1+
\overline W_{\rm pull})$ by (iii).  For the cost use [II,~Lemma~2.12(iv)]:
$p^*(f'-f_0)\le(1+\|U\|)(\|a'-a_0\|_1+C_fc(\delta))$ with $\delta:=\hat x'-\zh_0$;
by (b), $\sum_k\min(\lambda_k,|u_k(\delta)|)\le\sum_k\min(\lambda_k,4K_w\Pi)+
2\mathfrak t(N'')+W_{\rm lev}+W_{\rm pull}$, and [II,~(1)] bounds the first sum
by $C\Pi\log(e/\Pi)$; $\Delta_m(\delta)$ is bounded likewise.  (b) is
Lemma~\ref{lem:III-masses}(iii) with $\|e'-e_0\|\le\|e'-e_1\|+\Pi_X$ and
$\|e'-e_1\|\le2\|X_{\rm lev}\|/\nu_1\le2C_2r/\nu_1$, together with
Lemma~\ref{lem:III-DC}(ii).  (c) $|A'_m-A^0_m|\le\sum_k\lambda_k|u_k(\hat x')-
u_k(\zh_0)|$; $|C'-C|$ by the clamp equation through $k^\natural_m$
(Lemma~\ref{lem:III-dcons}(c) with $\Omega=\emptyset$, i.e.\ [I,~Lemma~7.14]'s
argument, on the ball of Lemma~\ref{lem:III-ball}); $q_0=1/(1+\sum_mA_m)$;
$\nu',e'$ by Lemma~\ref{lem:III-masses}.  Peak margins are Lipschitz in the
value and in $\bar\vartheta$ [I,~Lemma~1.7].  The gap clause is
Lemma~\ref{lem:III-DC}(iii), because $K_w\Pi\le C_{\min}\gap_{\min}/4$ when
$K_w\ge4/(C_{\min}\gap_{\min})$ times the constant of the $C$-bound; the rest is
Lemma~\ref{lem:III-dcons}(a).
(d) By [I,~Lemma~7.13(E5)], $0\le\mathrm{Bx}_m\le\ip{w'_m-w^0_m}{R_m(\hat x'-
\zh_0)}=\sum_k\lambda_k(w'_m(k)-w^0_m(k))u_k(\hat x'-\zh_0)$.  Coarse carriers
have $|w'-w^0|\le K_w\Pi$ (peaks: $|M'-M|$; $\Omega$: $|C'/C-1|$) and values
$\le2\Pi$; fine carriers have $|w'-w^0|\le2$ and contribute as displayed by (b).
(e) As in the proof of [I,~Lemma~7.17] with the perturbation size $\Pi$ in place
of $s_1$.  Let $K_*$ be the $f$-constant collecting the Lipschitz constants of
values, thresholds and the clamp formula.  By (b), (c), a peak of $f_0$ with
margin $>K_*\Pi$ lies in $S_1$; a carrier of $\Omega$ lies in $S_2$ as soon as
$\gap^0\ge2|C'-C|(1+(M-\gap^0)/C)$, because $w'=(C'/C)w^0$ on $\Omega$; a
strict non-peak with $\Phi\gap$ and $\gap$ $>K_*\Pi$ lies in $S_2$.  So
$A^{(1)}:=\{k\in A_m:t_k\le\Pi\}$ consists of carriers meeting lever
coordinates and of carriers with $\mu\le2K_*\Pi$ or $\Phi\gap\le2K_*\Pi$ or
$\gap\le2K_*\Pi$, whose $\Phi$-mass is $\le\mathrm{Scr}^{(0)}_m(2K_*\Pi)\le
C_S(2K_*\Pi)^2$ by (U5) ($2K_*\Pi\le x_0$).  As $\Pi\ge r\ge16K_JT^4s_1$,
$A^{(2)}\subseteq\{t_k>s_1T^4\}$, of $\Phi$-mass $\le s_1^2T^4$ by (iv).  With
$|w'-w|\le2$, $(C'-C)_+\le1$, $\Phi\le1$ and $\lambda_k=m\Phi_m(k)$ the sum over
$A_m$ is at most $4m$ times the $\Phi$-mass; carriers meeting lever coordinates
contribute $\le4\sum\lambda_k$.  Finally $\|D(w'-w^0)\|^2\le K_w\Pi^2\log(e/\Pi)$
by the last display of the proof of [I,~Lemma~7.14] with $K_Fs_1$ replaced by
$K_w\Pi$.
(f) Three changes: $\beta^\diamond-b^\diamond$ has $\ell_1$-norm $\le
\mathrm{tail}_W$ and $|h(y)-h(y')|\le\|U\|^2(\|y\|_1+\|y'\|_1)\|y-y'\|_1/\nu$;
$|\|P'^\perp y\|^2-\|P^\perp y\|^2|=|\ip ye^2-\ip y{e'}^2|\le2\|y\|^2\|e'-e\|$
with $\|y\|\le2\|U\|A_0/T$ and $\|e'-e_0\|\le K_w\Pi$; and $|c'/\nu'-q^0_0/\nu_0|\le K_w\Pi$; the constants in front of $A_0^2K_w\Pi/T^2$ and $A_0\mathrm{tail}_W/T$ are $f$-constants ($\|U\|$, $1/\nu_0\le2/\nu$, $1/q^0_0$).
(g) $g''_i-g_i=-b^\theta_i1_{(N_w,\infty)}+\sum_md^\theta_{i,m}R_m^*(w^0_m-w'_m)$
by $d$-consistency ($d'_m(\omega^\theta)=d_m(\omega^\theta)=d^\theta_{i,m}$), and
$R_m^*(w^0_m-w'_m)(\hat x')=\ip{w^0_m-w'_m}{R_m\hat x'}=-\mathrm{Bx}_m$.  So
$c_i=(g''_i-g_i)(\hat x')+g_i(\hat x'-\zh_0)$ (as $g_i(\zh_0)=0$), with
$\|\hat x'\|_\infty\le1+\|U\|$, $|d^\theta_{i,m}|\le\|D_m\omega^\theta\|_2\le
\max(2,A_2)\|\Phi_m\|_2/t_i$, and $g_i(\hat x'-\zh_0)=\ip{U^*g_i}{e'-e_0}+
\sum_j g_i(j)(z'_j-z_0(j))$: beyond $N''$ this is $\le\|g_i1_{(N'',\infty)}\|_1$;
at a pull, $|g_i(j_p)|\le K(\lambda_lv_l(j_p)+W_{\rm pull})/t_i\le
K(2^Gr+W_{\rm pull})/t_i$; at a (Z) coordinate the move is $\le r/v_l(j)$ and
$|g_i(j)|\le|b_i(j)|+K\lambda_lv_l(j)/t_i+K W_{\rm lev}/t_i$.  The bound for
$p^*(g'_i-\rho g_i)$ follows from $p^*(\cdot)\le(1+\|U\|)\|\cdot\|_1$,
$\|R_m^*(w^0-w')\|_1\le\sum_k\lambda_k|w^0(k)-w'(k)|\le K_w\Pi+2(c_{\rm fine}+
W_{\rm lev}+W_{\rm pull})$ and $q^*(a')=1$.

\emph{The form of $K_w$.}  Every estimate above, in Lemma~\ref{lem:III-DC} and in
the proof of Theorem~\ref{thm:III-UE} is a fixed finite product of the
quantities $n$, $A_0/T$, $1/\gap_{\min}$, $1/x_0$, $C_S$, $|\Omega|$ (entering
linearly through sums over levers), the lever constants and $f$-constants; no
estimate is iterated a window-dependent number of times (the Sherman--Morrison
inverse is rank one per block, the Poincar\'e--Miranda argument involves no
iteration, the clamp equation is used through one natural peak, and sums over
$\Omega$ enter as factors $|\Omega|$).  Hence $K_w$ has the stated form.
\end{proof}

\subsection{The uniform violation-tolerant per-piece bound}\label{subsec:III-UEthm}

\begin{theorem}[Uniform engineered bound]\label{thm:III-UE}
There are $f$-constants $t_1\in(0,1]$, $c^f_{\rm flat}\in(0,\frac1{16}]$,
$K_\sharp\ge1$, $\eta_1\in(0,\frac14]$ and $c_{\rm late}>0$, none of which
depends on $\gamma_B$, with the following property.  Put
\[
 c_{\rm flat}:=\min\big(c^f_{\rm flat},\ \gamma_B/(8\rho A_2)\big).
\]
Let $f_0$ be a base row and let a window family \textup{(U1)--(U7)} with scales
in $[T,t_1]$ be given at $f_0$, with window constant $K_w$ of
Lemma~\textup{\ref{lem:III-UE1}}.  Assume the \emph{window conditions}
\begin{enumerate}
\item[(W-a)] $c_{\rm fine}\le c_{\rm late}\delta T/\big(nA_0(1+C_2K_J)(1+
1/A_{\min})\big)$;
\item[(W-b)] $K_J2^{G+3}|\Omega|c_{\rm tiny}/A_{\min}\le\frac18$ and
$K_0\|Y\|_\infty\le\frac12$ at the pre-row;
\item[(W-c)] $N_w$ is chosen as in Definition~\textup{\ref{def:III-approx}(i)};
\end{enumerate}
and let $s_1$ satisfy
\begin{gather*}
 \textup{(VT$'$)}\ \ 64\rho\epsilon\le\delta s_1,\qquad
 \textup{(FINE)}\ \ W_{\rm lev}+\overline W_{\rm pull}\le\frac{c_{\rm late}\delta A_{\min}
 T^4s_1}{1+C_{\rm lev}},\\ \textup{(PULL)}\ \ r_{\rm pull}\le32K_JT^4s_1,\qquad
 \textup{(LATE)}\ \ s_1\le s_{\rm late}:=c_{\rm late}\Big(\frac{T\gap_{\min}x_0}
 {nA_0K_w}\Big)^7 .
\end{gather*}
Then the approximant $f'$ of Definition~\textup{\ref{def:III-approx}} \textup{(}for
every $N''$ as in (iv)\textup{)} is norm attaining and $d$-consistent with $f_0$
on $\Omega$, $\Pi\le K_ws_1$, $p^*(f'-f_0)\le K_w\big(\Pi\log(e/\Pi)+s_1\big)$,
and for every piece
\begin{gather}
 p^*(f'+\tau g'_i)\le1+\frac{\tau^2}2(1-\delta)\qquad(0<|\tau|\le c_{\rm flat}t_i),
 \label{eq:III-UE}\\
 p^*(g'_i-\rho g_i)\le K_w\big(K_ws_1+c_{\rm fine}+\mathrm{tail}_W\big)/t_i .
 \notag
\end{gather}
\end{theorem}

\begin{proof}
\emph{Step 0 (constants).}  The order of choices is: $f\to K_\sharp\to\eta_1\to$
transfer data and $r_f\to c^f_{\rm flat},t_1,c_{\rm late}\to$ window $\to N_w\to
s_1\to$ pre-row $\to r\to$ pulls $\to N''\to P^*$.  For rows in $B_f$
(Lemma~\ref{lem:III-ball}), (U1) gives $h^0(b^\diamond)\le4/q_0$ and
$H_m(\omega^\diamond)\le4/\sigma_m$.  Put $K_\sharp:=4\max(1/q_0,\max_m1/\sigma_m)
(1+\eta_0)+2$ and $\eta_1:=\min(\frac14,\delta/(192|I|K_\sharp))$; choose transfer
data of $f$ with this $\eta_1$ and $r_f:=r_f(\eta_1)$.  Let $a_{\min}:=\min_F|a_j|$,
$\gamma_0:=\min_m(M_m-\gamma_m)$, $\gamma_{\min}:=\min_m\gamma_m$,
$C_{\min}:=\min_mC_m$, $M_{\min}:=\min_mM_m$, $\Lambda_{\max}:=\max\Lambda^\varsigma_m$.
Let $K_W:=C_W(A_0+A_2+C_\Delta)$ with an $f$-constant $C_W$ such that
$\|W^\natural_m-w'_m\|_\infty+\|D_m(W^\natural_m-w'_m)\|_2\le(K_W/t_i)|\tau|$ in
Step~2, $K_A:=2(1+\|U\|)(2A_0+1)+2$, and $K'$ an $f$-constant times
$(A_0+A_2+C_\Delta+1)$ bounding the relative errors of [I,~Lemma~3.3(b)] and
[I,~Lemma~3.4(d)] per unit $c_{\rm flat}$.  Then $c^f_{\rm flat}\le\frac1{16}$ is
the largest number such that $c_{\rm flat}$ satisfies:
\begin{enumerate}
\item[(c1)] $\rho c_{\rm flat}A_0\le a_{\min}/32$;\quad (c2) $\rho c_{\rm flat}K_3
\le\frac1{16}$;\quad (c3) $\rho c_{\rm flat}(2A_2+C_\Delta)\le4$;
\item[(c4)] $\rho c_{\rm flat}\le\frac1{16}$, $\rho c_{\rm flat}(A_2+2)\le M_{\min}/8$,
$\rho c_{\rm flat}(A_2+2)\|\Phi_m\|_2\le C_{\min}/8$, and $8\rho c_{\rm flat}A_2\le
\gamma_B$;
\item[(c5)] $K_Wc_{\rm flat}\le\min(\gamma_0/8,C_{\min}/4)$;
\item[(c6)] $6|I|K_\sharp K_YK_Ac_{\rm flat}+48K_\sharp K_WK_yc_{\rm flat}/(C_{\min}/2)
\le\delta/32$;
\item[(c7)] $K'c_{\rm flat}(1+\rho^2)\le\delta/8$.
\end{enumerate}
Only the last inequality of (c4) involves $\gamma_B$, and every condition is an
upper bound on $c_{\rm flat}$; so $c_{\rm flat}=\min(c^f_{\rm flat},\gamma_B/
(8\rho A_2))$.  Then $t_1\in(0,1]$ is chosen with
\begin{gather*}
 6K_\sharp c_{\rm flat}^2t_1^2(2+\Lambda_{\max})\le\gamma_{\min}/2,\qquad
 6K_\sharp c_{\rm flat}^2t_1^2\le\gamma_0/8,\\
 6|I|K_\sharp K_Yc_{\rm flat}^2t_1^2\le1,\qquad
 36K_\sharp^2K_y^2c_{\rm flat}^2t_1^2/(C_{\min}/2)\le\delta/32
\end{gather*}
(these hold for $\gamma_B\le1$ when they hold with $c^f_{\rm flat}$).  Finally
$c_{\rm late}$ is a small $f$-constant fixed in Steps~3--5.

\emph{Consequences of (LATE), (W-a)--(W-c), (FINE).}  Lemma~\ref{lem:III-UE1}(a)
and the definition of $r$ give $\Pi\le K_ws_1$ (enlarging $K_w$).  As $K_w\ge1$
and $T,\gap_{\min},x_0\le1$, (LATE) with exponent $7$ implies:
$\Pi\le1/K_w$; $2K_w\Pi\le x_0$; $K_w^3s_1\log(e/s_1)\le c_{\rm late}\delta$;
$K_w^2s_1\le T^5$; $s_1\rho\max(2,A_2)/T\le\gap_{\min}/4$; $q^*(A(P^*))\le2$;
$p^*(f'-f)\le r_f$; and $\rho c_{\rm flat}t_i|c_i|\le\frac12$ (by
Lemma~\ref{lem:III-UE1}(g), (W-c) and (iv)).  So Lemma~\ref{lem:III-UE1}
applies, $f'$ lies in $B_f$, and Lemma~\ref{lem:III-ball} holds at $f'$.  By
Lemma~\ref{lem:III-UE1}(d), (e), $\Pi\le K_ws_1$, (W-a), (FINE), (LATE) and (iv):
\begin{equation}\label{eq:III-BxS}
 \mathrm{Bx}_m+\mathfrak S_m\le c'_{\rm late}\delta s_1\qquad(m\in I),
\end{equation}
with $c'_{\rm late}$ as small as we wish when $c_{\rm late}$ is small (for
$m\notin I_-$ only $\mathrm{Bx}_m$ is used).  Indeed
$K_w\Pi^2\log(e/\Pi)\le K_w^3s_1^2\log(e/s_1)$, $4c_{\rm fine}(\Pi_X+C_2r)\le
Cc_{\rm late}\delta s_1$ by (W-a) and the bounds of Definition~\ref{def:III-approx}
(iii), and $2C_{\rm lev}W_{\rm lev}r+4W_{\rm pull}+4(W_{\rm lev}+W_{\rm pull})\le
Cc_{\rm late}\delta s_1$ by (FINE).

\emph{Steps 1--2 (targets and decompositions).}  For $\diamond\in\{+,-,\theta\}$
put $B^\diamond:=\rho(\beta^\diamond-c_ia')$, $\Omega^\theta_m:=\rho(\omega^\theta_m-
d'_m(\omega^\theta_m)w'_m)$ and $\Omega^\pm_m:=\rho(\omega^\pm_m-d'_m(\omega^\theta_m)
w'_m-(d^\pm_m-d^\theta_m)w^0_m)$, $d^\diamond_m:=d_m(\omega^\diamond_m)$.  As in
Step~2 of the proof of [I,~Theorem~7.18], $f'+\tau g'_i=(a'+\tau B^\diamond)+
\sum_mR_m^*(w'_m+\tau\Omega^\diamond_m)$ for every $\tau$; the identity uses only
the two representations at $f_0$, so it holds for violated pairs
([II,~Lemma~2.30], proof).  By $d$-consistency (Lemma~\ref{lem:III-UE1}(c)),
$\kappa^\pm_m=0$ and
\[
 \Omega^\pm_m=\rho\big(\omega^\pm_m-d'_m(\omega^\pm_m)w'_m\big)+\rho(d^\pm_m-
 d^\theta_m)(w'_m-w^0_m).
\]
We use $\diamond=\theta$ for $0<|\tau|\le s_1$ and $\diamond=\sgn\tau$ for
$s_1<|\tau|\le c_{\rm flat}t_i$ (if $c_{\rm flat}t_i\le s_1$ only the first
regime occurs).  The anchors $V^{\rm an}_m$ ($=w^0_m$ if $\Delta_m\ge0$, $=y_{A,m}$
of [I,~Lemma~7.15] for $(w^0_m,w'_m)$ if $\Delta_m<0$), the vectors $W^\natural_m$
and the remainder $\mathcal Z:=\sum_{m\in I_-}r_mR_m^*Z_{A,m}$ are as there, with
$r_m:=\frac12|\tau|\rho|\Delta_m|\le|\tau|C_\Delta$.  Since
$|\omega^\diamond(k)|\le\max(2,A_2)/t_i$, $\|D\omega^\diamond\|_2\le\max(2,A_2)
\|\Phi_m\|_2/t_i$, $|d^\diamond|\le\|D\omega^\diamond\|_2$, $|\Delta_m|\le
C_\Delta$ and $\|V^{\rm an}_m-w'_m\|_\infty\le3$, the bound defining $K_W$ holds.

\emph{Step 3: blocks.}  As $\omega^\diamond_m$ vanishes on $P'_m$ (statuses of
$\Omega$ and peaks preserved), [I,~Lemma~3.2] at $f'$ gives
$\lin'_m(W^\natural_m)=r_m\ip{V^{\rm an}_m-w'_m}{R_mx'}/\sigma'_m$ (no $\kappa$
term).  Its modulus is at most $c'(\mathrm{Bx}_m+\mathfrak S_m)$ (for $y_{A,m}$
as in Step~3 of [I,~Theorem~7.18], with $|(R_mx')(k)|\le\lambda_kc'$), so by
\eqref{eq:III-BxS} $|\lin'_m|\le|\tau|C_\Delta(2/\sigma_m)c'_{\rm late}\delta s_1\le
\delta|\tau|s_1/(64|I|)$.  For the excess write $W_{0,m}-r_mw'_m=(1-r_m)
(w'_m+t'(\omega^\diamond_m-d'_m(\omega^\diamond_m)w'_m))$, $t':=\tau\rho/(1-r_m)$,
$|t'|\le2\rho c_{\rm flat}t_i$.  Put $\mathrm W:=w'+t'(\omega-d'(\omega)w')$.  For
$k$ of kind [1$'$]: $|t'\omega(k)|\le4\rho c_{\rm flat}\gap^0(k)\le\gap'(k)/2$,
since $\gap'\ge\gap^0/2$ (Lemma~\ref{lem:III-UE1}(c)) and $\rho c_{\rm flat}\le
\frac1{16}$.  Kind [2]: $|t'\omega(k)|\le2\rho c_{\rm flat}A_2\le\gamma_B/4\le
\gap'(k)/2$.  Kind [3]: with $\mathrm{vs}=\sgn w'(k)=\sgn w^0(k)$, the outward part
satisfies $\mathrm{vs}\,t'\omega(k)\le3\rho c_{\rm flat}\gap^0(k)\le\gap'(k)/2$, so
$\mathrm{vs}\,\mathrm W(k)\le(1-d't')(M'-\gap')+\gap'/2\le(1-d't')M'$ as
$|d't'|\le\frac12$ by (c4); and $-\mathrm{vs}\,\mathrm W(k)\le|t'\omega(k)|\le
2\rho c_{\rm flat}A_2\le(1-d't')M'$ by (c4) and $M'\ge\frac34M$.  Hence
$\|\mathrm W\|_\infty=(1-d't')M'$, the conclusion of [I,~Lemma~3.4(c)], whose
proof uses only these coordinatewise bounds; then the computation of
[I,~Lemma~3.4(d)] (which needs only that conclusion and $|d't'|M'/C'\le\frac12$,
given by (c4)) yields $N_m(\mathrm W)\le1+\frac{t'^2}2H'_m(\omega)(1+2|d't'|M'/C')$.
For $\diamond=\theta$ and $|\tau|\le s_1$ the radius condition
$|s_1\rho\omega^\theta(k)|\le\gap'(k)/2$ follows from (LATE).  Therefore, as in
[I,~Theorem~7.18],
\[
 G'_m(W^\natural_m)\le\hat G_m:=\frac{\rho^2\tau^2}2H'_m(\omega^\diamond_m)
 (1+K'c_{\rm flat})+r_m\mathfrak g_m,\qquad\mathfrak g_m\le K(\mathrm{Bx}_m+
 \mathfrak S_m),
\]
with $H'_m(\omega)=(C^0_m/C'_m)H_m(\omega)$ and $r_m\mathfrak g_m=0$ for $\theta$.

\emph{Step 3: base.}  By [I,~Lemma~7.1(b)] at $f'$, $G'_b(A_\tau)=\Exc'(\tau
B^\diamond+\mathcal Z)+\nu'\Psi'(U^*(\tau B^\diamond+\mathcal Z)/\nu')$ and
$\Exc'(\tau B^\diamond+\mathcal Z)\le\Exc'(\tau B^\diamond)+2\|\mathcal Z\|_1$.  We
claim
\begin{equation}\label{eq:III-base}
 \Exc'(\tau B^\pm)\le\tfrac12\rho|\tau|V_{>N''}+2\rho|\tau|\epsilon_i\ \ (s_1<|\tau|),
 \qquad\Exc'(\tau B^\theta)=0\ \ (0<|\tau|\le s_1).
\end{equation}
On $\supp a'$ the summand is $2(-\sgn(a'_j)\tau B_j-|a'_j|)_+$ with $a'_j+\tau
B_j=(1-\tau\rho c_i)a'_j+\tau\rho\beta_j$ and $1-\tau\rho c_i\ge\frac12$; it
vanishes if $|\tau\rho\beta_j|\le\frac12|a'_j|$ or if $\tau\rho\beta_j$ has the
sign of $a'_j$, and the map $\beta_j\mapsto$ summand is $2\rho|\tau|$-Lipschitz.
On $F$: $|a'_j|\ge a_{\min}/4$ and $|\tau\rho\beta_j|\le\rho c_{\rm flat}A_0\le
a_{\min}/32$ by (c1).  On $E$: the residue costs at most $2\rho|\tau||r^+_i(j)|$ in
the $\pm$ regimes and nothing in the $\theta$ regime ($r^\theta=0$); for the rest,
contact-like data have the sign of $a'_j$ on the designated side, carry a mass
$\mathsf m_j$ when $b^\theta(j)\ne0$ (so $|\tau\rho\beta^\theta_j|\le\mathsf m_j/4
\le|a'_j|/2$ for $|\tau|\le s_1$), and the alternative gives $|\tau\rho\tilde
b_j|\le\rho c_{\rm flat}K_3|a_0(j)|\le|a'_j|/8$ by (c2), since $|a'_j|\ge
|a_0(j)|/2$.  Window contacts with mass: for $\theta$, $|\tau\rho\beta^\theta_j|\le
\mathsf m_j/4\le|a'_j|/2$; for $\pm$ the summand is $\le2\rho|\tau|\mathrm{viol}(j)$
as in [II,~Lemma~2.30].  TU pulls: $|\tau\rho\beta^\pm_j|\le\rho c_{\rm flat}
(2A_2+C_\Delta)\lambda_lv_l(j_p)+\rho|\tau||r(j_p)|\le4\lambda_lv_l(j_p)+\rho|\tau|
|r(j_p)|$ by (c3), while $|a'_{j_p}|\ge\mathfrak p_l/q^*(A(P^*))\ge12\lambda_l
v_l(j_p)$: no flip from the main part, and the residue is paid as a violation
(only for $\pm$).  TU banks: contact-like data of sign $\epsilon_l=\sgn a'_{j'}$,
protected by a mass in the $\theta$ regime when $b^\theta(j')\ne0$ and $j'\le N_w$
(and $\beta^\theta_{j'}=0$ for $j'>N_w$).  (Z) coordinates: $\beta^\theta_j=0$, so
the $\theta$ summand vanishes whatever the state of $j$ at $f'$ (free, contact or
bank), and the $\pm$ summand is $\le2\rho|\tau||b^\pm_i(j)|$, included in
$\epsilon_i$.  Off $\supp a'$, at window contacts without mass, free coordinates
$j\le N_w$ (using (U4)), contacts in $(N_w,N'']$, free $j>N_w$ and contacts
$j>N''$: exactly as in the proof of [II,~Lemma~2.30].  Summing gives
\eqref{eq:III-base}.  Further $\|\mathcal Z\|_1\le|\tau|C_\Delta\sum_{m\in I_-}
\mathfrak S_m$ and $P'^\perp U^*a'=0$, so by [I,~Lemma~3.3]
\[
 G'_b(A_\tau)\le\hat G_b:=\frac{\rho^2\tau^2}2h'(\beta^\diamond)(1+K'c_{\rm flat})
 +\tfrac12\rho|\tau|V_{>N''}+2\rho|\tau|\epsilon_i+2|\tau|C_\Delta\sum_m\mathfrak S_m,
\]
the last three terms being absent for $\theta$.  For $|\tau|>s_1$:
$\frac12\rho V_{>N''}\le\delta s_1/128$ (iv), $2\rho\epsilon_i\le\delta s_1/32$
(VT$'$) and $2C_\Delta\sum_m\mathfrak S_m\le\delta s_1/128$ by \eqref{eq:III-BxS};
so these terms are $\le\frac\delta{16}|\tau|s_1\le\frac\delta{16}\tau^2$.  By
Lemma~\ref{lem:III-UE1}(c), (f), $H'\le2H$ and $h'\le2h^0$, hence
$\hat G_b,\hat G_m\le K_\sharp\tau^2$ for $|\tau|\le c_{\rm flat}t_i$.

\emph{Step 4: rebalancing.}  Apply [I,~Proposition~7.3] at $f'$ with $h:=\tau
g'_i$ ($h(x')=0$), the decomposition of Step~2, the bounds $\hat G_b,\hat G_m$,
and the transfer vectors of $f'$ (Lemma~\ref{lem:III-ball}(b)).  Then
$|\varepsilon_m|\le2(|\lin'_m|+\hat G_m+\hat\Gamma)\le6K_\sharp\tau^2$
($|\lin'_m|\le\tau^2$ for $|\tau|>s_1$, and $\lin'_m=0$ for $\theta$).  Condition
\textup{(R1)} follows from [I,~Lemma~7.2] with $\eta:=(K_W/t_i)|\tau|\le
K_Wc_{\rm flat}$, by (c5) and the conditions on $t_1$ (here $|M'_m-M_m|\le
\gamma_0/4$ is used), and $|\lambda|\le6|I|K_\sharp K_Y\tau^2\le1$.  The error
satisfies
\[
 E\le12|I|K_\sharp\eta_1\tau^2+6|I|K_\sharp K_Y|\tau|^3\frac{K_A}{t_i}+
 \frac{48K_\sharp K_WK_y|\tau|^3/t_i+36K_\sharp^2K_y^2\tau^4}{C_{\min}/2}
 \le\frac\delta8\tau^2,
\]
where $K_A/t_i$ bounds $(|q^*(A_\tau)-1|+q^*(A_\tau-a'))/|\tau|$ (as
$q^*(B^\diamond)\le(1+\|U\|)\rho(2A_0/t_i+|c_i|)$, $|c_i|\le1/t_i$ and
$\|\mathcal Z\|_1/|\tau|\le1$), by the choice of $\eta_1$, (c6), $|\tau|\le
c_{\rm flat}t_i$ and $\tau^4\le c_{\rm flat}^2t_1^2\tau^2$.

\emph{Step 5: levels.}  $\hat\Gamma=c'\hat G_b+\sum_m\sigma'_m\hat G_m\le
\frac{\rho^2\tau^2}2\big(c'h'(\beta^\diamond)+\sum_m\sigma'_mH'_m(\omega^\diamond_m)
\big)(1+K'c_{\rm flat})+(\text{linear terms})$.  By Lemma~\ref{lem:III-UE1}(c),
(f), the bracket is at most $\Gamma^{(0)}_\diamond+C_f(A_0^2K_w\Pi/T^2+A_0
\mathrm{tail}_W/T+K_w\Pi)\le\Gamma^{(0)}_\diamond+\delta/8$ by (LATE), (W-c) and
the choice of $c_{\rm late}$; and $\Gamma^{(0)}_\theta\le\max(\Gamma^{(0)}_+,
\Gamma^{(0)}_-)\le1+\eta_0/2$ by convexity of $\Gamma_w$ (which does not need
admissibility).  The linear terms are at most $\frac\delta{16}\tau^2$ (base) and
$\sum_m\sigma'_mr_m\mathfrak g_m\le|\tau|C_\Delta K\sum_m(\mathrm{Bx}_m+
\mathfrak S_m)\le\frac\delta{32}\tau^2$ (blocks) for $|\tau|>s_1$, and absent for
$\theta$.  With (c7) and $\rho^2(1+\eta_0/2)=1-2\delta$, for $0<|\tau|\le
c_{\rm flat}t_i$:
\[
 p^*(f'+\tau g'_i)\le1+\hat\Gamma+E\le1+\frac{\tau^2}2\Big(1-2\delta+\frac\delta8
 +\frac\delta8+\frac\delta8+\frac\delta{16}+\frac\delta4\Big)\le1+\frac{\tau^2}2
 (1-\delta).
\]
The remaining assertions are Lemma~\ref{lem:III-UE1}(a), (c), (g) with
$\Pi\le K_ws_1$ and (FINE).
\end{proof}

\begin{remark}[What changed relative to {[I,~Theorem~7.18]}]\label{rem:III-UE}
(a) The first-order junction mismatch is zero (exact $d$-consistency) instead of
$K(\omega)s_1$, so $K_\sharp$ is an $f$-constant.  (b) Every ``late stage''
statement is replaced by an explicit bound in the perturbation size $\Pi\le
K_ws_1$.  (c) The scrambling hypothesis is used in the explicit form (U5).
(d) Violations are tolerated as in [II,~Lemma~2.30], uniformly in the piece:
their cost $2\rho|\tau|\epsilon$ does not depend on $t_i$.  (e) Lever coordinates
are protected: zero data, contact-like data, or pull masses dominating
$c_{\rm flat}t_i$ times the data.  (f) The only interplay between $s_1$ and the
window is the interval $[64\rho\epsilon/\delta,s_{\rm late}]$; (VT$'$) bounds
$s_1$ from below by the violation mass, (LATE) from above by window quantities.
(g) The radius $c_{\rm flat}t_i$ is proportional to the scale; $c_{\rm flat}$
depends on the window only through $\gamma_B$.
\end{remark}

\subsection{Averaging at the norm-attaining approximant}\label{subsec:III-Eeng}

\begin{theorem}[Windowed recovery through $d$-consistent approximants]
\label{thm:III-Eeng}
Let $I$ be finite, $f\in S_{p^*}$ with $F$ finite, $g\in\Cset(f)$, $\rho\in(0,1)$,
$\eta_0\in(0,2]$ with $\rho^2(1+\eta_0)\le(1+\rho^2)/2$, and $\delta$ as above.
Let $f_j\to f$ be base rows, $\varepsilon_j:=p^*(f_j-f)$.  Suppose that for every
$j$ there are $T_j\in(0,t_1]$, $n_j\ge1$, a nonempty set $\mathcal S_j\subseteq
\{T_j2^{1-i}:1\le i\le n_j\}$ with $n'_j:=|\mathcal S_j|$, a number
$T'_j\in(0,\min\mathcal S_j]$, a number $\gamma_{B,j}\in(0,1]$, $K_j\ge1$, and a
window family at $f_j$ with scales $\mathcal S_j\subseteq[T'_j,t_1]$ \textup{(}window parameter $T:=T'_j$\textup{)} satisfying
\textup{(U1)--(U7)} with \emph{fixed} $A_0,A_2$ and with $\gamma_{B,j}$; let
$c_{\rm flat}(j):=\min(c^f_{\rm flat},\gamma_{B,j}/(8\rho A_2))$ and $K_{w,j}$
be the constants of Theorem~\textup{\ref{thm:III-UE}} and
Lemma~\textup{\ref{lem:III-UE1}} for this family.  Assume:
\begin{enumerate}
\item[(E-c)] $p^*(g-g_{j,t})\le K_jt$ for $t\in\mathcal S_j$;
\item[(E-d)] $K_jT_j\to0$ and $n'_jc_{\rm flat}(j)/K_j\to\infty$;
\item[(E-f)] the window conditions \textup{(W-a)--(W-c)} of
Theorem~\textup{\ref{thm:III-UE}} hold at $f_j$, and there is $s_{1,j}$ satisfying
\textup{(VT$'$)}, \textup{(FINE)} and \textup{(LATE)} at $f_j$, such that, with
$\Pi_j\le K_{w,j}s_{1,j}$ the resulting perturbation size,
\[
 K_{w,j}\big(K_{w,j}s_{1,j}+c_{{\rm fine},j}+\mathrm{tail}_{W,j}\big)\le K_j(T'_j)^2;
\]
\item[(E-e)] $\varepsilon_j+K_{w,j}\big(\Pi_j\log(e/\Pi_j)+s_{1,j}\big)\le
\theta_j(T'_j)^2$ with $\theta_j/c_{\rm flat}(j)^2\to0$.
\end{enumerate}
Then $(f,\rho g)\in\cl\NA((c_0,p),\ell_2^2)$.  If the hypotheses hold for every
$\rho<1$, then $(f,g)\in\cl\NA((c_0,p),\ell_2^2)$.
\end{theorem}

\begin{proof}
Fix $j$ large; write $n:=n'_j$, $K:=K_j$, $c:=c_{\rm flat}(j)$,
$\mathcal S:=\mathcal S_j$, $T':=T'_j$.  For large $j$, $p^*(f_j-f)\le r_f/2$.
Theorem~\ref{thm:III-UE} at $f_0:=f_j$ with $s_1:=s_{1,j}$ and $\mathrm{tail}_W$
as in (E-f) gives a norm-attaining $f'_j$ and functionals $g'_{j,t}$
($t\in\mathcal S$) with
\begin{enumerate}
\item[(A)] $p^*(f'_j+\tau g'_{j,t})\le1+\frac{\tau^2}2(1-\delta)$ for $|\tau|\le ct$;
\end{enumerate}
$\varepsilon'_j:=p^*(f'_j-f)\le\theta_jT'^2$ by (E-e); and $p^*(g'_{j,t}-\rho
g_{j,t})\le K_{w,j}(K_{w,j}s_1+c_{\rm fine}+\mathrm{tail}_W)/t\le KT'^2/t\le Kt$ by
(E-f), so $p^*(g'_{j,t}-\rho g)\le(1+\rho)Kt\le2Kt$.  Since $g\in\Cset(f)$,
[I,~Proposition~1.13(c)] gives
\begin{enumerate}
\item[(B)] $p^*(f'_j+rg'_{j,t})\le p^*(f+r\rho g)+\varepsilon'_j+|r|p^*(g'_{j,t}-\rho g)
\le s(\rho r)+\varepsilon'_j+2K|r|t$ for all real $r$.
\end{enumerate}
Put $\bar g_j:=\frac1n\sum_{t\in\mathcal S}g'_{j,t}$; by convexity
$p^*(f'_j+r\bar g_j)\le\frac1n\sum_tp^*(f'_j+rg'_{j,t})$.  Fix $r_0\in(0,1]$ with
$r_0^2\le\min(\delta,1-\rho^2)$.

\emph{Case $|r|\le r_0$.}  Let $\mathcal S_r:=\{t\in\mathcal S:ct<|r|\}$; as the
scales are distinct dyadic numbers, $\sum_{t\in\mathcal S_r}t<2|r|/c$.  If
$\mathcal S_r=\emptyset$, all terms obey (A), and $1+\frac{r^2}2(1-\delta)\le1+
\frac{r^2}2-\frac{r^4}8\le s(r)$ because $r^2\le\delta$ and $\sqrt{1+x}\ge1+
\frac x2-\frac{x^2}8$.  Otherwise $|r|>c\min\mathcal S\ge cT'$, so $\varepsilon'_j\le\theta_jr^2/c^2$
and
\[
 p^*(f'_j+r\bar g_j)\le\max\Big\{1+\frac{r^2}2(1-\delta),\,s(\rho r)\Big\}
 +\frac{\theta_j}{c^2}r^2+\frac{2K|r|}n\cdot\frac{2|r|}c .
\]
By (E-d), (E-e), for $j$ large $\theta_j/c^2+4K/(nc)\le\min(\delta/4,(1-\rho^2)/6)$,
and then the right side is $\le s(r)$, because $1+\frac{r^2}2(1-\delta)+
\frac{\delta r^2}4\le1+\frac{r^2}2-\frac{r^4}8$ (as $r^2\le\delta$) and
$s(\rho r)+(1-\rho^2)r^2/6\le s(r)$ by [I,~Lemma~3.6(a)] ($|r|\le1$).

\emph{Case $|r|\ge r_0$.}  By (B) for all $t$, $p^*(f'_j+r\bar g_j)\le s(\rho r)+
\varepsilon'_j+2K|r|\frac1n\sum_{t\in\mathcal S}t\le s(\rho r)+\varepsilon'_j+
4K|r|T_j/n$.  For $j$ large, $\varepsilon'_j\le(1-\rho^2)r_0^2/6$ and
$4KT_j/n\le4KT_j\le(1-\rho^2)r_0/6$; since $\min(r^2,|r|)\ge r_0|r|$,
[I,~Lemma~3.6(a)] gives $p^*(f'_j+r\bar g_j)\le s(r)$.

Hence $\bar g_j\in\Cset(f'_j)$ and, $f'_j$ being norm attaining,
$(f'_j,\bar g_j)\in\NA$ by [I,~Proposition~1.13(d)].  Finally
$\|(f'_j,\bar g_j)-(f,\rho g)\|\le\varepsilon'_j+p^*(\bar g_j-\rho g)\le
\varepsilon'_j+\frac1n\sum_t2Kt\le\varepsilon'_j+4KT_j/n\to0$.  The last sentence
follows since $(f,g)=\lim_{\rho\to1}(f,\rho g)$.
\end{proof}

\begin{remark}\label{rem:III-Eeng}
(a) Unlike [II,~Theorem~2.20] no exactness of the averaged data and no final
application of [I,~Corollary~7.21] are needed: the averaging is carried out at
the norm-attaining approximant itself, with the per-piece bounds (A).  This is
what makes violated data usable along a sequence of windows.  (b) The set of
scales enters only through $\sum_{\mathcal S_r}t<2|r|/c$, $\min\mathcal S\ge T'$
and $\frac1n\sum_{\mathcal S}t\le2T_j/n$; so any subset of a dyadic window may
be used, as in [II,~Lemma~5.14].  (c) $c_{\rm flat}(j)$ enters only through the
range of (A), the dyadic count and the threshold $|r|>c_{\rm flat}(j)T'_j$, all
tracked by (E-d), (E-e); a window-dependent $\gamma_{B,j}$ is therefore allowed.
(Numerically, in a model of inputs (A), (B) averaged over random half-subsets of a
dyadic window, the averaged bound depends on $c_{\rm flat}$ only through
$n'c_{\rm flat}/K$ and $\theta/c_{\rm flat}^2$; evidence only.)
\end{remark}

\subsection{The threshold comparison}\label{subsec:III-S2c}

The interval $[64\rho\epsilon/\delta,s_{\rm late}]$ of Theorem~\ref{thm:III-UE}
is nonempty only if the violation mass is small compared with the lateness
threshold.  The violation masses produced by fine carriers are of the size of the
first weight after the stage, and $s_{\rm late}$ is a power of window quantities;
the comparison is arranged by a design condition.

\begin{definition}[Design conditions (D-lev) and (W$_{\exp}$)]\label{def:III-designadd}
For the design $D^{U1'}$ of [II,~Definition~5.48] on $T_{\rm final}$:
\begin{enumerate}
\item[(D-lev)] for every stage $L$, with $\sigma_2(L):=\max_{l\le L}$ (second
element of $S_l\cap(s_{\max}(L),\infty)$),
\[
 \mathrm{Des}(L)\ge\max_{l\le L}\frac{n_l\,2^{\sigma_2(L)}}{(\mu^\star_{\sigma_2(L)})^2
 \delta_l} ;
\]
\item[(W$_{\exp}$)] for every \emph{main} stage $L$, the lower weight bound of stage
$L+1$ is chosen with $c^{\rm low}_{L+1}\le\exp(-1/T_{\rm lo}(L,M(L)))$, so that
$c_{L+1}\le\exp(-1/T_{\rm lo}(L,M(L)))$.
\end{enumerate}
\end{definition}

\begin{lemma}[The additions are harmless]\label{lem:III-designadd}
Both conditions are computable at the stage where they are imposed and do not
depend on $N$ or on $f$.  Imposing them keeps the recursion well founded, keeps
\textup{(T-a)--(T-d)}, \textup{(P1)}, \textup{(P2)} and $N$-independence, and every
result of Part~II valid for $D^{U1'}$ remains valid.  Moreover, under (D-lev)
$(\mu^\star_{j})^2v_l(j)\ge1/\mathrm{Des}(L)$ for every $l\le L$ and every
$j\in S_l\cap(s_{\max}(L),\sigma_2(L)]$.
\end{lemma}

\begin{proof}
(D-lev) is a lower bound on $\mathrm{Des}(L)$ by quantities fixed at stage $L$
(signature sets, base entries, $\delta_l$, $n_l$); (W$_{\exp}$) is an upper bound
on one weight per main stage, computed after all sub-windows of level $L$ are
fixed.  Imposed at main stages only, it does not conflict with the cluster weights
$c_{p+i}=c^{\rm low}_{p+1}4^{1-i}$ of $D^{U1'}$ (imposed at every stage it would,
since $\exp(-1/T_{\rm lo})\ll c^{\rm low}_{p+1}/4$).  The lower-bound uses
of weights in Part~II (the absorber length $R$, the least absorber entry
and the cluster capacity, [II,~Definition~5.48, Lemma~5.51]) are computed
from $c^{\rm low}$ after (W$_{\exp}$) is imposed, so they hold for the
modified weights; the absorber targets are produced by the same rule from
the new $c^{\rm low}$.  Enlarging $\mathrm{Des}$ enlarges $\bar D(w)$, and the
condition $Q(w)\ge\bar D(w)^3$ of [II,~Definition~5.48], on which the
class and cube pigeonhole rests, is imposed after $\mathrm{Des}(L)$ is fixed.
Every other window condition is monotone under enlarging $\mathrm{Des}$ and
shrinking weights ([II,~Lemma~3.22]); the other uses of weights in Part~II
are upper bounds.  For the last
claim: $v_l(j)=\delta_l2^{-j}/n_l$ on $S_l$, so $(\mu^\star_j)^2v_l(j)$ decreases
in $j$ and is $\ge(\mu^\star_{\sigma_2(L)})^2\delta_l2^{-\sigma_2(L)}/n_l\ge
1/\mathrm{Des}(L)$ for $j\le\sigma_2(L)$.
\end{proof}

\begin{proposition}[Threshold comparison]\label{prop:III-S2c}
Let $C_0\ge1$.  There is $T_*>0$, depending only on $C_0$ and
$f$-constants, such that the following holds for every window family at a base
row with scales in $[T,t_1]$ and $T\le T_*$.  Suppose
\begin{enumerate}
\item[(i)] $K_w$, $n$, $A_0/T$, $1/\gap_{\min}$, $1/x_0$ and $\Lambda_{\rm lev}$ are
at most $\exp(C_0\log^2(1/T))$, and \textup{(W-b)} holds;
\item[(ii)] $\epsilon$, $c_{\rm fine}$, $W_{\rm lev}+\overline W_{\rm pull}$ and
$r_{\rm pull}$ are at most $\exp(C_0\log^2(1/T))\,c_*$ for a number
$c_*\le\exp(-1/T)$;
\item[(iii)] $N_w$ is chosen with $\mathrm{tail}_W\le\exp(-1/T)$.
\end{enumerate}
Then $s_1:=\max(64\rho\epsilon/\delta,\exp(-1/(2T)))=\exp(-1/(2T))$ satisfies
\textup{(VT$'$)}, \textup{(FINE)}, \textup{(PULL)} and \textup{(LATE)};
\textup{(W-a)} and \textup{(W-c)} hold; and
\[
 K_w\big(K_ws_1+c_{\rm fine}+\mathrm{tail}_W\big)+K_w\big(\Pi\log(e/\Pi)+s_1\big)\le
 \exp\big(-1/(3T)\big).
\]
\end{proposition}

\begin{proof}
By (i), $s_{\rm late}=c_{\rm late}(T\gap_{\min}x_0/(nA_0K_w))^7\ge c_{\rm late}
\exp(-35C_0\log^2(1/T))$, which exceeds $\exp(-1/(2T))$ for $T$ small; and
\[
 64\rho\epsilon/\delta\le64\rho e^{C_0\log^2(1/T)}e^{-1/T}/\delta\le e^{-1/(2T)}
\]
for $T$ small.  So $s_1=e^{-1/(2T)}\in[64\rho\epsilon/\delta,s_{\rm late}]$.
(FINE) holds because
\[
 e^{C_0\log^2(1/T)}e^{-1/T}\le c_{\rm late}\delta A_{\min}T^4e^{-1/(2T)}/(1+C_{\rm lev})
\]
for $T$ small, as $1/A_{\min},C_{\rm lev}\le e^{C_0\log^2(1/T)}$; (PULL)
holds because $r_{\rm pull}\le e^{C_0\log^2(1/T)}e^{-1/T}\le32K_JT^4
e^{-1/(2T)}$ for $T$ small (note $K_J\ge2$).  (W-a) holds
likewise, since for $T$ small
\[
 \frac{c_{\rm late}\delta T}{nA_0(1+C_2K_J)(1+1/A_{\min})}\ge
 e^{-5C_0\log^2(1/T)}c_{\rm late}\delta T ;
\]
and (W-c) holds as $e^{-1/T}\le\min(T^3,c_{\rm late}\delta T/A_0)$.  Finally
$\Pi\le K_ws_1$, and every term on the left of the display of the proposition is
at most
\[
 3e^{3C_0\log^2(1/T)}e^{-1/(2T)}\log(e/s_1)\le e^{-1/(3T)}\qquad(T\text{ small}).
\]
\end{proof}

\begin{remark}[Where the hypotheses come from]\label{rem:III-S2c}
(a) By Lemma~\ref{lem:III-UE1}, $\log K_w\le C_{\rm abs}\log(C_fnA_0|\Omega|
\Lambda_{\rm lev}C_S/(T\gap_{\min}x_0))$; when $n$, $\mathrm{Des}$ and $|\Omega|$
are at most $\log_2(1/T)$ and $1/\gap_{\min},1/x_0\le T^{-C\,\mathrm{Des}}$, this is
$\le C\log^2(1/T)$: the design quantity $\mathrm{Des}$ enters the base of a power,
never the exponent of $K_w$.  (b) At companions of a clean sub-window $w$ of a
main stage $L$, the violation mass, $c_{\rm fine}$, $c_{\rm tiny}$, $W_{\rm lev}$
and $\overline W_{\rm pull}$ produced by carriers at stages $>L$ are at most
$C_f\mathrm{Des}(L)c_{L+1}$, and (W$_{\exp}$) with $T_{\rm lo}(L,M(L))\le
T_{\rm lo}(w)$ gives $c_{L+1}\le e^{-1/T}$; this is hypothesis (ii) with
$c_*:=c_{L+1}$.  (c) The weaker condition $c_{L+1}\le b(L,M(L))^2\le
T_{\rm lo}(w)^8$ of $T_{\rm final}$ does not suffice for the lateness threshold
obtained here, which contains $\gap_{\min}^7$; whether a sharper threshold would
make it suffice is not investigated.
\end{remark}

\subsection{Fine carriers enter only through their weight}\label{subsec:III-noB}

Hypothesis (U5) of Theorem~\ref{thm:III-UE} controls the scrambling of
\emph{all} carriers of the negative blocks.  It is used only in
Lemma~\ref{lem:III-UE1}(e).  When the fine carriers have total weight small
compared with $s_1$, it can be replaced by a condition on coarse carriers only.

\begin{lemma}[Coarse scrambling suffices]\label{lem:III-noB}
Replace \textup{(U5)} in Definition~\textup{\ref{def:III-window}} by
\begin{enumerate}
\item[(U5$'$)] every peak of $f_0$ that is not fine has margin $\ge x_0$, and every
$k\in\Omega_m$ has $\Phi_m(k)\gap^0(k)\ge x_0$ and $\gap^0(k)\ge x_0$,
\end{enumerate}
\textup{(}no condition on fine carriers and on the set $I_-$\textup{)}.  Then
Lemma~\textup{\ref{lem:III-UE1}} holds with \textup{(e)} replaced by
\[
 \mathfrak S_m\le K_w\Pi^2\log(e/\Pi)+4ms_1^2T^4+4c_{\rm fine}+4(W_{\rm lev}+
 W_{\rm pull})\qquad(m\in I),
\]
and Theorem~\textup{\ref{thm:III-UE}} holds with \textup{(U5)} replaced by
\textup{(U5$'$)} if in addition $4c_{\rm fine}\le c_{\rm late}\delta s_1$.
\end{lemma}

\begin{proof}
In the proof of Lemma~\ref{lem:III-UE1}(e), a coarse carrier with $t_k\le\Pi$ lies
in $S_1$ or $S_2$ of [I,~Lemma~7.15] once $2K_*\Pi\le x_0$ (margins $\ge x_0$;
$\Phi\gap$, $\gap\ge x_0$; for $\Omega$ also by $d$-consistency).  So
$A^{(1)}$ contains only fine carriers and carriers meeting lever coordinates, and
their contribution to $\mathfrak S_m$ is at most $\sum(\Phi_m(k)^2+\lambda_k(2+1))
\le4\sum\lambda_k$, as $\Phi_m(k)\le\lambda_k$.  $A^{(2)}$ and
$\|D_m(w'_m-w^0_m)\|_2^2$ are bounded as before (the clamp-formula bound of
[I,~Lemma~7.14] holds for every carrier, with value moves $\le2\Pi+t_k$ off the
lever coordinates and $|w'-w^0|\le2$ on them).  In the proof of
Theorem~\ref{thm:III-UE}, (U5) enters only through \eqref{eq:III-BxS}, which
holds again when $4c_{\rm fine}\le c_{\rm late}\delta s_1$.
\end{proof}

Under the hypotheses of Proposition~\ref{prop:III-S2c}, $4c_{\rm fine}\le
4e^{C_0\log^2(1/T)}c_*\le c_{\rm late}\delta e^{-1/(2T)}$ for $T$ small; so in that regime
Theorem~\ref{thm:III-UE} needs no information on the statuses of fine carriers.
This is not a statement that fine carriers are harmless in general: only their
total weight is used, against $s_1\ge e^{-1/(2T)}$, and this is where
(W$_{\exp}$) is essential.

\subsection{Activity classes in the violation-tolerant route}\label{subsec:III-mixed}

We now apply the results above at the companions of [II,~Subsection~5.6], for
every activity class, one-signed or mixed.  Throughout this subsection the
design is $D^{U1'}$ of [II,~Definition~5.48] on $T_{\rm final}$ with (D-lev) and
(W$_{\exp}$) (Definition~\ref{def:III-designadd}), $N\ge1$ is fixed, $f\in
S_{p_N^*}$ has finite base support, $w$ is a clean sub-window of a main stage
$L\ge l_f$, and $a$ is an activity class ([II,~Lemma~5.53]) satisfying
(KN$_{w,a}$) ([II,~Definition~5.55]; for a class without active shift this is
vacuous).  We use the notation of [II,~Subsection~5.6]: $A$ the active blocks,
$A_\mp:=\{m\in A:\sigma_m=\mp1\}$, $\Omega_{\rm act}$, $E_c(w)$, $J_{\rm fine}$,
$\mathcal V(\Delta\omega)$, $\mathcal L(\Delta',\gamma)$, the true shift $\Delta''$,
and $\mathcal S$ a set of scales of class $a$ in one cube of
[II,~Proposition~5.54(c)].  The new point compared with the proof of [II,~Theorem~5.68] is
that \emph{no absorber carries a switching coefficient}: residues on the coarse
coordinates and trace corrections are not cancelled but tolerated as violations.

\begin{lemma}[Completion without absorber switching]\label{lem:III-CMp}
Let $f^\#$ be obtained from $f$ by the moves of [II,~Definition~5.57] with the
following changes: \textup{(1a)} as there \textup{(}[II,~Proposition~5.56], or
[II,~Theorem~5.7] for a class without active shift\textup{)}; steps
\textup{(1b)}, \textup{(1c)} are omitted \textup{(}no absorber is tuned, and
after \textup{(1a)} no value of $\Omega_{\rm act}$ and no $\varkappa_m$ has to be
restored\textup{)}; in step \textup{(2)} the recursion of [II,~Lemma~5.60] is run
with owner candidates of types \textup{(a)}, \textup{(b)} and \textup{(d)} only,
where \textup{(d)} comprises every carrier at a stage $>L$ of an active block,
absorbers of $D^{U1'}$ included, with its forced coefficient $-\Delta''_m
\lambda_kw^\#_m(k)$; a type-\textup{(d)} carrier of a block of $A_-$ is
self-aligned and one of a block of $A_+$ is anti-aligned \textup{(}$z_j:=
\sgn(-\sigma_m\varsigma_ku_k(j))$ on $O(k)$, $\varsigma_k:=\sgn Y_k$\textup{)},
and unowned coordinates get $z_j:=0$.  For $t\in\mathcal S$ let $\Delta\omega(t)$
have the components of the data in the proof of [II,~Theorem~5.68] on $\Omega_{\rm act}$, $0$ on the
inactive $\Omega$-carriers, and \emph{no} absorber components; let
$\omega^+$ be the clamped decomposition used there, $\omega^-:=\omega^++
\Delta\omega(t)$, $\mathcal V:=\mathcal V(\Delta\omega(t))$, $\mathcal V_{\rm proj}:=
\mathcal L(\Delta''^*,\gamma(t))$ on $E_c(w)\setminus F^\#$ and $0$ elsewhere,
$\mathcal V_{\rm fine}:=\mathcal V-\mathcal V_{\rm proj}$ off $F^\#$ and $0$ on
$F^\#$, and
\[
 b^+:=\beta+\chi\mathcal V_{\rm proj}+\tfrac12\mathcal V_{\rm fine},\qquad
 b^-:=b^+-\mathcal V,
\]
with $\chi$ and $\beta$ \textup{(}supported in $F^\#$, $b^+(\xi^\#)=0$\textup{)} as in
the proof of [II,~Theorem~5.68].  Then for every $t\in\mathcal S$:
\begin{enumerate}
\item[(a)] every carrier at a stage $>L$ of a block of $A_-$ is a peak of $f^\#$ with
$|u_k(\zh^\#)|=|Y_k|+\mathrm{own}_k\ge\delta^\circ_k$ and margin $\ge q^\#_0
\delta^\circ_k/2$; hence $\mathrm{Scr}^\#_m(x)\le C(x/q^\#_0)^2$ for $m\in A_-$ and
$0<x\le x_0:=c_f\gap_{\min}\Phi_{\min}(L)$;
\item[(b)] $(b^\pm,\omega^\pm)$ represent a functional $g_t$ at $f^\#$ with
$p^*(g-g_t)\le Kt$, $K=C_f(\mathrm{Des}(L)/u(w))^C$, $\Gamma^\#_w(b^\pm,\omega^\pm)
\le1+\eta_0/2$, and the block parts are of kinds \textup{[1$'$]--[3]};
\item[(c)] the total violation satisfies $\epsilon(t)\le C_f\mathrm{Des}(L)c_{L+1}+
C_fc_{L+1}^2/T_{\rm lo}(w)$, and $b^\theta=0$ at every violated free coordinate;
\item[(d)] $p^*(f^\#-f)\le\theta_wT_{\rm lo}(w)^2$ with $\theta_w\to0$, and every
coarse carrier has the status and the gap of [II,~Proposition~5.56(b)].
\end{enumerate}
\end{lemma}

\begin{proof}
(a) At a type-(d) carrier $k$ of a block of $A_-$ the rule gives $z_j=
\varsigma_k\sgn u_k(j)$ on $O(k)\supseteq S_k$, and $\zh^\#_j=z_j$ on $J_{\rm fine}$
(diagonal base); so $u_k(\zh^\#)=Y_k+\varsigma_k\mathrm{own}_k$ and $|u_k(\zh^\#)|=
|Y_k|+\mathrm{own}_k\ge\delta^\circ_k$.  The peak criterion and the margin follow
from (W3) and [I,~Lemma~1.7] as in [II,~Lemma~5.60]; this computation does not
involve the other blocks.  Absorbers enter as carriers at stages $>L$; at
absorber stages the weight conditions of $D^{U1'}$ give $c_p\le(\delta^{\max}_p
H_p/2)^2$, so they are covered.  The scrambling bound is the proof of
[II,~Lemma~5.62]: coarse carriers do not contribute below $x_0$ (statuses and gaps
of [II,~Proposition~5.56(b)]; no absorber is a coarse carrier now), every other
carrier of a block of $A_-$ is a peak with margin $\ge q^\#_0\delta^\circ_k/2$,
$\Phi_k\le(\delta_kH_k)^2$ by (W3), and the super-exponential decay of
$\delta_kH_k$ gives $\sum_{\delta_kH_k\le X}(\delta_kH_k)^2\le2X^2$.

(b) $b^--b^+=-\mathcal V=-\sum_mR_m^*(\Delta\omega_m-\Delta''_mw^\#_m)$ and
$\omega^--\omega^+=\Delta\omega$, so both pairs represent the same functional
([II,~Lemma~2.22]).  Compared with the data of the proof of [II,~Theorem~5.68], the absorber
coefficients are removed and $\mathcal V_{\rm fine}$ is re-split; this changes
$g_t$ by at most $\|\mathcal V_{\rm fine}\|_1+\sum_a|\gamma_a|\le C_f\mathrm{Des}\,
c_{L+1}+C_fc_{L+1}^2/t\le Kt$ (the bound of $\|\mathcal V_{\rm fine}\|_1$ is
the one of (c) below).  The data change by a vector $y$ with $q_0h(y)$,
$\sigma_mH_m(y)=O(c_{L+1}^2)$, because $q_0h(y)\le\|U\|^2\|y\|_1^2/\nu$ and
$\sigma_mH_m(x)\le(\sum_k\lambda_k|x(k)|)^2/C_m$; since $\Gamma_w$ is quadratic,
the cross terms change $\Gamma_w$ by $O(c_{L+1})$, and $\sqrt{\Gamma_w}$ (a
seminorm) changes by $O(c_{L+1})$.  This is absorbed in the slack of the
proof of [II,~Theorem~5.68], which bounds $\sqrt{\Gamma^\#_w}$ by
$\sqrt{\Gamma_w(B_-,\Theta_-)}+Kt+C_fp\varepsilon$, since $c_{L+1}\le Kt$.  The remaining
claims are those of the proof of [II,~Theorem~5.68] (kinds: increments satisfy the
inward rows (X4), so near-threshold carriers are of kind [3]).

(c) At a contact $j\in E_c(w)\setminus F^\#$: $z^\#_j\mathcal V_{\rm proj}(j)\ge0$
(admissibility of $\mathcal L(\Delta''^*,\gamma)$, row (X1), [II,~Lemma~5.58(b)])
and $\chi\in[0,1]$ give $(z_jb^+_j)_-\le|\mathcal V_{\rm fine}(j)|/2$ and
$(z_jb^-_j)_+\le|\mathcal V_{\rm fine}(j)|/2$.  At a free coordinate of
$E_c(w)$, $\mathcal V_{\rm proj}(j)=0$ (row (X1)), so $b^\pm_j=\pm\frac12
\mathcal V_{\rm fine}(j)$, $b^\theta_j=0$ and $\mathrm{viol}(j)=|\mathcal V_{\rm
fine}(j)|$.  On $J_{\rm fine}$, $\mathcal V_{\rm proj}=0$ and $b^\pm=\pm\frac12
\mathcal V$: on coordinates owned by owners of types (a), (b) or by negative
type-(d) carriers the data are admissible by the dominance of
[II,~Lemma~5.61] (whose estimates do not use signs); unowned coordinates carry
$\mathcal V(j)=0$; on coordinates owned by positive type-(d) carriers
$\mathrm{viol}(j)\le|\mathcal V(j)|$.  On $E_c(w)$, $\mathcal V_{\rm fine}=
[\mathcal L(\Delta'',\gamma)-\mathcal L(\Delta''^*,\gamma)]+(\text{carriers at
stages}>L)$; the first part is the trace correction of [II,~Lemma~5.58(b)], of
$\ell_1$-norm $\le C_f\mathrm{Des}\,c_{L+1}+C_fc_{L+1}^2/t$, the second has
$\ell_1$-norm $\le C_f\sum_{l>L}\lambda_l$ (coefficients $\le C_f\lambda$,
$\|u\|_1\le1$); the same bound covers the coordinates owned by positive fine
carriers.  The residues at support coordinates are antisymmetric and are counted
in $\epsilon$ in the same way.

(d) [II,~Lemma~5.64]: moves (1b), (1c) are absent, and step (2) is as there.
\end{proof}

\begin{lemma}[Levers at the companion]\label{lem:III-LV}
In the setting of Lemma~\textup{\ref{lem:III-CMp}} let $\Omega$ be the set of
coarse strict non-peaks of $f^\#$, and for $l\in\Omega$ let $j(l)$ be the least
element of $S_l\cap(s_{\max}(L),\infty)\setminus F^\#$.  Then:
\begin{enumerate}
\item[(a)] every $l\in\Omega_{\rm act}$ has a TU lever with $\epsilon_l$ the sign
of its switching coefficient \textup{(}constant on $\mathcal S$\textup{)}, bank
coordinate $j':=j(l)$ and set of pull coordinates $\mathrm{Pull}_l:=S_l\cap
(j(l),s_{\rm far}(w)]\setminus F^\#$, all contacts of $f^\#$ with
$z^\#=\epsilon_l$; at $j'$ the data are contact-like of
sign $\epsilon_l$ up to the antisymmetric residue $\pm\frac12\mathcal V_{\rm
fine}(j')$, at a pull coordinate $t|b^\pm(j_p)|\le(2A_2+C_\Delta)\lambda_lv_l(j_p)$
up to such a residue, and, with $G:=6G(L)$, $G(L):=2^{L+1}$, and
$r_{\rm pull}:=2^{2G}\max_lv_l(s_{\rm far}(w))\le\mathrm{Des}(L)^2c_{L+1}^2$, for
every $r\in[r_{\rm pull},r_0]$ there is $j_p\in\mathrm{Pull}_l$ with
$v_l(j_p)\in[r,2^Gr]$;
\item[(b)] every $l\in\Omega$ with $\gamma_l(t)=0$ for $t\in\mathcal S$ \textup{(}the
inactive carriers, including those of class R and those held at value $0$\textup{)}
has a zero-data lever \textup{(Z)} at $j(l)$, where $b^\theta(j(l))=0$ and
$b^\pm(j(l))=\pm\frac12\mathcal V_{\rm fine}(j(l))$;
\item[(c)] the levers satisfy \textup{(SEP)}; the lever constants satisfy
\[
 \Lambda_{\rm lev}\le C_f\big(\mathrm{Des}(L)/(u(w)\gap_{\min})\big)^C;
\]
and $c_{\rm tiny}$,
$W_{\rm lev}$ and $\overline W_{\rm pull}$ are at most $C_f\mathrm{Des}(L)^Cc_{L+1}$.
\end{enumerate}
\end{lemma}

\begin{proof}
Beyond $s_{\max}(L)$ the only carrier $\le L$ whose vector is nonzero on $S_l$ is
$l$, with $u_l=v_l=\delta_l2^{-s}/n_l$ there (allowedness, [II,~Theorem~3.12(d)]);
this gives (SEP)(ii), and (SEP)(i) holds because the sets $S_l$ are pairwise
disjoint and $j(l)\notin F^\#$.  Inside $(s_{\max}(L),\sigma_2(L)]$ the support of
$f^\#$ meets $S_l$ at most in its first element (a bank of the assembly or of
(1a)), so $j(l)\le\sigma_2(L)$ and, by (D-lev) and Lemma~\ref{lem:III-designadd},
$\beta_{j(l)}=\nu/((\mu^\star_{j(l)})^2v_l(j(l)))\le\nu\,\mathrm{Des}(L)$.  On
$E_c(w)\cap S_l\cap(s_{\max}(L),\infty)$ the projected vector is $\mathcal
V_{\rm proj}=\gamma_lv_l$, so the data are $(\chi,\chi-1)\gamma_lv_l\pm\frac12
\mathcal V_{\rm fine}$.  (a) For $l\in\Omega_{\rm act}$, the sign of $\gamma_l$ is
fixed on the class ([II,~Proposition~5.54(e)]) and the coordinates of $S_l\cap
(s_{\max}(L),s_{\rm far}(w)]\setminus F^\#$ are contacts with $z^\#=\epsilon_l$
(closed by (C1) of [II,~Definition~4.16], unchanged by (1a) and by the release);
so the data at $j'$ are contact-like of sign $\epsilon_l$.  As $|\Delta\omega(l)|
\le2A_2/t$ and $|\Delta''|\le C_\Delta$, $|\gamma_l|\le\lambda_l(2A_2/t+C_\Delta)$,
which gives the bound at the pulls.  Consecutive elements of $S_l$ differ by
$2^{l+1}\le2^{L+1}$, so $v_l$ drops by at most $2^{G(L)}$ between them.  Hence
for $r\in[r_{\rm pull},r_0]$ (with $r_0$ decreased, if necessary, so that
$2^Gr_0\le\min_lv_l(j(l))$; this keeps the bound on $\Lambda_{\rm lev}$, as
$1/v_l(j(l))\le\mathrm{Des}(L)$) the interval $[r,2^Gr]$ lies in
$[v_l(s_{\rm far}(w)),v_l(j(l))]$ and contains at least six elements of
$S_l\cap(j(l),s_{\rm far}(w)]$, while the moves of (1a) (the tuning pull and
bank of [II,~Lemma~4.15], a hold or lever pull, and the assembly bank) make
at most five elements of each $S_l$ support coordinates of $f^\#$; so
$\mathrm{Pull}_l$ meets $[r,2^Gr]$.  Finally $2^{-s_{\rm far}(w)}\le c_{L+1}^2$
and $2^{2G}=2^{12G(L)}\le\mathrm{Des}(L)^2$ ([II,~Lemma~4.1(d)]) give
$r_{\rm pull}\le\mathrm{Des}(L)^2c_{L+1}^2$.  (b) Here
$\mathcal V_{\rm proj}(j(l))=\gamma_lv_l(j(l))=0$, so $b^\theta(j(l))=(\chi-\frac12)
\mathcal V_{\rm proj}(j(l))=0$.  (c) Lever coordinates of $l$ are met only by $l$
among carriers $\le L$; carriers at later stages meeting them have total weight
$\le C\sum_{k>L}\lambda_k\le Cc_{L+1}$ and meet them with $|u_k(s)|\le1$.  The
sum defining $c_{\rm tiny}$ runs over the (Z) and TU bank coordinates only,
which are the coordinates $j(l)$, and $1/v_l(j(l))\le\mathrm{Des}(L)$; the
pull coordinates enter only through $\overline W_{\rm pull}$, which is the
weight of carriers at stages $>s_{\max}(L)>L$ meeting $\mathrm{Pull}_l$.  The constants of Lemma~\ref{lem:III-DC} are
bounded by $C_f\mathrm{Des}(L)^C$ times powers of $1/u(w)$ and of $1/\gap_{\min}$,
which enter through $r_0$: the status radius of Lemma~\ref{lem:III-DC} is measured in
value units, $r_0\le c\min_{l\in\Omega}\Phi_m(l)\gap_m(l)A_{\min}/(mC_{\rm lev})$,
and $r_0\le c$ times the coarse peak margins $\ge(u/2)\bar\vartheta\Phi$ divided by
$C_{\rm lev}$.
\end{proof}

\begin{theorem}[Activity classes are recovered under (KN)]\label{thm:III-Cmix}
Let the design be $D^{U1'}$ with \textup{(D-lev)} and \textup{(W$_{\exp}$)}, $N\ge1$,
and $f\in S_{p_N^*}$ with finite base support.  Suppose that for infinitely many
main stages $L$ some clean sub-window $w$ of $L$ has at least $n(w)/\bar D(w)$
dyadic scales whose activity class $a$ \textup{(}without active shift, one-signed,
or mixed\textup{)} satisfies \textup{(KN$_{w,a}$)}.  Then $f\in\Rec$:
$(f,g)\in\cl\NA((c_0,p_N),\ell_2^2)$ for every $g\in\Cset(f)$.
\end{theorem}

\begin{proof}
Fix $g\in\Cset(f)$ and $\rho<1$.  For each such $(L,w)$ with $L\ge l_f$, the
pigeonhole of the proof of [II,~Theorem~5.68] gives a set $\mathcal S$ of at least
$n(w)/\bar D(w)^2\ge L2^{L^3}(\mathrm{Des}/u)^C$ scales of one class and one cube
(this count rests on the condition $Q(w)\ge\bar D(w)^3$ of
[II,~Definition~5.48], not on [II,~Lemma~3.22]).
Let $f^\#$ be the companion of Lemma~\ref{lem:III-CMp} and $\Omega$ the coarse
strict non-peaks of $f^\#$.  The data of Lemma~\ref{lem:III-CMp} form a window family
at $f_0:=f^\#$: (U1), (U2) by (b); (U3) as in the proof of [II,~Theorem~5.68] (at
the banked, pulled and support coordinates of (1a) the data are contact-like or
$t|b|\le K_3|a^\#|$, up to antisymmetric residues counted in $\epsilon$); (U4) by
(c); (U5) by (a); (U6) by Lemma~\ref{lem:III-LV}; (U7) by [II,~Lemma~2.23].  Take
$\gamma_B(w):=\min_mM_mu(w)/4$, so $c_{\rm flat}(w)\ge c_fu(w)$.

\emph{Threshold comparison.}  Use the window parameter $T:=T_{\rm lo}(w)\le
\min\mathcal S$ (scales lie in $[T,t_1]$).  Since $T_{\rm lo}(w)=2^{-n(w)}
T_{\rm hi}(w)$, the number of pieces is at most $n(w)\le\log_2(1/T)$; moreover
$L\le\mathrm{Des}(L)\le n(w)$, $|\Omega|\le2L$, $1/u(w)\le n(w)$
([II,~Proposition~3.14(a)]), $\gap_{\min}\ge c_f\eta(w)M\ge c_fT^4/(L\,\mathrm{Des})$
([II,~Proposition~5.56(b)] with $\eta(w)=T_{\rm lo}(w)^4/(L\,\mathrm{Des})$; robust
and kind-[1$'$] carriers have larger gaps), $x_0\ge c_f\gap_{\min}\Phi_{\min}(L)$
with $1/\Phi_{\min}(L)\le\mathrm{Des}(L)$, $C_S\le C/q_0^2$, and
$\Lambda_{\rm lev}\le C_f(\mathrm{Des}/(u\gap_{\min}))^C$ (Lemma~\ref{lem:III-LV}).  By
Lemma~\ref{lem:III-UE1}, $K_w\le T^{-C'}$ for an absolute $C'$ and $T$ small; in
particular hypothesis (i) of Proposition~\ref{prop:III-S2c} holds with an
absolute $C_0$.  By Lemma~\ref{lem:III-CMp}(c) and Lemma~\ref{lem:III-LV}(c),
hypothesis (ii) holds with $c_*:=c_{L+1}$, since $C_f\mathrm{Des}(L)^C\le
e^{C_0\log^2(1/T)}$, $C_fc_{L+1}^2/T\le c_{L+1}$ and $r_{\rm pull}\le
\mathrm{Des}(L)^2c_{L+1}^2\le c_{L+1}$ (Lemma~\ref{lem:III-LV}(a)); in particular
the radius satisfies $r\ge32K_JT^4s_1\ge r_{\rm pull}$; and $c_{L+1}\le
\exp(-1/T_{\rm lo}(L,M(L)))\le e^{-1/T}$ by (W$_{\exp}$), because
$T_{\rm lo}(L,M(L))\le T_{\rm lo}(w)$.  (W-b) follows from
Lemma~\ref{lem:III-LV}(c) and, for the $s_1$-dependent part of $Y$, from (LATE).
Choosing $N_w$ with $\mathrm{tail}_W\le e^{-1/T}$, for $L$ large
Proposition~\ref{prop:III-S2c} applies, with $s_1=e^{-1/(2T)}$.

\emph{Recovery.}  Apply Theorem~\ref{thm:III-Eeng} along the infinitely many
pairs $(L,w)$, with $f_j:=f^\#$, $\mathcal S_j:=\mathcal S$, $K_j:=K$ of
Lemma~\ref{lem:III-CMp}(b) and $s_{1,j}:=e^{-1/(2T)}$.  (E-c) is
Lemma~\ref{lem:III-CMp}(b).  (E-d): $KT_{\rm hi}(w)\to0$ by [II,~Lemma~3.22], and
$n'c_{\rm flat}/K\ge L2^{L^3}(\mathrm{Des}/u)^Cc_fu/(C_f(\mathrm{Des}/u)^{C'})\to
\infty$ for the constants of [II,~Theorem~5.68].  (E-f) and (E-e):
Proposition~\ref{prop:III-S2c} bounds the relevant quantities by $e^{-1/(3T)}\le
KT^2$, and $\varepsilon_j\le\theta_wT_{\rm lo}(w)^2\le\theta_wT^2$ by
Lemma~\ref{lem:III-CMp}(d); so (E-e) holds with $T'_j:=T$ and $\theta_j:=
\theta_w+e^{-1/(3T)}T^{-2}$, and $\theta_j/c_{\rm flat}^2\le C(\theta_w+e^{-1/(3T)}
T^{-2})n(w)^2\to0$
since $\theta_w\le C_f\mathrm{Des}\,T_{\rm lo}^2\log(1/T_{\rm lo})+C_fc_{L+1}
\log(1/c_{L+1})/T_{\rm lo}^2$ and $n(w)\le\log_2(1/T)$.  Hence $(f,\rho g)\in
\cl\NA$; let $\rho\to1$.
\end{proof}

\begin{corollary}[Step (2) is dispensable]\label{cor:III-noB}
In Lemma~\textup{\ref{lem:III-CMp}} the recursion of step \textup{(2)} may be
omitted: let $f^\#:=f^{(1a)}$ \textup{(}no release, $z$ unchanged on $J_{\rm
fine}$\textup{)} and put $b^\pm:=\pm\frac12\mathcal V$ on $J_{\rm fine}$, the data
being as in Lemma~\textup{\ref{lem:III-CMp}} elsewhere.  Then
\textup{(b)}--\textup{(d)} of Lemma~\textup{\ref{lem:III-CMp}} hold,
$\|\mathcal V1_{J_{\rm fine}}\|_1\le C_f\mathrm{Des}(L)c_{L+1}$, the violation
on $J_{\rm fine}$ is at most $\|\mathcal V1_{J_{\rm fine}}\|_1$ with $b^\theta=0$
there, and \textup{(U5$'$)} of Lemma~\textup{\ref{lem:III-noB}} holds with
$x_0:=c_f\gap_{\min}\Phi_{\min}(L)$.  Consequently
Theorem~\textup{\ref{thm:III-Cmix}} holds with this companion, and no
information on fine carriers \textup{(}ownership, self-alignment, release\textup{)}
is used.
\end{corollary}

\begin{proof}
A coordinate $j\in J_{\rm fine}$ lies in no coarse target and in no near part
$S_l\cap[1,s_{\rm far}(w)]$, so coarse carriers contribute to $\mathcal V(j)$ only
through $S_l\cap(s_{\rm far}(w),\infty)$, in total at most
$C_f\mathrm{Des}\sum_l\lambda_l(\delta_l/n_l)2^{1-s_{\rm far}(w)}\le
C_f\mathrm{Des}\,c_{L+1}^2$ ($2^{-s_{\rm far}(w)}\le c_{L+1}^2$); carriers at stages
$>L$ have forced coefficients $\le C_\Delta\lambda_k$ and $\|u_k\|_1\le1$, in total
$\le C_\Delta\sum_{k>L}\lambda_k\le Cc_{L+1}$.  At a free $j$, $b^\theta_j=0$
and $\mathrm{viol}(j)=|\mathcal V(j)|$; at a contact,
$\mathrm{viol}(j)=(z_j\mathcal V(j))_-\le|\mathcal V(j)|$.  The changes of $g_t$ and
$\Gamma_w$ against Lemma~\ref{lem:III-CMp}(b) are bounded as there.  (U5$'$):
by [II,~Proposition~5.56(b)] every coarse peak has relative margin $\ge u/2$, so
margin $\ge(u/2)\bar\vartheta\Phi\ge x_0$, and every coarse strict non-peak has
gap $\ge\gap_{\min}\ge x_0$ and $\Phi\gap\ge\Phi_{\min}\gap_{\min}\ge x_0$;
without step (2) there is no later drift.  In the proof of
Theorem~\ref{thm:III-Cmix} use Lemma~\ref{lem:III-noB}: $4c_{\rm fine}\le
Cc_{L+1}\le c_{\rm late}\delta e^{-1/(2T)}$ for $L$ large.
\end{proof}

\begin{remark}[Absorbers in the violation-tolerant route]\label{rem:III-abs}
The zero-value absorber pairs of $D^{U1'}$ make coarse data \emph{exact}
([II,~Proposition~5.59]).  They cannot be combined with the levers of
Lemma~\ref{lem:III-DC}: re-tuning an absorber by a mass at its tuning coordinate
has a range of order $m_0(\mu^\star_{p_0})^2$, which is not a design quantity of
level $L$; the coarse levers act at coordinates of $E_c(w)$ whose first absorber
pair has $|u_a(s)|\ge1/(8n_a)$, so a lever move, a pull flip or a window mass moves
the absorber's value far beyond its room, which is of the size of its weight;
untuned, the absorber becomes a peak with nonzero switching coefficient, which
costs a first-order term in the $\theta$ regime.  Therefore the statement that such
absorbers can be re-tuned by levers is false in this application and is not used,
and the mixed-class theorem with tuned absorbers is not claimed.  In the
violation-tolerant route absorbers carry no switching coefficient: they are
ordinary fine carriers, and residues on coarse coordinates become tolerated
violations of mass $O(\mathrm{Des}(L)c_{L+1})$
(Lemma~\ref{lem:III-CMp}(c)).  A finite-model computation illustrates the
obstruction (an absorber pushed $90$--$170$ times beyond its room by a coarse
lever move of size $10^{-6}$); it is evidence only.
\end{remark}

\begin{remark}[What is not proved here]\label{rem:III-open}
Theorem~\ref{thm:III-Cmix} needs (KN$_{w,a}$).  Without it the residual for finite
base support is the status-coherence problem of [II,~Remark~5.70]: rows at which,
at all large main stages and all clean sub-windows, almost all scales lie in
classes violating (KN), i.e.\ an active block with an active near-threshold
switching carrier and no inactive $\Omega$-carrier of robust relative position.
This section does not treat that problem, nor rows with infinite base support.
\end{remark}
\section{The operator $T_{\rm fin}$ and Master Theorem VI}\label{sec:III-MTVI}

This section assembles the results of Sections~\ref{sec:III-generic}
and~\ref{sec:III-UE} into a statement about one operator.  We define a single
operator $T_{\rm fin}$, independent of the number $N$ of blocks, and prove
that for the norms $p_N$ built with it every first row with finite base
support is recoverable, for every $N$, when $T_{\rm fin}$ is taken with
option~(a) of Definition~\ref{def:III-Tfindata} (with option~(b): for $N=1$,
and for $N\ge2$ at the rows with $\varrho_1(f)\ne1$)
(Theorem~\ref{thm:III-MTVI}).  By the reduction of [I,~Section~1] and
[II,~Proposition~6.48] the density of $\NA((c_0,p_N),\ell_2^2)$ is then
\emph{equivalent} (option~(a), every $N$; option~(b), $N=1$) to the approximation
statement (LSC-trunc)$_N$ at infinite base support
(Theorem~\ref{thm:III-VI2}); (LSC-trunc)$_N$ itself is not proved here.

The design is that of [II,~Definition~5.48] (rounds of absorber pairs around
main stages) on the ladder of [II,~Definition~3.5], with the changes listed
in Remark~\ref{rem:III-Tfincomp}, chiefly: the base is $\mu_s=2^{-s^2-1}$
instead of $\mu^\star_s$; there are no fresh-dominated stages; the weights
are chosen \emph{generically} (algebraically independent over a countable
field fixed in advance); the rate scheme and the design factor contain the
objects of Sections~\ref{sec:III-generic} and~\ref{sec:III-UE}; and the
window factor $Q(w)$ does not contain Part~II's $\bar D(w)^3$ and has
exponent $\omega(l)+l^2+20$.  Because the results of
Part~II are proved for Part~II's operator, we list the properties of the
operator that their proofs use and verify them for $T_{\rm fin}$
(Subsection~\ref{subsec:III-transfer}).  The recovery itself follows the
violation-tolerant route of Section~\ref{sec:III-UE}
(Subsection~\ref{sec:III-routeV}): the companion is the exactification of
Proposition~\ref{prop:III-KNtr}; the approximant is the $d$-consistent
engineered approximant of Definition~\ref{def:III-approx}, at which the block
scalars are frozen up to the weight of the fine carriers
(Lemma~\ref{lem:III-PM}); and the averaging is
Theorem~\ref{thm:III-Eeng}.

The verification status of the argument is recorded in
Remark~\ref{rem:III-cert}.  The route through exact data with zero-value
absorbers ([II,~Theorem~5.68]) is \emph{not} used: its absorber tuning
does not transfer to $T_{\rm fin}$ (Remark~\ref{rem:III-routeX}).

\subsection{The design}\label{sec:III-design}

Throughout, $\|\cdot\|$ on $\ell_1$-type objects and the notation of
[I,~Section~1] and [II,~Section~3] are used: a stage $l$ carries a target
$y_l$, a signature $\delta_lh_l$, the normalized target vector $u_l$, a
weight $c_l$, a block $m(l)$ and a position $k(l)$ in that block; $\Phi_l=
2^{-m(l)-k(l)}c_l$ and $\lambda_l=m(l)\Phi_l$.

\begin{definition}[Fixed data of $T_{\rm fin}$]\label{def:III-Tfindata}
\begin{enumerate}
\item[(i)] \emph{Base.}  $H=\ell_2$ with orthonormal basis $(k_s)_{s\ge1}$,
$\mu_s:=2^{-s^2-1}$, and $U:H\to c_0$ the diagonal operator $Uk_s:=\mu_se_s$.
Then $U$ is compact with dense range, $\|U\|=\mu_1=\frac14$, and
$q^*(a)=\|a\|_1+(\sum_s\mu_s^2a_s^2)^{1/2}$.
\item[(ii)] \emph{Signatures.}  $S_l:=\{2^l(2i+1):i\ge1\}$, $Z_0:=\N\setminus
\bigcup_lS_l$, $h_l$, $H_l$, $\delta^{\max}_l$ and $\iota_l=2^{-l}$ as in
[II,~Definition~3.4] (with $\|U\|=\frac14$).  The sets $S_l$ are pairwise
disjoint, consist of even numbers and have bounded gaps $G_l=2^{l+1}$.
\item[(iii)] \emph{Target pool.}  $Y_{\rm pool}:=\{x/q^*(x):0\ne x\in c_{00}$ with
rational entries$\}$, enumerated as $(y^{(i)})_{i\ge1}$ with $y^{(1)}:=
e^*_1/q^*(e^*_1)$; it is norm dense in $S_{q^*}$.  The \emph{absorber pool}
$A_{\rm pool}$ is the set of all $x/q^*(x)$ with $x=\beta e^*_s+e^*_{p_0}+
4\sum_{r=1}^R2^{-r}e^*_{p_r}$, $s\ge1$, $\beta=\pm1$, $R\ge1$ and odd integers
$p_0<\dots<p_R$.  Both pools are countable.
\item[(iv)] \emph{Schedules.}  A block sequence $(m_r)_{r\ge1}$ with $m_1:=2$ in
which every $m\ge1$ occurs infinitely often, and a sequence $(i_\nu)_{\nu\ge1}$
in which every positive integer occurs infinitely often.
\item[(v)] \emph{Normalization option.}  Under \emph{option (a)} the weight
$c_1$ is chosen generically in $[\frac12,1)$ like all other weights; under
\emph{option (b)} $c_1:=1$.
\item[(vi)] \emph{The field $k_0$.}  For $y\in Y_{\rm pool}\cup A_{\rm pool}$ and
$l\ge1$ let $\delta(y,l)$ be the value of $\delta_l$ produced by rule (GM) of
Definition~\ref{def:III-Tfinstage}(S2) for the target $y$ at a stage $l$
($\delta(y,1):=\delta^{\max}_1$), and $n(y,l):=q^*(y+\delta(y,l)h_l)$.  Then
$k_0$ is the subfield of $\R$ generated by all $y(j)$, $\delta(y,l)$, $n(y,l)$
\textup{(}$y\in Y_{\rm pool}\cup A_{\rm pool}$, $l,j\ge1$\textup{)}; it is countable.
\end{enumerate}
\end{definition}

The values $\delta(y,l)$ are well defined independently of the weights,
because (GM) involves only $y$, $H_l$ and $\delta^{\max}_l$; each of them is
$\delta^{\max}_l$ or an endpoint of an excluded interval, hence lies in
$\mathbb Q(y(j):j\ge1)$.  The numbers $n(y,l)$ may be transcendental and are
included by definition.

\begin{definition}[Rounds]\label{def:III-Tfinrounds}
The stages are grouped into rounds $r=1,2,\dots$; $l_r$ denotes the number of
stages before round $r$ ($l_1:=0$).
\begin{enumerate}
\item[(r.1)] \emph{Candidate.}  Let $m:=m_r$, $\nu:=\#\{r'\le r:m_{r'}=m\}$,
$y:=y^{(i_\nu)}$, $\Sigma(y):=\supp y\cap\bigcup_{l'}S_{l'}$ and $p:=l_r+
2|\Sigma(y)|+1$.  The candidate is \emph{allowed with margin} if
(a) $\supp y\cap S_{l'}=\emptyset$ for all $l'\ge p$ and (c$^+$) $\supp y\cap
S_{l'}\subseteq[1,p-2|\supp y|-1]$ for all $l'<p$.  If it is allowed, $Y_r:=y$;
otherwise $Y_r:=y^{(1)}$ and $p:=l_r+1$.  The stage $p_r:=p$ is the \emph{main
stage} of round $r$.
\item[(r.2)] \emph{PRE stages} $l_r+1,\dots,p_r-1$: for each $s\in\Sigma(Y_r)$, in
increasing order, two consecutive stages of block $1$ carrying absorber
targets for $(s,+1)$ and $(s,-1)$ (a \emph{biased pair} for $s$).
\item[(r.3)] \emph{MAIN stage} $p_r$: block $m_r$, target $Y_r$.
\item[(r.4)] \emph{CLUSTER stages} $p_r+1,\dots,p_r+R_{\rm cl}(p_r)$, where
$E^{\rm ch}(p):=T(p)\cup\bigcup_{l''\le p}(S_{l''}\cap[1,p])$, $T(p):=\bigcup_{l''
\le p}\supp y_{l''}$ and $R_{\rm cl}(p):=2|E^{\rm ch}(p)|$: a biased pair of block
$1$ for every $s\in E^{\rm ch}(p)$.  Then $l_{r+1}:=p_r+R_{\rm cl}(p_r)$.
\end{enumerate}
The absorber target of a PRE or CLUSTER stage $l$ for $(s,\beta)$ is
$x/q^*(x)\in A_{\rm pool}$ with $p_0<\dots<p_R$ the least odd integers
$>\max(9,\max T(l-1))$ and $R=R(l)$ the least integer $\ge1$ with
$2^{3-R}\le\mu_{p_0}^2{\rm lb}_l^2/4$, where ${\rm lb}_l$ is the a-priori lower
bound of Definition~\ref{def:III-Tfinstage}(S4).  For every stage $l$, $k(l):=1+
\#\{l'<l:m(l')=m(l)\}$, and $\mathfrak j(k,m)$ denotes the $k$-th stage of
block $m$.
\end{definition}

In round $1$ no candidate meeting a signature set is allowed (the bound in
(c$^+$) is $\le0$), so stage $1$ is the main stage of round $1$ and lies in
block $m_1=2$.  No fresh-dominated stages are used.

\begin{definition}[Stage $l$ of $T_{\rm fin}$]\label{def:III-Tfinstage}
Given all quantities of the stages $<l$, the following are formed in this
order.
\begin{enumerate}
\item[(S1)] \emph{Target} $y_l$ by Definition~\ref{def:III-Tfinrounds}; at an
absorber stage first ${\rm lb}_l$ (see (S4)), then $R(l)$, then $y_l$.
\item[(S2)] \emph{Signature size} $\delta_l$ by rule (GM) of
[II,~Definition~3.5(2)].
\item[(S3)] $n_l$, $u_l$, $v_l$, $\delta^\circ_l$, $\underline y_l$,
$\Lambda^\circ(l)$ as in [II,~Definition~3.5(3)].
\item[(S4)] \emph{Weight.}  At a PRE or MAIN stage, $c^{\max}_l$ is the minimum
of (W1)--(W3), (W5), (W7) of [II,~Definition~3.5(4)], of (W4$''$) of
[II,~Definition~5.48] (with $T_{\rm lo}(l-1,M(l-1))^3$), and, at MAIN stages
$l\ge2$, of (W6).  At the $i$-th CLUSTER stage of the main stage $p$,
$c^{\max}_{p+i}:=\min\{c^{\rm cl}(p)4^{1-i},\ c_{p+i-1}/4\}$, where $c^{\rm cl}(p)$
is the minimum of $c_p/4$, $b(p,M(p))^2$, $(\delta^{\max}_{p+1}H_{p+1}/2)^2$,
(W4$''$) for every $s\in E^{\rm ch}(p)$ (with $T_{\rm lo}(p,M(p))^3$), (W5) and
(W7) at stage $p+1$ with $\delta_{p+1}$ replaced by $\delta^{\max}_{p+1}/2$,
and
\[
 \textup{(W$_{\exp}$)}\qquad \exp\big(-1/T_{\rm lo}(p,M(p))\big).
\]
The a-priori lower bound is ${\rm lb}_l:=(1-2^{-l})c^{\rm low}_l$ at a PRE stage,
with $c^{\rm low}_l$ the minimum of $c_{l-1}/4$, $b(l-1,M(l-1))^2$,
$(\delta^{\max}_lH_l/2)^2$, (W4$''$) at $s$, (W5) with $\delta^{\max}_l/2$ and
(W7); and ${\rm lb}_{p+i}:=(1-2^{-(p+i)})c^{\max}_{p+i}$ at a CLUSTER stage.
\emph{Generic choice:} $c_l$ is a number in $[(1-2^{-l})c^{\max}_l,c^{\max}_l]$
that is transcendental over $k_0(c_1,\dots,c_{l-1})$ (for $l\ge2$, and for
$l=1$ under option (a) with $c^{\max}_1:=1$); under option (b) $c_1:=1$.  Then
$\Phi_l:=2^{-m(l)-k(l)}c_l$ and $\lambda_l:=m(l)\Phi_l$.
\item[(S5)] \emph{Level data.}  $T(l)$, $s_{\max}(l):=\max(T(l)\cup\{1\})$,
$\sigma(l):=\max_{l''\le l}\min(S_{l''}\cap(s_{\max}(l),\infty))$, $\sigma_2(l):=
\max_{l''\le l}$ (second element of $S_{l''}\cap(s_{\max}(l),\infty)$),
$B_\mu(l):=2^{\sigma(l)}/\mu_{\sigma(l)}^2$, $\delta_{\min}(l)$, $G(l)=2^{l+1}$,
$D_{\rm room}(l)$, and the design constants and rate objects of level $l$ of
[II,~Definitions~3.6 and~3.7] with $B^\star(l)$ replaced by $B_\mu(l)$.  In
addition the rate scheme of level $l$ contains, for every pattern $\kappa$ and
every class $a$ of level $l$ (patterns range over the carriers $l''\le l$ and
their blocks), the rate objects of the floor-form cone
$\Gamma^\flat(\kappa,a)$ of Definition~\ref{def:III-cone}, whose aggregated
signature rows carry the $s$-free coefficients $\alpha_{l''}=\|v_{l''}
1_{S_{l''}\cap(l,\infty)}\|_1$; and the constants $\mathrm{Lip}^*_\kappa$, $C_L(\kappa)$, $N_L(\kappa)$, $m_0$,
$C_Y$, $\alpha_Y$, $c_\kappa$, $\beta_\kappa$, $h_\kappa$ of
Definition~\ref{def:III-QCreq}, computed with the actual weights
$c_1,\dots,c_l$; under option (b) only patterns with $1\notin\mathcal N_T(\kappa)$
are used.  $\omega(l)$ is the number of rate objects.  The \emph{design factor}
$\mathrm{Des}(l)$ is the design factor of [II,~Definition~3.5(5)] (with $B_\mu$),
multiplied at main stages by $4^{R_{\rm cl}(l)}2^{k(l+R_{\rm cl}(l))}
\underline y_l^{-1}$, by $\max_{\kappa}\max\{\mathrm{Lip}^*_\kappa,C_L(\kappa),
N_L(\kappa),C_Y,1/m_0,1/c_\kappa,\beta_\kappa,1/h_\kappa,1\}$ (with $1/\infty:=0$),
and by $\max_{l''\le l}n_{l''}2^{\sigma_2(l)}/(\mu_{\sigma_2(l)}^2\delta_{l''})$
\textup{(}condition (D-lev) of Definition~\ref{def:III-designadd} for the base
$\mu$\textup{)}.
\item[(S6)] \emph{Sub-windows.}  $M(l):=\omega(l)+1$; for $w=(l,i)$, $1\le i\le
M(l)$, in lexicographic order: $u(w):=\frac14$ for $w=(1,1)$ and $u(w):=b(w^-)$
otherwise;
\[
 Q(w):=\Big(\frac{4\,\mathrm{Des}(l)}{u(w)}\Big)^{\omega(l)+l^2+20};
\]
$n(w)$, $T_{\rm hi}(w)$, $T_{\rm lo}(w)$ as in [II,~Definition~3.5(6)];
$\eta'(w):=T_{\rm lo}(w)^4/(l\,\mathrm{Des}(l))$; and $b(w):=\min\{b_{\rm II}(w),
b_{\rm QC}(w)\}$, where $b_{\rm II}(w)$ is the number $b(w)$ of
[II,~Definition~3.5(6)] with $\mathrm{Lip}(l)$ replaced by $\mathrm{Lip}^*(l)$, the
maximum of $\mathrm{Lip}(l)$ and of the numbers $\mathrm{Lip}^*_\kappa$, and $b_{\rm QC}(w)$ is the
largest $b$ satisfying (QC4) of Definition~\ref{def:III-QCreq}.
\item[(S7)] $T_{\rm fin}e_{k,m}:=c_{\mathfrak j(k,m)}u_{\mathfrak j(k,m)}$, extended
linearly and continuously to $\ell_1(\N\times\N)$.
\end{enumerate}
\end{definition}

\begin{remark}[Comparison with Part~II]\label{rem:III-Tfincomp}
(a) Relative to [II,~Definitions~3.5 and~5.48]: the base $\mu$ replaces
$\mu^\star$, and $B_\mu(l)$, formed at the first element beyond $s_{\max}(l)$,
replaces $B^\star(l)$, formed at the third; the fresh-dominated stages are
omitted; the absorber length $R$ is computed from ${\rm lb}_l$, and the tuning
coordinate is the least admissible odd integer (no condition of the form
$(\mu^\star_{p_0})^{2^{2^{p_0}}-1}<2^{-R}$); the cluster weights are generic
below $c^{\rm cl}(p)4^{1-i}$ instead of being equal to it; the factor
$\bar D(w)^3$ of [II,~Definition~5.48] does not enter $Q(w)$, and the
exponent of $Q(w)$ contains $l^2$ (this is used in
Lemma~\ref{lem:III-count}); the rate scheme, the design factor and $b(w)$
contain the objects of Definitions~\ref{def:III-QCreq}
and~\ref{def:III-designadd}.
(b) The aggregated signature rows of the level-$l$ rate objects carry the
coefficients $\alpha_{l''}$, which do not depend on $s_{\rm far}(w)$.  The cones
used at a sub-window $w$ of a main stage $L$ have the coefficients
$\|v_{l''}1_{S_{l''}\cap(L,s_{\rm far}(w)]}\|_1$, and $s_{\rm far}(w)$ depends on
$c_{L+1}$, which is chosen after the level-$L$ data.  The two matrices differ by
positive row scalings in $[1-2^{1+L-9\cdot2^L},1]$
(Lemma~\ref{lem:III-weightfree}(c)).  The scalings are $\le1$, so the zero
sets coincide, a rate object that is tiny for the rate scheme is tiny for
the cone, and one that is robust ($\ge u$) for the rate scheme has modulus
$\ge u/2$ for the cone, because the minors involve at most $R_\Gamma(L)$
scaled rows and $(1-2^{1+L-9\cdot2^L})^{R_\Gamma(L)}\ge\frac12$; Hoffman
\cite{Hoffman} and {\L}ojasiewicz constants agree up to a fixed factor.
Writing the rate objects with $\alpha_{l''}$ removes a forward reference to
$c_{L+1}$ from the level-$L$ data.
(c) The Lipschitz constant $\mathrm{Lip}^*$ in $b_{\rm II}$ and in (QC4) dominates
the Lipschitz constant of the floor-status function $F_\Phi$, which carries the
factor $1/\Phi_d$, and those of the floor statuses of all active coarse
carriers; the Lipschitz constant of the determinantal objects alone does not
dominate it.
\end{remark}

\begin{lemma}[Genericity of the weights]\label{lem:III-GENfin}
\textup{(a)} Every datum of $T_{\rm fin}$ that is not a weight lies in $k_0$ or is
an integer: $y_l\in Y_{\rm pool}\cup A_{\rm pool}$, $\delta_l=\delta(y_l,l)$,
$n_l=n(y_l,l)$, every value $u_l(j)$, every $\mu_s$ and $H_l$, every coefficient
$\alpha_l$ and every aggregated coefficient $\|v_l1_{S_l\cap(a,b]}\|_1$
\textup{(}$a<b\le\infty$ integers\textup{)}, and the integers $m(l)$, $k(l)$,
$R(l)$, $p_i$ and the stage types.

\textup{(b)} Under option \textup{(a)} the family $(c_l)_{l\ge1}$, under option
\textup{(b)} the family $(c_l)_{l\ge2}$, is algebraically independent over
$k_0$.  Consequently, for every finite set $D$ of stages \textup{(}not containing
$1$ under option \textup{(b)}\textup{)}, $(\Phi_d)_{d\in D}$ is algebraically
independent over $k_0(\Phi_l:l\notin D)$, and for every pattern $\kappa$ of
every level \textup{(}with $1\notin\mathcal N_T(\kappa)$ under option
\textup{(b)}\textup{)}, $(\Phi_d)_{d\in\mathcal N_T(\kappa)}\notin B_\kappa$.
\end{lemma}

\begin{proof}
(a) Main targets lie in $Y_{\rm pool}$ and absorber targets in $A_{\rm pool}$,
whatever the integers $s,\beta,R,p_i$ are; then $\delta_l=\delta(y_l,l)$ and
$n_l=n(y_l,l)$ by definition, $u_l(j)=(y_l(j)+\delta_l2^{-j}1_{S_l}(j))/n_l$, and
$\|v_l1_{S_l\cap(a,b]}\|_1=(\delta_l/n_l)\sum_{s\in S_l\cap(a,b]}2^{-s}$ with a
rational sum; $\mu_s$ and $H_l$ are rational.
(b) If $c_1,\dots,c_n$ are algebraically independent over $k_0$ and $c_{n+1}$
is transcendental over $k_0(c_1,\dots,c_n)$, then $c_1,\dots,c_{n+1}$ are
algebraically independent over $k_0$; induction gives the first claim, and
$\Phi_l$ is a nonzero rational multiple of $c_l$.  A subfamily of an
algebraically independent family is algebraically independent over the field
generated by the complementary subfamily.  For the last claim apply
Lemma~\ref{lem:III-GEN} with this $k_0$ (which contains every $u_l(j)$ and
every $\alpha_l$ by (a)) and $k:=k_0(\Phi_l:l\le L,\ l\notin\mathcal N_T)$.  By
Lemma~\ref{lem:III-weightfree}(a), $B_\kappa$ contains no datum of a row $f$:
a row only selects the pattern and the point $r(f)$.
\end{proof}

\begin{lemma}[Well-foundedness, positivity, $N$-independence]\label{lem:III-Wfin}
\textup{(a)} Every quantity of stage $l$ is a function of the fixed data, of the
quantities of the stages $<l$ and of the earlier steps of stage $l$.
\textup{(b)} Every set from which a choice is made is nonempty, and every
quantity used as a divisor is finite and positive; in particular $c_l>0$,
$R(l)$ exists, the constants of \textup{(S5)} are finite and positive, and
$b(w)>0$.
\textup{(c)} No rule refers to $N$.
\textup{(d)} \textup{(Weights.)}  $c_{l+1}\le c_l/4$ for every $l$.  For every
main stage $L$ and every sub-window $w$ of $L$:
\[
 \sum_{l>L}\lambda_l\le\frac{c_{L+1}}3,\qquad c_{L+1}\le b(w)^2\le T_{\rm lo}(w)^8,
 \qquad c_{L+1}\le\exp\big(-1/T_{\rm lo}(w)\big).
\]
\textup{(e)} \textup{(Lower bounds.)}  $c_l\ge c^{\max}_l/2$ for every $l\ge2$;
at a PRE stage $c_l/8\le{\rm lb}_l\le c_l$; at the $i$-th cluster stage of a main
stage $L$, $c_{L+i}\ge c^{\rm cl}(L)4^{1-i}/2$.
\textup{(f)} $T_{\rm fin}$ satisfies the QC requirements \textup{(QC1)--(QC5)} of
Definition~\textup{\ref{def:III-QCreq}} at every main stage, and \textup{(D-lev)}
of Definition~\textup{\ref{def:III-designadd}} with $\mu$ in place of $\mu^\star$.
\end{lemma}

\begin{proof}
(a) At stage $l$ the order is: target (at absorber stages ${\rm lb}_l$, which
uses $c_{l-1}$, $b(l-1,M(l-1))$, $T_{\rm lo}(l-1,M(l-1))$, $\delta^{\max}_l$, $H_l$,
$s$ and the (W7)-data of stage $l-1$, or, at a cluster stage, the data of the
main stage $p$ and $c_{p+i-1}$; then $R(l)$; then $y_l$), $\delta_l$, $n_l$,
$u_l$, $c^{\max}_l$, the generic $c_l$ (which uses only the field
$k_0(c_1,\dots,c_{l-1})$), the level data of (S5) (computed with $c_1,\dots,c_l$),
$\mathrm{Des}(l)$, and the sub-windows of (S6) in lexicographic order.  The number
$k(l+R_{\rm cl}(l))$ at a main stage $l$ is known at stage $l$, because the
stages $l+1,\dots,l+R_{\rm cl}(l)$ are block-$1$ cluster stages and $R_{\rm cl}(l)$
depends only on $T(l)$.  Requirements about later stages, namely (W2), (W7) and
(W$_{\exp}$), are imposed at those stages.  The aggregated rows of (S5) carry
$\alpha_{l''}$, which does not depend on later weights
(Remark~\ref{rem:III-Tfincomp}(b)).

(b) The set of admissible $\delta$ in (GM) is $[\frac12\delta^{\max}_l,
\delta^{\max}_l]$ minus at most $2\cdot3^{|\supp y_l|}$ open intervals of total
length $\le\delta^{\max}_l/8$, a nonempty finite union of closed intervals;
$c^{\max}_l$ is a minimum of finitely many positive numbers; the interval
$[(1-2^{-l})c^{\max}_l,c^{\max}_l]$ has positive length and $k_0(c_1,\dots,
c_{l-1})$ is countable, so a transcendental $c_l$ exists; $2^{3-R}\to0$, so
$R(l)$ exists.  The objects of [II,~Definition~3.5] are finite and positive as in
[II,~Lemma~3.10]; the constants of Definition~\ref{def:III-QCreq} are finite and
positive by Remark~\ref{rem:III-QCwell}(i), which applies because no pattern of
level $l$ is a coincidence (Lemma~\ref{lem:III-GENfin}(b); a pattern of level $l$
involves only $c_1,\dots,c_l$, chosen before the level-$l$ data).  $b_{\rm QC}(w)
>0$ by Remark~\ref{rem:III-QCwell}(ii), and $b_{\rm II}(w)>0$ trivially.

(c) The rounds, the schedules, the targets and the weights make no reference
to $N$, and the patterns range over all carriers $\le l$ and their blocks.  For
$p_N$ the carriers of blocks $m>N$ are absent; a pattern of $p_N$ is the design
pattern in which these blocks are inactive (they contribute no row), which the
design enumerates.

(d) $c_l\le c_{l-1}/4$ by (W1) at PRE and MAIN stages and by the cluster rule.
Hence $\sum_{l>L}c_l\le\frac43c_{L+1}$, and $\lambda_l\le c_l/4$
($m2^{-m}\le\frac12$, $2^{-k}\le\frac12$).  Stage $L+1$ is the first cluster
stage of $L$ ($R_{\rm cl}(L)\ge2$, as $T(L)\ne\emptyset$), so $c_{L+1}\le
c^{\rm cl}(L)\le b(L,M(L))^2$, and $b(L,M(L))\le b(w)$ because $b$ decreases
along the sub-windows of a level ($b(w)<u(w)=b(w^-)$); $b(w)\le\eta'(w)\le
T_{\rm lo}(w)^4$.  Finally $c_{L+1}\le\exp(-1/T_{\rm lo}(L,M(L)))\le\exp(-1/T_{\rm lo}
(w))$, since $T_{\rm lo}$ decreases along the sub-windows of $L$.

(e) $c_l\ge(1-2^{-l})c^{\max}_l\ge c^{\max}_l/2$.  At a PRE stage each entry of
$c^{\rm low}_l$ is at most the corresponding entry of $c^{\max}_l$ (because
$\delta_l\ge\delta^{\max}_l/2$) and at least a quarter of it, so $c_l/8\le
(1-2^{-l})c^{\max}_l/4\le{\rm lb}_l\le c_l$.  At cluster stages, by induction
from $c^{\max}_{L+i}=\min(c^{\rm cl}(L)4^{1-i},c_{L+i-1}/4)$,
$c_{L+i}\ge\prod_{j\le i}(1-2^{-(L+j)})c^{\rm cl}(L)4^{1-i}\ge(1-\sum_{n>L}2^{-n})
c^{\rm cl}(L)4^{1-i}\ge c^{\rm cl}(L)4^{1-i}/2$.

(f) (QC1) is Lemma~\ref{lem:III-GENfin}(b); (QC2) and (QC3) hold by (S5)
($\mathrm{Des}(L)$ contains $C_*(L)$ and $D_{\rm room}(L)\ge1/\Phi_{\min}(L)$ through
the factor of [II,~Definition~3.5(5)], and the QC constants by the additional
factor); (QC4) holds because $b(w)\le b_{\rm QC}(w)$ and $b(w)\le b_{\rm II}(w)\le
\eta(w)=\eta'(w)$; (QC5) is (d).  (D-lev) holds by the last factor of
$\mathrm{Des}(L)$.
\end{proof}

\begin{theorem}[Admissibility; Mart\'in's Lemma~B for $T_{\rm fin}$]\label{thm:III-Afin}
\textup{(a)} $q^*(T_{\rm fin}e_{k,m})=c_{\mathfrak j(k,m)}\le c_1$; under option
\textup{(a)} $c_1\in[\frac12,1)$, under option \textup{(b)} $c_1=1$, attained at
$(k,m)=(1,2)$.

\textup{(b)} $T_{\rm fin}$ satisfies \textup{(T-b)}, \textup{(T-c)} and
\textup{(T-d)} of [I,~Definition~1.3].  Under option \textup{(b)} it is admissible.
Under option \textup{(a)} it satisfies \textup{(T-a)} without its equality clause,
$\|T_{\rm fin}\|=c_1<1$.

\textup{(c)} The range $Y:=T_{\rm fin}(\ell_1(\N\times\N))$ is a separable operator
range with $Y\cap\NA((c_0,q),\R)=\{0\}$, and for every $m$ the normalized
vectors $\{u_{k,m}:k\ge1\}$ are norm dense in $S_{q^*}$.  Under option
\textup{(a)}, the norms $p_N$ and $p$ built with $T_{\rm fin}$ are Mart\'in-type
norms: they are given by Mart\'in's construction \cite{Martin} with an
injective operator of norm $c_1\le1$ having these properties.
\end{theorem}

\begin{proof}
(a) $q^*(u_l)=1$ and $c_l\le c_{l-1}/4$ (Lemma~\ref{lem:III-Wfin}(d)).

(b) \emph{Allowedness.}  Every target satisfies allowedness (a): main targets
by (r.1).  An absorber target for $(s,\beta)$ at stage $l$ has support
$\{s,p_0,\dots,p_R\}$ with $p_i$ odd, hence in $Z_0$, and $s$ either lies in
$Z_0$ or lies in some $S_{l''}$ with $l''<l$: for a PRE pair of round $r$,
$s\le p-2|\supp Y_r|-1\le l_r<l$ by (c$^+$), and $l''<\min S_{l''}\le s$; for a
cluster pair of the main stage $p<l$, either $s\in S_{l''}\cap[1,p]$, so $l''<
s\le p$, or $s\in T(p)\cap S_{l''}$, and then $l''$ is smaller than the stage of
the target containing $s$ by allowedness (a) of that target.  Consequently, for
$s\in S_{l'}$: $u_l(s)=0$ for $l<l'$, $u_{l'}(s)=\delta_{l'}2^{-s}/n_{l'}$, and
every target $y_l$ with $l>l'$ and $s\in\supp y_l$ satisfies
\[
 c_l\le\iota_{l'}2^{-2s}c_{l'}\delta_{l'}/2 :
\]
at PRE and MAIN stages by (W4$''$); at a cluster stage $p+i$ because
$c_{p+i}\le c^{\rm cl}(p)$, which contains (W4$''$) for every $s\in E^{\rm ch}(p)$,
and the coordinate $s$ of a cluster target lies in $E^{\rm ch}(p)$.  Also $n_l\in
[\frac34,\frac54]$, since $\delta_lq^*(h_l)\le\delta^{\max}_l(1+\|U\|)H_l\le\frac14$.

\emph{(T-b), (T-c).}  Let $x\in\ell_1(\N\times\N)$, indexed by stages, with
$Z:=T_{\rm fin}x=\sum_lx_lc_lu_l\in c_{00}$, and suppose $x_{l'}\ne0$.  For
$s\in S_{l'}$,
\[
 Z(s)=x_{l'}c_{l'}\delta_{l'}2^{-s}/n_{l'}+\sum_{l>l',\ s\in\supp y_l}x_lc_ly_l(s)/n_l .
\]
If the sum is nonempty with least index $L(s)$, then $\sum_{l\ge L(s)}c_l\le
\frac43c_{L(s)}\le\frac23\iota_{l'}2^{-2s}c_{l'}\delta_{l'}$, and $|y_l(s)|\le1$,
so
\[
 |Z(s)|\ge c_{l'}\delta_{l'}2^{-s}\big(|x_{l'}|/n_{l'}-\tfrac89\|x\|_\infty2^{-s}\big)>0
\]
for all large $s\in S_{l'}$.  This contradicts $Z\in c_{00}$.  So $x=0$, which is
(T-b), and $Y\cap c_{00}=\{0\}$, which is (T-c).

\emph{(T-d).}  Fix $m$ and $i$.  Block $m$ is the block of infinitely many
rounds, and $i_\nu=i$ for infinitely many $\nu$; at such a round the candidate
is $y^{(i)}$, and it is allowed with margin as soon as the main stage $p$
exceeds $1+\max\supp y^{(i)}+2|\supp y^{(i)}|$ (for (a): $S_{l'}\cap\supp y^{(i)}
\ne\emptyset$ forces $l'<\min S_{l'}\le\max\supp y^{(i)}<p$; for (c$^+$):
$\supp y^{(i)}\subseteq[1,p-2|\supp y^{(i)}|-1]$).  Hence $y_l=y^{(i)}$ for
infinitely many stages $l$ of block $m$, and $q^*(u_l-y^{(i)})\le|1/n_l-1|+
\delta_lq^*(h_l)/n_l\le4(1+\|U\|)2^{-l}H_l\to0$ along them.  Density of
$Y_{\rm pool}$ in $S_{q^*}$ gives (T-d).  Under option (b) (T-a) holds with
equality at $(1,2)$; under option (a) $\|T_{\rm fin}\|=\sup_lc_l=c_1<1$.

(c) $Y$ is the range of a bounded operator on a separable space; $\NA((c_0,q),\R)
=c_{00}$ [I,~Proposition~1.10(c)] (the proof uses only that $U$ is compact with
dense range), so (T-c) gives $Y\cap\NA((c_0,q),\R)=\{0\}$.  The density statement
is (T-d).  Mart\'in's construction [I,~Remarks~1.1 and~1.2] uses of the operator
only injectivity, the condition on the range, the density of the normalized
vectors in every block, and $\|T\|\le1$ (through $\Phi_m(k)\le2^{-m-k}$ and
$\|\Phi_m\|_1<1$ in Mart\'in's Lemma~A).
\end{proof}

\begin{remark}[The two options]\label{rem:III-options}
Option (b) keeps $T_{\rm fin}$ inside the class of [I,~Definition~1.3]; it costs
the genericity of one weight, $c_1=1$, and leads to the exception in
Theorem~\ref{thm:III-MTVI}(b).  Option (a) leaves that class only through the
equality clause of (T-a).  The results of Part~I used below remain valid,
because their proofs use the operator only through (T-b)--(T-d) and the upper
bound $q^*(Te_{k,m})\le1$ (Lemma~\ref{lem:III-partI}).  The normalized operator
$T_{\rm fin}/c_1$ of option (a) does not reduce option (b) to option (a): the
block norms are not homogeneous in the weights (the unit ball of $|\cdot|_m$
is $B_{\ell_1}+D_m(B_{\ell_2})$), so the coincidence sets of the two options are
unrelated.
\end{remark}

\subsection{Transfer of the tools to $T_{\rm fin}$}\label{subsec:III-transfer}

The results of [I] are stated for admissible operators, and those of [II],
Sections~3--5, for $T_{\rm final}$ and $D^{U1'}$.  We now show that every
result used below holds for $T_{\rm fin}$.  For [I] the point is the equality
clause of (T-a); for [II] we list the properties of the operator through which
the proofs of the results used below depend on it, and verify each for
$T_{\rm fin}$.

\begin{lemma}[Part~I under option (a)]\label{lem:III-partI}
The following results of [I] hold for every operator satisfying \textup{(T-b)},
\textup{(T-c)}, \textup{(T-d)} and $q^*(Te_{k,m})\le1$ for all $(k,m)$, in
particular for $T_{\rm fin}$ under both options: Propositions~1.9--1.13,
Lemmas~1.5--1.7 and~1.16, Definition~3.1 and Lemmas~3.2--3.4 and~3.6, the
results of Section~7 used in Section~\ref{sec:III-UE}
\textup{(}7.13--7.17, Theorem~7.18 and the lemmas used in its proof\textup{)},
Lemmas~8.8 and~8.10, and Lemma~8.27 \textup{(}equation~(16)\textup{)}.  The
same holds for the results of [I] used only in Section~\ref{sec:III-inf}:
Lemmas~1.6, 7.1, 8.7, 8.26 and~8.28 and equation~(15).
\end{lemma}

\begin{proof}
In [I] the operator enters through $\Phi_m(k)=2^{-m-k}q^*(Te_{k,m})$, through
$L^*w=T(\sum_m\sum_km2^{-m-k}w_m(k)e_{k,m})$ (an identity valid for every $T$),
and through (T-b)--(T-d).  The proofs of the listed results use, besides
(T-b)--(T-d), only the upper bounds $\Phi_m(k)\le2^{-m-k}$, $\lambda_{k,m}\le
m2^{-m-k}$, $\|\Phi_m\|_1\le2^{-m}<1$, $\|\Phi_m\|_2<\frac12$ and $\|R_m\|\le
m2^{-m}$, all implied by $q^*(Te_{k,m})\le1$.  The equality clause of (T-a)
and lower bounds on $q^*(Te_{k,m})$ occur in [I] only in
Definition~1.3, in the verifications that particular operators are
admissible, and in the resonance example of [I,~Section~6] and
[I,~Corollary~7.23]; none of these is used here.  Lemma~8.27 (equation~(16))
is derived from Lemma~8.8 without further use of the operator; Lemmas~8.7,
8.26 and~8.28 and equation~(15) are statements about the forced data and
the clamp formula whose proofs use $T$ only through (T-b)--(T-d) and
$q^*(Te_{k,m})\le1$ (via $\Phi_m(k)\le2^{-m-k}$ and
$\lambda_{k,m}\le m2^{-m-k}$).
\end{proof}

We write \emph{MS} for ``at the sub-windows of main stages''; windows are used
only there.

\begin{lemma}[The interface]\label{lem:III-interface}
$T_{\rm fin}$ has the following properties.
\begin{enumerate}
\item[(B1)] The base is diagonal.  Hence the identities of
[II,~Lemma~3.2(b)] and those of [II,~Lemma~4.14] hold exactly.
\item[(B2)] $\|U\|=\frac14$, $\mu_s\in(0,1]$ is decreasing, $\nu\le\|U\|$.
\item[(B3)] \textup{(Bank efficiency.)}  For every $l$, every $l''\le l$ and
$j':=\min(S_{l''}\cap(s_{\max}(l),\infty))$:
$\mu_{j'}^2v_{l''}(j')\ge\mathrm{Des}(l)^{-1/6}$, $(\mu_{j'}v_{l''}(j'))^{-2}\le
\mathrm{Des}(l)^{1/3}$ and $v_{l''}(j')\ge\mathrm{Des}(l)^{-1/6}$.
\item[(B4)] \textup{(Lever efficiency.)}  $\mu_j^2v_{l}(j)\ge1/\mathrm{Des}(L)$ for
$l\le L$ and $j\in S_l\cap(s_{\max}(L),\sigma_2(L)]$.
\item[(S1)] Allowedness (a) holds for every target, (c) with margin for main
targets; $u_{l'}(s)=0$ for $l'<l$ and $u_l=v_l$ on $S_l$ beyond $\max\supp y_l$.
\item[(S2)] $s_{\max}(l)\to\infty$.
\item[(W1)] MS: $\sum_{l'>L}\lambda_{l'}\le c_{L+1}/3\le b(L,M(L))^2/3$.
\item[(W2)] $\Phi_k\le c_k\le(\delta_kH_k)^2$ for every stage $k\ge2$.
\item[(W3)] \textup{(W4$''$)} holds at every target with $T_{\rm lo}(\cdot)^3$ read
at the preceding main stage at cluster stages.
\item[(W4)] $c_{l+1}\le c_l/4$; for consecutive stages $l<l'$ of block $1$,
$\lambda_l\ge8\lambda_{l'}$.
\item[(W5)] MS: $c_{L+1}\le\exp(-1/T_{\rm lo}(L,M(L)))$.
\item[(W6)] \textup{(Lower bounds.)}  The only lower bounds on weights used by
the results of Proposition~\textup{\ref{prop:III-transfer}} are: the least
entry $\eta_k$ of an absorber target satisfies $\eta_k\ge\mu_{p_0}^2c_k^2/2^{13}$;
and bounds computed from the actual weights after they are chosen
\textup{(}$\lambda_l\ge1/D(l)$, $\Phi_{\min}(L)\ge1/\mathrm{Des}(L)$\textup{)}.
\item[(D1)] $\mathrm{Des}(l)$ dominates every design constant used by the results
of Proposition~\textup{\ref{prop:III-transfer}} as in [II,~Lemma~4.1(c)] and
[II,~Proposition~3.16]: $l$, $2^{G(l)}$, $1+|T(l)|$, $D_{\rm room}(l)$, $G^{**}(l)$,
$H_{\rm tune}(l)$, $H_{\rm conf}(l)\le\mathrm{Des}(l)^{1/6}$, and $\mathrm{Des}(l)\ge
\mathrm{Lip}^*(l)$, $C_*(l)$, $C_L(l)$, $N_L(l)$, $2^{6l^3}$, $\Lambda^\circ(l)$.
\item[(D2)] MS: $Q(w)\ge(4\mathrm{Des}(l)/u(w))^{\omega(l)+20}$; $n(w)$, $T_{\rm hi}(w)$,
$T_{\rm lo}(w)$ are defined from $Q(w)$ as in [II,~Definition~3.5(6)];
$u(w^+)=b(w)<u(w)$; $b(w)$ satisfies every \emph{upper} bound imposed on $b(w)$
in [II].
\item[(D3)] \textup{(GW1)--(GW4)} of [II,~Lemma~3.22] hold MS.
\item[(A1)] \textup{(First touchers.)}  For a sub-window $w$ of a main stage $L$
and $s\in E_c(w)$, the first two carriers after $L$ whose targets contain $s$
form a biased pair, or there is no such carrier.
\end{enumerate}
\end{lemma}

\begin{proof}
(B1), (B2): Definition~\ref{def:III-Tfindata}(i); $\nu=\|U^*A\|$ with
$\|A\|_1\le1$.
(B3) $j'\le\sigma(l)$ and $n_{l''}\le\frac54$, so $v_{l''}(j')=\delta_{l''}
2^{-j'}/n_{l''}\ge\frac45\delta_{\min}(l)2^{-\sigma(l)}$, and, $\mu$ being
decreasing, $\mu_{j'}^2v_{l''}(j')\ge\frac45\delta_{\min}(l)/B_\mu(l)$.  The
bracket of $\mathrm{Des}(l)$ ([II,~Definition~3.5(5)] with $B_\mu$) contains
$l2^{l^3}B_\mu(l)\delta_{\min}(l)^{-1}$ to the sixth power, and $l2^{l^3}\ge2$,
so $\mathrm{Des}(l)^{1/6}\ge2B_\mu(l)/\delta_{\min}(l)$, which gives the three
inequalities ($B_\mu\ge2^\sigma$ for the third).
(B4) is the last factor of $\mathrm{Des}(L)$, as in Lemma~\ref{lem:III-designadd}.
(S1) is the allowedness argument in the proof of Theorem~\ref{thm:III-Afin}.
(S2) Every main stage $p$ is followed by $R_{\rm cl}(p)\ge2$ absorber stages, and
an absorber target at stage $l$ contains odd integers $p_R>\max T(l-1)$, so
$s_{\max}(l)>s_{\max}(l-1)$ at every absorber stage.  [This replaces the use of
fresh-dominated stages in [II,~Lemma~4.1(e)], the only use of those stages in
the results below.]
(W1), (W5): Lemma~\ref{lem:III-Wfin}(d).
(W2) At PRE and MAIN stages $k\ge2$ by (W3) of [II] (at stage $1$, which
carries $c_1\ge\frac12$, the inequality fails and is not claimed; (W2) is
used only for fine carriers $k>L\ge l_f$); at stage $p+1$ because $c^{\rm cl}(p)
\le(\delta^{\max}_{p+1}H_{p+1}/2)^2\le(\delta_{p+1}H_{p+1})^2$; at $p+i$, $i\ge2$:
$c_{p+i}\le b(p,M(p))^2\le2^{-8n(p,M(p))}$, while $(\delta_{p+i}H_{p+i})^2\ge
2^{-2(p+i)-4-6\cdot2^{p+i}}$, $p+i\le3p+2s_{\max}(p)$ and $n(p,M(p))\ge p2^{p^3}
\mathrm{Des}(p)\ge p2^{p^3}2^{12s_{\max}(p)^2}$ (as $B_\mu(p)^6\ge2^{12\sigma(p)^2}$).
(W3) by (S4).  (W4) by (W1) of [II] and the cluster rule; for consecutive
block-$1$ stages $k$ increases by $1$ and $c$ decreases by a factor $\ge4$.
(W6) Reading the proofs of the results listed in
Proposition~\ref{prop:III-transfer} shows that lower bounds on weights enter
only through the least absorber entry in [II,~Lemma~5.61(ii)] (owners inside
absorber blocks) and through the design-computed bounds; the absorber length
of the zero-value tuning [II,~Lemma~5.51] and the cluster capacity in
[II,~Proposition~5.59(c)] belong to the route not used here.  For the least
entry: $R$ is minimal, so $2^{4-R}>\mu_{p_0}^2{\rm lb}_k^2/4\ge\mu_{p_0}^2c_k^2/256$
by Lemma~\ref{lem:III-Wfin}(e); and $q^*(x)\le8$, so the least entry
$4\cdot2^{-R}/q^*(x)$ is $\ge2^{-R}/2\ge\mu_{p_0}^2c_k^2/2^{13}$.  The chain of
[II,~Lemma~5.61(ii)], $n(p'-1)\ge\mathrm{Des}(p'-1)\ge(B_\mu(p'-1)D(p'-1))^6\ge
(\eta_k\lambda_k)^{-6}$ for the first stage $p'$ after the absorber sequence,
survives the factor $2^{13}$: first, $\sigma(p'-1)>s_{\max}(p'-1)\ge p_R(k)=
p_0+2R(k)>R(k)$; second, the definition of $R(k)$ gives $2^{R(k)}\ge32/
(\mu_{p_0}^2{\rm lb}_k^2)\ge2^{11}/(\mu_{p_0}^2c_k^2)$, since ${\rm lb}_k\ge c_k/8$;
hence, with $1/\mu_\sigma^2\ge16$,
\[
 B_\mu(p'-1)=2^{\sigma(p'-1)}/\mu_{\sigma(p'-1)}^2\ge2^{15}/(\mu_{p_0}^2c_k^2),
\]
and, as $D(p'-1)\ge1/\Phi_k$, $\eta_k\ge\mu_{p_0}^2c_k^2/2^{13}$ and $\lambda_k\ge
\Phi_k$,
\[
 B_\mu(p'-1)D(p'-1)\ge\frac{2^{15}}{\mu_{p_0}^2c_k^2\Phi_k}\ge4(\eta_k
 \lambda_k)^{-1}.
\]
(Here $D(l):=D_{\rm room}(l)$ of [II,~Subsection~5.6]; [II,~Lemma~5.61(ii)]
needs only $n(p'-1)\ge\log_2(4C_f/(T_{\rm lo}\eta_k\lambda_k))$, so the
margin is large.)
(D1) The bracket of $\mathrm{Des}(l)$ contains these factors (with $B_\mu$ for
$B^\star$), and the remaining factors are products in $\mathrm{Des}(l)$.  The
results below use $G^{**}(l)$ ([II,~Proposition~4.24]) and not $H_{\rm conf}(l)$
itself, which enters only $\mathrm{Des}$ and the bound $G^{**}\le H_{\rm conf}$.
(D2) The exponent $\omega(l)+l^2+20\ge\omega(l)+20$ and $4\mathrm{Des}/u\ge1$.  The
bands are disjoint: $Q(w)\ge(4/u(w))^{20}$ gives $T_{\rm hi}(w)\le u(w)^{20}4^{-20}$
and $b(w)\le\eta'(w)\le T_{\rm hi}(w)^4<u(w)$.  $b(w)\le b_{\rm II}(w)$, and
$b_{\rm II}(w)$ is Part~II's $b(w)$ with a larger Lipschitz constant, hence
smaller; no result of [II] uses a lower bound on $b(w)$ (it enters later
quantities only through $c_{L+1}\le b^2$, $u(w^+)=b(w)$ and $1/u$ in $Q(w^+)$).
(D3) (GW1) follows from (D2) and $\mathrm{Des}\ge\Lambda^\circ$; (GW2) from (D2) and
$\mathrm{Des}(l)\ge(l2^{l^3}\Lambda^\circ(l))^6$; (GW3) from (W1) and the box bound
[I,~Lemma~8.10]; (GW4) from the definition of $T_{\rm hi}$; [II,~Lemma~3.22(b)]
uses (T-a) only through $q^*(Te_{k,m})\le1$.
(A1) Let $s\in E_c(w)$.  If $s\in T(L)$ or $s\in S_{l''}\cap[1,L]$, then $s\in
E^{\rm ch}(L)$, whose biased pair is in the cluster immediately after $L$, and
the absorber targets of that cluster are the first targets after $L$ that
contain $s$.  Otherwise $s\in S_{l''}\cap(L,s_{\rm far}(w)]$; later carriers
whose targets contain $s$ are PRE pairs of later main targets containing $s$
or cluster pairs of later main stages $p\ge s$, because absorber targets
contain, apart from their own $s$, only fresh odd integers, signature sets are
disjoint, and $y^{(1)}$ is supported in $Z_0$; the first such carriers form a
biased pair.
\end{proof}

\begin{proposition}[Transfer list]\label{prop:III-transfer}
The following results of [II] hold for $T_{\rm fin}$, with $\mathrm{Des}$ the design
factor of $T_{\rm fin}$, $\mu$ in place of $\mu^\star$ and ``level'' read as
``main stage'': Lemmas~2.9, 2.10, 2.12, 2.22, 2.23, 2.30, Corollary~2.11 and
Proposition~5.46 \textup{(}no property of the operator is used\textup{)};
Lemmas~5.1, 5.3, 5.14 and Theorem~5.7 \textup{(}algebra, {\L}ojasiewicz,
averaging, and \textup{(B3), (D1), (D2)}\textup{)}; Theorem~3.12, Proposition~3.14
and Lemma~3.22 \textup{(S1, W1, D2, D3)}; Theorem~3.19 and Corollary~3.20
\textup{(D2)}; Lemmas~4.7, 4.14, 4.15, 4.17--4.19, Definition~4.16 and
Proposition~4.24 \textup{(B1--B3, D1, D2, S1, S2, W1)}; Lemma~5.50 \textup{(A1)};
Definition~5.52, Lemma~5.53 and Proposition~5.54(a)--(c),(e) \textup{(D1, D2,
B1)}; Proposition~5.56 \textup{(B3, D1, D2)}; Definition~5.57(2), Lemmas~5.58,
5.60, 5.61, 5.62 and~5.64 \textup{(W1--W4, W6, D1, D2)}.  The counting bounds
``$\le C_f\bar D(w)$'' of Lemma~5.53 and Proposition~5.54(d) are replaced by
Lemma~\textup{\ref{lem:III-count}}.  Lemma~5.51 and Proposition~5.59(b)--(d)
\textup{(}zero-value absorbers\textup{)} are not claimed for $T_{\rm fin}$ and are
not used.
\end{proposition}

\begin{proof}
For each listed result, the proof in [II] uses the operator only through the
properties named in parentheses (and through the results of [I] covered by
Lemma~\ref{lem:III-partI}), and these hold for $T_{\rm fin}$ by
Lemma~\ref{lem:III-interface}.  We record the points where the two operators
differ quantitatively.  (i) The base enters Lemma~4.15 (tuning) and the
results built on it only through (B1)--(B3): with (B3), the contraction bound
$|M'|\le C_f\mathrm{Des}^{2/3}\varepsilon_x$ and the bound $\kappa_T\le C_f\mathrm{Des}\,
\varepsilon_x^2$ hold, the banks are positive and the fixed point exists for
$\varepsilon_x\le c_T=c_f\mathrm{Des}^{-1}$; this threshold contains the factor
$l^{-1}4^{-G(l)}(s'v')^2$ needed for positivity of the banks of shallow
carriers, because $\mathrm{Des}^{1/2}\ge l4^{G(l)}$ and $\mathrm{Des}^{1/3}\ge(s'v')^{-2}$.
The depth $s_{\rm t}(l)$ of Part~II (third element beyond $s_{\max}$) is used
there only for infinite base support; for finite base support the depths used
are the first element beyond $s_{\max}$ (tuning banks, peak pushes, closing
moves) and the second ($\sigma_2$, levers).  (ii) $\|U\|=\frac14$ instead of
$\frac1{16}$ changes the donor bound of Lemma~4.17(i) to $\mathrm{Des}^{1/3}
\Lambda/2$; Lemma~4.18 needs $N\mathrm{Des}^{1/3}\Lambda\le1$, true for main stages
$L\ge l_f$.  (iii) In Lemma~4.1, part (d) is (B3) and part (e) is (S2).
(iv) Part~II's $D^{U1'}$ is not transferred as a definition; its properties
used below are Lemma~\ref{lem:III-Wfin}, Theorem~\ref{thm:III-Afin} and (A1).
(v) The class and ray counts are the only window statements that are not of
the form treated by Lemma~3.22: their window factor has an exponent growing
with $L$, and Part~II handles them by $Q(w)\ge\bar D(w)^3$; for $T_{\rm fin}$
they are handled by Lemma~\ref{lem:III-count}.
\end{proof}

\begin{remark}[The exact route does not transfer]\label{rem:III-routeX}
The zero-value tuning of absorbers [II,~Lemma~5.51] needs $4e_1\le2^{-R}$ for
the perturbations $e_r:=\mu_{p_r}^2\theta_a/\nu$ of the binary coordinates;
Part~II obtains it from the doubly exponential base and its rule for $p_0$.
For $T_{\rm fin}$, with the first absorber of the cluster of a main stage $L$,
$2^{-R}$ is of the order $\mu_{p_0}^2c_{L+1}^2$ while $e_1$ is of the order
$\mu_{p_0+2}^2\theta_a/\nu$ with $\theta_a$ of the size of the residue to be
absorbed, which contains the true-shift correction and is not small compared
with $c_{L+1}$ in general.  So the intermediate-value sweep of
[II,~Lemma~5.51] can fail at the first used pair, and the route through exact
data with zero-value absorbers ([II,~Theorem~5.68]) is not available for
$T_{\rm fin}$ as defined.  Whether a different choice of the tuning coordinate
$p_0$ repairs this is not examined here.  The proof of Theorem~\ref{thm:III-MTVI}
tunes no absorber and is not affected.
\end{remark}

\subsection{Route V: the companion}\label{sec:III-routeV}

\subsubsection*{Standing hypotheses}
In this subsection and the next, $T_{\rm fin}$ is taken with option (a), or with
option (b) and only patterns $\kappa$ with $1\notin\mathcal N_T(\kappa)$; $N\ge1$;
$f\in S_{p_N^*}$ has finite base support $F$; $g\in\Cset(f)$ and $\rho\in(0,1)$.
The number $l_f=l_f(f,N,g,\rho)$ is a stage beyond which every $f$-constant is
absorbed; precisely, $l_f$ is so large that for every main stage $L\ge l_f$ and
every sub-window $w$ of $L$: the conclusions of Lemma~\ref{lem:III-alpha}(b) and
the hypotheses $L\ge l_f$ of Section~\ref{sec:III-generic} hold; every block
$m\le N$ has a carrier $\le L$ \textup{(}so the pattern of $f$ at $w$ is a pattern
of level $L$ of the design, and $J:=|\Omega|+N\le2L$\textup{)}; rows within
$C_fT_{\rm lo}(w)^2$ of $f$ lie in the ball $B_f$ of Lemma~\ref{lem:III-ball};
$T_{\rm hi}(w)\le t_1$ (Theorem~\ref{thm:III-UE}); $T_{\rm lo}(w)\le T_*$ of
Proposition~\ref{prop:III-S2c} for the constant $C_0$ of
Lemma~\ref{lem:III-S2ctr}; and $C_fN\le4\mathrm{Des}(L)/u(w)$ for the finitely
many $f$-constants $C_f$ named below.  Each condition holds for large $L$
because $T_{\rm lo}(w)\le2^{-L^3}$, $\mathrm{Des}(L)\to\infty$, and every block is
served in infinitely many rounds.

Let $L\ge l_f$ be a main stage, $w$ a sub-window of $L$ that is clean for $f$
([II,~Definition~3.18]), and $a$ an activity class of $w$ for the data of
$(g,\rho)$ ([II,~Lemma~5.53]: the class of a scale is read off the
decomposition data of $g$ at $f$ and is fixed before the companion is built),
with pattern $\kappa$ (Definition~\ref{def:III-pattern}).  We write
$T:=T_{\rm lo}(w)$, $b:=b(w)$, $u:=u(w)$, $\mathrm{Des}:=\mathrm{Des}(L)$, $\eta':=
\eta'(w)$, and
\begin{gather*}
 \mu_{\min}(w):=\min\{\mu_{\kappa'}(\eta'/3):\kappa'\in\mathfrak P(L),\
 \mathcal N_T(\kappa')\ne\emptyset\}\quad(:=1\text{ if there is none}),\\
 g_{\min}(w):=\tfrac18\min\{c_\flat\eta',\mu_{\min}(w)\},
\end{gather*}
where $c_\flat\in(0,1]$ is a constant depending only on $f$ such that the
gaps produced in blocks of kind (KN) and (U) by
Proposition~\ref{prop:III-KNtr}(b) are $\ge c_\flat\eta'M_m$: the minimum of
the numerical constant $1$ for the inward pushes in (U)-blocks and the
$f$-constant $c_f$ of [II,~Proposition~5.56(b)] for the inactive
near-threshold carriers of (KN)-blocks.  By Lemma~\ref{lem:III-Wfin}(f), Remark~\ref{rem:III-QCwell}(iii)
and (QC4), $8b\le\mu_{\min}(w)$.  A \emph{class set} is a set $\mathcal S$ of
dyadic scales of $w$ that all have the class $a$.

\begin{remark}[Route decision]\label{rem:III-routedecision}
For each $(L,w,a)$ one companion is built, for every kind of class: without
active shift, one-signed, or mixed.  The violation-tolerant route needs from
the coarse system only three things: the coarse statuses, the kind-[3]
property of the near-threshold switching carriers with the true shift, and
the exactness of the cone at the realized ratios.  Proposition~\ref{prop:III-KNtr}
supplies all three, and it treats each block separately, whatever the signs
of the shifts of the other blocks; so mixed classes need no separate
treatment.  Absorbers are not tuned: they are ordinary fine carriers.
(Tuned zero-value absorbers cannot follow the levers of
Lemma~\ref{lem:III-DC}, whose moves are much larger than the absorbers'
value rooms; Remark~\ref{rem:III-abs}.)
\end{remark}

\begin{lemma}[The route-V companion]\label{lem:III-CMtr}
Under the standing hypotheses there is a first row $f^\#$ with finite base
support $F^\#\supseteq F$, built in the order: class $a$, Step~A, Step~B
below, with the following properties.  Let $\Omega^\#$ be the set of coarse
strict non-peaks of $f^\#$.
\begin{enumerate}
\item[(a)] \textup{(Exactness.)}  The realized active ratios $r''$ \textup{(}the
active ratios of $f^{(1a)}$, i.e.\ $u_l(\zh^\#)/\varkappa'_m$ with the block
scalars $\varkappa'_m$ of $f^{(1a)}$\textup{)} lie in $Z_\kappa$; every robust
rate object has modulus $\ge u/4$ at $r''$; the
$\ell_1$-Hoffman constant of the cone $\Gamma^\#(\kappa,a)$ at $r''$ is
$\le C(\mathrm{Des}/u)^C$, and the bound $A_{\max}$ of [II,~Proposition~5.54] is
$\le C_f(\mathrm{Des}/u)^C$.
\item[(b)] \textup{(Statuses.)}  Every coarse carrier of $f^\#$ is either a
peak with relative margin $\ge u/4$, or a strict non-peak with $\gap^\#\ge
g_{\min}(w)$.  In every block of kind (T), every carrier $d\in\mathcal N(m)$ has
$\varrho_d(f^\#)\le1-\mu/4$ with $\mu:=\mu_\kappa(\eta'/3)$, and every carrier of
$\Omega_m\setminus\Omega_a(m)$ has $|u_l(\zh^\#)|\le4c_{L+1}^2$.
\item[(c)] \textup{(Data.)}  For every scale $t$ of a class set $\mathcal S$
there are pieces $(g_t;b^\pm_t,\omega^\pm_t)$ at $f^\#$, with
$\omega^\pm_t$ supported in $\Omega^\#$, such that $p^*(g-g_t)\le K_{\rm data}t$
with $K_{\rm data}$ as in Lemma~\textup{\ref{lem:III-count}},
$\Gamma^\#_w(b^\pm_t,\omega^\pm_t)\le1+\eta_0/2$, and every carrier of $\Omega^\#$
is of kind \textup{[1$'$]}, \textup{[2]} \textup{(}with $\gamma_B:=\gamma(w):=
\min_mM_mu/4$\textup{)} or \textup{[3]} of Definition~\textup{\ref{def:III-window}},
relative to $\gap^\#$.
\item[(d)] \textup{(Violations.)}  The total violation of these pieces is
$\epsilon(t)\le C_f\mathrm{Des}\,c_{L+1}+C_fc_{L+1}^2/T$, and $b^\theta_t(j)=0$ at
every free coordinate $j$ with a violation.
\item[(e)] \textup{(Scrambling.)}  Every fine carrier of a negative active
block is a peak with margin $\ge q^\#_0\delta^\circ_k/2$, and for every negative
active block $m$, $\mathrm{Scr}^\#_m(x)\le C(x/q^\#_0)^2$ for $0<x\le x_0(w):=
c_fg_{\min}(w)\Phi_{\min}(L)$.
\item[(f)] \textup{(Cost.)}  $p^*(f^\#-f)\le C_fT^3\log(1/T)$.
\end{enumerate}
\end{lemma}

\begin{proof}
\emph{Step A} is Proposition~\ref{prop:III-KNtr} applied to $(f,w,a)$; it
yields $f^{(1a)}$.  Its joint fixed point holds every non-active carrier of
$\Omega_m$ of a block of kind (T) at value exactly $0$, i.e.\ \emph{every}
nearly neutral coarse strict non-peak of such a block other than the active
ones (Remark~\ref{rem:III-KN1} shows that holding fewer carriers is not
enough), and all its pulls lie in the common window of
Lemma~\ref{lem:III-window} with $\eta_*=\max(\|x\|_\infty,c_{L+1})$, so every
bank is positive and every pull has mass $\ge24\lambda_lc_{L+1}$.  No absorber
is tuned.

\emph{Step B} is step (2) of Lemma~\ref{lem:III-CMp} (the fine-structure step
of [II,~Definition~5.57(2)] and [II,~Lemma~5.60] without absorber tuning),
applied to $f^{(1a)}$; steps (1b), (1c) of [II,~Definition~5.57] do not occur,
because the joint fixed point of Step~A already restores every $\varkappa_m$ and
every active value.  First the G1 release: the parts
$S_{l''}\cap(s_{\rm far}(w),\infty)$ of the signature sets of the coarse carriers
that do not own them are released into $J_{\rm fine}:=\N\setminus(F^{(1a)}\cup
E_c(w)\cup\{\text{coordinates fixed in Step A}\})$.  Then the static-ownership
recursion on $J_{\rm fine}$: the owner candidates are (a) the active carriers
of $\Omega$, (b) the coarse peaks of the active blocks, and (d) every carrier
at a stage $>L$ of an active block, absorbers included, with its forced
coefficient $-\Delta''_m\lambda_kw(k)$; a coordinate gets the $z$-value of
[II,~Lemma~5.60] for its first owner (negative fine owners self-aligned,
positive fine owners anti-aligned, coarse owners by the sign of their
coefficient), and unowned coordinates get $z:=0$.  The result is $f^\#$.  The
class enters Step~B only through the signs of the shifts and of the active
components, which are constant on the class ([II,~Lemma~5.53]); no common
shift value is used (Lemma~\ref{lem:III-nocube}).

(a) is Proposition~\ref{prop:III-KNtr}(a) (with the robust bound $u/4$), and
Step~B does not change the active values; the change of the block scalars
in Step~B is carried by the true shift $\Delta''_m=\Delta'_m\varkappa'_m/
\varkappa^\#_m$ in (c), so exactness is used only at $r''$.  The Hoffman bound is
[II,~Lemma~5.1] with the minors $\ge u/4$, and the bound on $A_{\max}$ is
[II,~Proposition~5.54].

(b) At $f^{(1a)}$: Proposition~\ref{prop:III-KNtr}(b) gives $\varrho_d\le1-\mu/2$
for $d\in\mathcal N(m)$ in blocks of kind (T), held values $0$, gaps $\ge
c_\flat\eta'M$ in blocks of kind (KN) and (U), and robust statuses moved by
$\le C\mathrm{Des}\,\eta'\ll u$.  Step~B changes the block vector $\zeta_m$ only
through the fine carriers ($\ell_1$-mass $\le\sum_{l>L}\lambda_l\le c_{L+1}/3$)
and through the released far parts of non-owners, which change a coarse value
by $\le2\delta_l\sum_{s>s_{\rm far}}2^{-s}/n_l\le4c_{L+1}^2$ (as $2^{-s_{\rm far}}\le
c_{L+1}^2$); it changes neither the base part $a$ (hence not $e$) nor any
active value.  By the clamp equation, as in Lemma~\ref{lem:III-dcons}(c) with
$\Omega=\emptyset$ and the natural peak of Lemma~\ref{lem:III-ball}, the block
scalar $C_m$ moves by $\le C_fc_{L+1}$; since $\varrho_k=\mathsf v_kC_m/(\Phi_k^2M_m)$
for strict non-peaks, every relative position moves by $\le C_f\mathrm{Des}\,
c_{L+1}$, and $\varkappa$ moves by $\le C_fc_{L+1}$.  By (QC5) and
Lemma~\ref{lem:III-drift}, $\varrho_d(f^\#)\le1-\mu/4$ for $d\in\mathcal N_T$; the
held values end at $\le4c_{L+1}^2$, which changes the protected statuses only
by $O(D^2c_{L+1}^4)\ll\mu$.  As $M_m\ge\frac12$, $\gap^\#(d)=M_m(1-\varrho_d)\ge
\mu/8\ge g_{\min}(w)$; robust margins and the gaps $\ge c_\flat\eta'M$ change by
$\le C_f\mathrm{Des}\,c_{L+1}\le C_f\mathrm{Des}\,b^2\ll g_{\min}(w)$.

(c), (d), (e) are Lemma~\ref{lem:III-CMp}(a)--(c) for the companion built with
Step~A in place of step (1a) there.  The proof of Lemma~\ref{lem:III-CMp}, and of
the results [II,~Lemmas~5.58, 5.60--5.62, 5.64] it invokes, uses step (1a) only
through: (i) the coarse statuses
(for the scrambling threshold $x_0$, for the coefficients $|w|=M$ of coarse
peak owners, and for the fact that in Step~B only carriers at stages $>L$
change status); (ii) the kind-[3] property of the near-threshold carriers of
$\Omega^\#$; (iii) the exactness of $\Gamma^\#(\kappa,a)$ at the realized ratios
(projection with error $K_{\rm data}t$ and the Hoffman bound).  Item (i) is
(b).  Item (iii) is (a).  For item (ii): in blocks of kind (T) the
decomposition data satisfy (X4$^0$) up to $K_{\rm data}t$
(Lemma~\ref{lem:III-floorrow}(b), applied at the realized ratios $r''$, whose
block coordinates satisfy $|r''^{(m)}-r(f)^{(m)}|\le2\eta'$ in blocks of kind
(T) by Proposition~\ref{prop:III-KNtr}(a)), and every point of the cone is inward at
$f^\#$ with the \emph{true} shift $\Delta''_m=\Delta'_m\varkappa'_m/\varkappa^\#_m$
([II,~Lemma~5.58(a)]; its middle term vanishes because Step~B does not move
active values): for a near-threshold carrier $l$, using (X4$^0$) at the cone
point in the form $\mathrm{vs}_l\gamma_l+\Delta'_m\lambda_l\varrho^0_l(r'')\ge0$
(Lemma~\ref{lem:III-floorrow}(a)),
\begin{align*}
 \lambda_l\mathrm{vs}_l\Delta\omega^\#(l)&=\mathrm{vs}_l\gamma_l+\Delta''_m\lambda_l
 \varrho^\#_l\\
 &\ge-|\Delta''_m-\Delta'_m|\lambda_l\varrho^\#_l-|\Delta'_m|\lambda_l
 |\varrho^\#_l-\varrho^0_l(r'')|\ge-C_f\lambda_l(c_{L+1}+D^Cc_{L+1}^2),
\end{align*}
which is $\ge-3\lambda_l\gap^\#(l)/t$ since $\gap^\#\ge M\mu/4$ and $c_{L+1}\le
b^2\le\mu^2/64$; on the $+$ side, $\mathrm{vs}_l\omega^+(l)\le\gap(f)/t$ with
$\gap(f)\le M_mb\le M_m\mu/8\le1.5\gap^\#(l)$.  In blocks of kind (KN) the row
(X4) of [II,~Definition~5.52] with the true shift gives $\mathrm{vs}_l
\Delta\omega^\#(l)\ge-|\Delta'^*_m|\gap^\#/M-C_fc_{L+1}\ge-3\gap^\#/t$; in blocks of
kind (U), $\Delta''_m=0$ and $\mathrm{vs}_l\Delta\omega=\mathrm{vs}_l\gamma_l/\lambda_l\ge0$.
Kinds of the other carriers of $\Omega^\#$: active robust carriers are of kind
[2]; held carriers and inactive nearly neutral carriers have $\Delta\omega=
\Delta''_mw^\#\approx0$ and gap close to $M$, kind [2]; the lever carrier of a
(KN)-block is of kind [2] in negative blocks (an active negative block has no
inactive near-threshold carrier, by [I,~Lemma~8.8] and [I,~(16)]) and of
kind [3] in positive blocks; the carriers of (U)-blocks pushed inward by $\eta'$
have $\Delta\omega=0$ and $|\omega|\le2\gap/t$, kind [1$'$] (their gap may be
smaller than $t^2$; Theorem~\ref{thm:III-UE} uses only $|\omega|\le2\gap^0/t$ for
them).
For the representation and the violation count of (c), (d): on a contact
of $E_c(w)$ the row (X1) at the data point gives $zV_{\rm proj}\ge0$, so the
violation is $\le|V_{\rm fine}|/2$; on a free coordinate of $E_c(w)$,
$V_{\rm proj}=0$, $b^\pm=\pm V_{\rm fine}/2$ and $b^\theta=0$; the trace correction
$-\sum_m(\Delta''_m-\Delta'_m)P_m$ on $E_c(w)$ has $\ell_1$-mass $\le C_fc_{L+1}$
and is a tolerated violation (it is not absorbed, since no absorber is
tuned); on $J_{\rm fine}$, coarse owners and negative fine owners dominate the
later coefficients ([II,~Lemma~5.61], including owners in absorber blocks
through the least-entry bound of Lemma~\ref{lem:III-interface}(W6)), unowned
coordinates have $V=0$, and positive anti-aligned owners contribute
violations of total mass $\le C_f\sum_{l>L}\lambda_l$.  With $\sqrt{\Gamma^\#_w}$ a
seminorm and $\|V_{\rm fine}\|_1\le C_f(\mathrm{Des}\,c_{L+1}+c_{L+1}^2/t)$, the
representation error and $\Gamma^\#_w\le1+\eta_0/2$ follow as in
[II,~Proposition~5.54(e)].  For (e): a negative fine carrier $k$ (any class,
absorbers included) is self-aligned, so $|u_k(\zh^\#)|\ge\delta^\circ_k\gg\bar\vartheta
\Phi_k$ by (W2), and $k$ is a peak with margin $\ge q^\#_0\delta^\circ_k/2$; coarse
peaks have margin $\ge q^\#_0\bar\vartheta\Phi_{\min}u/2>x_0$ and coarse strict
non-peaks have $\gap\ge g_{\min}$ and $\Phi\gap\ge\Phi_{\min}g_{\min}>x_0$ by (b);
so only fine carriers with $\delta^\circ_k\le2x/q^\#_0$ enter $\mathrm{Scr}_m(x)$, with
$\Phi_k\le(\delta_kH_k)^2$, and the super-exponential decay gives the bound
([II,~Lemma~5.62]).

(f) Step~A costs $\le C_f\mathrm{Des}^3\eta'\log(1/\eta')\le C_fT^3\log(1/T)$
(Proposition~\ref{prop:III-KNtr}(c)).  Step~B changes the row only on
$J_{\rm fine}$, where the coarse carriers move by $\le4c_{L+1}^2$ and the fine
carriers have total weight $\le c_{L+1}/3$; by [II,~Lemma~2.12] its cost is
$\le C_fc_{L+1}\log(1/c_{L+1})\le C_fT^4$ (Lemma~\ref{lem:III-Wfin}(d)).
\end{proof}

\begin{remark}[On Step~B]\label{rem:III-stepB}
The proof of Corollary~\ref{cor:III-noB} uses of the companion only the coarse
statuses and the bound on $\mathcal V$ on $J_{\rm fine}$; with (b) above it
applies to Lemma~\ref{lem:III-CMtr}.  So Step~B may be omitted if (U5) is
replaced by (U5$'$) of Lemma~\ref{lem:III-noB}: under (W$_{\exp}$) the fine
carriers then enter only through their total weight, and fine-origin misfits
are tolerated violations of mass $O(\mathrm{Des}\,c_{L+1})$.  We keep Step~B in
the main line because that is the version whose verification was carried out
twice.
\end{remark}

\subsection{Route V: base rows, frozen block scalars, threshold}\label{subsec:III-BRPM}

We keep the standing hypotheses of Subsection~\ref{sec:III-routeV}, let
$\mathcal S$ be a class set of $(L,w,a)$, and let $f^\#$ and the pieces
$(g_t;b^\pm_t,\omega^\pm_t)$, $t\in\mathcal S$, be those of
Lemma~\ref{lem:III-CMtr}.

\begin{proposition}[The companion is a base row]\label{prop:III-BR}
The pieces $(g_t;b^\pm_t,\omega^\pm_t)_{t\in\mathcal S}$ form a window family at
$f_0:=f^\#$ in the sense of Definition~\textup{\ref{def:III-window}}, with window
parameter $T=T_{\rm lo}(w)$, $\Omega:=\Omega^\#$, an $f$-constant $A_0$, $A_2=22$,
$\gamma_B:=\gamma(w)=\min_mM_mu(w)/4$, an $f$-constant $C_\Delta$, the levers of
Lemma~\textup{\ref{lem:III-LV}} with the pull window $[r,2^{G}r]$, $G:=6G(L)$, and
window quantities
\[
\begin{gathered}
 n\le\log_2(1/T),\quad\gap_{\min}\ge g_{\min}(w),\quad x_0=x_0(w),\quad
 C_S\le C/q_0^2,\quad c_{\rm fine}\le c_{L+1}/3,\\
 c_{\rm tiny}+W_{\rm lev}+\overline W_{\rm pull}\le C_f\mathrm{Des}(L)^Cc_{L+1},\qquad
 \Lambda_{\rm lev}\le C_f\big(\mathrm{Des}(L)/(u(w)g_{\min}(w))\big)^C .
\end{gathered}
\]
Moreover $c_{\rm flat}(w)\ge c_fu(w)$ for the number $c_{\rm flat}$ of
Theorem~\textup{\ref{thm:III-UE}}.
\end{proposition}

\begin{proof}
The scales of $\mathcal S$ lie in $\mathcal W(w)\subseteq[T,t_1]$ ($L\ge l_f$),
and $f^\#$ lies in $B_f$ with $F\subseteq F^\#$ (Lemma~\ref{lem:III-CMtr}(f) and
Lemma~\ref{lem:III-ball}(a)).

(U1), (U2): Lemma~\ref{lem:III-CMtr}(c); the bound $t\|b^\pm_t\|_1\le A_0$ and
the constant $A_2=22$ are those of [II,~Proposition~4.24(iii)], which apply by
Proposition~\ref{prop:III-transfer}.

(U3): The coordinates of $E:=F^\#\setminus F$ are created in Step~A (tuning
banks and pulls of active carriers and of held carriers, peak-push banks,
lever moves in blocks of kind (KN), and the inward pushes of the inactive
near-threshold carriers in blocks of kind (KN) and (U), realized by pulls in
the common window); Step~B creates none, and no absorber
coordinate is a support coordinate.  The data there are
$b^\pm=\beta+(\chi,\chi-1)\mathcal V_{\rm proj}\pm\frac12\mathcal V_{\rm fine}$.  The
last term is an antisymmetric residue $r^\pm$, counted in the total violation.
The remaining part is contact-like at banks (by (X2) the switching component
of an active carrier has the sign $\epsilon_l$ of its bank, and the
coefficient of a coarse peak has the sign of its trace), and at a pull $j_p$
it satisfies $t|\tilde b^\pm(j_p)|\le(2A_2+C_\Delta)\lambda_lv_l(j_p)$
([II,~Proposition~4.24(iv)], whose moves have the same types as those of
Step~A), while the pull mass is $|a^\#(j_p)|\ge c\lambda_lv_l(j_p)$; so
$t|\tilde b^\pm|\le K_3|a^\#(j_p)|$ with an $f$-constant $K_3$.  At a pull of a held
carrier $\gamma_l=0$, so $\tilde b^\pm(j_p)=0$; the residue there is
$\le C_fc_{L+1}$ in absolute value, while the pull mass is
$\ge24\lambda_lc_{L+1}$ by the common window (Lemma~\ref{lem:III-window}), so
even without splitting off the residue, $t|b^\pm(j_p)|\le K_3|a^\#(j_p)|$ for
$t\le T_{\rm hi}(w)$.  The same argument applies at the pulls of the inward
pushes: the pushed carriers are inactive in the class ($\gamma_l=0$,
Lemma~\ref{lem:III-LV}(b)), so the data there are the residue $\le C_fc_{L+1}$,
while the pull mass is $\ge24\lambda_lc_{L+1}$.

(U4): Lemma~\ref{lem:III-CMtr}(d).  (U5): Lemma~\ref{lem:III-CMtr}(e).

(U6): Lemma~\ref{lem:III-LV} at the companion of Lemma~\ref{lem:III-CMtr}.  Its
proof uses: allowedness and the disjointness of the signature sets; that the
support of $f^\#$ meets $S_l\cap(s_{\max}(L),\sigma_2(L)]$ at most in its first
element; that the coordinates of $S_l\cap(s_{\max}(L),s_{\rm far}(w)]\setminus F^\#$
are contacts with $z^\#=\epsilon_l$ (closing move (C1) of [II,~Definition~4.16],
unchanged by Step~A outside its own coordinates and by the release); the
constancy of the signs on the class; (D-lev); and (W1).  The first point holds
because Step~A places every bank at the first element of $S_l$ beyond
$s_{\max}(L)$, and its pulls have $v_l(j_p)\in[\eta_*,2^{G(L)}\eta_*)$ with
$\eta_*\le C_f\mathrm{Des}(L)\eta'$, far below $v_l$ at the second element, which is $\ge
1/\mathrm{Des}(L)$ by (D-lev) while $\eta'\le T^4$.  So the bank coordinate $j(l)$ of
Lemma~\ref{lem:III-LV} is at most the second element, and (B4) gives
$\beta_{j(l)}\le\nu\,\mathrm{Des}(L)$.  For a held carrier $l$ of class G,
$\gamma_l=0$, so it has the zero-data lever of Lemma~\ref{lem:III-LV}(b); for a
held carrier of class R the coordinate $j(l)$ is free and carries no data.
Step~A uses a bounded number of coordinates of each $S_l$ (at most a Step-0
push, one tuning pull and one bank, a lever or inward push, and one peak
push).  The pull coordinates of Lemma~\ref{lem:III-DC} have $v_l(j_p)$ of the
order of $r\le\Pi$; we choose them among the unused elements of $S_l$ in a
$v$-window of ratio $2^{6G(L)}$, which contains at least six elements of
$S_l$ because consecutive elements differ by $\le G(L)$; the factor
$2^{6G(L)}\le\mathrm{Des}(L)$ enters only $C_2$.  As $r\ge s_1T^4\ge
e^{-1/(2T)}T^4\gg c_{L+1}\ge2^{-s_{\rm far}(w)/2}$, these coordinates lie below
$s_{\rm far}(w)$.  Carriers outside $\Omega$ that meet lever coordinates are at
stages $>L$: beyond $s_{\max}(L)=\max T(L)$ no coarse target meets $S_l$.

(U7): $|\Delta_{t,m}|\le C_\Delta$ by [II,~Lemma~2.23], as the shifts of the
class are bounded by an $f$-constant.

Window quantities.  $n\le n(w)\le\log_2(1/T)$ since $T=2^{-n(w)}T_{\rm hi}(w)$.
$\gap_{\min}\ge g_{\min}(w)$ and the value of $x_0$ by
Lemma~\ref{lem:III-CMtr}(b), (e); $C_S$ by Lemma~\ref{lem:III-CMtr}(e).  A carrier
that is neither in $\Omega^\#$ nor a peak with margin $\ge x_0$ is fine
in the sense of Definition~\ref{def:III-window}; coarse peaks have margin
$\ge q^\#_0\bar\vartheta\Phi_{\min}(L)u/4\ge x_0$, so fine carriers are at stages
$>L$ and $c_{\rm fine}\le\sum_{l>L}\lambda_l\le c_{L+1}/3$.  The bounds on
$c_{\rm tiny}$, $W_{\rm lev}$, $\overline W_{\rm pull}$ and $\Lambda_{\rm lev}$ are
Lemma~\ref{lem:III-LV}(c) with $\gap_{\min}\ge g_{\min}(w)$ and the extra factor
$2^{6G(L)}$.  Finally $c_{\rm flat}=\min(c^f_{\rm flat},\gamma(w)/(8\rho A_2))\ge
c_fu(w)$, since $M_m\ge\frac12$ on $B_f$.
\end{proof}

Let now the window conditions of Theorem~\ref{thm:III-UE} hold for this
family (they are verified in Lemma~\ref{lem:III-S2ctr}), and let $f'$ be the
$d$-consistent engineered approximant of Definition~\ref{def:III-approx} at
$f_0=f^\#$, with lever parameter $P^*$; primes denote its data.

\begin{lemma}[The block scalars are frozen]\label{lem:III-PM}
In this situation, for $L\ge l_f$:
\begin{enumerate}
\item[(i)] $\mathrm{nv}'_k=\mathrm{nv}^\#_k$ for every $k\in\Omega^\#$, in particular for
every carrier of $\mathcal N_T$ and every held carrier;
\item[(ii)] $|C'_m-C^\#_m|\le C_fc_{L+1}$ for every block $m$;
\item[(iii)] for every $k\in\Omega^\#$ of block $m$,
\[
 \frac{\varrho'_k}{\varrho^\#_k}=\frac{C'_m(1-C^\#_m)}{C^\#_m(1-C'_m)},
\]
hence $|\varrho'_k-\varrho^\#_k|\le C_fc_{L+1}$, $\gap'(k)\ge\gap^\#(k)/2$, and
$\varrho_d(f')\le1-\mu/8$ for every $d\in\mathcal N_T$;
\item[(iv)] for every lever parameter $P$ with $|P|_\infty\le r$
\textup{(}not only at $P^*$\textup{)} no carrier of $\Omega^\#$ and no coarse peak of
$f^\#$ changes status; along the cube the carriers of $\Omega^\#$ move by at
most $\delta_1+C_{\rm lev}r$ and reach $\mathrm{nv}^\#$ at $P^*$; coarse peaks
of $f^\#$ meet no lever coordinate and move only at second order, by less
than their margins; among the carriers outside $\Omega^\#$, only fine
carriers \textup{(}stages $>L$, total weight $\le c_{L+1}/3$\textup{)} meet lever
coordinates.
\end{enumerate}
\end{lemma}

\begin{proof}
(i) is the $d$-consistency of $f'$ with $f^\#$ on $\Omega^\#$
(Theorem~\ref{thm:III-UE}).

(ii) Fix $m$ and put $C:=C^\#_m$.  By the clamp formula ([I,~Lemma~7.14];
Lemma~\ref{lem:III-dcons}(c)), $C^\#_m$ and $C'_m$ are the roots of
$F(c;v)=1$, $F(c;v')=1$ with $F(c;v)=\sum_k\min(\Phi_k(1-c)/c,v_k)^2$ and
$v_k=|\mathrm{nv}_k|/\Phi_k$.  Apply Lemma~\ref{lem:III-dcons}(c) with the natural
peak $k^\natural_m$ of Lemma~\ref{lem:III-ball}, which is a peak with margin
$\ge\mu^\natural_m/2$ at every row of $B_f$, so that its term is
$\Phi^2((1-c)/c)^2$ on an interval $[C-r',C+r']$ with $r'$ an $f$-constant.  In
the sum of Lemma~\ref{lem:III-dcons}(c) the terms of $\Omega^\#$ vanish by (i),
and so do those of the coarse peaks of $f^\#$: their relative margins are
$\ge u/4$, while their values move by $\le K_w\Pi\ll u$
(Lemma~\ref{lem:III-UE1}(b)) and $|C'_m-C|\le K_w\Pi$
(Lemma~\ref{lem:III-UE1}(c)), so they belong to $P^{\rm rob}_m$.  Every remaining
term belongs to a carrier at a stage $>L$ and is at most
$2\Phi_k(1-C)C^{-1}|v'_k-v_k|\le2(1-C)C^{-1}(|\mathrm{nv}'_k|+|\mathrm{nv}_k|)$, and
$|\mathrm{nv}_k|=\lambda_k|u_k(\zh)|/A_m\le(1+\|U\|)\lambda_k/A_{\min}$ at both rows.
Summing, $|C'_m-C|\le C_f\sum_{l>L}\lambda_l\le C_fc_{L+1}$ (Lemma~\ref{lem:III-Wfin}(d));
$C$ and $1-C$ are bounded below by $f$-constants on $B_f$.

(iii) For a strict non-peak $k$ of block $m$, by [I,~Lemma~1.7],
\[
 \varrho_k=\mathsf v_k/\bar\vartheta=|\mathrm{nv}_k|C_m/(\Phi_k^2M_m),
\]
and $\mathrm{nv}'_k=\mathrm{nv}^\#_k$,
which gives the identity and $|\varrho'_k-\varrho^\#_k|\le C_fc_{L+1}$.  By
Lemma~\ref{lem:III-dcons}(a), $\gap'(k)\ge\gap^\#(k)/2$ as soon as $|C'_m-C|\le
C\gap^\#(k)/(2(1-\gap^\#(k)))$; this holds because $\gap^\#(k)\ge g_{\min}(w)$,
$b\le\min(\eta',\mu_{\min}/8)\le8g_{\min}(w)/c_\flat$, and so $C_fc_{L+1}\le C_fb^2
\le64C_fc_\flat^{-2}g_{\min}(w)^2\le C^\#_mg_{\min}(w)/4$ for $L\ge l_f$ (as
$g_{\min}(w)\to0$ and $C^\#_m$ is bounded below on $B_f$).  For $d\in
\mathcal N_T$, Lemma~\ref{lem:III-CMtr}(b) gives $\varrho^\#_d\le1-\mu/4$, and
$C_fc_{L+1}\le C_fb^2\le\mu/8$.

(iv) By Lemma~\ref{lem:III-DC}(ii) and step (5) of its proof, along
$|P|_\infty\le r$ the carriers of $\Omega^\#$ move by at most
$\delta_1+C_{\rm lev}r$, and they reach $\mathrm{nv}^\#$ at $P^*$ by (i); a coarse
peak meets no lever coordinate, so it moves only through the normalization,
at second order; and the radius $r_0$ of
Lemma~\ref{lem:III-DC} is, in value units, below $c\min_{l\in\Omega}\Phi_l\gap_l
A_{\min}/(mC_{\rm lev})$ and below the coarse peak margins divided by
$C_{\rm lev}$; Theorem~\ref{thm:III-UE} uses $r\le r_0$ (proof of
Lemma~\ref{lem:III-UE1}).  The lever coordinates are private: by (SEP) they meet
no other carrier of $\Omega^\#$, and by (U6) the other carriers meeting them
are at stages $>L$.
\end{proof}

\begin{remark}[Why the freezing matters]\label{rem:III-PM}
The margin $\mu=\mu_\kappa(\eta'/3)$ protecting the carriers of $\mathcal N_T$
is only known to be $\ge\mathrm{Des}^{-1}(\eta'/3)^{\mathrm{Des}}$; it cannot be compared
with the size $r$ of the lever moves or with $K_w\Pi$ in general.  It does not
have to be: on $\Omega^\#$ the levers act exactly (i), and the only quantity
through which the other carriers act on a relative position of $\Omega^\#$ is
the block scalar $C_m$, which (ii) freezes up to the weight $c_{L+1}\le
\mu^2/64$ of the fine carriers.  This requires the levers on \emph{all} coarse
strict non-peaks: a single coarse strict non-peak left untuned changes $C_m$
at first order in the lever size.  (In finite block models, with all of
$\Omega$ tuned, $|C'-C|$ was at the working precision, and with one robust
strict non-peak left untuned the protected carrier moved by $10^{-7}$ to
$10^{-5}$ and in several instances became a peak; this is evidence only.)
\end{remark}

\begin{lemma}[Threshold comparison at the companion]\label{lem:III-S2ctr}
There is an absolute constant $C_0$ such that, under the standing hypotheses,
for the window family of Proposition~\textup{\ref{prop:III-BR}} with $T=T_{\rm lo}(w)$:
\begin{enumerate}
\item[(i)] $g_{\min}(w)\ge T^{5\mathrm{Des}(L)+2}\ge2^{-6\log_2^2(1/T)}$ and $x_0(w)\ge
T^{5\mathrm{Des}(L)+3}$;
\item[(ii)] hypotheses \textup{(i)--(iii)} of Proposition~\textup{\ref{prop:III-S2c}}
hold with this $C_0$ and $c_*:=c_{L+1}$, once $N_w$ is chosen with
$\mathrm{tail}_W\le e^{-1/T}$; in particular \textup{(W-b)} holds;
\item[(iii)] consequently $s_1:=e^{-1/(2T)}$ satisfies \textup{(VT$'$)},
\textup{(FINE)} and \textup{(LATE)}, \textup{(W-a)} and \textup{(W-c)} hold, and
\[
 K_w\big(K_ws_1+c_{\rm fine}+\mathrm{tail}_W\big)+K_w\big(\Pi\log(e/\Pi)+s_1\big)\le
 e^{-1/(3T)} .
\]
\end{enumerate}
\end{lemma}

\begin{proof}
(i) We have $L\le n(w)$ and $\mathrm{Des}(L)\le Q(w)\le n(w)\le\log_2(1/T)$, because
$T=2^{-n(w)}T_{\rm hi}(w)$.  Hence $1/T\ge2^{n(w)}\ge3n(w)^2\ge3L\,\mathrm{Des}(L)$
and $\eta'/3=T^4/(3L\,\mathrm{Des})\ge T^5$.  By Remark~\ref{rem:III-QCwell}(iii)
every $\mu_{\kappa'}(\eta'/3)$ is $\ge\mathrm{Des}^{-1}(\eta'/3)^{\mathrm{Des}}\ge
T^{5\mathrm{Des}+1}$ (as $1/\mathrm{Des}\ge T$), and $c_\flat\eta'\ge c_\flat T^5\ge T^6$,
so $g_{\min}(w)\ge\frac18T^{5\mathrm{Des}+1}\ge T^{5\mathrm{Des}+2}$ for $T\le\frac18$;
with $\mathrm{Des}\le\log_2(1/T)$ this is $\ge2^{-6\log_2^2(1/T)}$ when $\log_2(1/T)\ge2$.
Finally $x_0=c_fg_{\min}\Phi_{\min}(L)$ and $\Phi_{\min}(L)\ge1/\mathrm{Des}(L)\ge T$.

(ii) By Proposition~\ref{prop:III-BR} and (i), with natural logarithms:
$n\le\log_2(1/T)$; $A_0/T\le C_f/T$; $\log(1/\gap_{\min})\le9\log^2(1/T)$ and
$\log(1/x_0)\le11\log^2(1/T)$ for $T$ small; $1/u(w)\le n(w)$
([II,~Proposition~3.14(a)]); $|\Omega|\le2L\le2n$; $C_S\le C/q_0^2$; and
$\Lambda_{\rm lev}\le C_f(\mathrm{Des}/(ug_{\min}))^C\le e^{C\log^2(1/T)}$.  Since $K_w$ is
bounded by an $f$-constant times a power with \emph{absolute} exponent of
$nA_0|\Omega|\Lambda_{\rm lev}C_S/(T\gap_{\min}x_0)$ (Lemma~\ref{lem:III-UE1}), all
quantities in hypothesis (i) of Proposition~\ref{prop:III-S2c} are $\le
e^{C_0\log^2(1/T)}$ for an absolute $C_0$ and $T$ small.  The design quantity
$\mathrm{Des}(L)$ enters only the base $T$ of the power $T^{5\mathrm{Des}+2}$, and
$\mathrm{Des}(L)\le\log_2(1/T)$ turns it into a factor of the exponent of
$\log^2$.  Hypothesis (ii): by Lemma~\ref{lem:III-CMtr}(d) and
Proposition~\ref{prop:III-BR}, $\epsilon\le C_f\mathrm{Des}\,c_{L+1}+C_fc_{L+1}^2/T$,
$c_{\rm fine}\le c_{L+1}$ and $W_{\rm lev}+\overline W_{\rm pull}\le C_f\mathrm{Des}^Cc_{L+1}$, all
$\le e^{C_0\log^2(1/T)}c_{L+1}$; and $c_{L+1}\le e^{-1/T}$ by
Lemma~\ref{lem:III-Wfin}(d).  (W-b): $K_J2^{G+3}|\Omega|c_{\rm tiny}/A_{\min}\le
e^{2C_0\log^2(1/T)}e^{-1/T}\le\frac18$; the matrix $Y$ is nonzero only at bank
coordinates of levers that carry a mass of the pre-row, of size $\le4\rho s_1
\max_t|b^\theta_t(j')|\le4s_1A_0/T$, and $v_l(j')\ge1/\mathrm{Des}(L)$, so
$K_0\|Y\|_\infty\le C_f\mathrm{Des}\,s_1/T\le\frac12$.  Hypothesis (iii) is the choice
of $N_w$.

(iii) Proposition~\ref{prop:III-S2c}, for $T\le T_*$, which holds for $L\ge l_f$.
\end{proof}

\begin{remark}[The cost of the non-explicit margin]\label{rem:III-S2ctr}
In Theorem~\ref{thm:III-Cmix} the gaps of the switching carriers are $\ge
c_fT^4/(L\,\mathrm{Des})$, and $K_w\le T^{-C'}$.  Here the protected carriers have
gaps of the order of the Puiseux margin, so $K_w$ is only $\le
e^{C_0\log^2(1/T)}$, and $s_{\rm late}\ge e^{-35C_0\log^2(1/T)}$.  This is still far
above $e^{-1/(2T)}$, and (W$_{\exp}$) makes the fine weight $c_{L+1}\le e^{-1/T}$
negligible against it.  The estimate is robust: an exponent of $K_w$ depending
on $\mathrm{Des}(L)$ or on $|\Omega|\le2L$ would give $e^{C\log^3(1/T)}$, still
dominated.  No change of the design is needed for the Puiseux margin.
\end{remark}

\begin{lemma}[No common shift is needed]\label{lem:III-nocube}
The scales of $\mathcal S$ in Lemma~\textup{\ref{lem:III-CMtr}} and in
Proposition~\textup{\ref{prop:III-BR}} need not lie in one cube of the normalization of
[II,~Proposition~5.54(c)].  For $t\in\mathcal S$ let $(\Delta'(t),\gamma(t))\in
\Gamma^\#(\kappa,a)$ be the projection of the decomposition data
\textup{(}[II,~Proposition~5.54(a),(b)], [II,~Lemma~5.53]\textup{)}, let $\Delta''(t)
:=\Delta'(t)\varkappa'/\varkappa^\#$ be its true shift, and build the pieces as in
Lemma~\textup{\ref{lem:III-CMp}} with $\mathcal V_{\rm proj}:=\mathcal L(\Delta''(t),
\gamma(t))$ on $E_c(w)\setminus F^\#$, without normalization increment.  Then all
conclusions of Lemma~\textup{\ref{lem:III-CMtr}} and
Proposition~\textup{\ref{prop:III-BR}} hold for every class set, and
$\sqrt{\Gamma^\#_w(b^-,\omega^-)}\le\sqrt{\Gamma^\#_w(B_-,\Theta_-)}+K_{\rm data}t$,
where $(B_\pm,\Theta_\pm)$ is the decomposition of $g$ at scale $t$.
\end{lemma}

\begin{proof}
We go through every place where the construction could use that all scales
carry the same shift.  (1) Step~B: every ownership rule uses only the signs
$\sigma_m=\sgn\Delta''_m$ and $\epsilon_l$, which are constant on the class
([II,~Lemma~5.53]; the sign of $\Delta'_m$ is constant on the class cone), and
the quantities $Y_k$, which are computed from coordinates fixed before $k$ and
do not involve the data; the dominance of [II,~Lemma~5.61] is stated for every
shift and every component of the class.  The only construction built for one
shift value is the Schauder completion [II,~Lemma~5.63], which is not used.
(2) The true shift is defined per scale; the estimates of [II,~Lemma~5.58] are
per scale.  (3) The kind-[3] verification in the proof of
Lemma~\ref{lem:III-CMtr} holds for every point of the cone, and the data are
points of the cone.  (4) Conditions (U1)--(U7) are conditions on each piece
separately; (U5) concerns the set of negative active blocks, which is fixed by
the class; Theorem~\ref{thm:III-UE} uses the common set $\Omega$ and the
maxima of the masses over the pieces, and Theorem~\ref{thm:III-Eeng} averages
per-piece bounds.  (5) Lemma~\ref{lem:III-PM} and Lemma~\ref{lem:III-S2ctr} do
not involve the data.  (6) On contacts of $E_c(w)$ the row (X1) at the point
$(\Delta''(t),\gamma(t))$ gives $z\mathcal V_{\rm proj}\ge0$, as it did at the cube
corner.  For the last estimate: without the increment the $-$ side differs
from the decomposition only by the projection and the zeroing of the inactive
components, of size $\le K_{\rm data}t$, and $\sqrt{\Gamma_w}$ is a seminorm.
\end{proof}

\begin{lemma}[Window count]\label{lem:III-count}
Assume the standing hypotheses.  Let $B:=4\mathrm{Des}(L)/u(w)$, and let
$J:=|\Omega|+N$; then $J\le2L$.
\begin{enumerate}
\item[(a)] The number of activity classes of $w$ is $D_{\rm cl}(w)\le(J+1)3^J\le
B^{L+1}$.
\item[(b)] The class thresholds of [II,~Lemma~5.53] are $A_i=K_*^i$ with
$K_*\le B^{C_*}$, $C_*$ absolute, and the data of
Lemma~\textup{\ref{lem:III-CMtr}(c)} satisfy $p^*(g-g_t)\le K_{\rm data}t$ with
$K_{\rm data}:=CK_*^{J+2}K\le B^{C_*(2L+2)+C_K+2}$, where $K=C_f(\mathrm{Des}/u)^{C_K}$
is the constant of [II,~Proposition~5.54].
\item[(c)] Some class $a$ has a class set $\mathcal S$ with $|\mathcal S|\ge
n(w)/D_{\rm cl}(w)$, and there are absolute constants $C_1$, $L_0$ such that for
$L\ge\max(l_f,L_0)$, with $c_{\rm flat}(w)\ge c_fu(w)$,
\[
 \frac{|\mathcal S|\,c_{\rm flat}(w)}{K_{\rm data}}\ge L2^{L^3}\quad\text{and}\quad
 K_{\rm data}T_{\rm hi}(w)\le\frac{2^{-L^3}}L ,\qquad
 \log_B\frac{D_{\rm cl}(w)K_{\rm data}}{c_{\rm flat}(w)}\le C_1(L+1).
\]
\end{enumerate}
\end{lemma}

\begin{proof}
(a) A class records, for each of the $J$ components, active or inactive and,
if active, a sign, together with an index $i\le J$ of the threshold scale
([II,~Lemma~5.53]); this gives $\le(J+1)3^J$.  As $B\ge16$, $9^L\le B^L$, and
$2L+1\le B$.  (b) The thresholds satisfy $K_*\ge\max(2\mathrm{Des},4C_f\mathrm{Des}\,J
H(L))$ with the Hoffman bound $H(L)\le C(\mathrm{Des}/u)^C$, so one may take
$K_*\le B^{C_*}$ for $L\ge l_f$ ($C_fJ\le B$).  Zeroing the inactive components
of the data at scale $t$ costs $\le JA_{i(t)}Kt\le JK_*^{J+1}Kt$, and the
projection onto the cone moves the data by $\le\frac12A_{i(t)+1}Kt$; this gives
$K_{\rm data}$.  (The exponent of $K_{\rm data}$ grows linearly in $L$; a bound of
the form $C_f(\mathrm{Des}/u)^C$ with absolute $C$ is not available.)  (c) The
$n(w)$ dyadic scales of $w$ are distributed over the $D_{\rm cl}(w)$ classes;
by Lemma~\ref{lem:III-nocube} a class set needs no further pigeonhole.  With
$c_{\rm flat}\ge c_fu\ge B^{-2}$ for $L\ge l_f$, (a) and (b) give the bound on
$\log_B$ with $C_1:=2C_*+C_K+6$.  By Definition~\ref{def:III-Tfinstage}(S6),
$n(w)\ge L2^{L^3}Q(w)=L2^{L^3}B^{\omega(L)+L^2+20}$, so
$|\mathcal S|c_{\rm flat}/K_{\rm data}\ge L2^{L^3}B^{L^2+20-C_1(L+1)}\ge L2^{L^3}$ for
$L\ge L_0(C_1)$, and $K_{\rm data}T_{\rm hi}(w)\le B^{C_1(L+1)}2^{-L^3}/(LQ(w))\le
2^{-L^3}/L$.
\end{proof}

\begin{remark}[Two ways to count]\label{rem:III-count}
The count of [II,~Theorem~5.68] pigeonholes over classes and over the cubes of
the normalization; the cube count $(A_{\max}/\varepsilon)^p$ carries the number $p$
of extreme rays of the cone, which depends on the row and is bounded only by
$P(L)=2^{R_\Gamma(L)}$, and the data constant $K_{\rm data}$ has an exponent
growing with $L$.  So the claim that the window factor is $C_{f,N}(\mathrm{Des}/u)^C$
with an absolute exponent is not correct, and $T_{\rm fin}$ does not contain
Part~II's factor $\bar D(w)^3$.  Two corrections work; they are only
sketched here (\emph{sketch}).  (A) Pigeonhole first
over the support of a basic representation (at most $N$ extreme rays, by
Carath\'eodory; at most $(P(L)+1)^N$ choices) and then over $\varepsilon$-cubes in
the corresponding $N$ coordinates; with $P(L)\le\mathrm{Des}(L)^{1/2}$ and
$(2L+1)9^L\le B$ (as $\mathrm{Des}\ge2^{6L^3}$), $D_{\rm cls}(w)\le B^{1+N(2C+3)}$, the
increment $\sum_J(a^{\max}_j-a_j(t))r_j$ lies in the cone with $\ell_1$-norm
$\le N\varepsilon$, and the exponent $\log_B(D_{\rm cls}^2K_{\rm data}/c_{\rm flat})$ is
$O(L+N)$, absorbed by the term $L^2$ in the exponent of $Q(w)$ (or by
$\omega(L)\ge2^L-1$, which holds because the design enumerates the orthant
patterns of all subsets of $[1,L]$).  (B) Add $\bar D_{\rm fin}(w)^3$, with
$\bar D_{\rm fin}(w):=(2L+1)3^{2L}(\mathrm{Des}(L)/u(w))^{C_\Gamma(P(L)+2L+2)}$, to $Q(w)$
at main stages, as in Part~II.  We use neither: by Lemma~\ref{lem:III-nocube}
the cube is not needed in route V, and Lemma~\ref{lem:III-count} counts
classes only; correction (A) shows that the result does not depend on this
simplification.
\end{remark}

\begin{theorem}[Every activity class is recovered]\label{thm:III-Ctr}
Let $T_{\rm fin}$ be taken with option \textup{(a)}, or with option \textup{(b)}.
Let $N\ge1$, $f\in S_{p_N^*}$ with finite base support, $g\in\Cset(f)$ and
$\rho\in(0,1)$.  Suppose that for infinitely many main stages $L$ there are a
clean sub-window $w$ of $L$ and an activity class $a$ of $w$ \textup{(}of any
kind: without active shift, one-signed, or mixed\textup{)} with a class set of
at least $n(w)/D_{\rm cl}(w)$ scales, whose pattern satisfies $1\notin\mathcal
N_T(\kappa)$ if option \textup{(b)} is used.  Then $(f,\rho g)\in\cl\NA((c_0,p_N),
\ell_2^2)$.
\end{theorem}

\begin{proof}
Enumerate the given pairs $(L_j,w_j)$ with $L_j\ge\max(l_f,L_0)$, $L_j\to\infty$,
and let $\mathcal S_j$ be the class sets.  Let $f_j:=f^\#$ be the companion of
Lemma~\ref{lem:III-CMtr} at $(L_j,w_j,a_j)$, with the pieces of
Lemma~\ref{lem:III-nocube}.  By Proposition~\ref{prop:III-BR} they form a window
family at $f_j$ with fixed $A_0$, $A_2=22$ and $\gamma_{B,j}:=\gamma(w_j)$, and
$c_{\rm flat}(j)\ge c_fu(w_j)$.  We apply Theorem~\ref{thm:III-Eeng} with
$T_j:=T_{\rm hi}(w_j)$, $T'_j:=T_{\rm lo}(w_j)$, $n'_j:=|\mathcal S_j|$,
$K_j:=K_{\rm data}(w_j)$ and $s_{1,j}:=e^{-1/(2T'_j)}$.
\begin{itemize}
\item[(E-c)] Lemma~\ref{lem:III-CMtr}(c) with Lemma~\ref{lem:III-count}(b).
\item[(E-d)] Lemma~\ref{lem:III-count}(c).
\item[(E-f)] Lemma~\ref{lem:III-S2ctr}: the window conditions hold, $s_{1,j}$
satisfies (VT$'$), (FINE), (LATE), and $K_{w,j}(K_{w,j}s_{1,j}+c_{{\rm fine},j}+
\mathrm{tail}_{W,j})\le e^{-1/(3T'_j)}\le K_j(T'_j)^2$.
\item[(E-e)] $\varepsilon_j=p^*(f_j-f)\le C_fT'^3_j\log(1/T'_j)$
(Lemma~\ref{lem:III-CMtr}(f)) and $K_{w,j}(\Pi_j\log(e/\Pi_j)+s_{1,j})\le
e^{-1/(3T'_j)}$, so (E-e) holds with $\theta_j:=C_fT'_j\log(1/T'_j)+
e^{-1/(3T'_j)}T'^{-2}_j$.  As $1/u(w_j)\le n(w_j)\le\log_2(1/T'_j)$,
$\theta_j/c_{\rm flat}(j)^2\le C_f\theta_j\log_2^2(1/T'_j)\to0$.
\end{itemize}
Lemma~\ref{lem:III-PM} guarantees that the statuses used at $f_j$ (strict
non-peaks of $\Omega^\#$ with $\gap'\ge\gap^\#/2$, the protected carriers of
$\mathcal N_T$ with $\varrho\le1-\mu/8$, coarse peaks) persist at the
approximants.  Theorem~\ref{thm:III-Eeng} gives $(f,\rho g)\in\cl\NA((c_0,p_N),
\ell_2^2)$.
\end{proof}

\subsection{Master Theorem VI and the density criterion}\label{subsec:III-MTVI}

For $N\ge2$ and the operator $T_{\rm fin}$, stage $1$ is the carrier $(k,m)=(1,2)$;
for a first row $f\in S_{p_N^*}$ we write $\varrho_1(f)$ for the relative
position at $f$ of the coordinate $k=1$ of block $2$ ([I,~Lemma~1.7]).

\begin{theorem}[Master Theorem VI]\label{thm:III-MTVI}
Let $p_N$ be the norm on $c_0$ built with $T_{\rm fin}$ and the block set
$I=\{1,\dots,N\}$.
\begin{enumerate}
\item[(a)] \textup{(Option (a): $\|T_{\rm fin}\|=c_1<1$, a Mart\'in-type norm.)}  For
every $N\ge1$, every $f\in S_{p_N^*}$ with finite base support lies in
$\Rec(p_N)$: $(f,g)\in\cl\NA((c_0,p_N),\ell_2^2)$ for every $g\in\Cset(f)$.
\item[(b)] \textup{(Option (b): $T_{\rm fin}$ admissible, $c_1=1$.)}
\begin{enumerate}
\item[(b1)] For $N=1$, every $f\in S_{p_1^*}$ with finite base support lies in
$\Rec(p_1)$.
\item[(b2)] For $N\ge2$, every $f\in S_{p_N^*}$ with finite base support and
$\varrho_1(f)\ne1$ lies in $\Rec(p_N)$.
\end{enumerate}
\end{enumerate}
\end{theorem}

\begin{proof}
Fix $N$, $f$, $g\in\Cset(f)$ and $\rho\in(0,1)$, and let $l_f$ be as in the
standing hypotheses of Subsection~\ref{sec:III-routeV}.  Under option (b) and
$N\ge2$ with $\varrho_1(f)\ne1$, enlarge $l_f$ so that $b(w)<|\varrho_1(f)-1|$ for
every sub-window $w$ of every main stage $L\ge l_f$; this is possible because
$b(w)\le T_{\rm lo}(w)^4\le2^{-4L^3}$.

\emph{Step 1 (a clean sub-window at every main stage).}  Let $L$ be a main stage.
The rate objects of level $L$, evaluated at $f$, are $\omega(L)$ real numbers:
the relative positions and margins and the determinantal and shift objects of
[II,~Definitions~3.6 and~3.7], and the rate objects of every pattern and class
of level $L$ of Definition~\ref{def:III-cone}, evaluated at the ratio vector
of $f$ (QC2).  The $M(L)=\omega(L)+1$ bands $(b(w),u(w))$ of the sub-windows of
$L$ are pairwise disjoint (Lemma~\ref{lem:III-interface}(D2)), and each number
lies in at most one of them; so some band contains none of them, and the
corresponding sub-window $w(L)$ is clean: every rate object is $\le b(w(L))$ or
$\ge u(w(L))$.

\emph{Step 2 (a class set).}  Let $L\ge l_f$.  Every dyadic scale of $w(L)$ has
exactly one activity class for the data of $(g,\rho)$ ([II,~Lemma~5.53]), of one
of the three kinds (without active shift, one-signed, mixed); the
pigeonhole over the $D_{\rm cl}(w(L))$ classes gives a class $a$ with a class set
of at least $n(w(L))/D_{\rm cl}(w(L))$ scales (Lemma~\ref{lem:III-count}).  Under
option (b): if $N=1$, block $2$ is absent from $p_1$, so stage $1$ belongs to no
pattern of $p_1$; if $N\ge2$ and $\varrho_1(f)\ne1$, stage $1$ is not
near-threshold at $f$ at any clean sub-window of a main stage $L\ge l_f$, so
$1\notin\mathcal N_T(\kappa)$ for the pattern $\kappa$ of $(f,w(L),a)$.

\emph{Step 3.}  The hypothesis of Theorem~\ref{thm:III-Ctr} holds at every main
stage $L\ge l_f$, so $(f,\rho g)\in\cl\NA((c_0,p_N),\ell_2^2)$.  As $\rho<1$ is
arbitrary and $(f,g)=\lim_{\rho\to1}(f,\rho g)$, $(f,g)\in\cl\NA$; as $g\in\Cset(f)$
is arbitrary, $f\in\Rec(p_N)$.  The operator $T_{\rm fin}$ does not depend on $N$
(Lemma~\ref{lem:III-Wfin}(c)), and $N$ was arbitrary.
\end{proof}

The order of choices in the proof is: $T_{\rm fin}$ (independent of $N$ and of
every row), then $N$, $f$, $(g,\rho)$, $l_f$, a main stage $L\ge l_f$, the clean
sub-window $w(L)$, the class $a$ (read off the decomposition data of $g$ at
$f$), Step~A, Step~B, the data, $N_w$, $s_1$, the pre-row, the levers and the
approximant.  The weights are fixed before $f$; this is legitimate because the
coincidence sets $B_\kappa$ contain no datum of a row
(Lemma~\ref{lem:III-GENfin}(b)).

\begin{theorem}[Density criterion]\label{thm:III-VI2}
Let $T_{\rm fin}$ be taken with option \textup{(a)}, or with option \textup{(b)} and
$N=1$.  For every $N\ge1$ \textup{(}resp.\ for $N=1$\textup{)} the following are
equivalent:
\begin{enumerate}
\item[(i)] $\NA((c_0,p_N),\ell_2^2)$ is norm dense in $\mathcal L((c_0,p_N),\ell_2^2)$;
\item[(ii)] \textup{(LSC-trunc)$_N$}: for every $f\in S_{p_N^*}$, $g\in\Cset(f)$,
$\rho<1$ and $\varepsilon>0$ there is $f'\in S_{p_N^*}$ with finite base support
such that $p_N^*(f'-f)<\varepsilon$ and $\dist_{p_N^*}(\rho g,\Cset(f'))<
\varepsilon$.
\end{enumerate}
\end{theorem}

\begin{proof}
For $h_1,h_2\in(c_0,p_N)^*$ the operator $(h_1,h_2):x\mapsto(h_1(x),h_2(x))$ into
$\ell_2^2$ has norm $\le p_N^*(h_1)+p_N^*(h_2)$.
(ii)$\Rightarrow$(i): let $f\in S_{p_N^*}$, $g\in\Cset(f)$, $\rho<1$, $\varepsilon>0$,
and let $f'$ be given by (ii), with $g'\in\Cset(f')$ and $p_N^*(g'-\rho g)<\varepsilon$
($\Cset(f')$ is closed).  By Theorem~\ref{thm:III-MTVI}, $f'\in\Rec(p_N)$, so
$(f',g')\in\cl\NA$, and $\|(f',g')-(f,\rho g)\|<2\varepsilon$.  Hence $(f,\rho g)\in
\cl\NA$ for every $\rho<1$, so $(f,g)\in\cl\NA$ and $f\in\Rec(p_N)$; thus
$\Rec(p_N)=S_{p_N^*}$, which is (i) by [I,~Proposition~1.13] and
[I,~Definition~1.15].
(i)$\Rightarrow$(ii): by [I,~Lemma~1.16], density gives norm-attaining first rows
$f_n\to f$ with $\rho g\in\Li_n\Cset(f_n)$; a norm-attaining $f_n\in S_{p_N^*}$ has
finite base support ([I,~Propositions~1.10(c) and~1.12]).
The results of [I] used here hold for $T_{\rm fin}$ by Lemma~\ref{lem:III-partI}.
\end{proof}

\begin{corollary}[Mart\'in-type norm $p$]\label{cor:III-VI3}
Let $T_{\rm fin}$ be taken with option \textup{(a)}.  If \textup{(LSC-trunc)$_N$}
holds for infinitely many $N$, then $\NA((c_0,p),\ell_2^2)$ is dense in
$\mathcal L((c_0,p),\ell_2^2)$ for the Mart\'in-type norm $p$ \textup{(}$I=\N$\textup{)}
built with $T_{\rm fin}$.
\end{corollary}

\begin{proof}
By Theorem~\ref{thm:III-VI2}, $\NA((c_0,p_N),\ell_2^2)$ is dense for infinitely many
$N$; the norms $p_N$ and $p$ are built with the same operator, since $T_{\rm fin}$
does not depend on $N$.  By [I,~Remark~1.4] (Mart\'in tails, [I,~Lemma~1.5],
whose proof uses only $|R_mx|_m\le m2^{-m}q(x)$, valid for $\|T\|\le1$),
$\NA((c_0,p),\ell_2^2)$ is dense.
\end{proof}

\begin{remark}[What remains open]\label{rem:III-MTVIopen}
(a) \emph{Infinite base support.}  (LSC-trunc)$_N$ is not proved here, for any
$N$; by Theorem~\ref{thm:III-VI2} it is equivalent to the density of
$\NA((c_0,p_N),\ell_2^2)$ for the norms built with $T_{\rm fin}$.  So the density
of $\NA((c_0,p_N),\ell_2^2)$ and of $\NA((c_0,p),\ell_2^2)$ remains open.
(b) \emph{Option (b), $N\ge2$, $\varrho_1(f)=1$.}  At such rows stage $1$ may be an
active near-threshold carrier of a block of kind (T) at clean sub-windows of
infinitely many main stages, and its weight $\Phi_1=2^{-3}$ is rational, so
Lemma~\ref{lem:III-GEN} does not apply to it.  Whether $2^{-3}$ can be a critical
value of the floor function $\psi_1$ on a stratum of $Z_\kappa$, identically in the
transcendental parameters, is not known; normalizing the operator of option
(a) does not help (Remark~\ref{rem:III-options}).
(c) \emph{Other operators.}  Theorem~\ref{thm:III-MTVI} is proved for the designed
operator $T_{\rm fin}$ only, not for an arbitrary admissible operator of
Mart\'in's Lemma~B.
No counterexample to density is claimed, and none of the verifications
produced one.
\end{remark}

\begin{remark}[Verification]\label{rem:III-cert}
The results of Sections~\ref{sec:III-generic} and~\ref{sec:III-UE} on which
Theorem~\ref{thm:III-MTVI} rests were checked by one referee each and then
re-derived in a second, independent verification organized in four clusters
(the route-V companion; the generic-weight chain of
Section~\ref{sec:III-generic}; the engineered approximant and the averaging
of Section~\ref{sec:III-UE}; the transfer to $T_{\rm fin}$ and the window
bookkeeping); the combination in this section was checked by two referees
and re-derived in the same verification.  The corrections introduced in the
second verification (the common pull window, the transfer interface of
Lemma~\ref{lem:III-interface} and Proposition~\ref{prop:III-transfer}, the
$s$-free coefficients, the gap criterion of Lemma~\ref{lem:III-dcons}(a), the
bounds of Lemmas~\ref{lem:III-DC} and~\ref{lem:III-UE1}(g)) were each made
and checked once; only the corrected window count was checked twice.
Proposition~\ref{prop:III-J}, Lemma~\ref{lem:III-noB},
Corollary~\ref{cor:III-noB} and Lemma~\ref{lem:III-nocube} were checked by
one referee.  The following corrections were found and are incorporated.
\begin{enumerate}
\item[(1)] \emph{Companion.}  Carriers pushed inward in blocks of kind (U) are
of kind [1$'$] (Lemma~\ref{lem:III-CMtr}(c)); the companion is built in the order class,
Step~A, Step~B; every non-active nearly neutral coarse strict non-peak of a
(T)-block is held exactly (Remark~\ref{rem:III-KN1}); all pulls of the joint
fixed point lie in one common window, without which a held carrier can
require a negative bank (Lemma~\ref{lem:III-window},
Remark~\ref{rem:III-window}); the Puiseux point is chosen with $F_\Phi\le
1-\frac34\mu$; the cost carries $\mathrm{Des}^3$; robust rate objects are only
$\ge u/4$ after the rescaling of (U)-blocks; the inequality ``$Mb\le\gap^\#/4$''
is replaced by $\gap(f)\le1.5\gap^\#$ (proof of Lemma~\ref{lem:III-CMtr}).
\item[(2)] \emph{Base row.}  At the supports created by the companion the
side condition (U3) holds up to antisymmetric fine residues, which are
tolerated violations; Step~A uses a bounded number of coordinates of each
$S_l$, and the pull coordinates of the levers are taken among the unused ones
in a window of ratio $2^{6G(L)}$ (Proposition~\ref{prop:III-BR}).
\item[(3)] \emph{Generic weights.}  The Lipschitz constant in $b(w)$ and in (QC4)
dominates the floor-status function and the floor statuses of all active
coarse carriers; the level data use $s$-free aggregated coefficients (no
forward reference to $c_{L+1}$); normalizations of the {\L}ojasiewicz constants
are fixed; the $\alpha$-mass identity (Lemma~\ref{lem:III-alpha}) supplies the
natural peak and $\sigma_{\rm rob}>0$.
\item[(4)] \emph{Approximant.}  The neighbourhood lemma includes the margin of
the natural peak and $|M''-M|$ (Lemma~\ref{lem:III-ball}); the bound on
$|F(0)|$ in Lemma~\ref{lem:III-DC} contains the second-order effect of the fixed
pulls and banks; the gap criterion is $|C'-C|\le C\gap/(2(1-\gap))$, not
$|C'-C|\le\gap_{\min}/4$ (Lemma~\ref{lem:III-dcons}(a)); the coefficients $c_i$ are
bounded through the Bregman identity (Lemma~\ref{lem:III-UE1}(g)); the averaging
theorem allows a window-dependent $\gamma_B$ (Theorem~\ref{thm:III-Eeng}).  The
exponents of $K_w$ were confirmed to be absolute.
\item[(5)] \emph{Transfer and bookkeeping.}  The patterns range over the blocks
of stages $\le l$, and $l_f$ exceeds the stage after which every block $\le N$
has a carrier; the inspection claim that every result of Part~II survives for
$T_{\rm fin}$ is replaced by Lemma~\ref{lem:III-interface} and
Proposition~\ref{prop:III-transfer} (with $s_{\max}\to\infty$ proved without
fresh-dominated stages, the bank efficiency proved for the base $\mu$, and the
complete list of lower-bound uses of weights, including the least absorber
entry); the window count is corrected (Lemma~\ref{lem:III-count},
Remark~\ref{rem:III-count}); Part~I is used under option (a) after reading the
proofs, not only by locating the equality clause of (T-a)
(Lemma~\ref{lem:III-partI}); the route with zero-value absorbers does not
transfer to $T_{\rm fin}$ and is not used (Remark~\ref{rem:III-routeX}).
\end{enumerate}
The proof given here uses Lemma~\ref{lem:III-nocube} (in
Lemma~\ref{lem:III-count}(c) and in the proof of Theorem~\ref{thm:III-Ctr}).
Without it, Theorem~\ref{thm:III-Ctr} and Theorem~\ref{thm:III-MTVI} hold with
sets of scales of one class and one cube, the data built at the cube corner,
and the count of Remark~\ref{rem:III-count}(A), which is only sketched here;
that version was verified twice in the sources (the count by two
independent checks, the order class--Step~A--cube--Step~B by two
independent checks), so Theorem~\ref{thm:III-MTVI} does not depend on
Lemma~\ref{lem:III-nocube}.  A further round of three referees read the
present text; their corrections (among them the moduli of rate objects,
the reduced matrix of the cone, the constant $C_{**}$, the a priori pull
weight $\overline W_{\rm pull}$, the pull range $[r_{\rm pull},r_0]$, the realized
ratios in Lemma~\ref{lem:III-CMtr}(a), and the chain in (W6) of
Lemma~\ref{lem:III-interface}) are incorporated and were checked once.  Finite
numerical models were used as sanity checks of several identities (the clamp
equation, the $\alpha$-mass identity, the bank equations of the tuning lemma);
they are evidence only.
\end{remark}
\section{Infinite base support: truncation and transport}\label{sec:III-inf}

By [II,~Proposition~6.48], density of $\NA((c_0,p_N),\ell_2^2)$ is
equivalent to Lemma~Z at the rows with finite base support together with
\textup{(LSC-trunc)}$_N$: for every $f\in S_{p_N^*}$, $g\in\Cset(f)$,
$\rho\in(0,1)$ and $\varepsilon>0$ there is a first row $f'$ with finite base
support, $p_N^*(f'-f)<\varepsilon$ and $\dist(\rho g,\Cset(f'))<\varepsilon$.
For the design of Section~\ref{sec:III-design} with option (a) the first half
is Theorem~\ref{thm:III-MTVI}; so for that design density for $p_N$ is
equivalent to \textup{(LSC-trunc)}$_N$.  From
Subsection~\ref{subsec:III-cushion} on, the operator is $T^\infty_{\rm fin}$
(Definition~\ref{def:III-designinf}); for it the first half of the
equivalence is expected but not verified (Remark~\ref{rem:III-infdep}).  This
section treats \textup{(LSC-trunc)}.  It proves four things.
\begin{enumerate}
\item[(1)] Truncating the base support is asymptotically free at finitely
many scales fixed beforehand (its cost $2\Lambda_M+\varepsilon''_M$ tends to
$0$; lowering transfer, Lemma~\ref{lem:III-LT}); inside the
violation-tolerant engineered approximant of Section~\ref{sec:III-UE} the
truncated support becomes a set of contacts whose two-sided use is a
tolerated violation.  Consequently \textup{(LSC-trunc)} at $(f,g,\rho)$
follows from window families with \emph{side cushion domination} and exact
$d$-consistency at base rows of arbitrary support
(Theorems~\ref{thm:III-UEinf} and~\ref{thm:III-Eenginf}).
\item[(2)] For the design, at deep-raised rows the route-V data of
Section~\ref{sec:III-routeV} are cushion dominated, on the whole support of
the companion, with the absolute constant $C_+=5$
(Proposition~\ref{prop:III-cushion}).
\item[(3)] After a cheap preparation of the base row, the $d$-consistency
system is uniformly solvable at every row (Theorem~\ref{thm:III-SysDC}),
while Step~A of the companion construction is solvable whenever a cone
condition holds; the obstructions to the cone condition have a rigid
structure (Lemma~\ref{lem:III-cone}, Theorem~\ref{thm:III-Aobs},
Corollary~\ref{cor:III-Asuff}).
\item[(4)] The transport theorem (Theorem~\ref{thm:III-transport}):
\textup{(LSC-trunc)}$_N$ holds at every $(f,g,\rho)$ satisfying a joint
push condition \textup{(H1$^*$)} and the cone condition \textup{(H2)} at
infinitely many main stages.  Rows not covered satisfy, at almost all main
stages, one of the conditions \textup{(W$^*$)}, \textup{(A$^*$/U$^*$)} of
Subsection~\ref{subsec:III-residual}; these are
necessary conditions for not being covered, nothing more.
\end{enumerate}
\textup{(LSC-trunc)}$_N$ for every $f$, and hence density of
$\NA((c_0,p_N),\ell_2^2)$ and of $\NA((c_0,p),\ell_2^2)$, is not proved here.

\subsection{Contact truncations and lowering transfer}\label{subsec:III-trunc}

\subsubsection*{Conventions}
Throughout this section $I$ is finite, $p=p_N$, and the base is diagonal,
$Uk_s=\mu^\star_se_s$ with $\mu^\star_s>0$ \emph{non-increasing} and
$(k_s)$ orthonormal.  A first row $f\in S_{p^*}$ has forced data
$(\xi,q_0,a,w,F,z,\zh,e,\nu)$ [I,~Proposition~1.9]: $f=a+L^*w$,
$q^*(a)=\|a\|_1+\nu=1$, $\nu=\|U^*a\|_H$, $e=U^*a/\nu$, $F=\supp a$
\emph{arbitrary} (finite or infinite), $z=\sgn a$ on $F$, $\|z\|_\infty\le1$,
$\zh=z+Ue$.  For $A\in\ell_1\setminus\{0\}$ and an admissible $z$ ($z=\sgn A$
on $\supp A$) we write $f[A,z]$ for the first row with base vector
$A/q^*(A)$ and sign vector $z$, as in Section~\ref{sec:III-UE}.  Put
$s_j:=\sgn a_j$ for $j\in F$, $\alpha(M):=\sum_{j\in F,\,j>M}|a_j|$,
$a_{>M}:=a1_{(M,\infty)}$, $a_{\le M}:=a1_{[1,M]}$ and $s(t):=(1+t^2)^{1/2}$.

\begin{lemma}[Contact truncations]\label{lem:III-ctrunc}
Let $f\in S_{p^*}$ and $M$ with $a_{\le M}\ne0$.  Put $c_M:=q^*(a_{\le M})$,
$a^M:=a_{\le M}/c_M$ and $f^M:=f[a_{\le M},z]$, with forced data
$(a^M,z,e^M,\zh^M,w^M,\dots)$.  The coordinates $j\in F$, $j>M$, are contacts
of $f^M$ with $z_j=s_j$, and:
\begin{enumerate}
\item[(a)] $|1-c_M|\le\alpha(M)+\|U^*a_{>M}\|$;
\item[(b)] $\|e^M-e\|\le2\|U^*a_{>M}\|/\nu$ and $\zh^M-\zh=U(e^M-e)$;
\item[(c)] $f^M\to f$ in norm as $M\to\infty$, and all forced data converge as
in [I,~Proposition~1.11];
\item[(d)] $\|U^*a_{>M}\|\le\mu^\star_{M+1}\alpha(M)$.
\end{enumerate}
\end{lemma}

\begin{proof}
$z$ is admissible for $a_{\le M}$, and $z_j=s_j$ with $a^M_j=0$ for $j\in F$,
$j>M$.  (a) $1=\|a_{\le M}\|_1+\alpha(M)+\|U^*a\|$ and $c_M=\|a_{\le M}\|_1+
\|U^*a_{\le M}\|$; the two Hilbert norms differ by at most $\|U^*a_{>M}\|$.
(b) $\|x/\|x\|-y/\|y\|\|\le2\|x-y\|/\|y\|$ with $x=U^*a_{\le M}$, $y=U^*a$;
since $z$ is unchanged, $\zh^M-\zh=U(e^M-e)$.  (c) By (a), $a^M\to a$ in
$\ell_1$.  The proof of the implication (ii)$\Rightarrow$(i) of
[I,~Proposition~1.12] uses only $\ell_1$-convergence of the base vectors and
coordinatewise convergence of the sign vectors, not norm attainment; with
the constant sign vector $z$ it gives $f^M\to f$, and [I,~Proposition~1.11]
gives the convergence of the forced data.  (d) $\|U^*a_{>M}\|^2=\sum_{j>M}
(\mu^\star_j)^2a_j^2\le(\mu^\star_{M+1})^2\alpha(M)^2$ because $\mu^\star$ is
non-increasing.
\end{proof}

\begin{lemma}[Lowering transfer]\label{lem:III-LT}
Let $f\in S_{p^*}$, $M$ as in Lemma~\textup{\ref{lem:III-ctrunc}}, $h\in X^*$,
$r\ne0$, and let $f+rh=A+L^*W$ with $q^*(A)\le\pi$ and $\|W\|_{V^*}\le\pi$.
Put $B:=(A-a)/r$ and
\[
 \Lambda_M(r,B):=\sum_{j\in F,\ j>M}\min\big(|a_j|,(-s_jrB_j)_+\big)\qquad
 \big(\le\min(\alpha(M),|r|\,\|B1_{(M,\infty)}\|_1)\big).
\]
\begin{enumerate}
\item[(a)] $q^*(a^M+rB)\le q^*(A)+2\Lambda_M(r,B)+2\|U^*a_{>M}\|$, and
\[
 p^*(f^M+rh)\le\pi+2\Lambda_M(r,B)+2\|U^*a_{>M}\|+p^*\big(L^*(w^M-w)\big).
\]
\item[(b)] \textup{(Partial lowering.)}  Let $\Delta a$ be supported in $F$ with
$s_j\Delta a_j\le0$ and $|\Delta a_j|\le|a_j|$, let $f^{\rm low}:=f[a+\Delta a,
z]$ with normer $w^{\rm low}$, and put $y_j:=(-s_jrB_j)_+$,
$\alpha'_j:=|a_j+\Delta a_j|$.  Then
\[
 p^*(f^{\rm low}+rh)\le\pi+2\sum_{j}\min\big(|\Delta a_j|,(y_j-\alpha'_j)_+
 \big)+2\|U^*\Delta a\|+p^*\big(L^*(w^{\rm low}-w)\big).
\]
\end{enumerate}
\end{lemma}

\begin{proof}
(a) $f^M+rh=A'+L^*W+L^*(w^M-w)$ with $A':=A+a^M-a=a^M+rB$.  By
[I,~Lemma~1.6] and the triangle inequality it suffices to show
$q^*(A')\le q^*(A)+2\Lambda_M+2\|U^*a_{>M}\|$.  Write
$A'=(a_{\le M}+rB)+(1/c_M-1)a_{\le M}$; the second term has $q^*$-norm
$|1-c_M|$.  \emph{$\ell_1$-part.}  At $j\le M$ and at $j\notin F$ the
coordinates of $a_{\le M}+rB$ and of $A=a+rB$ agree.  For $j\in F$, $j>M$,
put $\alpha:=|a_j|$, $x:=rB_j$, $y:=(-s_jx)_+$; checking the cases
$s_jx\ge0$, $0<y\le\alpha$ and $y>\alpha$ gives
\[
 |s_j\alpha+x|-|x|=\alpha-2\min(\alpha,y),
\]
i.e.\ $|x|=|A_j|-|a_j|+2\min(|a_j|,y)$.  Summing,
$\|a_{\le M}+rB\|_1=\|A\|_1-\alpha(M)+2\Lambda_M$.  \emph{Hilbert part.}
$\|U^*(a_{\le M}+rB)\|=\|U^*(A-a_{>M})\|\le\|U^*A\|+\|U^*a_{>M}\|$.  With
Lemma~\ref{lem:III-ctrunc}(a),
$q^*(A')\le q^*(A)-\alpha(M)+2\Lambda_M+\|U^*a_{>M}\|+\alpha(M)+
\|U^*a_{>M}\|$.

(b) Put $c:=q^*(a+\Delta a)$ and $A':=(a+\Delta a)/c+rB$.  Since the lowering
keeps signs, $\|a+\Delta a\|_1=\|a\|_1-\|\Delta a\|_1$, so $|1-c|\le\|\Delta
a\|_1+\|U^*\Delta a\|$.  At $j\in\supp\Delta a$, with $x:=rB_j$,
$\alpha:=|a_j|$ and $\alpha'=\alpha-|\Delta a_j|$, the identity above at
$\alpha$ and at $\alpha'$ gives
$|s_j\alpha'+x|-|s_j\alpha+x|=-|\Delta a_j|+2\min(|\Delta a_j|,(y_j-\alpha'_j)_+)$.
Hence $\|a+\Delta a+rB\|_1=\|A\|_1-\|\Delta a\|_1+2\sum_j\min(|\Delta a_j|,
(y_j-\alpha'_j)_+)$, $\|U^*(a+\Delta a+rB)\|\le\|U^*A\|+\|U^*\Delta a\|$, and
the conclusion follows as in (a).
\end{proof}

The $\ell_1$-mass $\alpha(M)$ of the truncation cancels exactly, as the raise
mass does in the raise transfer [II,~Lemma~2.15].  What remains is the
\emph{cushion loss} $\Lambda_M$ (anti-sign usage of the decomposition on the
far support, capped at the cushion $|a_j|$), the Hilbert footprint
$2\|U^*a_{>M}\|\le2\mu^\star_{M+1}\alpha(M)$ and the normer change, which is
bounded by [II,~Lemma~2.12(iv)] in terms of $\zh^M-\zh=U(e^M-e)$.  In (b)
only usage beyond the \emph{new} cushion $\alpha'_j$ is lost.

\begin{remark}[The bound is sharp and scale limited]\label{rem:III-LTsharp}
(a) The factor $2$ in front of $\Lambda_M$ in the base inequality of (a)
cannot be lowered.  Let $j\in F$, $j>M$, $0<r\le|a_j|$, $B:=-s_je^*_j$ and
$A:=a+rB$.  Then $\Lambda_M(r,B)=r$, $q^*(a+rB)\le\|a\|_1-r+\nu+
\mu^\star_jr=1-r+\mu^\star_jr$ and $q^*(a^M+rB)\ge\|a^M\|_1+r+\|U^*a^M\|-
\mu^\star_jr=1+r-\mu^\star_jr$, so
$q^*(a^M+rB)-q^*(a+rB)\ge2(1-\mu^\star_j)\Lambda_M(r,B)$.

(b) The cushion loss is not uniformly small as $r\to0$.  Fix $M$ and $r_0>0$
and put $h:=-a_{>M}/r_0$.  For $0<r\le r_0$, $f+rh=(a+rh)+L^*w$ with
$q^*(a+rh)\le\|a\|_1-(r/r_0)\alpha(M)+\nu+(r/r_0)\|U^*a_{>M}\|\le1$ (as
$\|U^*a_{>M}\|\le\alpha(M)$), so $\pi=1$, while $\Lambda_M(r,h)=
(r/r_0)\alpha(M)$ is linear in $r$.  At $f^M$ the coordinates $j>M$ of $F$
are contacts with $z_j=s_j$, and the base excess of $rh$ there,
$\sum_{j>M}(|rh_j|-s_jrh_j)$ in the sense of [I,~Lemma~7.1(b)], equals
$2r\alpha(M)/r_0$.  Hence truncation is free only at finitely many scales
fixed before $M$ (Corollary~\ref{cor:III-LTwin}), and in the engineered
approximant below the truncation point is chosen after the smallest scale
$s_1$ of the construction.
\end{remark}

\begin{corollary}[Truncation at finitely many scales]\label{cor:III-LTwin}
Let $f\in S_{p^*}$, let $\mathcal S$ be a finite set of nonzero reals, and
for $r\in\mathcal S$ let $f+rh_r=(a+rB_r)+L^*W_r$ with $\max(q^*(a+rB_r),
\|W_r\|_{V^*})\le\pi_r$.  For every $\eta>0$ there is $M_0$ such that for all
$M\ge M_0$ and $r\in\mathcal S$
\[
 p^*(f^M+rh_r)\le\pi_r+\eta r^2+\varepsilon''_M,\qquad\varepsilon''_M:=
 2\|U^*a_{>M}\|+p^*\big(L^*(w^M-w)\big)\xrightarrow[M\to\infty]{}0 .
\]
\end{corollary}

\begin{proof}
$\Lambda_M(r,B_r)\le|r|\,\|B_r1_{(M,\infty)}\|_1\to0$ for each of the finitely
many $r$, since $B_r\in\ell_1$; apply Lemma~\ref{lem:III-LT}(a).
$\|U^*a_{>M}\|\to0$, and $p^*(L^*(w^M-w))\to0$ by
Lemma~\ref{lem:III-ctrunc}(c) and [I,~Proposition~1.11].
\end{proof}

\begin{remark}[Heuristic: why exact local data are needed at the approximant]
\label{rem:III-trunc-heur}
This remark is heuristic and is not used.  For $f'$ near $f$ with
$p^*(f'-f)=d>0$, scales $|r|\gg(d/(1-\rho^2))^{1/2}$ are covered by the slack of
$\rho g$, but for $|r|\ll d^{1/2}$ the inequality $p^*(f'+rg')\le s(r)$ is a
statement about the second-order structure of $p^*$ at $f'$ alone.  The
normer changes for every $f'\neq f$ (Lemma~\ref{lem:III-ctrunc}(b)), so
statuses and levels of deep carriers at $f'$ differ from those at $f$ at
arbitrarily small scales, and the mate must carry exact local data at $f'$.
This is why the truncation is carried out inside an engineered approximant
whose local data are controlled, rather than by approximating $f$ first.
\end{remark}

\subsection{The engineered approximant at arbitrary base support}
\label{subsec:III-UEinf}

We now show that finite base support of the approximant costs nothing in the
violation-tolerant route of Section~\ref{sec:III-UE}.  Fix $f\in S_{p^*}$
($F$ arbitrary), $\rho\in(0,1)$, $\eta_0$, $A_0,A_2\ge1$, $\gamma_B\in(0,1]$ and
$\delta$ as in Section~\ref{sec:III-UE}; $f$-constants are as there.

\begin{lemma}[Uniform neighbourhoods at arbitrary support]\label{lem:III-ballinf}
Let $f\in S_{p^*}$ have arbitrary base support, and fix $\gamma_m$,
$k^\natural_m$ and transfer data as in Lemma~\textup{\ref{lem:III-ball}}.  There
is $r_f(\eta_1)>0$ such that every $f''\in S_{p^*}$ with $p^*(f''-f)\le
r_f(\eta_1)$ satisfies \textup{(a)} of Lemma~\textup{\ref{lem:III-ball}} without
the clause ``$F\subseteq\supp a''$ with $|a''_j|\ge a_{\min}/2$ on $F$'', and
\textup{(b)} of Lemma~\textup{\ref{lem:III-ball}}.
\end{lemma}

\begin{proof}
The proof of Lemma~\ref{lem:III-ball} uses [I,~Proposition~1.11] (valid for
every $F$), the continuity of the strict threshold inequality of
[I,~Lemma~1.7] and [I,~Lemma~7.5], whose proof uses only
[I,~Proposition~1.11], the threshold lemma and dominated convergence in the
blocks.  The removed clause is the only place where $F$ enters.
\end{proof}

\begin{definition}[Window families at arbitrary support]\label{def:III-windowinf}
Let $C_+\ge1$.  A \emph{base row} is $f_0\in S_{p^*}$ with $p^*(f_0-f)\le
r_f(\eta_1)/2$ (Lemma~\ref{lem:III-ballinf}) and forced data
$(a_0,z_0,e_0,\nu_0,\zh_0,w^0,\dots)$, $F_0:=\supp a_0$ arbitrary, contact set
$K_0$, $s_j:=\sgn a_0(j)$ on $F_0$.  A \emph{window family with cushion
constant $C_+$} at $f_0$ consists of pieces $(g_i;b^\pm_i,\omega^\pm_i)$,
$i=1,\dots,n$, at scales $t_i\in[T,t_1]$ (Definition~\ref{def:III-pieces}),
supported in a finite set $\Omega$ of strict non-peaks of $f_0$, satisfying
\textup{(U1)}, \textup{(U2)}, \textup{(U5)}, \textup{(U7)} of
Definition~\ref{def:III-window} and:
\begin{enumerate}
\item[(U3$^{\inf}$)] \emph{side cushion domination}: for every $i$ and
$j\in F_0$ there is $r_i(j)\ge0$ with
\[
 (s_jb^+_i(j))_-\le\frac{C_+|a_0(j)|}{t_i}+r_i(j),\qquad
 (s_jb^-_i(j))_+\le\frac{C_+|a_0(j)|}{t_i}+r_i(j);
\]
\item[(U4$^{\inf}$)] $b^\theta_i(j)=0$ at every free coordinate $j\notin
F_0\cup K_0$ with $\mathrm{viol}_i(j)>0$, and the total violation
$\epsilon_i:=\sum_{j\notin F_0}\mathrm{viol}_i(j)+2\sum_{j\in F_0}r_i(j)$ plus
the masses $|b^\pm_i(j)|$ at lever coordinates of type \textup{(PR$\alpha$)}
below is at most $\epsilon$;
\item[(U6$^{\inf}$)] a \emph{lever system} (Definition~\ref{def:III-levsys}) is
given.
\end{enumerate}
Nothing is required of the sign-direction usage on $F_0$ or of $b^\theta$ on
$F_0$.  (Contact-like data on $F_0$, and data with $t_i|b^\pm_i(j)|\le
|a_0(j)|$, satisfy (U3$^{\inf}$) with $r_i(j)=0$ and $C_+=1$; antisymmetric
residues as in (U3) of Definition~\ref{def:III-window} are absorbed in
$r_i(j)$.)  Alternatively (U5) may be replaced by (U5$'$) of
Lemma~\ref{lem:III-noB}.
\end{definition}

The pre-row and the approximant are defined first; the lever system acts on
the pre-row.

\begin{definition}[The truncated $d$-consistent approximant]\label{def:III-approxinf}
Given a window family at $f_0$ with lever coordinates $\mathrm{Lev}$
(Definition~\ref{def:III-levsys}), choose in this order:
\begin{enumerate}
\item[(i)] $s_1\in(0,1]$;
\item[(ii)] $N_w\ge\max\mathrm{Lev}$ such that
\begin{gather*}
 \mathrm{tail}_W:=\max_i\big(\|b^\pm_i1_{(N_w,\infty)}\|_1+\|v_i1_{(N_w,\infty)}
 \|_1\big)\le\min\Big(T^3,\frac{c_{\rm late}\delta T}{A_0},\frac{\delta s_1}
 {64\rho}\Big),\\
 \|U^*a_01_{(N_w,\infty)}\|\le s_1T^4,\qquad \alpha_0(N_w):=\sum_{j\in F_0,\,
 j>N_w}|a_0(j)|\le\theta_{\rm tr},
\end{gather*}
for a prescribed truncation tolerance $\theta_{\rm tr}$ with
\[
 0<\theta_{\rm tr}\le\min\Big(s_1T^4,\ \frac{r_f(\eta_1)}{8(1+\|U\|)^2}\Big);
\]
\item[(iii)] the \emph{mass set} $W$ of contacts $j\in[1,N_w]\setminus F_0$
of $f_0$ with $\max_i|b^\theta_i(j)|>0$, with masses $\mathsf m_j:=4\rho s_1
\max_i|b^\theta_i(j)|$ of sign $z_0(j)$, and the \emph{raise set}
$W_F:=\{j\in F_0\cap[1,N_w]:|a_0(j)|<32\rho s_1\max_i|b^\theta_i(j)|\}$, where
$a_0(j)$ is raised, with its sign, to modulus $32\rho s_1\max_i|b^\theta_i(j)|$;
\item[(iv)] the pre-row $f_1:=f[A_1,z_1]$, where $A_1$ is $a_01_{[1,N_w]}$ with
the raises of $W_F$ and the masses of $W$, and $z_1:=z_01_{[1,N'']}$ ($N''$
from (vi)); put $\Pi_{\rm tr}:=2\|U^*a_01_{(N_w,\infty)}\|/\nu_0$,
$\Pi_X:=4\|U^*(A_1-a_01_{[1,N_w]})\|/\nu_0$ (neither depends on $N''$);
\item[(v)] the \emph{lever radius} $r:=4K_{\rm lev}(1+1/A_{\min})(\Pi_X+
\Pi_{\rm tr}+2T^4s_1)$;
\item[(vi)] $N''>N_w$ with the conditions (iv) of
Definition~\ref{def:III-approx};
\item[(vii)] $P^*$ by the lever system, $f':=f[A(P^*),z(P^*)]$ with data
$a',\hat x',c',w'_m,d'_m,\nu'$;
\item[(viii)] $g''_i:=b^\theta_i1_{[1,N_w]}+\sum_mR_m^*(\omega^\theta_{i,m}-
d'_m(\omega^\theta_{i,m})w'_m)$, $c_i:=g''_i(\hat x')$, $g'_i:=\rho(g''_i-c_ia')$.
\end{enumerate}
The perturbation size is $\Pi:=\Pi_X+\Pi_{\rm tr}+r+\mathfrak t(N'')^{1/2}$.
In the decompositions of the proof we use $\beta^\theta:=b^\theta1_{[1,N_w]}$
and $\beta^\pm:=b^\pm1_{[1,N_w]}\pm\frac12v1_{(N_w,\infty)}$, so that
$\beta^+-\beta^-=v$ and $\frac12(\beta^++\beta^-)=\beta^\theta$.
\end{definition}

\begin{definition}[Lever systems]\label{def:III-levsys}
Let $f_0$ be a base row and $\Omega$ a finite set of strict non-peaks of
$f_0$.  A \emph{lever system} consists of a finite set $\mathrm{Lev}$ of
coordinates, a finite index set $\Lambda_P$, numbers $r_0\in(0,1]$ and
$K_{\rm lev},C_2,C_{\rm lev}\ge1$, and, for every pre-row $f_1=f[A_1,z_1]$ of
Definition~\textup{\ref{def:III-approxinf}}, a continuous family
$P\mapsto(A(P),z(P))$, $P\in[-r_0,r_0]^{\Lambda_P}$, of admissible pairs with
$(A(0),z(0))=(A_1,z_1)$ and $A(P)-A_1$, $z(P)-z_1$ supported in
$\mathrm{Lev}$, such that:
\begin{enumerate}
\item[(DC)] if $r\le r_0$ and $K_{\rm lev}|F(0)|_\infty\le r/2$, where $F$ is
the map of Lemma~\ref{lem:III-DC} for this family, then $F$ has a zero $P^*$
with $|P^*|_\infty\le K_{\rm lev}r$; at $f(P^*)$ every $l\in\Omega$ is a
strict non-peak, the peaks of $f_0$ whose margin in value units exceeds
$C_{\rm lev}(K_{\rm lev}r+\delta_1)$ plus the change of the threshold remain
peaks with the same sign, the lever masses satisfy $\|X_{\rm lev}\|\le C_2r$,
a carrier meeting no lever coordinate moves by at most $|u_k(\hat x_1)-
u_k(\zh_0)|+\|e(P^*)-e_1\|$ and any other by at most that plus $C_{\rm lev}
K_{\rm lev}r$;
\item[(PR)] every $j\in\mathrm{Lev}$ is of one of two types:
\textup{(PR$\alpha$)} $b^\theta_i(j)=0$ for all $i$ (the masses $|b^\pm_i(j)|$
are counted in $\epsilon$); \textup{(PR$\beta$)} $j\in F_0$, and for
$|P|_\infty\le K_{\rm lev}r$ with $r\le r_0$, $A(P)(j)$ has the sign $s_j$ and
$|A(P)(j)|\ge|A_1(j)|/4$.
\end{enumerate}
The lever constants are $\Lambda_{\rm lev}:=\max(K_{\rm lev},C_2,C_{\rm lev},
1/r_0,1/A_{\min},1/C_{\min})$; $W_{\rm lev}$ is the total weight
$\sum\lambda_k$ of the \emph{fine} carriers $k$ meeting $\mathrm{Lev}$, and
$W_{\rm pull}:=0$.  Coarse carriers outside $\Omega$ that meet $\mathrm{Lev}$ (peaks
whose signature sets contain lever coordinates) are not counted in $W_{\rm lev}$;
by (DC) they move by at most $C_{\rm lev}K_{\rm lev}r$ more than the other
carriers.
\end{definition}

The levers of Lemma~\ref{lem:III-DC} with (SEP) form a lever system when no TU
lever is used; the system used below is that of
Theorem~\ref{thm:III-SysDC}.

\begin{theorem}[Engineered bound at arbitrary base support]\label{thm:III-UEinf}
Let $t_1$, $c^f_{\rm flat}$, $K_\sharp$, $\eta_1$, $c_{\rm late}$ be the
$f$-constants of Theorem~\textup{\ref{thm:III-UE}}, with condition \textup{(c1)}
of its proof deleted.  For $C_+\ge1$ put
\[
 c_{\rm flat}:=\min\Big(c^f_{\rm flat},\ \frac{\gamma_B}{8\rho A_2},\
 \frac1{32\rho C_+}\Big).
\]
Let $f_0$ be a base row of arbitrary support and let a window family with
cushion constant $C_+$ and scales in $[T,t_1]$ be given at $f_0$, with window
constant $K_w$ of Lemma~\textup{\ref{lem:III-UE1}} computed with the lever
constants of \textup{(U6$^{\inf}$)}.  Assume \textup{(W-a)} of
Theorem~\textup{\ref{thm:III-UE}} and let $s_1$ satisfy \textup{(VT$'$)},
\textup{(FINE)} and \textup{(LATE)} of Theorem~\textup{\ref{thm:III-UE}}, and
$K_w s_1\le r_0$.  Then for every choice of $N_w$, $N''$ as in
Definition~\textup{\ref{def:III-approxinf}} the approximant $f'$ is norm
attaining, has finite base support, is $d$-consistent with $f_0$ on $\Omega$,
$\Pi\le K_ws_1$,
\[
 p^*(f'-f_0)\le K_w\big(\Pi\log(e/\Pi)+s_1\big)+2(1+\|U\|)^2\alpha_0(N_w),
\]
and for every piece $i$
\begin{gather*}
 p^*(f'+\tau g'_i)\le1+\frac{\tau^2}2(1-\delta)\qquad(0<|\tau|\le c_{\rm flat}
 t_i),\\
 p^*(g'_i-\rho g_i)\le K_w\big(K_ws_1+c_{\rm fine}+\mathrm{tail}_W\big)/t_i .
\end{gather*}
\end{theorem}

\begin{proof}
The proof is that of Theorem~\ref{thm:III-UE} and Lemma~\ref{lem:III-UE1}
with the following modifications; everything not mentioned is verbatim.

\emph{(M0) Order of choices.}  $f\to$ $f$-constants $\to$ window family $\to
s_1\to N_w\to$ pre-row $\to r\to N''\to P^*$.  $s_1$ may precede $N_w$ because
$K_w$, $s_{\rm late}$, (VT$'$) and (FINE) do not involve $N_w$.

\emph{(M1) Neighbourhoods.}  Lemma~\ref{lem:III-ballinf} replaces
Lemma~\ref{lem:III-ball}; the removed clause was used only in condition
(c1), which is replaced by $32\rho c_{\rm flat}C_+\le1$.

\emph{(M2) Pre-row and levers.}  By Lemma~\ref{lem:III-ctrunc}(b) applied at
$f_0$ (that is, by the inequality $\|x/\|x\|-y/\|y\|\|\le2\|x-y\|/\|y\|$, valid
for all nonzero $x,y$; Lemma~\ref{lem:III-masses}(ii) is stated only for
finitely many mass coordinates), truncating $a_0$ beyond $N_w$ moves $e$ by at
most $\Pi_{\rm tr}$ and moves no $z$, and the raises and masses move $e$ by at
most $\Pi_X$; so
$\|e_1-e_0\|\le\Pi_X+\Pi_{\rm tr}$, and by Lemma~\ref{lem:III-masses}(iii)
every value moves by at most $\Pi_X+\Pi_{\rm tr}+t_k(N'')$.  As
$\|U^*(A_1-a_01_{[1,N_w]})\|\le32\rho s_1\mu^\star_1\sum_i\|b^\theta_i\|_1$,
$\Pi_X\le128\rho s_1\mu^\star_1nA_0/(T\nu_0)$.  Since $F_l(0)=\epsilon_l[(u_l
(\hat x_1)-u_l(\zh_0))-((A_m(f_1)-A_m(f_0))/A_m(f_0))u_l(\zh_0)]$ and
$|A_m(f_1)-A_m(f_0)|\le\sum_k\lambda_k|u_k(\hat x_1)-u_k(\zh_0)|$, we get
$|F(0)|_\infty\le(1+1/A_{\min})(\Pi_X+\Pi_{\rm tr}+T^4s_1)$, hence
$K_{\rm lev}|F(0)|_\infty\le r/4$; and $r\le K_ws_1\le r_0$.  So (DC)
applies.  In Lemma~\ref{lem:III-UE1}, $\Pi$ is enlarged by $\Pi_{\rm tr}$,
$K_J$ is replaced by $K_{\rm lev}$, and $W_{\rm pull}=0$.  The truncation also
changes $\nu$: $a'=A(P^*)/q^*(A(P^*))$ with $q^*(A(P^*))$ lowered by about
$\alpha_0(N_w)$, so the two-sided bound $|\nu'-\nu_0|\le K_w\Pi$ of
Lemma~\ref{lem:III-UE1}(c) is replaced by the lower bound $\nu'\ge\nu_0-
K_w\Pi$; this holds because the truncation lowers $q^*$ and changes
$\|U^*\cdot\|$ by at most $\|U^*a_01_{(N_w,\infty)}\|\le s_1T^4\le\Pi$, and
$\alpha_0(N_w)\le\theta_{\rm tr}\le s_1T^4\le\Pi$ changes $q^*$ by at most
$\Pi$.  Only this lower bound is used downstream (in
Lemma~\ref{lem:III-UE1}(f) and in the relative error of [I,~Lemma~3.3] in
Step~3), and the other clauses of Lemma~\ref{lem:III-UE1} hold verbatim,
with $K_w$ enlarged by an $f$-constant.  Finally $f'$ lies in $B_f$: as in the
step ``Consequences of (LATE)'' of the proof of Theorem~\ref{thm:III-UE},
$p^*(f'-f)\le r_f(\eta_1)/2+K_w(\Pi\log(e/\Pi)+s_1)+2(1+\|U\|)^2\alpha_0(N_w)$,
where the middle term is $\le r_f(\eta_1)/4$ by (LATE) and the last is
$\le r_f(\eta_1)/4$ by the choice of $\theta_{\rm tr}$; so Lemma~\ref{lem:III-ballinf}
holds at $f'$ and Step~4 applies.  A coarse carrier meeting lever
coordinates moves by at most $2(\Pi_X+\Pi_{\rm tr}+C_2r)+t_k(N'')+C_{\rm lev}
K_{\rm lev}r\le K_w\Pi$, so in Lemma~\ref{lem:III-UE1}(c)--(e) it is treated like
a carrier meeting no lever coordinate (its margin or gap exceeds $2K_*\Pi$ for
$K_w\Pi\le x_0$); only fine carriers meeting $\mathrm{Lev}$ enter through
$W_{\rm lev}$.  In the cost bound (a) the truncation
adds $q^*(a_01_{(N_w,\infty)})\le(1+\|U\|)\alpha_0(N_w)$ to $q^*(A(P^*)-a_0)$;
with the normalization and [II,~Lemma~2.12(iv)] this gives the extra term
$2(1+\|U\|)^2\alpha_0(N_w)$.  The base form (f) uses only $\|\beta^\diamond-
b^\diamond\|_1\le\mathrm{tail}_W$, which holds.

\emph{(M3) Step 3, base.}  We verify \eqref{eq:III-base} with the extra term
$\rho|\tau|\mathrm{tail}_W$ on the right of the first inequality.  Recall that on
$\supp a'$ the summand is $2(-\sgn(a'_j)\tau B_j-|a'_j|)_+$ with
$a'_j+\tau B_j=(1-\tau\rho c_i)a'_j+\tau\rho\beta_j$ and $1-\tau\rho c_i\ge
\frac12$, and that only the part of $\tau\rho\beta_j$ of sign opposite to
$a'_j$ can create excess.  New classes of coordinates:
\begin{enumerate}
\item[(1)] $j\in F_0\cap[1,N_w]$ not of type (PR$\alpha$).  Put $A''_j:=A_1(j)$
(the raised value if $j\in W_F$).  By (PR$\beta$) for lever coordinates and
trivially otherwise, $|A(P^*)(j)|\ge|A''_j|/4$ with the sign $s_j$; as
$q^*(A(P^*))\le2$, $|a'_j|\ge|A''_j|/8\ge|a_0(j)|/8$.  $\theta$-piece
($|\tau|\le s_1$): $|\tau\rho\beta^\theta_j|\le s_1\rho|b^\theta_i(j)|\le
|A''_j|/32\le|a'_j|/4$ by the definition of $W_F$: summand $0$.  $\pm$-pieces
($s_1<|\tau|\le c_{\rm flat}t_i$, $\diamond=\sgn\tau$): by
(U3$^{\inf}$) the anti-sign part of $\tau\rho\beta^\diamond_j$ is at most
$\rho c_{\rm flat}C_+|a_0(j)|+\rho|\tau|r_i(j)\le|a'_j|/4+\rho|\tau|r_i(j)$;
the first part creates no excess and the second costs at most
$2\rho|\tau|r_i(j)$, which is counted in $2\rho|\tau|\epsilon_i$.
\item[(2)] $j\in F_0\cap(N_w,N'']$: a contact of $f'$ with $z'_j=s_j$, $a'_j=0$,
$\beta^\theta_j=0$, $\beta^\pm_j=\pm\frac12v_j$.  The summand is $0$ for
$\theta$ and at most $\rho|\tau||v_j|$ for $\pm$; in total at most
$\rho|\tau|\mathrm{tail}_W\le\delta s_1|\tau|/64\le\delta\tau^2/64$ for
$|\tau|>s_1$: a tolerated violation.
\item[(3)] $j\in F_0$, $j>N''$: $z'_j=0$; inside the cut-off term
$\frac12\rho|\tau|V_{>N''}$ exactly as in [II,~Lemma~2.30].
\item[(4)] Lever coordinates of type (PR$\alpha$): $\beta^\theta_j=0$, so the
$\theta$ summand vanishes in every state of $j$ at $f'$, and the $\pm$ summand is
at most $2\rho|\tau||b^\pm_i(j)|$, counted in $\epsilon_i$.
\end{enumerate}
Coordinates $j\notin F_0$ are treated exactly as in the proof of
Theorem~\ref{thm:III-UE}.  Hence $\Exc'(\tau B^\pm)\le\frac12\rho|\tau|
V_{>N''}+2\rho|\tau|\epsilon_i+\rho|\tau|\mathrm{tail}_W$ and
$\Exc'(\tau B^\theta)=0$.

\emph{(M4) Step 5.}  The extra base term adds at most $\delta/32$ to the
bracket of Step~5 of the proof of Theorem~\ref{thm:III-UE}, and
$1-2\delta+\frac\delta8+\frac\delta8+\frac\delta8+\frac\delta{16}+
\frac\delta4+\frac\delta{32}\le1-\delta$.

\emph{(M5) The remaining steps.}  Steps 1, 2 and 4, Lemma~\ref{lem:III-UE1}(d),
(e), and the (U5$'$)-variant of Lemma~\ref{lem:III-noB} use the base only
through $q^*(B^\diamond)\le(1+\|U\|)\rho(2A_0/t_i+|c_i|)$, the base form
(f), and Lemma~\ref{lem:III-masses} at the truncated pre-row; they hold
verbatim.  $f'$ is norm attaining with finite base support because
$\supp A(P^*)\subseteq[1,N_w]\cup\mathrm{Lev}$ is finite and $z'$ vanishes
beyond $N''$ [I,~Proposition~1.10(c)].
\end{proof}

\begin{remark}\label{rem:III-UEinf}
(a) The far support $F_0\cap(N_w,\infty)$ of the base row becomes, inside the
approximant, a set of contacts; its two-sided use by the side data is a
violation of mass at most $\rho\,\mathrm{tail}_W\le\delta s_1/64$ per unit
$|\tau|$, which the $\pm$ regime tolerates, while the $\theta$ regime does not
see it ($\beta^\theta=0$ there).  The truncation is made after $s_1$, as
Remark~\ref{rem:III-LTsharp}(b) requires.  (b) On the kept support no lower
bound on $|a_0(j)|$ is used: (U3$^{\inf}$) and the radius $c_{\rm flat}t_i$
replace the finite-support condition $\rho c_{\rm flat}A_0\le a_{\min}/32$.
(c) $c_{\rm flat}$ depends on the window through $\gamma_B$ and $C_+$ only.
\end{remark}

\begin{theorem}[Windowed recovery at arbitrary base support, with raise transfer]
\label{thm:III-Eenginf}
Let $f\in S_{p^*}$ have arbitrary base support, $g\in\Cset(f)$, $\rho\in(0,1)$,
$\eta_0\in(0,2]$ with $\rho^2(1+\eta_0)\le(1+\rho^2)/2$, and $\delta$ as above.
Let $f_j$ be base rows of arbitrary support with $p^*(f_j-f)\to0$, and let
$\lambda_j\ge1$, $\lambda_j\to1$, $\varepsilon'_j\ge0$ satisfy
\begin{enumerate}
\item[(a)] $p^*(f_j+r\rho g)\le1+\big(s(\lambda_j\rho r)-1\big)/\lambda_j+
\varepsilon'_j$ for all $r\in\R$.
\end{enumerate}
Suppose that for every $j$ there are $T_j\in(0,t_1]$, $n_j\ge1$, a nonempty set
$\mathcal S_j\subseteq\{T_j2^{1-i}:1\le i\le n_j\}$ with $n'_j:=|\mathcal S_j|$,
$T'_j\in(0,\min\mathcal S_j]$, $\gamma_{B,j}\in(0,1]$, $C_{+,j}\ge1$, $K_j\ge1$,
and a window family with cushion constant $C_{+,j}$ at $f_j$ with scales
$\mathcal S_j$ (window parameter $T:=T'_j$), fixed $A_0,A_2$, and $\gamma_{B,j}$;
let $c_{\rm flat}(j):=\min(c^f_{\rm flat},\gamma_{B,j}/(8\rho A_2),
1/(32\rho C_{+,j}))$ and $K_{w,j}$ be the constants of
Theorem~\textup{\ref{thm:III-UEinf}} for this family.  Assume:
\begin{enumerate}
\item[(E-c)] $p^*(g-g_{j,t})\le K_jt$ for $t\in\mathcal S_j$;
\item[(E-d)] $K_jT_j/n'_j\to0$ and $n'_jc_{\rm flat}(j)/K_j\to\infty$;
\item[(E-f)] \textup{(W-a)} holds at $f_j$, and there is $s_{1,j}$ satisfying
\textup{(VT$'$)}, \textup{(FINE)}, \textup{(LATE)} and $K_{w,j}s_{1,j}\le r_{0,j}$
at $f_j$ such that, with $N_{w,j}$ chosen by
Definition~\textup{\ref{def:III-approxinf}} and $\Pi_j\le K_{w,j}s_{1,j}$,
\[
 K_{w,j}\big(K_{w,j}s_{1,j}+c_{{\rm fine},j}+\mathrm{tail}_{W,j}\big)\le
 K_j(T'_j)^2 ;
\]
\item[(E-e)] $\varepsilon'_j+K_{w,j}\big(\Pi_j\log(e/\Pi_j)+s_{1,j}\big)+
2(1+\|U\|)^2\alpha_{0,j}(N_{w,j})\le\theta_j(T'_j)^2$ with
$\theta_j/c_{\rm flat}(j)^2\to0$.
\end{enumerate}
Then $(f,\rho g)\in\cl\NA((c_0,p),\ell_2^2)$, along pairs $(f'_j,\bar g_j)$
whose first rows $f'_j$ are norm attaining with finite base support.  In
particular \textup{(LSC-trunc)} holds at $(f,g,\rho)$.
\end{theorem}

\begin{proof}
Fix $j$ large and write $n:=n'_j$, $K:=K_j$, $c:=c_{\rm flat}(j)$,
$\lambda:=\lambda_j$, $T':=T'_j$, $\mathcal S:=\mathcal S_j$.  Choose $N_{w,j}$
with $\alpha_{0,j}(N_{w,j})$ so small that (E-e) holds (possible since
$\theta_{\rm tr}$ in Definition~\ref{def:III-approxinf}(ii) is arbitrary).
Theorem~\ref{thm:III-UEinf} at $f_j$ with $s_1:=s_{1,j}$ gives a norm-attaining
$f'_j$ with finite base support and $g'_{j,t}$ ($t\in\mathcal S$) such that
\begin{enumerate}
\item[(A)] $p^*(f'_j+\tau g'_{j,t})\le1+\frac{\tau^2}2(1-\delta)$ for
$0<|\tau|\le ct$,
\end{enumerate}
$p^*(f'_j-f_j)+\varepsilon'_j\le\theta_jT'^2$ by (E-e), and
$p^*(g'_{j,t}-\rho g_{j,t})\le KT'^2/t\le Kt$ by (E-f), hence
$p^*(g'_{j,t}-\rho g)\le2Kt$ by (E-c).  By (a) and
[II,~Corollary~2.17], for every $r$,
\begin{enumerate}
\item[(B)] $p^*(f'_j+rg'_{j,t})\le p^*(f_j+r\rho g)+p^*(f'_j-f_j)+|r|\,
p^*(g'_{j,t}-\rho g)\le s(\rho r)+(\lambda-1)\min\big(\frac{\lambda^2r^2}2,1
\big)+\theta_jT'^2+2K|r|t$.
\end{enumerate}
Let $\bar g_j:=\frac1n\sum_{t\in\mathcal S}g'_{j,t}$; by convexity
$p^*(f'_j+r\bar g_j)\le\frac1n\sum_tp^*(f'_j+rg'_{j,t})$.  Fix $r_0\in(0,1]$
with $r_0^2\le\min(\delta,1-\rho^2)$.  For $j$ large,
\begin{gather*}
 \frac{(\lambda-1)\lambda^2}2+\frac{\theta_j}{c^2}+\frac{4K}{nc}\le
 \min\Big(\frac\delta4,\frac{1-\rho^2}6\Big),\\
 \theta_jT'^2+(\lambda-1)\le\frac{(1-\rho^2)r_0^2}{6},\qquad
 \frac{4KT_j}n\le\frac{(1-\rho^2)r_0}6,
\end{gather*}
by (E-d), (E-e), $\lambda_j\to1$ and $T'\le1$.

\emph{Case $|r|\le r_0$.}  Let $\mathcal S_r:=\{t\in\mathcal S:ct<|r|\}$;
since the scales are distinct dyadic numbers, $\sum_{t\in\mathcal S_r}t<
2|r|/c$.  If $\mathcal S_r=\emptyset$, every term obeys (A), and
$1+\frac{r^2}2(1-\delta)\le1+\frac{r^2}2-\frac{r^4}8\le s(r)$ as
$r^2\le\delta$.  Otherwise $|r|>cT'$, so $\theta_jT'^2\le\theta_jr^2/c^2$, and
\[
 p^*(f'_j+r\bar g_j)\le\max\Big\{1+\frac{r^2}2(1-\delta),\ s(\rho r)+
 \frac{(\lambda-1)\lambda^2r^2}2\Big\}+\frac{\theta_j}{c^2}r^2+
 \frac{2K|r|}n\cdot\frac{2|r|}c\le s(r),
\]
because $1+\frac{r^2}2(1-\delta)+\frac\delta4r^2\le1+\frac{r^2}2-\frac{r^4}8$
and $s(\rho r)+\frac{1-\rho^2}6r^2\le s(r)$ by [I,~Lemma~3.6(a)] ($|r|\le1$).

\emph{Case $|r|\ge r_0$.}  By (B) for all $t$, $p^*(f'_j+r\bar g_j)\le s(\rho r)+
(\lambda-1)+\theta_jT'^2+4K|r|T_j/n\le s(\rho r)+\frac{1-\rho^2}3r_0|r|\le
s(r)$, by [I,~Lemma~3.6(a)] and $\min(r^2,|r|)\ge r_0|r|$.

So $\bar g_j\in\Cset(f'_j)$, and $(f'_j,\bar g_j)$ attains its norm
[I,~Proposition~1.13(d)].  Finally $\|(f'_j,\bar g_j)-(f,\rho g)\|\le
p^*(f_j-f)+p^*(f'_j-f_j)+\frac1n\sum_t2Kt\le p^*(f_j-f)+\theta_jT'^2+
4KT_j/n\to0$.  As $f'_j$ has finite base support and
$\dist(\rho g,\Cset(f'_j))\le p^*(\rho g-\bar g_j)\to0$, (LSC-trunc) holds at
$(f,g,\rho)$.
\end{proof}

\begin{remark}[What the two theorems reduce \textup{(LSC-trunc)} to]
\label{rem:III-reduction}
(a) With $\lambda_j=1$ and $\varepsilon'_j=p^*(f_j-f)$, (a) holds trivially;
raised base rows ([II,~Lemma~2.15]) satisfy (a) with $\varepsilon'_j=
2\|U^*\Delta a_j\|_H+p^*(L^*(w^{(j)}-w))$ by [II,~Corollary~2.17], and
truncated rows $f^M$ (Lemma~\ref{lem:III-LT}) satisfy (a) for all $r$ with
$\lambda_j=1$ and $\varepsilon'_j=2\alpha(M)+\varepsilon''_M$, since
$\Lambda_M\le\alpha(M)$ (apply Lemma~\ref{lem:III-LT}(a) to an optimal
decomposition of $f+r\rho g$).  The better bound $2\max\Lambda_M+
\varepsilon''_M$ holds only at the finitely many scales of
Corollary~\ref{cor:III-LTwin} (Remark~\ref{rem:III-LTsharp}(b)).  (b) $c_{\rm flat}(j)$ may depend on the
window through $\gamma_{B,j}$ and $C_{+,j}$; the theorem needs exactly
$n'_jc_{\rm flat}(j)/K_j\to\infty$ and $\theta_j/c_{\rm flat}(j)^2\to0$.  (c)
Consequently \textup{(LSC-trunc)} at $(f,g,\rho)$ follows from the
\emph{existence} of window families with side cushion domination
(U3$^{\inf}$), a lever system (U6$^{\inf}$) giving exact $d$-consistency, and
(E-c)--(E-f), at base rows of arbitrary support converging to $f$ in the
sense of (a).  Producing them is the infinite-support exactification problem
treated in the rest of this section.
\end{remark}

\subsection{The designed operator: cushion domination at deep-raised rows}
\label{subsec:III-cushion}

From now on the operator is the design of Section~\ref{sec:III-design} over the
doubly exponential diagonal base $\mu^\star_s:=2^{-2^{2^s}}$, enlarged as in
Definition~\ref{def:III-designinf} below (all enlargements only increase
$\mathrm{Des}(L)$ and the rate scheme, Lemma~\ref{lem:III-designinf}).  We use
the notation of Subsection~\ref{subsec:III-gen-setting}, of
[II,~Subsection~5.6] and of the route-V companion of
Section~\ref{sec:III-routeV}: $N\ge1$, $f\in S_{p_N^*}$ with \emph{arbitrary}
base support $F$, $g\in\Cset(f)$, $\rho\in(0,1)$, $l_f=l_f(f,N,g,\rho)$, a main
stage $L\ge l_f$, a sub-window $w$ of $L$ clean for $f$, $T:=T_{\rm lo}(w)$,
$b:=b(w)$, $u:=u(w)$, $\eta':=T^4/(L\,\mathrm{Des}(L))$; $D(L)$ is the design
quantity of [II,~Subsection~5.6] with $\lambda_k\ge1/D(L)$ for coarse
carriers.  $T(L)$ is the union of the supports of the coarse targets and
$s_{\max}(L):=\max T(L)$.  For a coarse carrier $l$, $S_l$ is its signature
set and $v_l:=u_l$ on $S_l$; beyond $s_{\max}(L)$ the only coarse carrier
whose vector is nonzero on $S_l$ is $l$, and $v_l(s)=\delta_l2^{-s}/n_l>0$
there (proof of Lemma~\ref{lem:III-LV}).  The \emph{near coordinates} $N_l$ of
$l$ are the first six elements of $S_l\cap(s_{\max}(L),\infty)$, and
\[
 d(L):=\max_{l\le L}\max N_l,\qquad H(L):=\max_{l\le L}\max_{s\in N_l}
 \frac1{v_l(s)} .
\]
Sums over coarse and fine carriers satisfy $\sum_{k>L}\lambda_k\le Cc_{L+1}$,
and (W$_{\exp}$) gives $c_{L+1}\le e^{-1/T}$.

The \emph{box function} of level $L$ is $\mathrm{Bx}_L(j):=\sum_{l\le L}
\lambda_l|u_l(j)|$; $\sum_j\mathrm{Bx}_L(j)\le1$.

\subsubsection*{Route-V data}
For a scale $t$ of the window let $(B_\pm,\Theta_\pm)$ be the two-sided
decomposition at scale $t$ of the mate at the row considered, and
$\Delta\omega(t)$ the route-V switching on $\Omega$ (all coarse strict
non-peaks of the companion).  The switching vector
$V:=\sum_mR_m^*(\Delta\omega_m-\Delta_mw_m)$ (Definition~\ref{def:III-pieces};
$R_m^*e_k=\lambda_ku_k$) is $V=\sum_lc_lu_l$ over all carriers, with
$c_l=\lambda_l(\Delta\omega(l)-\Delta_mw_m(l))=:\gamma_l(t)$ for $l\in\Omega$
and $c_k=-\Delta_m\lambda_kw_m(k)$ for every other carrier $k$ of block $m$
(coarse peaks and fine carriers), since $\Delta\omega_m$ vanishes off
$\Omega$; $V_c:=\sum_{l\le L}c_lu_l$ and $V_f:=V-V_c$.  By (U7),
$|c_k|\le C_\Delta\lambda_k$ for $k\notin\Omega$, so $\|V_f\|_1\le
C_\Delta\sum_{k>L}\lambda_k\le Cc_{L+1}$.  Activity classes are those of
[II,~Lemma~5.53] read off the decomposition data, with the thresholds
$A_i:=K_*(w)^i$ defined below, and the projected class data
$(\Delta'(t),\gamma(t))$ lie in the class cone $\Gamma^\#(\varkappa,a)$
[II,~Definition~5.52].  The next lemma restates Lemma~\ref{lem:III-nocube} in the present setting and adds the count of class sets.

\begin{lemma}[Route V needs no shift normalization]\label{lem:III-nocubeinf}
In the route-V construction of Section~\textup{\ref{sec:III-routeV}} the cube
pigeonhole and the normalization increment of
[II,~Proposition~5.54(c)--(e)] may be omitted: if for every scale $t$ of a set
$\mathcal S$ of one activity class the data are built from the projected class
data $(\Delta'(t),\gamma(t))\in\Gamma^\#(a)$ without increment, then all
conclusions of the route-V companion lemma hold, with the $\Gamma_w$-bound
improved to $\Gamma^\#_w(b^-,\omega^-)^{1/2}\le\Gamma^\#_w(B_-,\Theta_-)^{1/2}+
K_{\rm data}t$, with $K_{\rm data}$ as in Proposition~\ref{prop:III-cushion},
and the pigeonhole over classes alone gives $|\mathcal S|\ge n(w)/D_{\rm cl}(w)$
with $D_{\rm cl}(w):=(J+1)3^J$, $J:=|\Omega|+N$.
\end{lemma}

\begin{proof}
We check every place where a common shift of all scales of $\mathcal S$
could be used.  (1) The anti-alignment and ownership rules of the route-V
companion use only the signs $\sgn\Delta''_m$ and $\epsilon_l$, which are
fixed by the class ([II,~Lemma~5.53]; the sign of $\Delta'_m$ is constant on
the class cone), and quantities computed from coordinates fixed before the
carrier; the dominance estimates of the completion are stated for every
shift and every $\gamma$ of the class.  (2) The relation between true and
reference shift ([II,~Lemma~5.58]) and the trace-correction bound are per
scale.  (3) Inwardness (kind [3]) holds at every point of the cone; the
increment was one more point of the cone.  (4) The conditions (U1)--(U7) of a
window family concern each piece separately, (U5) only through the set of
negative active blocks, which is fixed by the class; the approximant uses the
masses $\max_i|b^\theta_i|$ and one common $\Omega$; the averaging uses
convexity only.  (5) The status lemmas at the approximant do not involve the
data.  (6) Every later use of $|\mathcal S|$ is a lower bound.  Without the
increment the $-$ side differs from the decomposition's $-$ side only by the
projection and zeroing errors of [II,~Lemma~5.53]; the zeroing error is
$A_{i(t)}Kt$ with $A_i=K_*(w)^i\le K_*(w)^{J+1}$, which is why the constant is
$K_{\rm data}$ and not the projection constant $K$.  This gives the
$\Gamma_w$-bound since $\Gamma_w^{1/2}$ is a seminorm in the data.
\end{proof}

\begin{remark}[Why the increment must be dropped; heuristic]\label{rem:III-cube}
This remark is heuristic and not used.  The increment is a nonnegative
combination of extreme rays of the class cone with coefficients up to the
cube side, which is a fixed power of $u/\mathrm{Des}$ and not proportional to
$t$.  At an anti-sign support coordinate $j\in S_l\cap F$ below the raise
depth it may add anti-sign switching of order (cube side)$\cdot v_l(j)$, which
exceeds the cushion $|a_j|/t$ when $|a_j|$ is small compared with
$t\,v_l(j)$; Lemma~\ref{lem:III-nodeep} controls only the decomposition's own
switching.
\end{remark}

\begin{lemma}[Deep raise]\label{lem:III-deepraise}
Let $J>\max(d(L),s_{\rm t}(L))$, where $s_{\rm t}(L)$ is the depth of
[II,~Lemma~6.45], and put
\[
 R(J):=\{j\in F:j\ge J,\ |a_j|<8\mathrm{Bx}_L(j)\},\qquad
 \Delta a:=\sum_{j\in R(J)}s_j\big(8\mathrm{Bx}_L(j)-|a_j|\big)e^*_j,
\]
$\lambda_r:=q^*(a+\Delta a)$ and $f^r:=f[a+\Delta a,z]$.  Then:
\begin{enumerate}
\item[(a)] $\|\Delta a\|_1\le\psi_L(J):=8\sum_{j\ge J}\mathrm{Bx}_L(j)\to0$ as
$J\to\infty$, $\|U^*\Delta a\|\le\mu^\star_J\psi_L(J)\le8\mu^\star_J$,
$1\le\lambda_r\le1+2\|\Delta a\|_1$, and $|a^r_j|\ge4\mathrm{Bx}_L(j)$ for every
$j\in F$, $j\ge J$, once $\psi_L(J)\le\frac12$;
\item[(b)] for every $r\in\R$, $p^*(f^r+r\rho g)\le1+(s(\lambda_r\rho r)-1)/
\lambda_r+\varepsilon'_r$ with $\varepsilon'_r:=2\|U^*\Delta a\|+
p^*(L^*(w^r-w))\le C_f\mu^\star_J\log(e/\mu^\star_J)$;
\item[(c)] $F$, $z$, the contacts and the rooms are unchanged; every value
$u_k(\zh)$ moves by at most $2\|U^*\Delta a\|/\nu$.  Among the rate objects of
level $L$ \textup{(}[II,~Lemma~6.45]\textup{)}: the rooms and the objects
\textup{(O1)}, \textup{(O2)}, \textup{(O6)}, \textup{(O11)} are unchanged; the
support moduli \textup{(}\textup{(O10)} and the moduli $|a_s|$ of
\textup{(I3)}\textup{)} are divided by $\lambda_r\in[1,1+2\psi_L(J)]$; the minors
and cofactor polynomials of the enlarged class cones of \textup{(I3)} depend
only on design data and on the pattern, which is unchanged; and every other
rate object of level $L$ moves by at most $M_f(L)\|U^*\Delta a\|$.
\end{enumerate}
Let $J(w)$ be the least $J>\max(\bar d(L),s_{\rm t}(L))$
\textup{(}$\bar d(L)$ from Definition~\textup{\ref{def:III-designinf}}\textup{)}
with $C_f\mu^\star_J\log(e/\mu^\star_J)\le\theta_wT^2$ and $M_f(L)\mu^\star_J
\psi_L(J)\le b^2$, where $\theta_w:=T$ \textup{(}so the first condition reads
$C_f\mu^\star_J\log(e/\mu^\star_J)\le T^3$\textup{)}.  Then at $f^r$ every object of the rate scheme that is tiny at
$f$ is at most $2b$ and every robust one is at least $u/2$, and
$2^{J(w)}\le\mathrm{Des}(L)+2\log_2(Cn(w))$.
\end{lemma}

\begin{proof}
(a) The raise has the sign of $a$, so $\|a+\Delta a\|_1=\|a\|_1+\|\Delta a\|_1$
and $\|U^*(a+\Delta a)\|\ge\nu-\|U^*\Delta a\|$; as $\mu^\star$ is non-increasing,
$\|U^*\Delta a\|\le\mu^\star_J\|\Delta a\|_1\le\|\Delta a\|_1$, so $\lambda_r\ge1$,
and $\lambda_r\le1+2\|\Delta a\|_1$.  On $R(J)$, $|a_j+\Delta a_j|=8\mathrm{Bx}_L(j)$,
elsewhere on $F\cap[J,\infty)$, $|a_j|\ge8\mathrm{Bx}_L(j)$; divide by
$\lambda_r\le2$.  (b) is [II,~Corollary~2.17] for the raise; the normer
change is bounded by [II,~Lemma~2.12(iv)] with $\delta=U(e^r-e)$, $\|e^r-e\|
\le2\|U^*\Delta a\|/\nu$ (by the inequality in the proof of
Lemma~\ref{lem:III-ctrunc}(b); the raise set $R(J)$ may be infinite) and $c(\delta)\le
C\epsilon\log(e/\epsilon)$ for $\epsilon:=\|U\|\|e^r-e\|$ [II,~(1)].  (c) The
raise acts only on $F\cap[J,\infty)$ and keeps signs; the rest is
[II,~Lemma~6.45] (whose hypothesis $J>s_{\rm t}(L)$ holds), which distinguishes
the unchanged objects, the support moduli (divided by $\lambda_r$) and the
objects moved by $M_f(L)\|U^*\Delta a\|$; the moduli $|a_s|$ of (I3) are
support moduli, and the class-cone minors of (I3) do not involve the row.
The consequence for tiny and robust objects is the last part of
[II,~Lemma~6.45], since $\|\Delta a\|_1\le\frac12$ for $L\ge l_f$: a tiny
modulus stays $\le b$ and a robust one stays $\ge u/2$ because
$\lambda_r\le2$, and the moved objects change by at most
$M_f(L)\mu^\star_J\psi_L(J)\le b^2$.  For the bound on $J(w)$: both
conditions hold as soon as $\mu^\star_J\le\epsilon_w$ with $\log_2(1/\epsilon_w)
\le C\,\mathrm{Des}(L)\log_2(1/T)$, because $b$, $\theta_w$ and
$1/M_f(L)$ are at least design powers of $T$ (Section~\ref{sec:III-design});
since $\mu^\star_J=2^{-2^{2^J}}$ this holds when $2^J\ge\log_2(C\,\mathrm{Des}
(L)\log_2(1/T))$; with $\theta_w=T$, $\log_2(1/(\theta_wT^2))=3\log_2(1/T)$.
With $\log_2(1/T)\le3n(w)$ for $L$ large, $2^{\bar d(L)+1}\le\mathrm{Des}(L)$
(Definition~\ref{def:III-designinf}) and $2^{s_{\rm t}(L)+1}\le\mathrm{Des}(L)$
(the factor $B^\star(L)=2^{s_{\rm t}(L)}/(\delta_{\min}(\mu^\star_{s_{\rm t}})^2)$ of
$\mathrm{Des}(L)$ in Part~II), the
least admissible $J$ satisfies $2^{J(w)}\le\mathrm{Des}(L)+2\log_2(Cn(w))$.
\end{proof}

\subsubsection*{Classes with pinning thresholds}
For a target coordinate $j\in T(L)\cap F$ the clean sub-window gives
$|a_j|\le b$ (\emph{tiny support}) or $|a_j|\ge u$ (\emph{robust support}),
because the moduli $|a_s|$, $s\le\bar d(L)$, belong to the rate scheme
(Definition~\ref{def:III-designinf}).  The \emph{enlarged class cone} is
$\Gamma^\#(a)$ of [II,~Definition~5.52] with, for every tiny support
coordinate $j\in T(L)$, the contact row
\[
 \text{(X1$^F$)}\qquad s_j\Big(\sum_{l\in\Omega}\gamma_lu_l(j)-\sum_m\Delta'_m
 \sum_{k\in P^{\rm c}_m}\mathrm{vs}_k\lambda_ku_k(j)\Big)\ge0 .
\]
Its coefficients are non-weight design data, so the weight-freeness and
genericity arguments of Section~\ref{sec:III-generic} apply verbatim to the
enlarged cones, and their minors and cofactor polynomials are rate objects
(Definition~\ref{def:III-designinf}).  Let $C_q:=\frac83(4/q_0+8)$ be the
constant of [II,~Lemma~6.37] at rows of the ball ($q^r_0\ge q_0/2$), $H(L)$ the
Hoffman bound and $C_\Gamma$ the row bound of [II,~Lemma~5.53],
$K=C_f(\mathrm{Des}/u)^C$ the projection constant, $\delta_{\min}(L):=
\min_{l\le L}\delta_l$, and
\begin{gather*}
 K_P(w):=\frac{8C_q2^{J(w)}D(L)}{\delta_{\min}(L)},\\
 K_*(w):=\max\Big\{K_*,\ 6\big(H(L)(C_\Gamma J+1)+1\big),\ \big(K_P(w)/K
 \big)^{1/2}\Big\},
\end{gather*}
with $K_*$ of [II,~Lemma~5.53].  Classes are read off the decomposition data
at $f^r$ with the thresholds $A_i:=K_*(w)^i$.  Write $c_l$ for the
coefficients of $V_c$ built from the projected class data, and
$c^{\rm act}_l:=-\Delta\theta_l$ for the actual coefficients of the
decomposition $(B_\pm,\Theta_\pm)$ at $f^r$ ($\Delta B=-\sum_l\Delta\theta_lu_l$,
[I,~(15)]).

\begin{lemma}[Active coefficients track the decomposition]\label{lem:III-track}
For every $t\in\mathcal S$ and every coarse carrier $l$, either $c_l=0$, or
$|c_l-c^{\rm act}_l|\le|c_l|/5$ and $|c_l|\ge K_*(w)^2Kt/(4D(L))$.
\end{lemma}

\begin{proof}
For $l\in\Omega$, $c_l=\gamma_l(t)$.  An inactive component is set to $0$.  An
active one has $|\gamma_{{\rm dec},l}|\ge A_{i(t)+1}Kt=K_*(w)A_{i(t)}Kt$; zeroing
the inactive components and projecting onto the enlarged cone moves it by at
most $H(L)(1+C_\Gamma JA_{i(t)})Kt\le H(L)(C_\Gamma J+1)A_{i(t)}Kt\le
|\gamma_{{\rm dec},l}|/6$ (proof of [II,~Lemma~5.53]), so $|\gamma_l|\ge
\frac56|\gamma_{{\rm dec},l}|$ and $|\gamma_l-\gamma_{{\rm dec},l}|\le|\gamma_l|/5$;
also $|\gamma_l|\ge\frac56K_*(w)^2Kt$.  For a coarse peak $k$ of block $m$,
$c_k=-\Delta''_m\lambda_k\mathrm{vs}_k$ with $\Delta''_m:=M_m\Delta_m$
($w_m(k)=M_m\mathrm{vs}_k$ at a peak), and by [I,~Lemma~8.27] at $f^r$,
$c^{\rm act}_k=-\Delta'_{{\rm dec},m}\lambda_k\mathrm{vs}_k+e_k$ with $|e_k|\le
C_fD(L)t/u$ at peaks of relative margin $\ge u/2$.  The shift components obey
the same zeroing and projection bounds as the $\gamma$-components, and the
true shift differs from the reference shift only by the trace correction of
[II,~Lemma~5.58], which is $\le C_f\mathrm{Des}(L)c_{L+1}\ll Kt$; with
$|\Delta'_m|\ge K_*(w)^2Kt$ for an active shift, $\lambda_k\ge1/D(L)$,
$M_m\ge\frac12$, and $K_*(w)\ge K_*\ge C_fD(L)/u$, the same arithmetic gives the
claim.  Inactive shifts give $c_k=0$.
\end{proof}

\begin{lemma}[No deep anti-sign set at the raised row]\label{lem:III-nodeep}
Let $t\in\mathcal S$ and put, for $j\in T(L)$,
$\mathrm{tr}_j:=|\sum_m(\Delta''_m-\Delta'_m)\sum_{k\in P^{\rm c}_m}\mathrm{vs}_k
\lambda_ku_k(j)|$, and $\mathrm{tr}_j:=0$ otherwise.  Then
$\sum_j\mathrm{tr}_j\le C_f\mathrm{Des}(L)c_{L+1}+C_fc_{L+1}^2/t$, and for every
$j\in F$
\[
 s_jV_c(j)\ge-\frac{4|a^r_j|}t-\mathrm{tr}_j .
\]
\end{lemma}

\begin{proof}
The bound on $\sum\mathrm{tr}_j$ is the trace-correction estimate of
[II,~Lemma~5.58(b)], applied per scale.  \emph{Coefficient bound.}  For coarse
$l$, $|c_l|\le C_fD(L)\lambda_l$: for $l\in\Omega$, $c_l=\lambda_l(\Delta\omega(l)
-\Delta_mw_m(l))$ with $|\Delta\omega(l)|\le\|D_m\Delta\omega_m\|_2/\Phi_m(l)\le
C_fD(L)$ by (U1) and Proposition~\ref{prop:III-J}(a), and $|\Delta_m|\le
C_\Delta$; for peaks $|c_k|\le C_\Delta\lambda_k$.  Hence $|V_c(j)|\le C_fD(L)
\mathrm{Bx}_L(j)$.  Four cases.

(D) $j\ge J(w)$: $|a^r_j|\ge4\mathrm{Bx}_L(j)$ by Lemma~\ref{lem:III-deepraise}(a),
so $4|a^r_j|/t\ge16\mathrm{Bx}_L(j)/t\ge|V_c(j)|$ since $t\le T_{\rm hi}(w)\le
16/(C_fD(L))$.

(T) $j\in T(L)$.  If the support is robust, $|a^r_j|\ge u/2$ and $4|a^r_j|/t\ge
2u/T_{\rm hi}(w)\ge C_fD(L)\ge|V_c(j)|$, since $T_{\rm hi}(w)\le u/(C\,
\mathrm{Des}(L))$.  If it is tiny, the projected data satisfy (X1$^F$) exactly,
and $V_c(j)$ differs from the left side of (X1$^F$) (divided by $s_j$) by the
trace term, so $s_jV_c(j)\ge-\mathrm{tr}_j$.

(S) $j<J(w)$, $j\notin T(L)$, $j\in S_l$ for a coarse $l$.  Among coarse
carriers only $l$ lives at $j$, so $V_c(j)=c_lv_l(j)$.  Suppose $c_l\ne0$ and
$s_jc_lv_l(j)<-4|a^r_j|/t$.  By Lemma~\ref{lem:III-track}, $c^{\rm act}_l$ has the
sign of $c_l$ and $|c^{\rm act}_l|\ge\frac45|c_l|$; so with $\tau_l:=\epsilon_l
c^{\rm act}_l$ the coordinate $j$ is of the type opposite to the switching
direction and $|\tau_l|v_l(j)>\frac{16}5|a^r_j|/t$.  As $j$ is unraised and not
in $T(L)$, [II,~Lemma~6.37] at $f^r$ (its term $\sum|r_s|\le8t^2$ requires
$\sum_{l''>L}\lambda_{l''}\le t^3$, which holds) gives $|c^{\rm act}_l|\le
C_qt/v_l(j)\le\frac54C_q2^{J(w)}t/\delta_{\min}(L)$, since $v_l(j)=\delta_l2^{-j}
/n_l$ and $n_l\le\frac54$.  Hence $|c_l|\le\frac{25}{16}C_q2^{J(w)}t/
\delta_{\min}(L)$, while Lemma~\ref{lem:III-track} gives $|c_l|\ge K_*(w)^2Kt/
(4D(L))\ge K_P(w)t/(4D(L))=2C_q2^{J(w)}t/\delta_{\min}(L)$: a contradiction.

(O) Otherwise no coarse carrier lives at $j$ (a coordinate of $S_l$ with $l>L$
is not in $T(L)$), and $V_c(j)=0$.
\end{proof}

\subsubsection*{The deep assignment on the support}
Let $f^\#$ be the companion of $f^r$ produced by Step~0, the preparation
(Definition~\ref{def:III-prep}) and Step~A (Lemma~\ref{lem:III-cone}) of the
route-V construction, with the moves of Subsection~\ref{subsec:III-levers}
(no step~(2), Lemma~\ref{lem:III-noBinf}).  Let $F_{\rm low}\subseteq F$ be the
set of support coordinates lowered by these moves: far pull coordinates
$j\in F$ with $s_j=\epsilon_l$, which are lowered to contacts and then pulled
(so they carry, at $f^\#$, the pull mass $24\lambda_lv_l(j)$ of sign
$-\epsilon_l$).  Let $Z$ be the set of coordinates of type ZD
(Lemma~\ref{lem:III-prep}(d)): coordinates $s\le\bar d(L)$, private to a
zero-data coarse carrier, which are not thick; at $f$ they are off $F$ or thin.
We use the following properties of $f^\#$, established in
Lemmas~\ref{lem:III-prep} and~\ref{lem:III-cone}:
(P$^\#$1) $|a^\#_j|\ge|a^r_j|/2$ and $\sgn a^\#_j=s_j$ for $j\in F\setminus
(F_{\rm low}\cup Z)$; (P$^\#$2) the support coordinates of $f^\#$ outside $F$ are banks (of
(P$_2$) and of the tuning pre-placements), at which the route-V data are
those of a contact of the bank's sign, i.e.\ contact-like up to the fine
residue and the base term, and pulled coordinates (pre-placed tuning pulls
and Step-0 pulls at pullable far contacts), which carry the pull mass
$24\lambda_lv_l(j)$ of sign $-\epsilon_l$; (P$^\#$3) the coarse
statuses of the route-V companion.

For $t\in\mathcal S$ let $(B_\pm,\Theta_\pm)$ be the decomposition of the mate
$g^r$ at $f^r$ at scale $t$ (as in [II,~Subsection~6.7]), $V$ the route-V
switching vector and $A_j:=2|a^r_j|/t$ ($j\in F$).  On $F\setminus(F_{\rm low}
\cup Z)$
define $b^+_1$ by the deep assignment of [II,~Lemma~6.28] with $X:=V_c|_F$
and the trace slack: with $x_j:=s_jB_+(j)$, $y_j:=s_j(B_+(j)-V_c(j))$, let
$\sigma_j$ be the point of $[y_j-A_j-\mathrm{tr}_j,x_j+A_j]$ closest to $0$ and
$b^+_1(j):=B_+(j)-s_j\sigma_j$; the interval is nonempty because $x_j-y_j=
s_jV_c(j)\ge-2A_j-\mathrm{tr}_j$ (Lemma~\ref{lem:III-nodeep}).  On $F_{\rm low}$
put $b^+_1(j):=\chi_jV_c(j)$, where $\chi_j$ is the projection of
$s_jB_+(j)/(s_jV_c(j))$ onto $[0,1]$ if $s_jV_c(j)>0$ and $\chi_j:=0$ otherwise.
On $Z\cap F$ put $b^+_1(j):=\frac12V(j)$ ($=\frac12V_f(j)$, since only a
zero-data coarse carrier lives at $j$).  Off $F$ the $+$ side is that of the
route-V companion (at coordinates of $Z\setminus F$ it is $\frac12V_{\rm fine}(j)$,
as in Lemma~\ref{lem:III-LV}(b)).  With $a^\#_Z:=a^\#1_{\N\setminus Z}$ put
\[
 b^+:=b^+_1+b^+_{\rm off}-\varkappa\frac{a^\#_Z}{a^\#_Z(\zh^\#)},\quad
 \varkappa:=(b^+_1+b^+_{\rm off})(\zh^\#),\qquad b^-:=b^+-V;
\]
$a^\#_Z(\zh^\#)\ge\frac14$ because the masses on $Z$ are at most $2\bar d(L)b$.
and $\omega^\pm$ are those of the route-V companion.

\begin{proposition}[Cushion domination at the deep-raised companion]
\label{prop:III-cushion}
For every $t\in\mathcal S$ the pairs $(b^\pm,\omega^\pm)$ represent one
functional $g_t$ at $f^\#$, and together with the other scales of $\mathcal S$
they satisfy \textup{(U1)}, \textup{(U2)}, \textup{(U4$^{\inf}$)}, \textup{(U7)} and
\begin{enumerate}
\item[(i)] for $j\in F\setminus(F_{\rm low}\cup Z)$:
$(s_jb^+(j))_-\le5|a^\#_j|/t$ and $(s_jb^-(j))_+\le5|a^\#_j|/t+|V_f(j)|+
\mathrm{tr}_j$;
\item[(ii)] for $j\in F_{\rm low}$: at a pulled coordinate $t|b^\pm(j)|\le
|a^\#_j|$;
\item[(iii)] for $j\in Z$: $b^\pm(j)=\pm\frac12V(j)$ (with $V(j)=V_f(j)$ on
$F$), so $b^\theta(j)=0$ and $\mathrm{viol}(j)\le|V(j)|$;
\item[(iv)] at banked coordinates $s\in\supp a^\#\setminus F$ \textup{(}banks
of \textup{(P$_2$)} and the base banks of the tuning pre-placements of
Step~A\textup{)}: $b^\pm(s)$ are contact-like of sign $z^\#_s$ up to the fine
residue, plus the base term $-\varkappa a^\#_s/a^\#_Z(\zh^\#)$ of modulus
$\le4Kt|a^\#_s|\le|a^\#_s|/t$;
\item[(v)] at pulled coordinates $j\in\supp a^\#\setminus F$: $t|b^\pm(j)|\le
|a^\#_j|$;
\end{enumerate}
so \textup{(U3$^{\inf}$)} holds on all of $\supp a^\#$ with $C_+:=5$ and excess
$r(j):=|V_f(j)|+\mathrm{tr}_j$, of total mass at most $C_f\mathrm{Des}(L)c_{L+1}+
C_fc_{L+1}^2/t$, which is counted in $\epsilon$.  Moreover
\begin{gather*}
 p^*(g-g_t)\le K_{\rm data}t,\qquad K_{\rm data}:=(JK_*(w)^{J+1}+C_f)K,\\
 \|B_+1_F-b^+_1\|_1\le\frac{(1+\eta')t}{q^r_0}+Kt+\frac{7Cc_{L+1}}t .
\end{gather*}
\end{proposition}

\begin{proof}
(i) \emph{Deep assignment.}  By [I,~Lemma~8.28] at $f^r$ (with the budget of the
window), $s_jB_+(j)\ge-|a^r_j|/t-\mathfrak f^+_j$ and $s_jB_-(j)\le|a^r_j|/t+
\mathfrak f^-_j$, with $\sum_F(\mathfrak f^+_j+\mathfrak f^-_j)\le(1+\eta')t/
q^r_0$.  The point of $[y-A-\mathrm{tr},x+A]$ closest to $0$ has modulus
$(-x-A)_++(y-A-\mathrm{tr})_+\le\mathfrak f^+_j+\mathfrak f^-_j+|(\Delta B-V_c)_j|$,
because $B_+-V_c=B_-+(\Delta B-V_c)$.  Hence $|B_+(j)-b^+_1(j)|\le\mathfrak f^+_j+
\mathfrak f^-_j+|(\Delta B-V_c)_j|$, $s_jb^+_1(j)=x_j-\sigma_j\ge-A_j$ and
$s_j(b^+_1-V_c)(j)=y_j-\sigma_j\le A_j+\mathrm{tr}_j$; also $|b^+_1(j)|\le A_j+
|V_c(j)|+\mathrm{tr}_j$.  On $Z\cap F$, $V_c(j)=0$, and the flip bounds give
$|B_+(j)|\le A_j/2+\mathfrak f^+_j+\mathfrak f^-_j+|(\Delta B-V_c)_j|$, so
$|B_+(j)-b^+_1(j)|\le A_j+\mathfrak f^+_j+\mathfrak f^-_j+|(\Delta B-V_c)_j|+
|V_f(j)|$; since $|a^r_j|\le2b$ on $Z$, $\sum_{Z\cap F}A_j\le4\bar d(L)b/t\le Kt$.
On $F_{\rm low}$, $V_c(j)=c_lv_l(j)$ has the sign
$s_j=\epsilon_l$ (a lowered far coordinate of a data-carrying carrier); the
flip bounds give
$-A_j/2-\mathfrak f^+_j\le s_jB_+(j)\le s_jV_c(j)+A_j/2+\mathfrak f^-_j+|(\Delta B-
V_c)_j|$, so $|B_+(j)-\chi_jV_c(j)|\le A_j+\mathfrak f^+_j+\mathfrak f^-_j+
|(\Delta B-V_c)_j|$; and $\sum_{F_{\rm low}}A_j\le Kt$ because $|F_{\rm low}|\le
2|\Lambda|$, the lowered moduli are at most $\max(2b,x'_{\max})$ and
$x'_{\max}\le T^3$ (Definition~\ref{def:III-resources}).

(ii) \emph{$\ell_1$-bound and closeness.}  $\|\Delta B-V_c\|_{\ell_1(F)}\le
\sum_{l\le L}|c^{\rm act}_l-c_l|\,\|u_l\|_1+\|\sum_{l>L}\Delta\theta_lu_l\|_1
\le Kt+6\sum_{l>L}\lambda_l/t$, by Lemma~\ref{lem:III-track} (whose errors are
the zeroing and projection errors of [II,~Lemma~5.53], summing to at most $Kt$)
and [I,~Lemma~8.10] for the fine carriers.  This gives the $\ell_1$-bound.  The
functional $g_t$ differs from the one built in Section~\ref{sec:III-routeV}
from $B_+1_F$ by $b^+_1-B_+1_F$ and the corresponding $\varkappa$-terms.  The
closeness of that functional to $g$ is bounded by the zeroing cost
$JA_{J+1}Kt\le JK_*(w)^{J+1}Kt$ of [II,~Lemma~5.53] plus $C_fKt$, and
$(1+\|U\|)\|B_+1_F-b^+_1\|_1\le C_fKt$ because $c_{L+1}\le e^{-1/T}\le t^2$; so
$p^*(g-g_t)\le K_{\rm data}t$ after enlarging $C_f$.  Both pairs represent $g_t$ because $b^+-b^-=V$ and
$\omega^--\omega^+=\Delta\omega$ ([II,~Lemma~2.22]), and $b^+(\xi^\#)=0$ by the
choice of $\varkappa$.

(iii) \emph{Cushions.}  $|\varkappa|\le Kt$: $b^+_1$ differs from $B_+1_F$ by
$O(Kt)$ in $\ell_1$, $|B_+(\zh^r)|\le t/(2q^r_0)$ by [I,~Lemma~8.7(a)], and
$|\ip{U^*B_+}{e^\#-e^r}|\le\|B_+\|_1\|e^\#-e^r\|\le t$ since the companion
moves $e$ by at most $C\,\mathrm{Des}(L)T^3$.  On $F\setminus(F_{\rm low}\cup Z)$,
by (i) and (P$^\#$1), $(s_jb^+(j))_-\le A_j+4|\varkappa||a^\#_j|\le
4|a^\#_j|/t+4Kt|a^\#_j|\le5|a^\#_j|/t$, and $(s_jb^-(j))_+=(s_j(b^+_1-V_c-V_f)(j)
-\varkappa s_ja^\#_j/a^\#_Z(\zh^\#))_+\le A_j+\mathrm{tr}_j+|V_f(j)|+4Kt|a^\#_j|$,
as $4Kt^2\le1$.  On $Z$ the $\varkappa$-term vanishes, so $b^+(j)=\frac12V(j)$
(and $b^+(j)=\frac12V_{\rm fine}(j)$ off $F$), $b^-(j)=-\frac12V(j)$, $b^\theta(j)=0$.
On $F_{\rm low}$: at a pulled coordinate
$t|b^\pm(j)|\le t|V(j)|\le24\lambda_lv_l(j)=|a^\#_j|$, since $t\le T_{\rm hi}(w)$
and $|c_l|\le C_fD(L)\lambda_l$.  The bound on $\sum(|V_f(j)|+\mathrm{tr}_j)$ is
Lemma~\ref{lem:III-nodeep} and $\|V_f\|_1\le Cc_{L+1}$.
At a banked coordinate $s\notin F$ the route-V data are those of a contact
of the bank's sign up to the fine residue (as in Proposition~\ref{prop:III-BR}),
and the base term has modulus $|\varkappa||a^\#_s|/a^\#_Z(\zh^\#)\le4Kt|a^\#_s|
\le|a^\#_s|/t$ because $4Kt^2\le1$; this gives (iv).  At a pulled coordinate
$j\notin F$ (a pre-placed tuning pull, or a Step-0 pull at a pullable far
contact), the argument for $F_{\rm low}$ applies verbatim, since it does not
use $j\in F$: $t|b^\pm(j)|\le t|V(j)|\le24\lambda_lv_l(j)=|a^\#_j|$ for
$t\le T_{\rm hi}(w)$; this gives (v).

(iv$'$) \emph{The other conditions.}  Only the base part on $F$ changed; $\omega^\pm$
and the data off $F$ are those of the route-V companion, so (U2) and (U7) hold
as there.  (U1): $t\|b^\pm\|_1\le\sum_FtA_j+2t\|V\|_1+t\sum\mathrm{tr}_j+
(\text{off-}F)\le4\|a^r\|_1+C_f(1+D(L)t)\le A_0$; the $\Gamma_w$-bound changes
by $O(Kt)$, because $\Gamma_w^{1/2}$ is a seminorm and $q_0h(y)\le\|U\|^2
\|y\|_1^2/\nu$.  (U4$^{\inf}$): the violations of the route-V companion, plus
the excesses of (i), which sit at support coordinates.
\end{proof}

\begin{lemma}[Fine carriers enter only through their weight]\label{lem:III-noBinf}
Lemma~\textup{\ref{lem:III-noB}} and Corollary~\textup{\ref{cor:III-noB}} hold for
base rows of arbitrary support, with $J_{\rm fine}:=\N\setminus(F^\#\cup
E_c(w))$: step~\textup{(2)} of the route-V companion is omitted, the data on
$J_{\rm fine}$ are $b^+:=\chi V$, $b^-:=-(1-\chi)V$ with $\chi=\frac12$ at free
coordinates and $\chi\in[0,1]$ given by [I,~Lemma~8.26] at contacts,
$\|V1_{J_{\rm fine}}\|_1\le C_f\mathrm{Des}(L)c_{L+1}$, the violation on
$J_{\rm fine}$ is at most $\|V1_{J_{\rm fine}}\|_1$ with $b^\theta=0$ at
violated free coordinates, and \textup{(U5$'$)} holds with
$x_0:=c_f\gap_{\min}\Phi_{\min}(L)$.  In Theorem~\textup{\ref{thm:III-UEinf}}
hypothesis \textup{(U5)} may then be replaced by \textup{(U5$'$)} provided
$4c_{\rm fine}\le c_{\rm late}\delta s_1$, which holds under \textup{(W$_{\exp}$)}
for $s_1\ge e^{-1/(2T)}$ and $T$ small.
\end{lemma}

\begin{proof}
The proofs of Lemma~\ref{lem:III-noB} and Corollary~\ref{cor:III-noB} do not
use finiteness of $F$: a coordinate of $J_{\rm fine}$ lies in no coarse target
and in no near part $S_l\cap[1,s_{\rm far}(w)]$, so coarse carriers contribute
to $V(j)$ only through $S_l\cap(s_{\rm far}(w),\infty)$, at most $C_f D(L)
c_{L+1}^2$ in total, and fine carriers at most $C_\Delta\sum_{k>L}\lambda_k$.
At a free $j$, $|b^+_j|+|b^-_j|=|V(j)|$ and $b^\theta_j=0$; at a contact,
$(z_jb^+_j)_-+(z_jb^-_j)_+=(z_jV(j))_-$.  (U5$'$) follows from the coarse
statuses at $f^\#$ (Lemma~\ref{lem:III-cone}(c)); without step~(2) there is no
later drift.  The replacement of (U5) is Lemma~\ref{lem:III-noB}, whose proof
is local in the scrambled sets and does not involve $F$.  Finally
$4c_{\rm fine}\le4C\sum_{k>L}\lambda_k\le Ce^{-1/T}\le c_{\rm late}\delta
e^{-1/(2T)}$ for $T$ small.
\end{proof}

\begin{remark}\label{rem:III-noB}
At arbitrary base support the static-ownership recursion of the finite-support
construction would require the signature sets of fine carriers to avoid $F$,
which cannot be arranged.  Lemma~\ref{lem:III-noBinf} removes the question: only
the total weight $c_{L+1}$ of the fine carriers is used, against $s_1\ge
e^{-1/(2T)}$, and this is where (W$_{\exp}$) is essential.
\end{remark}

\subsection{Levers inside the support and the preparation of the base row}
\label{subsec:III-levers}

At finite base support every lever of the route-V construction acts at
coordinates off the support.  At arbitrary support the near and far
coordinates of a carrier may lie in $F$, and the moves available there have
first-order side effects on all values (common modes).  This subsection
collects the first-order calculus and the preparation of the base row.

\subsubsection*{Depth}
Let $\Lambda_L$ be the set of coarse carriers, $U_L:=\bigcup_{k\in\Lambda_L}
\supp u_k$, and let $\bar d(L)$ be the least $D\ge d(L)$ with
\[
 \max_{k\in\Lambda_L}\|u_k1_{(D,\infty)}\|_1\le2^{-L^2}\big(|\Lambda_L|H(L)\big)^{-4}
\]
(it exists since the $u_k$ are finitely many vectors of $\ell_1$), and
$v_{\min}(L):=\min\{|u_k(s)|:k\in\Lambda_L,\ s\in\supp u_k\cap[1,\bar d(L)]\}>0$.
Coordinates $s\le\bar d(L)$ are said to be at \emph{design depth}.  By
Definition~\ref{def:III-designinf}, $\mathrm{Des}(L)\ge2^{\bar d(L)}\bar d(L)
H(L)/((\mu^\star_{\bar d(L)})^2v_{\min}(L))$, so a mass change at $s\le\bar d(L)$
has normalized efficiency $(\mu^\star_s)^2/\nu\ge1/\mathrm{Des}(L)$ at rows with
$\nu\le1$.

\begin{definition}[Resources]\label{def:III-resources}
Fix a row $f_1$ (the input row of a linear system: of Step~A, or the
pre-row of Definition~\ref{def:III-approxinf}), a set $\Lambda$ of coarse
carriers to be tuned and a scale $x$ of the system.  For a coarse carrier
$l$ put $\mathrm{val}_l:=\epsilon_lu_l(\hat x_1)$, with $\epsilon_l$ the sign of
the switching of $l$ if $l$ carries data and a fixed sign otherwise, and
$g_k:=\epsilon_k\ip{U^*u_k}{e_1}$.  Let
\[
 x'_{\max}:=\frac{u^2T^3}{\mathrm{Des}(L)^2|\Lambda|},\qquad
 \mathrm{Far}_l(x):=\{j\in S_l:x\le v_l(j)\le x'_{\max}\} .
\]
A far $j$ is \emph{pullable} for $l$ if $j\notin F$ and $\epsilon_lz_1(j)>-\frac12$,
or $j\in F$, $s_j=\epsilon_l$ and $|a_j|\le x'_{\max}$ (then $j$ is first lowered
to a contact).  A coarse carrier is \emph{zero-data} if its coefficient $c_l(t)$
vanishes for all $t\in\mathcal S$ (inactive $\Omega$-carriers, held carriers,
pushed carriers, coarse peaks of inactive blocks), and \emph{data-carrying}
otherwise.  The \emph{resources} of $l$ are:
\begin{enumerate}
\item[(P1)] $l$ zero-data, and a near $s\in N_l$ off $F$ or thin
($|a_s|\le b$): the composite move at $s$ (Lemma~\ref{lem:III-actions}(iii));
\item[(P2)] (TU) $l$ data-carrying, a near bank coordinate $s\in N_l$ which is a
contact of sign $\epsilon_l$ or a thin support coordinate with $s_s=
\epsilon_l$, and a pullable $j_p\in\mathrm{Far}_l(x)$ ([II,~Lemma~4.15]);
\item[(P3)] (Hilbert pair) two thick $s,s'\in N_l$ ($|a_s|,|a_{s'}|\ge u$) with
robust spread $|1-\mathfrak a_s/\mathfrak a_{s'}|\ge u$, $\mathfrak a:=a/v_l$
([II,~Lemma~6.46(b)]);
\item[(P4)] (single mass) a thick $s\in N_l$: a mass change at $s$;
\item[(P5)] (anchor) a thick $p\in F\cap[1,\bar d(L)]$ with $p\notin U_L$.
\end{enumerate}
$l$ is \emph{privately tunable} if it has (P1), (P2) or (P3).
\end{definition}

The choice of $x'_{\max}$ (rather than a larger far depth) is what makes the
costs of far moves $o(u^2T^2)$, as required by
Theorem~\ref{thm:III-Eenginf}(E-e) with $c_{\rm flat}\ge c_fu$
(Lemma~\ref{lem:III-costs}).

\begin{lemma}[Exact identity and signed dichotomy]\label{lem:III-identity}
Let $f_1$ be a row with data $(a,z,e,\nu,\hat x)$ and $F=\supp a$.  For $u\in
\ell_1$ put $\varkappa_u:=\ip{U^*u}{e}/\nu$, $D_u(s):=u(s)-\varkappa_ua_s$ and
$c_s(u):=D_u(s)\sgn a_s$ for $s\in F$, $c_s(u):=u(s)z_s$ for $s\notin F$.  Then
for every $D\ge1$, with $\alpha(D):=\sum_{s\in F,\,s>D}|a_s|$,
\begin{equation}\label{eq:III-identity}
 u(\hat x)=\varkappa_u\big(1-\alpha(D)\big)+\sum_{s\le D}c_s(u)+\sum_{s>D}
 u(s)z_s,\qquad1-\alpha(D)\ge\nu .
\end{equation}
Consequently: if $\varkappa_u\ge0$, then $\sum_{s\le D}(c_s(u))_-\ge
\ip{U^*u}{e}-u(\hat x)-\|u1_{(D,\infty)}\|_1$; if $\varkappa_u\le0$, then
$\sum_{s\le D}(c_s(u))_+\ge|\ip{U^*u}{e}|+u(\hat x)-\|u1_{(D,\infty)}\|_1$.
\end{lemma}

\begin{proof}
$u(\hat x)=u(z)+\ip{U^*u}{e}=\sum_Fu(s)\sgn a_s+\sum_{F^c}u(s)z_s+
\varkappa_u\nu$, and $\sum_Fu(s)\sgn a_s=\varkappa_u\|a\|_1+\sum_FD_u(s)\sgn a_s$;
$\|a\|_1+\nu=1$.  For $s\in F$, $s>D$, write $D_u(s)\sgn a_s=u(s)z_s-
\varkappa_u|a_s|$; then $\varkappa_u(\nu+\|a\|_1-\alpha(D))=\varkappa_u(1-
\alpha(D))$ and $1-\alpha(D)=\|a1_{[1,D]}\|_1+\nu$.  If $\varkappa_u\ge0$,
$\varkappa_u(1-\alpha(D))\ge\varkappa_u\nu=\ip{U^*u}{e}$, and
\eqref{eq:III-identity} with $|z|\le1$ gives $\sum_{s\le D}c_s(u)\le u(\hat x)-
\ip{U^*u}{e}+\|u1_{(D,\infty)}\|_1$, whence the bound on the negative parts;
the case $\varkappa_u\le0$ is symmetric.
\end{proof}

\begin{lemma}[First-order actions]\label{lem:III-actions}
Let $f_1$ be a row as in Lemma~\textup{\ref{lem:III-identity}} and $u\in\ell_1$.
Up to remainders at most $C\|U^*u\|(\mu^\star_s|\mathrm dm|/\nu)^2$ for
$|\mathrm dm|\le\nu/(2\mu^\star_s)$, and after renormalization by $q^*$ (which
changes neither $e$ nor $z$):
\begin{enumerate}
\item[(i)] a change $\mathrm dm$ of $a_s$ at $s\in F$ \textup{(}sign kept\textup{)}
changes $u(\hat x)$ by $((\mu^\star_s)^2/\nu)D_u(s)\,\mathrm dm$;
\item[(ii)] a new mass $\mathrm dm\,z_s$ ($\mathrm dm\ge0$) at a contact
$s\notin F$ changes $u(\hat x)$ by $((\mu^\star_s)^2/\nu)u(s)z_s\,\mathrm dm$;
\item[(iii)] a change $\mathrm dz$ of $z_s$ at $s\notin F$ changes $u(\hat x)$ by
exactly $u(s)\,\mathrm dz$.  At a zero-data contact $s$ the \emph{composite
move} \textup{(}for $P>0$ a bank of mass $(\nu/(\mu^\star_s)^2)Pz_s$, for
$P\le0$ the move $z_s\mapsto z_s(1+P)$\textup{)} is $C^{1,1}$ in $P$ with slope
$u(s)z_s$ on both sides.
\end{enumerate}
In particular, for a single mass \textup{(P4)} at a thick $s\in N_l$, normalized so
that a unit step changes $\mathrm{val}_l$ privately by $1$, every
coarse $k$ satisfies $\mathrm d\mathrm{val}_k=(\delta_{kl}-g_kh_l)
\,\mathrm dP$ with the signed coefficient $h_l:=\epsilon_la_s/(\nu v_l(s))$,
$|h_l|\le H(L)/\nu$; for an anchor \textup{(P5)} the same holds without the
private term.  Every move of the size $x$ at design depth moves $e$ by at most
$C\,\mathrm{Des}(L)x$ and no $z$ outside its own coordinate.
\end{lemma}

\begin{proof}
$U^*a$ changes by $\mu^\star_s\mathrm dm\,k_s$, and the derivative of
$y\mapsto y/\|y\|$ at $U^*a$ is $(I-ee^T)/\nu$; with $e_s:=\ip{k_s}{e}=
\mu^\star_sa_s/\nu$ we get $\mathrm de=(\mu^\star_s\mathrm dm/\nu)(k_s-e_se)$ and
$\ip{U^*u}{k_s-e_se}=\mu^\star_sD_u(s)$.  The second derivative of $y\mapsto
y/\|y\|$ is at most $3\|\mathrm dy\|^2/\|y\|^2$.  (ii) is (i) with $a_s=0$;
(iii): $\hat x=z+Ue$ and $e$ does not depend on $z$.  For the composite move both
branches have derivative $u(s)z_s$ at $P=0$, and the bank branch is smooth.
For (P4): for $k\ne l$ coarse, $u_k(s)=0$ ($s$ lies beyond $s_{\max}(L)$ in
$S_l$); by (i) a mass change $\mathrm dm$ changes $\epsilon_k u_k(\hat x)$ by
$\epsilon_k((\mu^\star_s)^2/\nu)(u_k(s)-\varkappa_{u_k}a_s)\mathrm dm$; with
$\mathrm dP:=\epsilon_l((\mu^\star_s)^2/\nu)v_l(s)\mathrm dm$ this is
$(\delta_{kl}-g_kh_l)\mathrm dP$.  The last claim is
Lemma~\ref{lem:III-masses}(ii) and the efficiency bound above.
\end{proof}

The other resources act as follows (proved in the cited places): (P1) has
slope exactly $1$ on both sides and affects no other coarse carrier
(Lemma~\ref{lem:III-DC}); (P2) shifts $\mathrm{val}_l$ down by at
least $v_l(j_p)/2\ge x/2$ through the pull, and the bank provides every upward
correction continuously, with second-order cross effects [II,~Lemma~4.15];
(P3) is private and two-sided with efficiency $\ge u/\mathrm{Des}(L)$ and cross
effects $O(\mathrm{Des}^2u^{-4}x^2)$ [II,~Lemma~6.46(b)].  A far coordinate
gives only a \emph{one-sided} private move: a contact is at the end of its range,
so $z$-moves there lower $\mathrm{val}_l$ only, and a bank there has
negligible efficiency.

\begin{lemma}[Offsets do not create two-sidedness]\label{lem:III-offset}
Let $N$ be a square matrix with a left null vector $y$, $c$ a column with
$y^Tc\ne0$, and $F^0$ a vector.  If $(P,P_0)$ solves the linear system
$F^0+NP+cP_0=0$, then $P_0=-y^TF^0/(y^Tc)$.  In particular, for a one-sided
lever with column $c$ \textup{(}a bank, admissible only for $P_0\ge0$ measured
from zero lever mass\textup{)} the total position is forced and does not
depend on any offset placed along the lever in the same row; if
$y^TF^0/(y^Tc)>0$, the linearized system has no admissible solution.
\end{lemma}

\begin{proof}
Apply $y^T$.  An offset $m_0$ placed in the row from which $F^0$ is computed
is a move along $c$, so it changes $F^0$ by $c\,m_0$ and the remaining
position by $-m_0$; the total is unchanged.
\end{proof}

So a bank placed \emph{with an offset} at the row where the linear system is
solved, as a block-norm lever along a degenerate direction, is one-sided and
cannot replace a two-sided lever.  The remedy below is different: the offsets
are placed in the \emph{base row}, whose values are the targets of the system,
so that the forced position is the offset itself up to $O(r)$.

\begin{lemma}[Exact drift identity]\label{lem:III-driftid}
Let $f^\circ$ be a row with data $(a^\circ,z^\circ,F^\circ,\nu^\circ,e^\circ)$, and
let $u=\sum_kc_ku_k$ be a finite combination of carrier vectors with
$u|_{F^\circ}=\varkappa a^\circ|_{F^\circ}$.  Let $f^1$ be a row with data
$(a^1,z^1)$ such that $z^1=z^\circ$ on $\supp u$ and
$(\supp a^1\setminus F^\circ)\cap\supp u=\emptyset$.  Then
\[
 u(\zh^1)-u(\zh^\circ)=-\frac{\varkappa\nu^\circ}2\|e^1-e^\circ\|^2 .
\]
\end{lemma}

\begin{proof}
$u(\zh)=u(z)+\ip{U^*u}{e}$ and $u(z^1)=u(z^\circ)$.  $U^*u=\varkappa U^*(a^\circ
1_{F^\circ})+U^*(u1_{(F^\circ)^c})=\varkappa\nu^\circ e^\circ+\sum_{s\notin F^\circ}
\mu^\star_su(s)k_s$, while $e^1=\sum_{s\in\supp a^1}\mu^\star_sa^1(s)k_s/\nu^1$; so
$\ip{U^*u}{e^1}=\varkappa\nu^\circ\ip{e^\circ}{e^1}$ (the cross terms vanish by
hypothesis) and $\ip{U^*u}{e^\circ}=\varkappa\nu^\circ$.  The difference is
$\varkappa\nu^\circ(\ip{e^\circ}{e^1}-1)=-(\varkappa\nu^\circ/2)\|e^1-e^\circ\|^2$.
\end{proof}

\begin{remark}[Why the original lever set can be obstructed]\label{rem:III-drift}
Lemma~\ref{lem:III-driftid} allows arbitrary raises and lowerings on $F^\circ$,
truncations to contacts of the same sign, and masses and $z$-moves outside
$\supp u$.  So, at a row where a combination $u$ of carriers is proportional to
$a^\circ$ on the support, support moves change $u(\zh)$ only at second order and
only in one direction (it decreases if $\varkappa>0$), and the only two-sided
first-order resources for $u$ are moves on $\supp u\setminus F^\circ$.  Applied
to the left null vector of the $d$-consistency system built from single masses
on the support and block-norm levers (which never act on $\supp u\setminus
F^\circ$), the same computation, keeping the cross term, shows that the
corresponding combination of the equations equals $-(\varkappa\nu_0/2)\|e-e_0\|^2$
plus the first-order effect of the window masses on $\supp u\setminus F_0$ plus
the cut-off beyond $N''$; when the window term is of first order in $s_1$,
exact $d$-consistency with those levers is impossible for small $s_1$.
Theorem~\ref{thm:III-SysDC} avoids this by preparing the base row: the
coordinates of $\supp u\setminus F_0$ at design depth then carry banks and
become two-sided resources.
\end{remark}

\begin{definition}[Preparation]\label{def:III-prep}
Let $\theta_0=\theta_0(w):=b/\mathrm{Des}(L)^{C_P}$ (Definition~\ref{def:III-designinf}).
The \emph{prepared row} $f^P$ is obtained from the row after the deep raise
and Step~0, and before the \L ojasiewicz step of Step~A, by:
\begin{enumerate}
\item[(P$_1$)] closing to a contact every free $s\le\bar d(L)$, $s\in U_L$,
with tiny room $1-|z_s|\le b$;
\item[(P$_2$)] placing a bank of mass $\theta_0z_s$ at every contact $s\le\bar
d(L)$, $s\in U_L$;
\item[(P$_3$)] raising, with its sign, every $s\in F\cap U_L$, $s\le\bar d(L)$,
with $|a_s|<\theta_0$, to $|a_s|=\theta_0$.
\end{enumerate}
\end{definition}

\begin{lemma}[Properties of the preparation]\label{lem:III-prep}
Let $L\ge l_f$ and $w$ clean.
\begin{enumerate}
\item[(a)] $p^*(f^P-f^r)\le C_f\mathrm{Des}(L)\,b\log(1/b)$.
\item[(b)] Every value moves by at most $b+C\bar d(L)\theta_0\le2b$ \textup{(}the
closing moves of \textup{(P$_1$)} by at most $b$, the banks and raises by at most
$C\bar d(L)\theta_0$\textup{)}, every rate object by at most
$2\,\mathrm{Lip}\cdot b$; the clean classification survives with the factor $2$.
\item[(c)] At a banked coordinate $s$ the route-V data are those of a contact of
sign $z_s$: $b^+=\chi V'+\beta_s$, $b^-=-(1-\chi)V'+\beta_s$ with $z_sV'\ge0$ up
to the fine residue, and $|\beta_s|\le4Kt|a^\#_s|$; hence
\textup{(U3$^{\inf}$)} holds there with $C_+=5$.  At raised thin
coordinates the data of Proposition~\textup{\ref{prop:III-cushion}} satisfy
\textup{(U3$^{\inf}$)} a fortiori.
\item[(d)] At the companion and at every pre-row of
Definition~\textup{\ref{def:III-approxinf}} built at it, every $s\le\bar d(L)$
with $s\in U_L$, or $s\in F$ with $|a_s|\ge u/2$, is of one of the types:
\textup{TS} (thick support, $|a_s|\ge u/2$; an \emph{anchor} if $s\notin U_L$);
\textup{S$_{\rm th}$} (support with $|a_s|\ge\theta_0/2$: banked contacts, raised
thin coordinates); \textup{FR} (free with room $\ge u/2$; zero data or tolerated
violations with $b^\theta=0$); \textup{ZD} (a coordinate of a zero-data carrier,
in any state).  Every other $s\le\bar d(L)$ is off $U_L$ and either off $F$ or
thin, $|a_s|\le2b$.
\end{enumerate}
\end{lemma}

\begin{proof}
(a) Banks off $F$ and same-sign raises: the costs of [II,~Lemma~2.15] and
[II,~Lemma~2.12(iv)], with value moves bounded as in (b) and
$c(\delta)\le C\epsilon\log(e/\epsilon)$; there are at most $2\bar d(L)$ such
moves of size $\le\theta_0$.  The closing moves of (P$_1$) are $z$-moves of
size up to $b$ at up to $\bar d(L)$ coordinates; by [II,~Lemma~2.12(iv)] their
normer change costs $C_fc(\delta)$ with value moves of order $b$, that is at
most $C_f(|\Lambda_L|+\bar d(L))b\log(1/b)$.  As $|\Lambda_L|+\bar d(L)\le
\mathrm{Des}(L)$, (a) follows.  (b) Lemma~\ref{lem:III-actions}: each move
changes $u_k(\hat x)$ by at most $\theta_0\|u_k\|_1\cdot C$ (masses: efficiency
$\le1$; closing a room: $|u_k(s)|(1-|z_s|)\le b|u_k(s)|$, so the closing moves
change each value by at most $b\|u_k\|_1\le b$ in total).  The rate objects are
Lipschitz in the values, moduli and rooms; the clean dichotomy at depth
$\bar d(L)$ (Lemma~\ref{lem:III-depthclean}) is preserved with factor $2$.
(c) A banked $s$ was a contact of $f^r$ in the support of a coarse carrier.  If
the carrier carries data, the route-V data at $s$ are the contact row of the
class cone (row (X1), or (X1$^F$)), so $V'$ is $z_s$-signed up to the fine
residue; the companion support outside $\supp a^r$ is treated as a set of
contacts in the transplant of the route-V companion, which gives the stated
form with the base term $\beta_s=-\varkappa a^\#_s/a^\#_Z(\zh^\#)$, with the
normalization of Proposition~\ref{prop:III-cushion}: $|\varkappa|\le Kt$ and
$a^\#_Z(\zh^\#)\ge\frac14$.  Then $(z_sb^+)_-\le|\beta_s|$ and $(z_sb^-)_+\le
|\beta_s|$, and $|\beta_s|\le4Kt|a^\#_s|\le|a^\#_s|/t$ as $4Kt^2\le1$; this is
Proposition~\ref{prop:III-cushion}(iv).  Zero-data coordinates carry only
fine residues.  (d) By Lemma~\ref{lem:III-depthclean}, at $f^r$ every $s\in U_L\cap
[1,\bar d(L)]$ is thick, thin, a contact, free with robust room, or free with
tiny room; (P$_1$)--(P$_3$) turn thin coordinates, contacts and tiny rooms into
S$_{\rm th}$ (zero-data ones are also ZD).  Step~A moves TS and FR coordinates
by at most $\mathrm{Des}^Cx'_{\max}\ll u$ and uses S$_{\rm th}$ coordinates with
data only through raises (Lemma~\ref{lem:III-cone}); the pre-row moves
design-depth coordinates by at most $K_w\Pi\ll\theta_0$ (Lemma~\ref{lem:III-late});
truncation and deep raise act beyond $\bar d(L)$.
\end{proof}

\subsection{Uniform solvability of the $d$-consistency system}
\label{subsec:III-SysDC}

Let $f_0$ be the prepared companion after Step~A, $\Omega$ the set of all
coarse strict non-peaks of $f_0$ (no absorbers in route V), and $f_1$ a
pre-row of Definition~\ref{def:III-approxinf} at $f_0$.  The
\emph{$d$-consistency system} is $F(P)=0$ with $F_l(P):=\epsilon_l[u_l(\hat x(P))-
(A_m(P)/A^0_m)u_l(\zh_0)]$ ($l\in\Omega_m$), $A^0_m:=A_m(f_0)$, as in
Lemma~\ref{lem:III-DC}.  Put $p_l:=\epsilon_lu_l(\zh_0)/A^0_m$, $q_k:=
\epsilon_k\lambda_kw^0_m(k)$ ($k,l\in\Omega_m$), $g_k:=\epsilon_k\ip{U^*u_k}{e_1}$,
$G_m:=\ip{U^*R_m^*w^0_m}{e_1}=\sum_k\lambda_kw^0_m(k)\ip{U^*u_k}{e_1}$ (all
carriers $k$ of block $m$) and $\tilde g:=g-\sum_mG_mp_m$; here $p_m,q_m$ are
the restrictions to $\Omega_m$.  Then $q_m^Tp_m\le C^0_m\le\frac12$
(proof of Lemma~\ref{lem:III-DC}), $\|q_m\|_2\le1$ and $\|p_m\|_2\le
2|\Omega|^{1/2}/A_{\min}$.

\begin{lemma}[Left combinations]\label{lem:III-leftDC}
For $y\in\R^\Omega$ put $\beta_m:=y^Tp_m$ and $\Psi_y:=\sum_{l\in\Omega}y_l
\epsilon_lu_l-\sum_m\beta_mR_m^*w^0_m\in\ell_1$.  Then:
\begin{enumerate}
\item[(a)] $\Psi_y(\zh_0)=0$;
\item[(b)] for every move $\delta$ of Lemma~\textup{\ref{lem:III-actions}} at $f_1$
\textup{(}or a composite of such moves\textup{)}, with $c(\delta)$ the first-order
change of $F$ and $l_\delta(u)$ the first-order change of $u(\hat x)$,
$|y^Tc(\delta)-l_\delta(\Psi_y)|\le C|\delta|\,\|y\|(K_w\Pi+c_{L+1})$;
\item[(c)] $\ip{U^*\Psi_y}{e_1}=y^T\tilde g$.
\end{enumerate}
\end{lemma}

\begin{proof}
(a) $\sum_{l\in\Omega_m}y_l\epsilon_lu_l(\zh_0)=A^0_m\beta_m$, and
$(R_m^*w^0_m)(\zh_0)=\ip{w^0_m}{R_m^{**}\zh_0}=A^0_m$ because $w^0_m$ norms
$R_m^{**}\zh_0$.  (b) $\mathrm dF_l=\epsilon_l\mathrm du_l-p_l\mathrm dA_m$ and,
the block norm being differentiable at $R_m^{**}\hat x_1$ with norming functional
$w^1_m$, $\mathrm dA_m=l_\delta(R_m^*w^1_m)+O(|\delta|^2)$; so $y^Tc(\delta)=
l_\delta(\Psi_y)+\sum_m\beta_ml_\delta(R_m^*(w^0_m-w^1_m))$, and
$\|R_m^*(w^0_m-w^1_m)\|_1\le K_w\Pi+2c_{L+1}$ (Lemma~\ref{lem:III-UE1}, with the
fine carriers bounded by their weight).  (c) is the definition of $\tilde g$.
\end{proof}

\begin{theorem}[The $d$-consistency system is uniformly solvable]
\label{thm:III-SysDC}
Let $L\ge l_f$, $w$ clean, $f_0$ the prepared companion after Step~A and
$f_1$ a pre-row at $f_0$ with $\Pi\le1/K_w$.  Let $\mathsf C$ be the column
matrix of the $d$-consistency system at $f_1$ formed by:
\begin{enumerate}
\item[(i)] for every $l\in\Omega$ one two-sided lever: if $l$ is zero-data with
an off-$F$ or thin near coordinate, the composite move at it (type ZD); if
$l$ carries data and has an off-$F$ or thin near coordinate, a mass change at
it (it is banked or raised by the preparation: type S$_{\rm th}$); if all near
coordinates of $l$ are thick, (P3) if a pair of robust spread exists and (P4)
otherwise;
\item[(ii)] one elementary column at every $s\le\bar d(L)$ of type TS
(anchors included) or S$_{\rm th}$ (mass change), FR ($z$-move), ZD
(composite move).
\end{enumerate}
Then
\[
 s_{\min}(\mathsf C):=\min_{\|y\|_2=1}\|\mathsf C^Ty\|_2\ \ge\
 \varepsilon_{\rm DC}:=\frac{\nu_1(\mu^\star_{\bar d(L)})^2}{16\,\bar d(L)|\Omega|H(L)
 \big(1+8N|\Omega|^{1/2}/A_{\min}\big)},
\]
so $\mathsf C$ has the right inverse $R:=\mathsf C^T(\mathsf C\mathsf C^T)^{-1}$ with
$\|R\|_{\infty\to\infty}\le K_{\rm DC}:=|\Omega|^{1/2}/\varepsilon_{\rm DC}\le
C_{f,N}\mathrm{Des}(L)^{C_0}$.  Consequently the moves of \textup{(i)},
\textup{(ii)} form a lever system \textup{(}Definition~\textup{\ref{def:III-levsys}}%
\textup{)} with $K_{\rm lev}=K_{\rm DC}$, $W_{\rm pull}=0$ and no hypothesis on
Sherman--Morrison denominators, bordered scalars, spreads or anchors: if
\textup{(L-K2)} of Lemma~\textup{\ref{lem:III-late}} holds, $F$ has a zero
$P^*$ with $|P^*|_\infty\le K_{\rm DC}r$, every carrier of $\Omega$ is a strict
non-peak of $f(P^*)$ with $\mathrm{nv}=\mathrm{nv}^0$, the coarse peaks keep
their statuses, the S$_{\rm th}$ masses stay $\ge\theta_0/4$, and the TS and FR
coordinates stay robust.
\end{theorem}

\begin{proof}
Let $y$ be a unit vector with $\|\mathsf C^Ty\|_2<\varepsilon$; we show
$\varepsilon\ge\varepsilon_{\rm DC}$.  Put $\gamma:=y^T\tilde g$.

\emph{Step 1 (shape of the lever columns).}  By Lemma~\ref{lem:III-actions}
and the linearization of Lemma~\ref{lem:III-DC}, the lever of $l$, with its
parameter scaled so that its private slope is $1$, has the column
$e_l-h_l\tilde g-p_{m(l)}q_l+E_l$, with $h_l$ the signed common-mode coefficient ($h_l=0$ for (P3) and
ZD levers, $|h_l|\le H(L)/\nu_1$ always) and $\|E_l\|\le\varepsilon_E:=C_f
\mathrm{Des}(L)^C(K_w\Pi+c_{L+1})\le\varepsilon_{\rm DC}/8$ by (L-K2).  Hence
$y_l=\beta_{m(l)}q_l+\gamma h_l+e'_l$ with $e'_l:=(\mathsf C^Ty)_l-y^TE_l$,
$|e'_l|\le|(\mathsf C^Ty)_l|+\varepsilon_E$.

\emph{Step 2 ($\gamma$ is robust).}  Since $q_m^Tp_{m'}=0$ for $m\ne m'$, taking
the scalar product with $p_m$ gives $\beta_m(1-q_m^Tp_m)=\gamma h^Tp_m+e'^Tp_m$
(restrictions to $\Omega_m$), so $|\beta_m|\le4(|\gamma|\|h\|_2+\|e'\|_2)\|p_m\|_2$, and
\begin{align*}
 1=\|y\|_2&\le2\Big(\sum_m|\beta_m|\|q_m\|_2+|\gamma|\|h\|_2+\|e'\|_2\Big)\\
 &\le2\big(|\gamma|\|h\|_2+\|e'\|_2\big)\big(1+8N|\Omega|^{1/2}/A_{\min}\big).
\end{align*}
If $\|e'\|_2\ge1/(4(1+8N|\Omega|^{1/2}/A_{\min}))$, then $\varepsilon\ge
\varepsilon_{\rm DC}$ already.  Otherwise $h\ne0$ and $|\gamma|\ge\gamma_{\min}:=
1/(4\|h\|_2(1+8N|\Omega|^{1/2}/A_{\min}))\ge\nu_1/(4|\Omega|H(L)(1+8N|\Omega|^{1/2}
/A_{\min}))$; replacing $y$ by $-y$ we may assume $\gamma>0$.

\emph{Step 3 (the combination).}  By Lemma~\ref{lem:III-leftDC}, $\Psi_y(\zh_0)=0$
and $\ip{U^*\Psi_y}{e_1}=\gamma$.  Split $\Psi_y=\Psi^c+\Psi^f$, $\Psi^f$ the
fine-carrier part of $-\sum_m\beta_mR_m^*w^0_m$, with $\|\Psi^f\|_1\le
\sum_m|\beta_m|\sum_{k>L}\lambda_k\le C(\gamma/\gamma_{\min})c_{L+1}$.  The
coefficients of $\Psi^c$ are at most $C\gamma/\gamma_{\min}$, so by the choice
of $\bar d(L)$, $\|\Psi^c1_{(\bar d(L),\infty)}\|_1\le\gamma/16$ for $L\ge l_f$.
At $f_1$, $|\Psi^c(\hat x_1)|\le|\Psi_y(\hat x_1)-\Psi_y(\zh_0)|+\|\Psi^f\|_1\le
\|U^*\Psi_y\|\|e_1-e_0\|+\|\Psi_y1_{(N'',\infty)}\|_1+\|\Psi^f\|_1\le\gamma/16$
(the pre-row changes $z$ only beyond $N''$), and $\ip{U^*\Psi^c}{e_1}\ge
\frac{15}{16}\gamma$.

\emph{Step 4 (a resource).}  Lemma~\ref{lem:III-identity} at $f_1$ with
$u:=\Psi^c$ and $D:=\bar d(L)$ gives some $s\le\bar d(L)$ with
$|c_s(\Psi^c)|\ge\gamma/(2\bar d(L))$.  If $s\notin U_L$, then $\Psi^c(s)=0$, so
$s\in F_1$ and $|c_s|=\varkappa|a_s|$ with $\varkappa=\ip{U^*\Psi^c}{e_1}/\nu_1
\le2\gamma/\nu_1$; hence $|a_s|\ge\nu_1/(4\bar d(L))\gg2b$, and $s$ is TS (an
anchor).  (Thin $s\notin U_L$ contribute at most $2\varkappa b$ each, in total
$\ll\gamma/16$.)  If $s\in U_L$, $s$ has type TS, S$_{\rm th}$, FR or ZD by
Lemma~\ref{lem:III-prep}(d).  In all cases $\mathsf C$ contains the elementary
column at $s$.

\emph{Step 5 (contradiction).}  By Lemma~\ref{lem:III-leftDC}(b) and
Lemma~\ref{lem:III-actions}, $|y^Tc_s|\ge\mathrm{eff}_s|c_s(\Psi_y)|-\mathrm{err}$
with $\mathrm{eff}_s\ge(\mu^\star_{\bar d(L)})^2$ for masses ($\nu_1\le1$) and
$\mathrm{eff}_s=1$ for $z$-moves and composite moves, and $|c_s(\Psi_y)-
c_s(\Psi^c)|$ negligible.  So $\varepsilon>|y^Tc_s|\ge(\mu^\star_{\bar d(L)})^2
\gamma_{\min}/(4\bar d(L))\ge\varepsilon_{\rm DC}$.

\emph{Step 6 (nonlinear solution).}  Put $G(y):=F(Ry)=F(0)+y+\mathrm{Rem}(Ry)$
on the cube $|y|_\infty\le r$.  On the face $y_l=r$, $G_l\ge r-|F(0)|_\infty-
C_2K_{\rm DC}^2r^2\ge0$ when $K_{\rm DC}|F(0)|_\infty\le r/2$ and $C_2K_{\rm DC}^2r
\le\frac14$, with $C_2$ a design constant (second derivatives of
Lemma~\ref{lem:III-actions} at depth $\le\bar d(L)$ and of the levers, the
composite moves being $C^{1,1}$); symmetrically on $y_l=-r$; the
Poincar\'e--Miranda theorem gives a zero.  The S$_{\rm th}$ masses move by at most
$\mathrm{Des}(L)K_{\rm DC}r\le\theta_0/4$, TS and FR coordinates by $\ll u$;
values move by at most $C_{\rm lev}K_{\rm DC}r$, below every gap and margin
(Lemma~\ref{lem:III-late}).  The protection (PR) holds: TS and S$_{\rm th}$
coordinates are of type (PR$\beta$) since their masses stay above half their
pre-row values, and FR and ZD coordinates are of type (PR$\alpha$)
(Proposition~\ref{prop:III-cushion}(iii) and Lemma~\ref{lem:III-prep}(d)).
\end{proof}

\begin{remark}[The degenerate configurations are bypassed]\label{rem:III-SysDC}
(a) Steps 2--5 are a uniform dichotomy: a near-null left vector of the system
has a robust common mode $\gamma$, and then the identity of
Lemma~\ref{lem:III-identity} forces a robust first-order two-sided resource at
design depth.  No rate object for Sherman--Morrison denominators, bordered
scalars or spreads is needed.  The \emph{square} lever system (one lever per
carrier) may still be degenerate; it is the enlarged column system that is
uniformly onto.  (b) The resource of Step~4 is never the lever coordinate of an
$\Omega$ carrier, since there $|c_s(\Psi_y)|\le|y^T(\text{lever})|/\mathrm{eff}_s+
\mathrm{err}$ is small by Step~1.  (c) No TU lever and no pull at scale $r$ is
used.  (d) Without the preparation the resource of Step~4 may be a
data-carrying contact, where only a bank is available, and a bank placed at the
pre-row is one-sided along the degenerate direction
(Lemma~\ref{lem:III-offset}); after (P$_2$) the offset lies in the base row,
whose values are the targets of the system, and the forced position is
$\theta_0+O(K_{\rm DC}r)>0$.
\end{remark}

\begin{lemma}[Tolerance variant]\label{lem:III-tol}
In Theorem~\textup{\ref{thm:III-UEinf}} replace exact $d$-consistency on
$\Omega$ by
\begin{gather*}
 \textup{(TOL)}\qquad|\mathrm{mis}_{i,m}|\le s_1\ \text{ for every piece }i
 \text{ and block }m,\\ \mathrm{mis}_{i,m}:=(d'_m-d_m)(\Delta\omega_{i,m}),
\end{gather*}
and add \textup{(LATE-tol)} $C_f|\Omega|A_2^2K_w^2s_1\le c_{\rm late}\delta T^2$.
Then all conclusions of Theorem~\textup{\ref{thm:III-UEinf}} hold, with $K_w$
enlarged by an $f$-constant factor.  In particular an approximate zero of the
$d$-consistency system with $|F|_\infty\le c_fs_1/|\Omega|^{1/2}$ suffices.
\end{lemma}

\begin{proof}
$d$-consistency is used in exactly four places.  (u1) Step~2: without it,
$\Omega^\diamond_m=\rho(\omega^\diamond_m-d'_m(\omega^\diamond_m)w'_m)+
\kappa^\diamond_mw'_m+\rho(d^\diamond_m-d^\theta_m)(w'_m-w^0_m)$ with
$\kappa^\pm_m=\mp\frac\rho2\mathrm{mis}_m$ and $\kappa^\theta_m=0$.  (u2) Step~3:
$W_0=(1+\tau\kappa)(w'+t''(\omega^\diamond-d'(\omega^\diamond)w'))$ with
$t''=\tau\rho/(1+\tau\kappa)$ and $|\tau\kappa|\le\frac1{16}$, so the radius
conditions hold with $c_{\rm flat}$ replaced by $\frac{16}{15}c_{\rm flat}$
(absorbed by shrinking $c^f_{\rm flat}$), and $\lin'_m=\tau\kappa^\diamond_m+
r_m\ip{V^{\rm an}_m-w'_m}{R_mx'}/\sigma'_m$, so $|\lin'_m|\le\frac\rho2|\tau|s_1+
\frac\delta{64|I|}|\tau|s_1\le\tau^2$ for $|\tau|>s_1$; for $|\tau|\le s_1$ the
$\theta$ piece is used and $\lin'_m=0$.  Step~4 uses $\lin'_m$ only through
$|\lin'_m|\le\tau^2$.  (u3) In Lemma~\ref{lem:III-UE1}(c), $|d'_m(\omega)-
d_m(\omega)|\le|\Omega|(A_2/t)K_w\Pi$ and $|H'_m(\omega)-(C_m/C'_m)H_m(\omega)|\le
C_f|\Omega|A_2^2K_w\Pi/t^2$; in Step~5 this is a relative error at most
$C_f|\Omega|A_2^2K_w^2s_1/T^2\le\delta/32$ by (LATE-tol); gaps move by at most
$CK_w\Pi\le\gap/2$.  (u4) In Lemma~\ref{lem:III-UE1}(g), $g''_i-g_i$ acquires
$\sum_m(d_m-d'_m)(\omega^\theta)R_m^*w'_m$, of $\ell_1$-norm at most
$N|\Omega|A_2K_w\Pi/t_i$.  Finally $|\mathrm{mis}_{i,m}|\le(m/A'_m)\|D_m\Delta
\omega_{i,m}\|_2\|F|_{\Omega_m}\|_2$ because $\mathrm{nv}'_l-\mathrm{nv}_l=\epsilon_l
\lambda_lF_l/A'_m$.
\end{proof}

Theorem~\ref{thm:III-SysDC} gives exact zeros, so Lemma~\ref{lem:III-tol} is a
robustness statement and is not used below.

\subsection{Step A: cone realization and the structure of its obstructions}
\label{subsec:III-SysA}

Step~A of the route-V companion (Section~\ref{sec:III-routeV}) sets the ratios
of the active switching carriers to the \L ojasiewicz targets on the exact set
$Z_\varkappa$, holds the nearly neutral inactive carriers of (T)-blocks
exactly at $0$, and keeps the other coarse strict non-peaks at their ratios.
Downstream only \emph{ratios} are used: the class cone in ratio variables
((X3) of [II,~Definition~5.52]), the statuses (Lemma~\ref{lem:III-K}) and the
true shift.  So the block scalars may float, and we tune the ratios of all
coarse strict non-peaks (this is necessary as soon as a common-mode resource
is used, since a common mode moves every value at first order).

\begin{definition}[The Step-A system]\label{def:III-SysA}
Let $f_1$ be the row after the deep raise, Step~0 (in the variant described
before Definition~\ref{def:III-H1}) and the preparation.  Let
$\Lambda_A$ be the set of coarse strict non-peaks of $f_1$ together with the
carriers pushed inward by Step~A; for a block $m$ let $\varkappa_m:=
\varkappa(\hat\zeta_m;\Omega_m)$ be the block scalar of
Subsection~\ref{subsec:III-gen-setting}.  The \emph{target ratios}
$\mathrm{rt}_l$ are: $\epsilon_lr''_l$ for active carriers ($r''$ the targets
on $Z_\varkappa$), $0$ for the held nearly neutral carriers of (T)-blocks, the
ratio at $f_1$ for the other held carriers, and the pushed ratio for pushed
carriers.  The \emph{Step-A system} is $\tilde F(P)=0$,
$\tilde F_l:=\mathrm{val}_l-\mathrm{rt}_l\varkappa_{m(l)}$ ($l\in\Lambda_A$),
$\mathrm{val}_l=\epsilon_lu_l(\hat x)$; at $f_1$, $|\tilde F(f_1)|_\infty\le x$,
where $x:=C\,\mathrm{Des}(L)\eta'$ is the Step-A scale.
\end{definition}

While the peak set of block $m$ is fixed, $\varkappa_m$ is a $C^1$ function of
the block vector $\zeta_m=(\lambda_ku_k(\hat x))_k$, positively homogeneous of
degree $1$ (the threshold equation is invariant under $\zeta\mapsto c\zeta$,
$\bar\vartheta\mapsto c\bar\vartheta$).  Put $K_m:=\sum_k(\partial\varkappa_m/
\partial\zeta_m(k))\lambda_ku_k\in\ell_1$ (derivatives at the row considered),
$G^A_m:=\ip{U^*K_m}{e_1}$ and $\tilde g^A:=g-(\mathrm{rt}_lG^A_{m(l)})_{l}$.

\begin{lemma}[Euler combination]\label{lem:III-Euler}
For $y\in\R^{\Lambda_A}$ put $\beta^A_m:=\sum_{l\in\Lambda_A,\,m(l)=m}y_l
\mathrm{rt}_l$ and $\Psi^A_y:=\sum_ly_l\epsilon_lu_l-\sum_m\beta^A_mK_m$.  Then
\textup{(a)} $\Psi^A_y(\hat x_1)=y^T\tilde F(f_1)=:\varphi_y$, $|\varphi_y|\le
x\|y\|_1$; \textup{(b)} for every move $\delta$ of
Lemma~\textup{\ref{lem:III-actions}}, the first-order change of $y^T\tilde F$ is
$l_\delta(\Psi^A_y)$, up to $O(|\delta|^2)$; \textup{(c)}
$\ip{U^*\Psi^A_y}{e_1}=y^T\tilde g^A$.
\end{lemma}

\begin{proof}
Euler's identity gives $\varkappa_m(f_1)=\sum_k\zeta_m(k)\partial\varkappa_m/
\partial\zeta_m(k)=K_m(\hat x_1)$, whence (a); (b) is the chain rule, as
$\mathrm d\tilde F_l=\epsilon_l\mathrm du_l-\mathrm{rt}_l\,l_\delta(K_m)$; (c)
follows from the definitions.
\end{proof}

\begin{definition}[Generators, A-regularity, near-certificates]\label{def:III-Areg}
The \emph{generators} at $f_1$ are the normalized columns of the Step-A system
(their action on $y^T\tilde F$ is $l_\delta(\Psi^A_y)$ by
Lemma~\ref{lem:III-Euler}):
\emph{two-sided}: the levers (P1), (P3), (P4) of tuned carriers; TU levers
whose pull and base bank are pre-placed; one two-sided resource of a peak per
block (optional); one elementary move at every design-depth coordinate of type
TS (mass change), FR ($z$-move) and ZD (composite move);
\emph{one-sided} (allowed direction only): raises of S$_{\rm th}$ coordinates
carrying data.  Each two-sided generator contributes its column $g$ and $-g$
to the set $\mathcal G$, each one-sided one its column.  A \emph{pre-placement}
$\mathcal P$ is a set of discrete far moves, at most one per tuned carrier, at
the smallest admissible depth in $[x,x'_{\max}]$: the pulls of TU levers with
their compensating base banks.  Put
\begin{gather*}
 b:=-\tilde F(f_1+\mathcal P),\qquad W:=\mathrm{Des}(L)^{C_A}u^{-3}\max(x,\|b\|_2),\\
 \rho_0:=C_2W^2,\qquad Q:=\mathrm{conv}(\mathcal G\cup\{0\}),
\end{gather*}
with $C_2$ the second-order constant of the generators (a design constant).
The stage $(L,w,a)$ is \emph{A-regular} if $B(b,\rho_0)\subseteq(W/2)Q$ for some
$\mathcal P$.  A \emph{near-certificate} of size $\varepsilon$ is a unit
$y$ with $\max_{g\in\mathcal G}(-y^Tg)\le\varepsilon$; $\varphi_y:=
y^T\tilde F(f_1+\mathcal P)$.
\end{definition}

\begin{lemma}[Cone realization of Step A]\label{lem:III-cone}
Let $L\ge l_f$.
\begin{enumerate}
\item[(a)] If $\varepsilon_{\rm PS}:=\min_{\|y\|_2=1}\max_{g\in\mathcal G}(-y^Tg)\ge
u^3\mathrm{Des}(L)^{-C_1}$ and $C_A\ge C_1+2$, the stage is A-regular with
$\mathcal P=\emptyset$, for every right-hand side.
\item[(b)] If the stage is not A-regular for $\mathcal P$, there is a
near-certificate $y$ of size $\varepsilon<4u^3\mathrm{Des}(L)^{-C_A}$ with
$\varphi_y>-\rho_0$.
\item[(c)] If the stage is A-regular, Step~A is realized: there is a companion
$f^{(1a)}$ with the conclusions of Step~A of Section~\textup{\ref{sec:III-routeV}}
\textup{(}active ratios exactly $r''$ on $Z_\varkappa$, held values exactly
$0$, the threshold protection with the Puiseux margin, robust statuses
kept\textup{)}, all moves at most $W$, and $p^*(f^{(1a)}-f_1)\le
C_f\mathrm{Des}(L)^CW\log(1/W)$.  The coarse statuses at $f^{(1a)}$ are those
used in Lemma~\textup{\ref{lem:III-noBinf}}, and support signs are kept.
\end{enumerate}
\end{lemma}

\begin{proof}
(a) If some $b'$ with $\|b'\|<\varepsilon_{\rm PS}$ were outside $Q$, a separating
unit $y$ would give $y^Tb'<\min_gy^Tg\le-\varepsilon_{\rm PS}$, which is
impossible; so $B(0,\varepsilon_{\rm PS})\subseteq Q$.  Every $b'\in B(b,\rho_0)$
then lies in $((\|b\|+\rho_0)/\varepsilon_{\rm PS})Q\subseteq(W/2)Q$, since
$\|b\|/\varepsilon_{\rm PS}\le\mathrm{Des}^{C_1}u^{-3}\|b\|\le W/4$ and
$\rho_0/\varepsilon_{\rm PS}=C_2W^2\mathrm{Des}^{C_1}u^{-3}\le W/4$; the latter
holds because $\max(x,\|b\|)\le\mathrm{Des}^Cx'_{\max}$, so $W\le\mathrm{Des}^C
T^3/u$, and $T\le(u/(4\,\mathrm{Des}(L)))^{20}$ (Section~\ref{sec:III-design})
gives $4C_2W\mathrm{Des}^{C_1}\le u^3$ for $L$ large.

(b) Pick $b''\in B(b,\rho_0)$ outside the compact convex set $(W/2)Q$ and
separate: a unit $y$ with $y^Tb''<(W/2)\min(0,\min_gy^Tg)$.  Then $y^Tb<\rho_0-
(W/2)\max(0,\max_g(-y^Tg))$, so $\max_g(-y^Tg)<2(\rho_0+\|b\|)/W\le2C_2W+
2u^3\mathrm{Des}^{-C_A}\le4u^3\mathrm{Des}^{-C_A}$ for $L$ large, and
$\varphi_y=-y^Tb>-\rho_0$.

(c) \emph{Frozen statuses.}  Compute the block scalars with the coarse peak set
\emph{frozen} at $f_1+\mathcal P$: the threshold equations of [I,~Lemma~1.7]
with a fixed peak set form a smooth system whose Jacobian is bounded below by
$\Phi_{P_m}^2\ge D(L)^{-2}$, so the frozen map is $C^{1,1}$ on the parameter set;
fine carriers are computed truly (their status changes cost at most
$Cc_{L+1}\|t\|\ll\rho_0$).  This is necessary: at a threshold crossing of an
$\Omega$ carrier the derivative of the true block scalar jumps by an amount of
order one, so the quadratic remainder bound fails without freezing.
\emph{Brouwer.}  Write the frozen system as $\tilde F(f_1+\mathcal P)+
\mathcal Gt+\mathrm{Rem}(t)$ with $t\ge0$ indexed by $\mathcal G$ and
$\|\mathrm{Rem}(t)\|\le C_2\|t\|^2$.  For $b'\in\mathrm{cone}(\mathcal G)$ let
$S(b'):=\mathrm{argmin}\{\|s\|_2:\mathcal Gs=b',\ s\ge0\}$; $S$ is continuous on
$\mathrm{cone}(\mathcal G)$ (a strictly convex parametric quadratic program whose
feasible sets depend Hausdorff-continuously on $b'$ by Hoffman's bound).  On
$\mathcal K:=\{t\ge0:\|t\|_2\le W\}$ the map $\mathcal T(t):=S(b-\mathrm{Rem}(t))$
is defined and continuous, since $\|\mathrm{Rem}(t)\|\le\rho_0$ puts
$b-\mathrm{Rem}(t)$ into $B(b,\rho_0)\subseteq(W/2)Q$, which also gives
$\|\mathcal T(t)\|_2\le W/2$.  By Brouwer's theorem $t^*=\mathcal T(t^*)$ solves
the frozen system with $t^*\ge0$.  \emph{Statuses.}  At $t^*$ the ratios are the
targets, so the frozen relative positions computed by Lemma~\ref{lem:III-K}
satisfy $\varrho_d\le1-\mu/2$ for the near-threshold active carriers and the
held values vanish; robust statuses move by at most $\mathrm{Des}^CW/u\ll u$.
Hence the frozen statuses are consistent, the frozen block scalars solve the
true threshold equations, and by uniqueness of the threshold solution
[I,~Lemma~1.7] they are the true ones.  The Puiseux protection of
Section~\ref{sec:III-routeV} uses only the exactness of the realized ratios and
held values, not the size of the moves.  \emph{Costs.}  The moves have size at
most $W\le\mathrm{Des}^Cx'_{\max}u^{-3}$ in value units; the cost bound is
[II,~Lemma~2.12(iv)] together with the costs of the pre-placements
(Lemma~\ref{lem:III-costs}).  One-sided generators are used only in their allowed
direction, so supports keep their signs.
\end{proof}

\begin{lemma}[Bordered rank-one systems]\label{lem:III-border}
Let $g,h,w\in\R^n$, $\beta>0$ and $\mathsf N:=I-gh^T$.  If $|1-h^Tg|\ge\beta$ or
$|h^Tw|\ge\beta$, then every $b\in\R^n$ has a preimage under $(P,s)\mapsto
\mathsf NP+ws$, linear in $b$, with $\|P\|_\infty+|s|\le(1+\|g\|_\infty
\|h\|_1/\beta)(1+\|w\|_\infty)(1+\|h\|_1/\beta)\|b\|_\infty$.  Every unit
vector $y$ with $\|y^T\mathsf N\|_2\le\beta'\le\frac12$ is within $2\beta'$ of
$h/\|h\|_2$ or of $-h/\|h\|_2$.
\end{lemma}

\begin{proof}
If $|1-h^Tg|\ge\beta$: $s:=0$ and $P:=b+g\,h^Tb/(1-h^Tg)$ (Sherman--Morrison).
Otherwise $|h^Tw|\ge\beta$: put $s:=h^Tb/(h^Tw)$ and $P:=b-ws$; then $h^TP=0$ and
$\mathsf NP+ws=P+ws=b$.  The bounds are immediate.  For the last claim write
$y=th/\|h\|_2+y_\perp$ with $y_\perp\perp h$; then $y^T\mathsf N=y^T-(y^Tg)h^T$, whose
component orthogonal to $h$ is $y_\perp^T$, so $\|y_\perp\|\le\beta'$,
$|t|\ge(1-\beta'^2)^{1/2}\ge1-\beta'^2$, and $y$ is within $\beta'+\beta'^2\le
2\beta'$ of $\sgn(t)h/\|h\|_2$.
\end{proof}

\subsubsection*{The structure of Step-A obstructions}
Fix a unit $y\in\R^{\Lambda_A}$, work at $f_1+\mathcal P$ (all derivatives taken
there) and write $\Psi:=\Psi^A_y$.  Its coarse part is $\sum_{k\le L}
\mathrm{coef}_ku_k$, with $\mathrm{coef}_l=\epsilon_ly_l-\beta^A_{m}\lambda_l
\partial\varkappa_m/\partial\zeta_m(l)$ for a tuned carrier $l$ of block $m$ and
$\mathrm{coef}_k=-\beta^A_m\lambda_k\partial\varkappa_m/\partial\zeta_m(k)$ for an
untuned coarse carrier (a peak); its fine part has $\ell_1$-norm at most
$Cc_{L+1}$.  For $s\in S_k$ beyond $s_{\max}(L)$ (private to $k$ among coarse
carriers) $\Psi(s)=\mathrm{coef}_kv_k(s)$ up to the fine part.  Put
$\gamma:=\ip{U^*\Psi}{e_1}$ and $\kappa_\Psi:=\gamma/\nu_1$, and let $h$ be the
vector of signed (P4) coefficients ($h_l:=0$ for carriers without (P4)).  In
route V, $\Lambda_A$ contains every coarse strict non-peak, so the
derivative of $\varkappa_m$ with respect to a strict non-peak coordinate is
small ([II,~Subsection~5.6]) and the non-peak Euler share
$\tilde X_m:=\sum_{l\in\Lambda_A,\,m(l)=m}\mathrm{rt}_l\epsilon_l\lambda_l
\partial\varkappa_m/\partial\zeta_m(l)$ satisfies $|\tilde X_m|\le\frac14$.  Since
$\sum_{l}\mathrm{rt}_l\epsilon_l\mathrm{coef}_l=\beta^A_m(1-\tilde X_m)$ (sum over
the tuned carriers of block $m$), one gets $\|\mathrm{coef}_{\rm tuned}\|_2\ge
c_{\rm A}:=1/(2+CN|\Lambda_A|^{1/2}/A_{\min})$.  We write $\mathrm{tail}_D:=C_f2^{-L^2}
(|\Lambda_L|H(L))^{-4}+Cc_{L+1}$.

\begin{theorem}[Structure of Step-A obstructions]\label{thm:III-Aobs}
Let $L\ge l_f$ and let $y$ be a near-certificate of size $\varepsilon\le
\varepsilon_*:=u^3\mathrm{Des}(L)^{-C}$ at $f_1+\mathcal P$, $C$ a large absolute
constant.  Then:
\begin{enumerate}
\item[(1)] \textup{(no positive common mode)} if $\gamma>0$, then $\gamma\le
2(|\varphi_y|+\bar d(L)\varepsilon\mathrm{Des}(L)+\mathrm{tail}_D)$, where
$|\varphi_y|\le|\Lambda_A|^{1/2}\max(x,\|b\|_2)$;
\item[(2)] \textup{(action of the generators)} at every $s\le\bar d(L)$,
$|c_s(\Psi)|\le\varepsilon\mathrm{Des}(L)$ if $s$ carries a two-sided generator,
and $c_s(\Psi)\ge-\varepsilon\mathrm{Des}(L)$ if it carries a one-sided one.
Hence, for every coarse carrier $k$:
\begin{enumerate}
\item[(2a)] if $S_k$ has a private FR or ZD coordinate in $[1,\bar d(L)]$, then
$|\mathrm{coef}_k|\le\varepsilon\mathrm{Des}(L)^2$; if $k$ has a pre-placed TU lever
whose base bank has mass $\mathsf m_k$, then $|\mathrm{coef}_k|\le(\varepsilon
\mathrm{Des}(L)+|\kappa_\Psi|\mathsf m_k)\mathrm{Des}(L)$;
\item[(2b)] at every thick private $p\in S_k\cap[1,\bar d(L)]$:
$|\mathrm{coef}_kv_k(p)-\kappa_\Psi a_p|\le\varepsilon\mathrm{Des}(L)$;
\item[(2c)] at every S$_{\rm th}$ private $s'\in S_k\cap[1,\bar d(L)]$ carrying
data: $\mathrm{coef}_kz_{s'}\ge-(\varepsilon\mathrm{Des}(L)+2|\kappa_\Psi|b)
\mathrm{Des}(L)$;
\end{enumerate}
\item[(3)] \textup{(robustly negative common mode)} $\gamma\le-c_*$, where
\[
 c_*:=c_{\rm A}\min\{1/(4\|h\|_2),\ \nu_1/(8H(L))\}
\]
\textup{(}with $1/(4\|h\|_2):=\infty$ if $h=0$\textup{)};
\item[(4)] \textup{(proportional thick support)} at every $s\le\bar d(L)$ with
$|a_s|\ge u/2$: $|\Psi(s)-\kappa_\Psi a_s|\le\varepsilon\mathrm{Des}(L)$ while
$|\kappa_\Psi a_s|\ge c_*u/(2\nu_1)$; so $\Psi(s)\ne0$, $s\in U_L$ \textup{(}there is
no anchor\textup{)}, and $s$ lies in the support of a coarse carrier with a robust
coefficient;
\item[(5)] $|\Psi(s)|\le\varepsilon\mathrm{Des}(L)$ at every FR or ZD coordinate
$s\le\bar d(L)$;
\item[(6)] if $\varphi_y>-\rho_0$, the sum of $c_s(\Psi)$ over the design-depth
coordinates carrying one-sided generators is at least $c_*/2$.
\end{enumerate}
\end{theorem}

\begin{proof}
(2) By Lemmas~\ref{lem:III-Euler}(b) and~\ref{lem:III-actions}, a generator at
$s$ acts on $y^T\tilde F$ by $\mathrm{eff}_s$ times $D_\Psi(s)$ (mass change at a
support coordinate), $\Psi(s)$ ($z$-move at a free coordinate) or $\Psi(s)z_s$
(composite move), times its orientation, with $\mathrm{eff}_s\ge1/\mathrm{Des}(L)$;
in each case the action is $\mathrm{eff}_s$ times $c_s(\Psi)$ or $\Psi(s)$ with
$|c_s(\Psi)|\le|\Psi(s)|$ off $F$.  A near-certificate has $y^Tg\ge-\varepsilon$
for all $g\in\mathcal G$, which gives the two inequalities.  (2a): at a private
FR or ZD coordinate $s'$, $|\mathrm{coef}_k|v_k(s')=|\Psi(s')|\le\varepsilon
\mathrm{Des}(L)$ and $v_k(s')\ge v_{\min}(L)\ge1/\mathrm{Des}(L)$; a pre-placed
TU bank is a private two-sided mass generator at a support coordinate of mass
$\mathsf m_k$, where $D_\Psi=\mathrm{coef}_kv_k-\kappa_\Psi\mathsf m_k\,\sgn$.
(2b) is the TS generator at $p$.  (2c): $c_{s'}(\Psi)=(\mathrm{coef}_kv_k(s')-
\kappa_\Psi a_{s'})z_{s'}\ge-\varepsilon\mathrm{Des}(L)$, with $|a_{s'}|\le2b$ at
S$_{\rm th}$ coordinates of $f_1+\mathcal P$ and $v_k(s')\ge1/\mathrm{Des}(L)$,
$v_k>0$ on $S_k$.

(1) Apply \eqref{eq:III-identity} at $f_1+\mathcal P$ with $u:=\Psi$ and
$D:=\bar d(L)$: $\varphi_y=\kappa_\Psi(1-\alpha(D))+\sum_{s\le D}c_s(\Psi)+
\sum_{s>D}\Psi(s)z_s$.  By (2), every design-depth coordinate carrying a
generator has $c_s\ge-\varepsilon\mathrm{Des}(L)$; the remaining ones (the last
kind of Lemma~\ref{lem:III-prep}(d)) have $c_s=0$ up to the fine part, or
$c_s=-\kappa_\Psi|a_s|\ge-2|\kappa_\Psi|b$; the deep sum is at least
$-\mathrm{tail}_D$ (the coefficients of $\Psi$ are at most $C_f$).  If $\gamma>0$,
then $\kappa_\Psi(1-\alpha(D))\ge\kappa_\Psi\nu_1=\gamma$ and $2\bar d(L)
\kappa_\Psi b\le\gamma/2$, whence (1).  The bound on $|\varphi_y|$ is
Lemma~\ref{lem:III-Euler}(a) at $f_1+\mathcal P$.

(3) Put $\gamma_0:=c_{\rm A}/(4\|h\|_2)$ and suppose $\gamma>-\gamma_0$; by (1),
$|\gamma|\le\gamma_0$.  After the preparation, the off-$F$ or thin near
coordinates of zero-data carriers are ZD and those of data-carrying carriers
are S$_{\rm th}$ of the data sign $\sigma_l$ (Lemma~\ref{lem:III-prep}(c),
Lemma~\ref{lem:III-nodeep} case (S)).  So every tuned $l$ is of one of the kinds:
(i) $N_l$ contains a ZD coordinate, or $l$ has a pre-placed TU lever: then
$|\mathrm{coef}_l|\le\varepsilon\mathrm{Des}^2+|\kappa_\Psi|\mathsf m_l\mathrm{Des}$
by (2a); (ii) $S_l\cap[1,\bar d(L)]$ contains two thick private coordinates
$s,s'$ of robust spread: by (2b) at both, $|\kappa_\Psi||\mathfrak a_s-\mathfrak a_{s'}|
\le2\varepsilon\mathrm{Des}^2$, so $|\kappa_\Psi|\le8\varepsilon\mathrm{Des}^3/u^2$;
(iii) otherwise, $N_l$ contains a thick $p$: by (2b), $\mathrm{coef}_l=\kappa_\Psi
a_p/v_l(p)+O(\varepsilon\mathrm{Des}^2)=\epsilon_l\gamma h_l+O(\varepsilon
\mathrm{Des}^2)$; (iv) otherwise every near coordinate of $l$ is S$_{\rm th}$ of
the data sign, and $\mathrm{coef}_l\sigma_l\ge-(\varepsilon\mathrm{Des}+2|\kappa_\Psi|
b)\mathrm{Des}$ by (2c).  If a kind-(ii) carrier exists, $\kappa_\Psi$ is tiny,
which contradicts the conclusion $\kappa_\Psi\le-c_U/2$ reached below; so then no
near-certificate with $\gamma>-\gamma_0$ exists.  Otherwise
$\|\mathrm{coef}_{\rm (i)}\|_2+\|\mathrm{coef}_{\rm (iii)}\|_2\le|\gamma|\|h\|_2+
\mathrm{small}\le c_{\rm A}/4+\mathrm{small}$, hence $\|\mathrm{coef}_{\rm (iv)}\|_2\ge
c_{\rm A}/2$.  At an S$_{\rm th}$ near coordinate $s$ of a kind-(iv) carrier $l$,
$c_s(\Psi)\ge\mathrm{coef}_l\sigma_l/H(L)-2|\kappa_\Psi|b$ when $\mathrm{coef}_l
\sigma_l\ge0$, and $\sum_l\max(\mathrm{coef}_l\sigma_l,0)\ge c_{\rm A}/2-\mathrm{small}$.
With (2) for all other design-depth coordinates,
\[
 \sum_{s\le\bar d(L)}c_s(\Psi)\ge\frac{c_{\rm A}}{2H(L)}-\bar d(L)\big(\varepsilon
 \mathrm{Des}^2+2|\kappa_\Psi|b\big)-|\Lambda_A|\,\mathrm{small}\ge\frac{c_{\rm A}}{4H(L)}
 =:c_U .
\]
By \eqref{eq:III-identity}, $\kappa_\Psi(1-\alpha(D))\le|\varphi_y|-c_U+
\mathrm{tail}_D\le-c_U/2$, so $\kappa_\Psi\le-c_U/2$ and $\gamma=\kappa_\Psi\nu_1
\le-c_{\rm A}\nu_1/(8H(L))$.  In both cases $\gamma\le-c_*$.

(4) The TS generator at $s$ is two-sided: $|c_s(\Psi)|=|\Psi(s)-\kappa_\Psi a_s|\le
\varepsilon\mathrm{Des}$, and $|\kappa_\Psi|\ge c_*/\nu_1$ by (3).  If $s\notin U_L$,
$\Psi(s)$ is the fine part, which is tiny.  (5) FR and ZD generators are
two-sided with action $\Psi(s)$ times the orientation.  (6) By
\eqref{eq:III-identity}, $\sum_{s\le D}c_s(\Psi)=\varphi_y-\kappa_\Psi(1-\alpha(D))-
(\text{deep})\ge-\rho_0+|\kappa_\Psi|\nu_1-\mathrm{tail}_D\ge c_*-\rho_0-
\mathrm{tail}_D$; two-sided coordinates contribute at most $\bar d(L)\varepsilon
\mathrm{Des}$ in absolute value and the last kind of
Lemma~\ref{lem:III-prep}(d) at most $2\bar d(L)|\kappa_\Psi|b$; the rest is the
one-sided sum.
\end{proof}

\begin{corollary}[Sufficient conditions for A-regularity]\label{cor:III-Asuff}
Let $L\ge l_f$, $w$ clean, $a$ a class, after the preparation, with $\mathcal P$
the TU pre-placements.  The generators positively span with $\varepsilon_{\rm PS}
>\varepsilon_*$ -- so the stage is A-regular for every right-hand side by
Lemma~\textup{\ref{lem:III-cone}(a)} -- if one of the following holds:
\begin{enumerate}
\item[(a)] every tuned carrier is privately tunable;
\item[(b)] there is an anchor: a thick $s\le\bar d(L)$ outside $U_L$;
\item[(c)] some coarse carrier $k$ \textup{(}tuned, or a peak\textup{)} has, in
$S_k\cap[1,\bar d(L)]$ and private to $k$, a thick coordinate together with an FR or
ZD coordinate \textup{(}or $k$ has a pre-placed TU lever\textup{)};
\item[(d)] some coarse carrier $k$ has, in $S_k\cap[1,\bar d(L)]$ and private to
$k$, a thick coordinate $p$ and an S$_{\rm th}$ coordinate $s'$ carrying data,
with $\sgn a_p=z_{s'}$;
\item[(e)] some coarse carrier $k$ has two thick private coordinates in
$S_k\cap[1,\bar d(L)]$ of robust spread.
\end{enumerate}
Moreover the stage is A-regular for the given right-hand side if
\begin{enumerate}
\item[(f)] every tuned carrier is of kind (i) or (iii) of the proof of
Theorem~\textup{\ref{thm:III-Aobs}}, and $|1-h^T\tilde g^A|\ge u$, or
$|(1-\tilde X_m)\lambda_ch^T\mathrm{rt}^{(m)}+h_ch^T\tilde g^A|\ge u$ for a two-sided
peak resource $c$ of a block $m$ \textup{(}$h_c$ its common-mode coefficient,
$\mathrm{rt}^{(m)}$ the target vector of block $m$\textup{)}, or
$\varphi_{y_+}\ge2\rho_0+C\mathrm{Des}(L)^{-C_A}\|\tilde F(f_1)\|_2$ for
$y_+:=h/\|h\|_2$.
\end{enumerate}
\end{corollary}

\begin{proof}
Suppose a near-certificate of size $\varepsilon\le\varepsilon_*$ exists;
Theorem~\ref{thm:III-Aobs} applies.  (a) The private two-sided columns are
$e_l$ up to small errors and span.  (b) $\Psi(s)$ is tiny at an anchor, against
(4).  (c) $|\mathrm{coef}_k|$ is tiny by (2a), so $\Psi(p)=\mathrm{coef}_kv_k(p)$
is tiny at the thick $p$, against (4).  (d) By (2b) and (3), $\mathrm{coef}_k
v_k(p)=\kappa_\Psi a_p+O(\varepsilon\mathrm{Des})$ with $|\kappa_\Psi a_p|$ robust
and $\kappa_\Psi<0$, so $\sgn\mathrm{coef}_k=-\sgn a_p$ and $|\mathrm{coef}_k|$ is
robust; (2c) then forces $z_{s'}=\sgn\mathrm{coef}_k=-\sgn a_p$.  (e) This is
kind (ii) in the proof of (3): no near-certificate exists.  So in (a)--(e)
$\varepsilon_{\rm PS}>\varepsilon_*$.  (f) Without kind-(iv) carriers the proof of
(3) shows that every near-certificate has $\gamma\le-\gamma_0$.  The two-sided
columns form the square system $I-\tilde g^Ah^T$ (up to small errors) bordered
by the peak columns $-(1-\tilde X_m)\lambda_c\mathrm{rt}^{(m)}-h_c\tilde g^A$, and a
near-certificate $y$ is a near-null left vector of all of them:
$\|y^T(I-\tilde g^Ah^T)\|_2\le C\varepsilon_*|\Lambda_A|^{1/2}$.  By
Lemma~\ref{lem:III-border}, $y$ is within $C\varepsilon_*|\Lambda_A|^{1/2}$ of
$\pm y_+$.  Now $y_+^T(I-\tilde g^Ah^T)=(1-h^T\tilde g^A)h^T/\|h\|_2$ and
$-\|h\|_2\,y_+^T\mathrm{col}_c=(1-\tilde X_m)\lambda_ch^T\mathrm{rt}^{(m)}+
h_ch^T\tilde g^A$; since $\|h\|_2\le|\Lambda_A|^{1/2}H(L)/\nu_1\le\mathrm{Des}(L)^C$,
either scalar being $\ge u$ gives a two-sided column $g$ with $|y^Tg|\ge
u\,\mathrm{Des}(L)^{-C}-C\varepsilon_*\mathrm{Des}(L)^C>\varepsilon_*$, which is
impossible.  So no near-certificate exists.  Otherwise $+y_+$ has $\gamma=h^T\tilde g^A/
\|h\|_2>0$, excluded by (3), so the only candidate is close to $-y_+$, and it is
an obstruction (Lemma~\ref{lem:III-cone}(b)) only if $\varphi_{-y_+}>-\rho_0$ up to
its deviation from $-y_+$, i.e.\ only if $\varphi_{y_+}<\rho_0+C\mathrm{Des}(L)^{
-C_A}\|\tilde F(f_1)\|_2$.
\end{proof}

\begin{remark}[Further cases and limits of the corollary]\label{rem:III-Asuff}
(a) \emph{Reduction (not a sufficient condition by itself).}  If every block
containing kind-(iv) carriers has a private two-sided peak resource, no
kind-(iii) carrier with nonzero target, and kind-(iv) carriers with
$|\mathrm{rt}_l|\ge u$ and $\mathrm{rt}_l\epsilon_l\sigma_l$ of one sign, then every
near-certificate has $\gamma\le-\gamma_0$: in such a block the peak coefficient
is $O(\varepsilon\mathrm{Des}^2)$, hence $\beta^A_m=O(\varepsilon\mathrm{Des}^2)$, and
the one-signed terms $\mathrm{rt}_l\epsilon_l\mathrm{coef}_l$ are each
$O(\varepsilon\mathrm{Des}^2/u)$, so $\|\mathrm{coef}_{\rm (iv)}\|$ is small.  A-regularity
then still needs the criterion of (f).
(b) \emph{Held carriers in a degenerate direction.}  A carrier held at value $0$
is zero-data.  If its near set contains an off-$F$ or thin coordinate, it is ZD
after the preparation and has kind (i); if its near set is thick it has kind
(ii) (then there is no near-certificate) or (iii).  So held carriers in a
degenerate direction are not a separate obstruction: they are either realized
or produce a near-certificate of Theorem~\ref{thm:III-Aobs}.
(c) \emph{Far moves against a certificate (sketch).}  This item is a sketch and
is not used.  A far placement at $j\in S_k$, $x\le v_k(j)\le x'_{\max}$, of
support of sign $\sigma$ acts on $\Psi(\hat x)$ by about $\mathrm{coef}_kv_k(j)
(\sigma-z_j)$; pulls (placements of sign $-\epsilon_k$) break a certificate when
$\epsilon_k\mathrm{coef}_k>0$, and placements of sign $+\epsilon_k$ would break it
when $\epsilon_k\mathrm{coef}_k<0$, provided such placements are admissible (the
mirror image of the pulls) and no new certificate appears after an overshooting
pre-placement.  Neither point is proved here.
\end{remark}

\subsection{The design for arbitrary base support}\label{subsec:III-designinf}

\begin{definition}[The design $T^\infty_{\rm fin}$]\label{def:III-designinf}
$T^\infty_{\rm fin}$ is the design $T_{\rm fin}$ of Section~\ref{sec:III-design}
(option (a)) over the diagonal base $\mu^\star_s=2^{-2^{2^s}}$, with the
following insertions at every stage $L$.
\begin{enumerate}
\item[(I1)] After the carriers of stages $\le L$, $s_{\max}(L)$ and
$\delta_{\min}(L)$ are fixed: $d(L)$, $H(L)$, $\Lambda_L$, $U_L$, $\bar d(L)$ and
$v_{\min}(L)$ as in Subsections~\ref{subsec:III-cushion}
and~\ref{subsec:III-levers}.
\item[(I2)] $\mathrm{Des}(L)$ is replaced by $\mathrm{Des}(L)\cdot
\big(2^{\bar d(L)}\bar d(L)H(L)|\Lambda_L|/((\mu^\star_{\bar d(L)})^2v_{\min}(L))
\big)^{C_D}$ with an absolute $C_D\ge2$.
\item[(I3)] The rate scheme $\omega(L)$ is enlarged by: the moduli $|a_s|$ and
the rooms $1-|z_s|$ for $s\le\bar d(L)$; the minors and cofactor polynomials
of the enlarged class cones (with the contact rows (X1$^F$)) for all patterns
and classes of level $L$.
\item[(I4)] For every sub-window $w$ of $L$, $\theta_0(w):=b(w)/\mathrm{Des}(L)^{C_P}$
with an absolute $C_P\ge C_D+4$.
\item[(I5)] In the conditions (D1), (D2) of Section~\ref{sec:III-design} that
define the admissible interval $(0,b_{X1}(w)]$ of $b(w)$, the number $b$ is
replaced by $2\mathrm{Des}(L)^Cb/u(w)$ in (D1) and by $2b$ in (D2).
\end{enumerate}
\end{definition}

\begin{lemma}[The insertions are admissible]\label{lem:III-designinf}
\textup{(a)} Every inserted quantity of stage $L$ is a function of quantities
fixed earlier in the recursion of Section~\textup{\ref{sec:III-design}}, so the
recursion stays well founded.  \textup{(b)} All inserted quantities are finite
and positive.  \textup{(c)} No rule refers to $N$ or to a row; the number
$2\bar d(L)$ of new rate objects is $N$-free.  \textup{(d)} The insertions only
enlarge $\mathrm{Des}(L)$, $\omega(L)$, $Q(w)$ and $n(w)$, and only lower
$T_{\rm lo}(w)$, $\eta'(w)$ and $b(w)$; every window condition used in Parts~II
and~III is monotone under these changes, the $\omega(L)+1$ sub-windows of a
main stage still have pairwise disjoint bands, and for every row some
sub-window of every main stage $L\ge l_f$ is clean.  \textup{(e)}
$2^{\bar d(L)+1}\le\mathrm{Des}(L)$, $1/v_{\min}(L)\le\mathrm{Des}(L)$, and a mass
change at any $s\le\bar d(L)$ has normalized efficiency at least
$1/\mathrm{Des}(L)$.
\end{lemma}

\begin{proof}
(a) $d(L)$, $H(L)$, $\bar d(L)$, $U_L$ and $v_{\min}(L)$ depend only on the
carriers of stages $\le L$ and on $\delta_{\min}(L)$; $\mathrm{Des}(L)$ on these;
$\omega(L)$ on $\bar d(L)$ and the patterns of level $L$; $\theta_0(w)$ on $b(w)$
and $\mathrm{Des}(L)$; and the admissible interval of $b(w)$ on the same data as
before together with $u(w)=b(w^-)$, which is fixed before $b(w)$.  Both sides of
the modified (D1), (D2) are monotone in $b$, so the admissible $b$ still form an
interval $(0,b_{X1}(w)]$.  (b) $\bar d(L)<\infty$ because finitely many $\ell_1$
vectors have tails tending to $0$; $v_{\min}(L)>0$ is a minimum of finitely many
nonzero numbers; $\mu^\star_{\bar d(L)}>0$.  (c) is clear; for $p_N$ the objects are
evaluated on the blocks present.  (d) [II,~Lemma~3.22] and the remarks after it:
every window condition is an upper bound on rate objects, weights and windows or a
lower bound on $\mathrm{Des}$ and $Q(w)$; the disjointness of the bands
($u(w^+)=b(w)<u(w)$) and the existence of clean sub-windows
([II,~Theorem~3.19]) do not depend on the size of $\omega(L)$.  (e) follows from
(I2), since $\nu\le1$.
\end{proof}

\begin{lemma}[Clean dichotomy at depth $\bar d(L)$]\label{lem:III-depthclean}
Let $w$ be clean for $f$.  For every $s\le\bar d(L)$, at $f$: $|a_s|\le b$ or
$|a_s|\ge u$, and $1-|z_s|\le b$ or $1-|z_s|\ge u$.  Along the construction
$f\to f^r\to$ Step~0 $\to$ preparation $\to$ Step~A $\to$ pre-row $\to$
approximant, every $s\le\bar d(L)$ keeps its type in the sense of
Lemma~\textup{\ref{lem:III-prep}(d)}.
\end{lemma}

\begin{proof}
The moduli and rooms at depth $\le\bar d(L)$ are rate objects (I3), so
cleanness gives the dichotomy.  The deep raise acts beyond $J(w)>\bar d(L)$ and
divides the moduli at depth $\le\bar d(L)$ by $\lambda_r\le2$
(Lemma~\ref{lem:III-deepraise}).  Step~0 (in the variant described before
Definition~\ref{def:III-H1}) closes rooms and pushes weak peaks,
moving values by at most $\mathrm{Des}^Cb/u\ll u$ (Definition~\ref{def:III-H1}).
The preparation moves values by at most $2b$ (Lemma~\ref{lem:III-prep}(b)).  Step~A moves TS and FR
coordinates by at most $\mathrm{Des}^CW\ll u$ and moves S$_{\rm th}$ coordinates
with data only upward (Lemma~\ref{lem:III-cone}).  The pre-row moves
design-depth coordinates by at most $K_w\Pi\ll\theta_0$ and the levers by at most
$\mathrm{Des}(L)K_{\rm DC}r\le\theta_0/4$ (Lemma~\ref{lem:III-late}).
\end{proof}

\begin{lemma}[Lateness conditions]\label{lem:III-late}
Let $s_1:=\max(64\rho\epsilon/\delta,e^{-1/(2T)})$, as in
Proposition~\textup{\ref{prop:III-S2c}}.  For $T$ small \textup{(}i.e.\ $L$
large\textup{)} the following hold, and they are the conditions \textup{(L-K2)}
used above:
\begin{enumerate}
\item[(L1)] $s_1\le\theta_0T/(64\rho A_0)$ \textup{(}so no S$_{\rm th}$
coordinate is in the raise set $W_F$ of Definition~\textup{\ref{def:III-approxinf}}\textup{)};
\item[(L2)] $\mathrm{Des}(L)K_{\rm DC}r\le\theta_0/4$;
\item[(L3)] $C_{\rm lev}K_{\rm DC}r\le\frac12\min(\gap_{\min},\text{robust
margins})$;
\item[(L4)] $\varepsilon_E\le\varepsilon_{\rm DC}/8$ and $C_2K_{\rm DC}^2r\le\frac14$;
\item[(L5)] $K_{\rm DC}s_1\le r_0$ and $K_w s_1\le r_0$.
\end{enumerate}
\end{lemma}

\begin{proof}
$b(w)$ is bounded below by a power of $\eta'(w)=T^4/(L\,\mathrm{Des}(L))$ whose
exponent and constants are bounded by powers of $\mathrm{Des}(L)$ (the
definition of $b_{X1}(w)$ through (D1), (D2) and the modulus of the threshold
protection), and $\mathrm{Des}(L)\le n(w)\le\log_2(1/T)$; hence $\log(1/b)\le
C\log^3(1/T)$ and $\theta_0\ge\exp(-C\log^3(1/T))$.  On the other side
$r\le\exp(C\log^2(1/T))s_1$, $K_{\rm DC}\le C_{f,N}\log^{C}(1/T)$, $C_{\rm lev},
K_w\le\exp(C\log^2(1/T))$, $\varepsilon_E\le C_f\mathrm{Des}^C(K_w\Pi+c_{L+1})$
with $\Pi\le K_ws_1$ and $c_{L+1}\le e^{-1/T}$ by (W$_{\exp}$), while
$\varepsilon_{\rm DC}\ge\mathrm{Des}(L)^{-C}$ and $\gap_{\min}$, the robust margins
and $r_0$ are at least $\exp(-C\log^2(1/T))$.  Since $s_1\le(1+64\rho C/\delta)
e^{-1/(2T)}\le\exp(-C'\log^3(1/T))$ for $T$ small, all five conditions hold.
\end{proof}

\begin{lemma}[Costs]\label{lem:III-costs}
Let $c_{\rm flat}(w)\ge c_fu(w)$.  The companion moves satisfy, with
$T=T_{\rm lo}(w)$:
\begin{enumerate}
\item[(a)] deep raise: $\varepsilon'_r\le\theta_wT^2$ (choice of $J(w)$);
\item[(b)] Step~0 and the preparation: at most $C_f\mathrm{Des}(L)^C(b/u)
\log(1/b)+C_f\mathrm{Des}(L)b\log(1/b)$ \textup{(}Lemma~\ref{lem:III-prep}(a)\textup{)};
\item[(c)] Step~A: at most $C_f\mathrm{Des}(L)^C\,W\log(1/W)\le
C_f\mathrm{Des}(L)^C\,T^3\log(1/T)/u$;
\item[(d)] the pre-placements and the lowering of far coordinates in $F$:
at most $C\,\mathrm{Des}(L)|\Lambda|x'_{\max}\log(1/T)$, i.e.\ at most
$Cu^2T^3\log(1/T)/\mathrm{Des}(L)$.
\end{enumerate}
Writing the sum as $\theta'_wT^2$, $\theta'_w/c_{\rm flat}(w)^2\to0$ along
main stages.
\end{lemma}

\begin{proof}
(a) is Lemma~\ref{lem:III-deepraise}.  (b)--(d): the cost of a family of moves
is bounded by [II,~Lemma~2.12(iv)] (base change plus $C_fc(\delta)$ with
$c(\delta)\le C\epsilon\log(e/\epsilon)$ for value moves $\epsilon$), the raises
by [II,~Lemma~2.15]; a pull at depth $v$ costs at most $C_f\lambda_lv\log(e/v)$
for the $z$-flip, $24\lambda_lv$ for its mass and $2\,\mathrm{Des}(L)\nu v$ for
the compensating bank, and $\log(e/v)\le6\log(1/T)$ since $v\ge x\ge T^5$;
pulls are taken at the smallest pullable depth in $[x,x'_{\max}]$, and lowering
an in-$F$ far coordinate costs at most $2(1+\|U\|)|a_j|\le2(1+\|U\|)x'_{\max}$.
Dividing by $c_f^2u^2T^2$: (a) gives $\theta_w/(c_fu)^2\le T^{9/10}/c_f^2\to0$ since
$\theta_w=T$ and $u\ge T^{1/20}$; (b) gives at most $\mathrm{Des}^C
T^2\log^3(1/T)/u^3$; (c) at most $\mathrm{Des}^CT\log(1/T)/u^3$; (d) at most
$CT\log(1/T)/c_f^2$; all tend to $0$ because $b\le T^4$ and $T\le
(u/(4\,\mathrm{Des}(L)))^{20}$.
\end{proof}

\begin{remark}\label{rem:III-costs}
A far depth larger than $x'_{\max}$ (for instance $T^2/(L\,\mathrm{Des}(L)|\Lambda|)$)
is not admissible: the zeroth-order cost of pulls and lowerings would then be of
order $T^2\log(1/T)/L$, which is not $o(u^2T^2)$, and $\theta'_w/c_{\rm flat}^2$
would diverge.  The window $[x,x'_{\max}]$ still has ratio at least $2^{6G(L)}$,
so it contains at least three elements of every $S_l$.
\end{remark}

\begin{lemma}[Window bookkeeping]\label{lem:III-bookkeeping}
Let $C_+:=5$ and $\gamma_B(w):=\min_mM_m\,u(w)/4$, so that
$c_{\rm flat}(w)\ge c_f\,u(w)$.  Then
\[
 K_{\rm data}(w)\le C_f^L\big(\mathrm{Des}(L)/u(w)\big)^{CL}(\log_2n(w))^{L+1},
 \qquad D_{\rm cl}(w)\le(2L+1)3^{2L},
\]
and for a class
set $\mathcal S$ with $|\mathcal S|\ge n(w)/D_{\rm cl}(w)$,
\[
 \frac{|\mathcal S|\,c_{\rm flat}(w)}{K_{\rm data}(w)}\to\infty,\qquad
 \frac{K_{\rm data}(w)T_{\rm hi}(w)}{|\mathcal S|}\to0
\]
along main stages $L\to\infty$.
\end{lemma}

\begin{proof}
$J\le2L$ for $L\ge l_f$ [II,~Lemma~5.53], and $K_*(w)\le C_f(\mathrm{Des}/u)^C
(\mathrm{Des}(L)+2\log_2(Cn(w)))^{1/2}$ by Lemma~\ref{lem:III-deepraise}; so
$K_{\rm data}=(JK_*(w)^{J+1}+C_f)K$ has the stated bound.  By the window count
of Section~\ref{sec:III-design}, $n(w)\ge L2^{L^3}Q(w)$ with $Q(w)\ge
(4\,\mathrm{Des}(L)/u(w))^{\omega(L)+L^2+20}$.  For $L$ large,
$(\mathrm{Des}/u)^{CL+1}\le Q(w)^{1/4}$ (as $L^2\ge4(CL+1)$), $C_f^LD_{\rm cl}(w)\le
2^{L^3/2}$, and $(\log_2n)^{L+1}\le n^{1/4}$ because $\log_2n\ge L^3$.  Hence
$|\mathcal S|c_{\rm flat}/K_{\rm data}\ge c_fu(w)n(w)^{3/4}/(2^{L^3/2}Q(w)^{1/4})
\ge c_fn(w)^{1/4}\to\infty$ (as $n(w)\ge2^{L^3}Q(w)$ and $Q(w)^{1/4}\ge1/u(w)$) and
$K_{\rm data}T_{\rm hi}(w)/|\mathcal S|\le K_{\rm data}D_{\rm cl}/n(w)\le n(w)^{-1/2}
\to0$.
\end{proof}

\subsection{The transport theorem}\label{subsec:III-transport}

\emph{The variant of Step~0 used at arbitrary base support.}  The
companion of Proposition~\ref{prop:III-KNtr} does not push the active weak
peaks (Step~0 there consists of the closing moves only, and the active weak
peaks enter the pattern's $\Omega$ under the floor convention).  At arbitrary
base support we use the variant of Step~0 in which, after the closing moves,
every active weak peak $k$ of a (T)-block ($\varrho_k\in[1,1+b]$) is pushed
to $\varrho_k\le1-b$ by moves of total size $\le\mathrm{Des}(L)^Cb/u$
(condition (H1$^*$) below), so that it becomes a strict non-peak.  Steps~1--3
of the proof of Proposition~\ref{prop:III-KNtr} hold for this variant: at the
pushed row one has $P\supseteq P_{\rm rob}$ and $\Omega\subseteq Q$, so
Lemma~\ref{lem:III-floor} gives $F_\Phi(r_1)\le1+C_{**}b+\mathrm{Des}^{C}b/u$,
which is absorbed in $\delta_0$ after the enlargement of the radius of (D1)
by (I5) of Definition~\ref{def:III-designinf}.  The pushes are needed here,
and not in Proposition~\ref{prop:III-KNtr}, for two reasons of this section:
Lemma~\ref{lem:III-cone}(c) computes the block scalars with the coarse
statuses frozen at $f_1+P$, which is consistent only if the near-threshold
active carriers are strict non-peaks there; and Theorem~\ref{thm:III-Aobs}
uses $|\tilde X_m|\le\frac14$, which needs every tuned carrier to be a strict
non-peak of $f_1$ (the derivative of $\varkappa_m$ in a peak coordinate is of
order one).

\emph{Orientation.}  The push must lower $\mathrm{vs}_ku_k(\hat x)=
\mathrm{vs}_k\epsilon_k\mathrm{val}_k$.  By [I,~Lemma~8.27], at an active weak
peak $k$ the route-V data $\gamma_k=\lambda_k(\Delta\omega(k)-\Delta_mw(k))$,
with $w(k)=M\mathrm{vs}_k$ and $\Delta'_m=\Delta_mM$, satisfy $\mathrm{vs}_k
\Delta\omega(k)=e_k:=|\omega_+(k)|+|\omega_-(k)|\ge0$; hence
$\mathrm{vs}_k\gamma_k=\lambda_k(e_k-\Delta'_m)$, that is,
$\epsilon_k=\mathrm{vs}_k\sgn(e_k-\Delta'_m)$.  (The formula
$\epsilon_k=-\sgn(\Delta''_m)\mathrm{vs}_k$ holds at robust peaks, where
$\Delta\omega=0$; it can fail at a weak peak when $0<\Delta'_m<e_k$.)  A far
pull at a contact of sign $\epsilon_k$ lowers $\mathrm{val}_k$, and a near
bank raises $\mathrm{val}_k$.  Hence a far pull pushes inward iff
$\epsilon_k=\mathrm{vs}_k$ ($e_k>\Delta'_m$), a near bank pushes inward iff
$\epsilon_k=-\mathrm{vs}_k$ ($e_k<\Delta'_m$), and if $e_k=\Delta'_m$ the peak
carries no data and the composite move (P1) is available.

\begin{definition}[The push condition \textup{(H1$^*$)}]\label{def:III-H1}
At a main stage $L$, a clean sub-window $w$ and a class $a$, condition
\textup{(H1$^*$)} holds if there are admissible push moves -- pulls at pullable
far coordinates of depth in $[2b\bar\vartheta_m\Phi_k/m,\mathrm{Des}(L)^Cb/u(w)]$,
near banks of weak peaks with $\epsilon_k=-\mathrm{vs}_k$, resources (P1), (P3),
(P4), and moves of Lemma~\ref{lem:III-actions} at design depth admissible for
their data -- of total size at most $\mathrm{Des}(L)^Cb/u(w)$, after which
every active weak peak of every (T)-block has $\varrho\le1-b$, the block
scalars being computed with the coarse statuses frozen as in
Lemma~\ref{lem:III-cone}(c).  In particular, (H1$^*$) holds if the Step-0
system, with the needed decreases as right-hand side, is A-regular in the
sense of Definition~\ref{def:III-Areg} with $x$ replaced by $2b$ and
$W\le\mathrm{Des}(L)^Cb/u(w)$ (Brouwer argument of
Lemma~\ref{lem:III-cone}(c), statuses frozen); the converse need not hold,
since a feasible inequality system need not contain a ball around its
right-hand side.
\end{definition}

A condition for each weak peak separately does not suffice: a (P4) push of a
weak peak $k$ changes every value at first order by $-g_{k'}h_k$ per unit, with
$|h_k|$ up to $H(L)/\nu$, so with several such pushes the system on the
weak-peak rows is $I-g_Wh_W^T$, which is singular when $1-h_W^Tg_W=0$; already
for one peak a (P4) push moves its own value by $(1-g_kh_k)$ per unit.  If the
left null vector of $I-g_Wh_W^T$ has components of one sign in the inward
orientation, the push inequalities are infeasible although every peak has its own
resource.

\begin{lemma}[Sufficient conditions for \textup{(H1$^*$)}]\label{lem:III-H1suff}
\textup{(H1$^*$)} holds in each of the following cases:
\begin{enumerate}
\item[(i)] every active weak peak $k$ of a (T)-block has a \emph{private} push
resource of the correct orientation: if $\epsilon_k=\mathrm{vs}_k$ (i.e.\
$e_k>\Delta'_m$), a pullable far coordinate $j\in S_k$ with $v_k(j)\in
[2b\bar\vartheta_m\Phi_k/m,\mathrm{Des}(L)^Cb/u(w)]$; if $\epsilon_k=
-\mathrm{vs}_k$ ($e_k<\Delta'_m$), a near bank; if $e_k=\Delta'_m$, the composite
move (P1); or, in each case, (P3);
\item[(ii)] the active weak peaks $W_4$ that are pushed by (P4) at distinct
private thick near coordinates satisfy $|1-h_{W_4}^Tg_{W_4}|\ge u$; the
active weak peaks $W_s$ pushed by shared design-depth moves are such that
the matrix of first-order actions of the chosen moves on the rows of $W_s$
(private and common-mode parts) has a right inverse of norm
$\le\mathrm{Des}(L)^C/u(w)$; and all other active weak peaks have private push
resources as in (i).
\end{enumerate}
\end{lemma}

\begin{proof}
(i) By the orientation discussion before Definition~\ref{def:III-H1}, each
listed resource moves $\mathrm{vs}_ku_k(\hat x)$ downwards.  The needed
decrease is at most $2b\bar\vartheta_m\Phi_k/m$ in value units (from
$\varrho_k\le1+b$ to $\varrho_k\le1-b$).  A far pull at depth $v$ moves the
value by $2v$, so the pull of depth $v_k(j)\in[2b\bar\vartheta_m\Phi_k/m,
\mathrm{Des}(L)^Cb/u]$ realizes the push with a move of size $\le\mathrm{Des}
(L)^Cb/u$; near banks and (P1), (P3) have continuous size and are used with
size $\le2b\bar\vartheta_m\Phi_k/m$.  The pull mass $24\lambda_kv_k(j)$ acts on
the other values only through $e$, with efficiency $(\mu^\star_j)^2$, which
is negligible for $j>s_{\max}(L)$; banks and raises act on other values only
at second order (Lemma~\ref{lem:III-masses}(i) and [II,~Lemma~4.15]), by at
most $C\,\mathrm{Des}^C(b/u)^2\ll b$ in relative position.  So pushing each
peak to $\varrho\le1-2b$ realizes Step~0 with moves of total size
$\le\mathrm{Des}(L)^Cb/u$, which keeps the {\L}ojasiewicz start inside the
radius of (D1) as enlarged by (I5) of Definition~\ref{def:III-designinf}; an
overshoot is bounded by the same size, so floor statuses and tiny minors at
$r_1$ move by $O(\mathrm{Des}^Cb/u)$, which (I5) accounts for.  (ii) For the
peaks of $W_4$ the column of a (P4) push at a private thick near coordinate
is $e_k-h_kg$ (Lemma~\ref{lem:III-actions}), so their rows form the system
$I-g_{W_4}h_{W_4}^T$ up to second-order terms; by the Sherman--Morrison
formula $\|(I-g_{W_4}h_{W_4}^T)^{-1}\|\le1+\|g_{W_4}\|\|h_{W_4}\|/u$.  For the
peaks of $W_s$ the system is not rank one (a shared coordinate has
first-order private entries for every peak whose support contains it), and
the assumed right inverse is used instead.  In both cases the rows are
solved exactly with moves of size $\le\mathrm{Des}^Cb/u$; the private pushes of
the other peaks, oriented as in (i), then absorb the first-order cross
effects of these moves one-sidedly (the system is block triangular).
\end{proof}

\begin{theorem}[Transport]\label{thm:III-transport}
Let the design be $T^\infty_{\rm fin}$, $N\ge1$, $f\in S_{p_N^*}$ with base
support finite or infinite, $g\in\Cset(f)$ and $\rho\in(0,1)$.  Suppose that
for infinitely many main stages $L$ there are a clean sub-window $w$ of $L$ and
an activity class $a$ \textup{(}thresholds $K_*(w)^i$, no cube\textup{)} with at
least $n(w)/D_{\rm cl}(w)$ scales such that
\begin{enumerate}
\item[(H1$^*$)] the push condition of Definition~\textup{\ref{def:III-H1}} holds,
and
\item[(H2)] Step~A is A-regular after Step~0 and the preparation
\textup{(}Definition~\textup{\ref{def:III-Areg})}.
\end{enumerate}
Then $(f,\rho g)\in\cl\NA((c_0,p_N),\ell_2^2)$, along pairs whose first rows are
norm attaining with finite base support; in particular
\textup{(LSC-trunc)}$_N$ holds at $(f,g,\rho)$.  If the hypothesis holds for every
$g\in\Cset(f)$ and every $\rho\in(0,1)$, then $f\in\Rec(p_N)$.
\end{theorem}

\begin{proof}
Fix $N,f,g,\rho$ and $l_f$ (all $f$-constants absorbed; companions and
approximants lie in the ball of Lemma~\ref{lem:III-ballinf}).  The order of
choices at a stage $L\ge l_f$ of the hypothesis is: $w\to J(w),f^r\to$ classes
$\to$ Step~0 $\to$ preparation $\to$ \L ojasiewicz targets $\to\mathcal P\to$
Step~A $\to$ data $\to s_1\to N_w\to N''\to$ $d$-consistency levers $\to$
approximant.  Put $T:=T_{\rm lo}(w)$.

\emph{Step 1 (clean sub-window).}  By Lemma~\ref{lem:III-designinf}(d) the
sub-windows of $L$ have disjoint bands for the enlarged rate scheme, and $w$
is clean for $f$.

\emph{Step 2 (deep raise).}  $f\to f^r$ beyond $J(w)$
(Lemma~\ref{lem:III-deepraise}): hypothesis (a) of
Theorem~\ref{thm:III-Eenginf} holds at $f^r$ with $\lambda_r$ and
$\varepsilon'_r\le\theta_wT^2$; tiny and robust objects stay tiny and robust
with factor $2$; the budget lemmas hold at $f^r$ for the mate
$g^r:=g-(g(\xi^r)/q^r_0)a^r$ on the window ([II,~Subsection~6.7]).

\emph{Step 3 (classes).}  Two-sided decompositions of $g^r$ at $f^r$ at every
dyadic scale of $w$, classes with the thresholds $K_*(w)^i$; by
Lemma~\ref{lem:III-nocubeinf} a class $a$ with a set $\mathcal S$ of at least
$n(w)/D_{\rm cl}(w)$ scales suffices, and Lemmas~\ref{lem:III-track}
and~\ref{lem:III-nodeep} hold on $\mathcal S$.

\emph{Step 4 (companion).}  The variant of Step~0 described before
Definition~\ref{def:III-H1}, under (H1$^*$); the preparation
(Lemma~\ref{lem:III-prep}); Step~A by Lemma~\ref{lem:III-cone}(c) under (H2).
The resulting companion $f^\#:=f^{(1a)}$ has the Step-A conclusions of
Section~\ref{sec:III-routeV} for this variant (exact ratios on $Z_\varkappa$, exact held zeros,
threshold protection, robust statuses), the types of Lemma~\ref{lem:III-prep}(d),
and $p^*(f^\#-f^r)\le\theta'_wT^2$ with $\theta'_w/c_{\rm flat}(w)^2\to0$
(Lemma~\ref{lem:III-costs}).  No step~(2) (Lemma~\ref{lem:III-noBinf}).

\emph{Step 5 (data).}  By Proposition~\ref{prop:III-cushion},
Lemma~\ref{lem:III-prep}(c) and Lemma~\ref{lem:III-noBinf}, the route-V data with
the deep assignment on $F$ form, for $t\in\mathcal S$, a window family at
$f_0:=f^\#$ with cushion constant $C_+=5$: (U1), (U2), (U3$^{\inf}$)
(support excesses counted in $\epsilon$), (U4$^{\inf}$), (U5$'$), (U7), and
$p^*(g-g_t)\le K_{\rm data}t$.  At companion support coordinates outside $F$,
(U3$^{\inf}$) holds by Proposition~\ref{prop:III-cushion}(iv) at banks (the
data are those of contacts of the bank's sign, plus the base term) and by
Proposition~\ref{prop:III-cushion}(v) at pulled coordinates.

\emph{Step 6 (levers).}  Theorem~\ref{thm:III-SysDC} provides a lever system
(U6$^{\inf}$) with $K_{\rm lev}=K_{\rm DC}\le C_{f,N}\mathrm{Des}(L)^{C_0}$, under
the conditions of Lemma~\ref{lem:III-late}.

\emph{Step 7 (approximant).}  We use $T$ as the window parameter, and we put
$s_1:=\max(64\rho\epsilon/\delta,\,e^{-1/(2T)})$.  Every new window constant ($K_{\rm DC}$,
$K_*(w)\le(\mathrm{Des}/u)^C(\log n(w))^{1/2}$, $C_+$) is at most a design
power times a power of $\log(1/T)$, so $K_w\le\exp(C\log^2(1/T))$ and
hypothesis (i) of Proposition~\ref{prop:III-S2c} holds; the violation mass,
$c_{\rm fine}$ and $W_{\rm lev}$ are at most $C_f\mathrm{Des}(L)^Cc_{L+1}+
C_fc_{L+1}^2/T$ with $c_{L+1}\le e^{-1/T}$ by (W$_{\exp}$), which is hypothesis
(ii); choose $N_w>\bar d(L)$ after $s_1$ with $\mathrm{tail}_W\le e^{-1/T}$ and
$\alpha_0(N_w)\le\theta_{\rm tr}$, where $\theta_{\rm tr}$ satisfies the bound
of Definition~\ref{def:III-approxinf}(ii) and is as small as (E-e) requires.
Proposition~\ref{prop:III-S2c} gives (VT$'$),
(FINE), (LATE), (W-a), and Lemma~\ref{lem:III-late} the conditions (L-K2) and
$K_ws_1\le r_0$.  In (W-a) of Theorem~\ref{thm:III-UE} the constant $K_J$ is
$K_J:=K_{\rm lev}=K_{\rm DC}$; hypothesis (i) of Proposition~\ref{prop:III-S2c}
contains ``(W-b) holds'', a condition on Lemma~\ref{lem:III-DC}, which in this
section is replaced by (DC) of the lever system (Theorem~\ref{thm:III-SysDC});
there is no pull, so $r_{\rm pull}=0$ and (PULL) is void.  Theorem~\ref{thm:III-UEinf} then yields a norm-attaining
approximant $f'_L$ with finite base support, $d$-consistent with $f^\#$ on
$\Omega$, with the per-piece bounds.  The statuses of coarse carriers at $f'_L$
are those of the status lemma at the approximant in Section~\ref{sec:III-routeV},
whose proof uses only $\mathrm{nv}'=\mathrm{nv}^\#$ on $\Omega$, the coarse peaks
being clamped at both rows (their margins are at least $u/2$, far above
the value moves $C_{\rm lev}K_{\rm DC}r$ and the common modes
$\mathrm{Des}^Cr$), and the fine
carriers through their weight; none of these uses privacy of the levers.

\emph{Step 8 (averaging).}  Apply Theorem~\ref{thm:III-Eenginf} along the
infinitely many stages $L_j$, with $f_j:=f^\#(L_j)$, $\lambda_j:=\lambda_r(L_j)$,
$\varepsilon'_j:=\varepsilon'_r+p^*(f^\#-f^r)\le(\theta_w+\theta'_w)T^2$,
$\mathcal S_j:=\mathcal S$, $K_j:=K_{\rm data}(w)$, $T'_j:=T$, $s_{1,j}:=s_1$.
(a) holds at $f^\#$ by Step~2 and the triangle inequality (the extra
$p^*(f^\#-f^r)$ is in $\varepsilon'_j$).  (E-c) is Step~5.  (E-d) is
Lemma~\ref{lem:III-bookkeeping}.  (E-f) and (E-e): by
Proposition~\ref{prop:III-S2c} the quantities there are at most $e^{-1/(3T)}\le
K_jT^2$, and with Lemma~\ref{lem:III-costs}, $\theta_j:=\theta_w+\theta'_w+
e^{-1/(3T)}T^{-2}$ satisfies $\theta_j/c_{\rm flat}(w)^2\to0$.  Hence $(f,\rho g)
\in\cl\NA$ along the norm-attaining rows $f'_{L_j}$, which have finite base
support.  The last sentence follows by letting $\rho\to1$.
\end{proof}

\begin{corollary}[Explicit sufficient conditions]\label{cor:III-transport}
\textup{(LSC-trunc)}$_N$ holds at $(f,g,\rho)$ if, for infinitely many main
stages $L$, at a clean sub-window and a class as in
Theorem~\textup{\ref{thm:III-transport}}, one of the conditions of
Lemma~\textup{\ref{lem:III-H1suff}} holds and one of the conditions
\textup{(a)}--\textup{(f)} of Corollary~\textup{\ref{cor:III-Asuff}} holds.  In
particular this is the case if, for infinitely many main stages $L$, the support
$F$ meets no signature set $S_l$, $l\le L$, beyond $s_{\max}(L)$; then
$f\in\Rec(p_N)$.
\end{corollary}

\begin{proof}
The first statement is Theorem~\ref{thm:III-transport} with
Lemma~\ref{lem:III-H1suff}, Corollary~\ref{cor:III-Asuff} and
Lemma~\ref{lem:III-cone}(a).  For the second: if $F\cap S_l\cap(s_{\max}(L),
\infty)=\emptyset$ for all $l\le L$, every near coordinate of a coarse carrier
is off $F$, and every far coordinate of $S_l$ in $[x,x'_{\max}]$ is off $F$ and,
for a data-carrying carrier, a contact with $z=\epsilon_l$ after the closing
moves of Step~0 (as in Lemma~\ref{lem:III-LV}(a)), hence pullable.  So every
tuned carrier is privately tunable ((P1) or (P2)), which is
Corollary~\ref{cor:III-Asuff}(a).  Every active weak peak $k$ has a private
push of the right orientation, which is Lemma~\ref{lem:III-H1suff}(i): if
$\epsilon_k=\mathrm{vs}_k$, a far pull, since its far coordinates are contacts
of sign $\epsilon_k$ and, consecutive elements of $S_k$ having $v$-ratio
$2^{-2^{k+1}}\ge\mathrm{Des}(L)^{-1}$, the smallest one with $v_k(j)\ge2b\bar
\vartheta_m\Phi_k/m$ satisfies $v_k(j)\le\mathrm{Des}(L)2b\bar\vartheta_m
\Phi_k/m\le\mathrm{Des}(L)^Cb/u$; if $\epsilon_k=-\mathrm{vs}_k$, a near bank,
since its near coordinates are off $F$ and contacts of sign $\epsilon_k$; and
(P1) if it carries no data.  This holds for every $(g,\rho)$.
\end{proof}

\begin{remark}[Relation to a modulus condition]\label{rem:III-LR}
If at a stage the two-sided Step-A columns have a right inverse of norm at most
$\mathrm{Des}(L)^{C_0}/u$, then $\varepsilon_{\rm PS}\ge u\,\mathrm{Des}(L)^{-C_0}
|\Lambda_A|^{-1/2}\ge u^3\mathrm{Des}(L)^{-C_1}$, and (H2) holds by
Lemma~\ref{lem:III-cone}(a) once $C_A\ge C_1+2$.  So (H2) is weaker than such a
modulus condition, and the $d$-consistency half of any such condition is not
needed at all (Theorem~\ref{thm:III-SysDC}).  A condition requiring a push
resource for each weak peak separately does not imply (H1$^*$) (the discussion
before Lemma~\ref{lem:III-H1suff}); (H1$^*$) is the correct hypothesis.
\end{remark}

\subsection{The residual classes}\label{subsec:III-residual}

\begin{definition}[Residual conditions]\label{def:III-residual}
Let $L\ge l_f$ be a main stage, $w$ a clean sub-window of $L$ and $a$ an
activity class with at least $n(w)/D_{\rm cl}(w)$ scales.  With
$\gamma_0:=c_{\rm A}/(4\|h\|_2)$ and $c_*$ as in Theorem~\ref{thm:III-Aobs}:
\begin{enumerate}
\item[(W$^*$)] condition (H1$^*$) fails at $(L,w,a)$;
\item[(A$^*$/U$^*$)] (H1$^*$) holds, and for every pre-placement $\mathcal P$
there is a near-certificate $y$ of size $<4u^3\mathrm{Des}(L)^{-C_A}$ with
$\varphi_y>-\rho_0$ and $\gamma\le-c_*$; the certificate is of \emph{type}
\textup{(A$^*$)} if $\gamma\le-\gamma_0$ and of \emph{type} \textup{(U$^*$)} if
$-\gamma_0<\gamma\le-c_{\rm A}\nu_1/(8H(L))$.
\end{enumerate}
The type may depend on the pre-placement; this is why the two types form one
condition.
\end{definition}

\begin{proposition}[Necessary conditions for not being covered]
\label{prop:III-residual}
Let the design be $T^\infty_{\rm fin}$, $N\ge1$, $f\in S_{p_N^*}$, $g\in\Cset(f)$
and $\rho\in(0,1)$.  If the hypothesis of
Theorem~\textup{\ref{thm:III-transport}} fails for $(f,g,\rho)$, then for all but
finitely many main stages $L$, at every clean sub-window $w$ of $L$ and every
class $a$ with at least $n(w)/D_{\rm cl}(w)$ scales, \textup{(W$^*$)} or
\textup{(A$^*$/U$^*$)} holds.  In the case \textup{(A$^*$/U$^*$)} the
near-certificates have all the properties
\textup{(1)}--\textup{(6)} of Theorem~\textup{\ref{thm:III-Aobs}}, none of the
conditions \textup{(a)}--\textup{(e)} of Corollary~\textup{\ref{cor:III-Asuff}}
holds, and:
\begin{enumerate}
\item[(i)] the thick support at design depth is approximately proportional
to the certificate's carrier combination (error $\varepsilon\mathrm{Des}(L)/
|\kappa_\Psi|$; here $a_s$ are the moduli at $f_1+P$, i.e.\ those of $f$ at
design depth divided by $\lambda_r$, up to the Step-0 and preparation
moves): $|a_s-\Psi_y(s)/\kappa_\Psi|\le
\varepsilon\mathrm{Des}(L)/|\kappa_\Psi|$ for every $s\le\bar d(L)$ with
$|a_s|\ge u/2$, with one common $\kappa_\Psi\le-c_*/\nu_1$, and every such $s$
lies in $U_L$;
\item[(ii)] for a certificate of type \textup{(A$^*$)} the coefficients of the
carriers with only (P4) resources are $\epsilon_l\gamma h_l$ up to
$O(\varepsilon\mathrm{Des}^2)$; for a certificate of type \textup{(U$^*$)} the
coefficients of the carriers of kind \textup{(iv)}
\textup{(}all near coordinates S$_{\rm th}$ of the data sign\textup{)} have
$\ell_2$-norm at least $c_{\rm A}/2$ and oriented signs $\mathrm{coef}_l\sigma_l\ge
-O(\varepsilon\mathrm{Des}^2)$;
\item[(iii)] the one-sided generators carry a total action at least $c_*/2$.
\end{enumerate}
\end{proposition}

\begin{proof}
If the hypothesis of Theorem~\ref{thm:III-transport} fails, then for all but
finitely many main stages, at every clean sub-window and admissible class,
(H1$^*$) or (H2) fails.  If (H2) fails, the stage is not A-regular for any
$\mathcal P$, and Lemma~\ref{lem:III-cone}(b) gives for every $\mathcal P$ a
near-certificate of size $<4u^3\mathrm{Des}^{-C_A}\le\varepsilon_*$ (for
$C_A\ge C+2$) with $\varphi_y>-\rho_0$.  Theorem~\ref{thm:III-Aobs} applies:
(3) gives $\gamma\le-c_*$, so (A$^*$/U$^*$) holds, the type of each
certificate being (A$^*$) or (U$^*$) according to whether $\gamma\le-\gamma_0$; (4) gives (i); the proof of (3) gives (ii) (in case
$\gamma>-\gamma_0$ the kind-(iv) coefficients carry $\ell_2$-mass $\ge c_{\rm A}/2$, and
in case $\gamma\le-\gamma_0$ the kind-(iii) coefficients are $\epsilon_l\gamma
h_l+O(\varepsilon\mathrm{Des}^2)$); (6) gives (iii).  If one of
Corollary~\ref{cor:III-Asuff}(a)--(e) held, there would be no near-certificate of
size $\le\varepsilon_*$.
\end{proof}

\begin{remark}[Status of the residual]\label{rem:III-residual}
(a) Proposition~\ref{prop:III-residual} gives necessary conditions only; the
converse is not claimed (Lemma~\ref{lem:III-cone}(b) is one-directional, and a
certificate for one pre-placement need not obstruct another).  It is not shown
that rows satisfying (W$^*$) or (A$^*$/U$^*$) at all large main stages exist,
nor that such rows have a mate $g$ with $(f,\rho g)\notin\cl\NA$.  No
counterexample to density is known.
(b) The certificates of (A$^*$/U$^*$) are sign (cone) obstructions, not rank obstructions: the
$d$-consistency system has no such obstruction after the preparation
(Theorem~\ref{thm:III-SysDC}), but the preparation cannot help Step~A, because the
targets of Step~A are ratios fixed by the \L ojasiewicz step at the displacement
scale $x\ge\mathrm{Des}(L)b$, not values of the row; moves that keep the tiny
objects tiny have size at most $b$, so data-carrying contacts and thin
coordinates give only one-sided resources at the scale $x$.
(c) Item (i) is an infinite-depth form of the structure of
[II,~Lemma~6.39]: at every large stage the thick part of $a$ at design depth is
approximately proportional (error $\varepsilon\mathrm{Des}(L)/|\kappa_\Psi|$) to a fixed combination of carrier vectors with a robustly negative
common mode.
(d) \emph{Heuristic directions (not used).}  If the exact set $Z_\varkappa$ of
the \L ojasiewicz step has positive dimension at the selected point, the targets
might be chosen so that $\varphi_y\le0$ for the finitely many near-certificates,
which would remove (A$^*$/U$^*$); placements of support of the sign
$+\epsilon_k$ at far coordinates (the mirror image of the pulls) might break
certificates as in Remark~\ref{rem:III-Asuff}(c); and for (W$^*$), raising another
robust peak of the same block raises $\bar\vartheta_m$ and lowers $\varrho_k$,
an additional push resource.  None of this is proved.
(e) Statements of earlier drafts that are \emph{not} used and are not true as
stated: a bank placed with an offset at the row where the $d$-consistency system
is solved is not a two-sided lever along a degenerate direction
(Lemma~\ref{lem:III-offset}); a per-peak push condition does not imply (H1$^*$);
and descriptions of the non-covered class by carrier types are necessary
conditions only, not characterizations.
\end{remark}

\begin{remark}[Dependence on earlier results]\label{rem:III-infdep}
Theorems~\ref{thm:III-UEinf}, \ref{thm:III-Eenginf}, \ref{thm:III-SysDC},
\ref{thm:III-transport} and the results of
Subsections~\ref{subsec:III-cushion}--\ref{subsec:III-SysA} are proved here
modulo the cited results of Parts~I and~II and of
Sections~\ref{sec:III-generic}--\ref{sec:III-routeV}, in particular the route-V
companion of Section~\ref{sec:III-routeV} (Step~0, the \L ojasiewicz step and its
threshold protection, the status lemma at the approximant) and the window count
of Section~\ref{sec:III-design}.  These results of
Sections~\ref{sec:III-generic}--\ref{sec:III-MTVI} are used for
$T^\infty_{\rm fin}$ through the base independence of the window theorems
([II,~Lemma~3.22]; the base enters only through the bank factor in
$\mathrm{Des}$); this was not re-verified line by line for the base
$\mu^\star$.

\emph{Sketch, not verified.}  We expect Theorem~\ref{thm:III-MTVI} to hold for
$T^\infty_{\rm fin}$.  The insertions are covered by
Lemma~\ref{lem:III-designinf}(d).  The change of base $\mu\to\mu^\star$
enters Section~\ref{sec:III-MTVI} only through (B3), (W2), (W6) of
Lemma~\ref{lem:III-interface} and the tuning lemma behind
Proposition~\ref{prop:III-transfer}; each has a $\mu^\star$-analogue in
Part~II (where, however, [II,~Lemma~4.15] is formed with $B^\star$ at the
third element beyond $s_{\max}$, while $T_{\rm fin}$'s $\mathrm{Des}$ uses the
first), but this transfer has not been written out or verified.  If, in
addition, Theorem~\ref{thm:III-MTVI} holds for $T^\infty_{\rm fin}$, then,
combined with [II,~Proposition~6.48], the results of this section give: if
every $(f,g,\rho)$ with infinite base support satisfies the hypothesis of
Theorem~\ref{thm:III-transport}, then $\NA((c_0,p_N),\ell_2^2)$ is dense for
the design $T^\infty_{\rm fin}$; this hypothesis is not verified for all
rows.
\end{remark}

\section{Status and open problems}\label{sec:III-status}

\subsection{What is proved}
We collect the main results with their hypotheses.
\begin{enumerate}
\item[(1)] \emph{Finite base support} (Theorem~\ref{thm:III-MTVI}).  For the
operator $T_{\rm fin}$ of Section~\ref{sec:III-design}, which does not
depend on $N$, satisfies (T-b)--(T-d) and, with option (b), the conclusion
of Mart\'in's Lemma~B, while with option (a) it has norm $c_1<1$
(Theorem~\ref{thm:III-Afin}): with option (a), for every
$N\ge1$ every first row $f\in S_{p_N^*}$ with finite base support lies in
$\Rec(p_N)$.  With option (b) the same holds for $N=1$, and for $N\ge2$ under
$\varrho_1(f)\ne1$.
\item[(2)] \emph{Density criterion} (Theorem~\ref{thm:III-VI2},
Corollary~\ref{cor:III-VI3}).  With option (a), for every $N\ge1$:
$\NA((c_0,p_N),\ell_2^2)$ is dense if and only if (LSC-trunc)$_N$ holds; and
(LSC-trunc)$_N$ for infinitely many $N$ implies density of
$\NA((c_0,p),\ell_2^2)$ for the Mart\'in-type norm $p$ built with
$T_{\rm fin}$.
\item[(3)] \emph{Tools valid beyond $T_{\rm fin}$.}  Status coherence for
generic weights (Lemmas~\ref{lem:III-GEN} and~\ref{lem:III-QC}); exactification
with threshold protection without (KN), for every design satisfying the QC
requirements (Proposition~\ref{prop:III-KNtr});
exact $d$-consistency by levers (Lemma~\ref{lem:III-DC}); the uniform
engineered bound and averaging at the norm-attaining approximant
(Theorems~\ref{thm:III-UE} and~\ref{thm:III-Eeng}); the lowering transfer
(Lemma~\ref{lem:III-LT}); and the engineered bound and windowed recovery at
arbitrary base support (Theorems~\ref{thm:III-UEinf}
and~\ref{thm:III-Eenginf}).  Each is stated with the hypotheses under which
it is proved.
\item[(4)] \emph{Infinite base support} (Theorem~\ref{thm:III-transport},
Corollary~\ref{cor:III-transport}, Proposition~\ref{prop:III-residual}).  For
$T^\infty_{\rm fin}$, (LSC-trunc)$_N$ holds at every $(f,g,\rho)$ satisfying
(H1$^*$) and (H2) at a clean sub-window and a class with at least
$n(w)/D_{\rm cl}(w)$ scales, at infinitely many main stages, and the triples
not covered satisfy (W$^*$) or (A$^*$/U$^*$) at almost all main stages
(each result checked by one referee).
\end{enumerate}

\subsection{Open problems}
\begin{problem}[Density for the designed operator]\label{prob:III-LSC}
Does \textup{(LSC-trunc)$_N$} hold for $T_{\rm fin}$ \textup{(}option
\textup{(a))} and some, or every, $N\ge1$?  Equivalently
\textup{(}Theorem~\textup{\ref{thm:III-VI2})}: is $\NA((c_0,p_N),\ell_2^2)$ dense?
\end{problem}

\begin{problem}[The Step-A residual]\label{prob:III-Astar}
For $T^\infty_{\rm fin}$, are there triples $(f,g,\rho)$ that satisfy
\textup{(A$^*$/U$^*$)} of
Definition~\textup{\ref{def:III-residual}} at every clean sub-window of all
large main stages?  If so, is $(f,\rho g)\in\cl\NA((c_0,p_N),\ell_2^2)$ for
such triples?  A positive answer to the second question for all such triples,
together with a positive answer to Problem~\textup{\ref{prob:III-Wstar}}, would give
\textup{(LSC-trunc)$_N$} for $T^\infty_{\rm fin}$.
\end{problem}

\begin{problem}[The push residual]\label{prob:III-Wstar}
Does the push condition \textup{(H1$^*$)} of
Definition~\textup{\ref{def:III-H1}} hold at every row, or can the weak-peak
pushes of Step~$0$ be dispensed with?  By the discussion before
Lemma~\textup{\ref{lem:III-H1suff}} a condition for each weak peak separately
does not suffice; Lemma~\textup{\ref{lem:III-H1suff}} gives sufficient
conditions.
\end{problem}

\begin{problem}[Option (b) at the first stage]\label{prob:III-optb}
For $T_{\rm fin}$ with option \textup{(b)} and $N\ge2$, is every first row with
finite base support and $\varrho_1(f)=1$ recoverable?  The question is whether
the rational weight $\Phi_1=2^{-3}$ can be a critical value of the floor
function of stage~$1$ on a stratum of an exact zero set, identically in the
transcendental parameters \textup{(}Remark~\textup{\ref{rem:III-MTVIopen}(b))}.
\end{problem}

\begin{problem}[Other operators]\label{prob:III-other}
Does Master Theorem~VI hold for an arbitrary admissible operator, in
particular for the operators produced by Mart\'in's Lemma~B in
\cite{Martin}?  The proofs here use genericity of the weights, the window
structure and the design factor of $T_{\rm fin}$ in an essential way.
\end{problem}

\begin{problem}[Property B]\label{prob:III-B}
Does $\ell_2^2$ have Lindenstrauss' property~B \cite{Lindenstrauss}?  By
[I,~Section~9], a negative answer to Problem~\textup{\ref{prob:III-LSC}} for one
$N$ would give a negative answer here.
\end{problem}

\subsection{Remarks on the status}
\begin{remark}[The two designs]\label{rem:III-twodesigns}
Section~\ref{sec:III-MTVI} uses $T_{\rm fin}$ over the base
$\mu_s=2^{-s^2-1}$, and Section~\ref{sec:III-inf} uses $T^\infty_{\rm fin}$,
which is $T_{\rm fin}$ with option (a) over the doubly exponential base
$\mu^\star_s=2^{-2^{2^s}}$ of [II,~Definition~3.1] and with the insertions of
Definition~\ref{def:III-designinf}.  The doubly exponential base is needed
at infinite base support for the reason given in [II,~Remark~3.3]; the base
$\mu_s$ is used at finite base support because the tuning and bank estimates
of the transfer list were re-derived for it.  Whether Master Theorem~VI, and
hence the density criterion, hold for $T^\infty_{\rm fin}$ is not proved
here.  The monotonicity argument of Remark~\ref{rem:III-infdep} covers the
insertions (I1)--(I5).  It does not cover the change of base from $\mu$ to
$\mu^\star$, which enters the proof of Theorem~\ref{thm:III-MTVI} through
(B3), (W2), (W6) of Lemma~\ref{lem:III-interface} and the tuning lemma
behind Proposition~\ref{prop:III-transfer}, all verified for $\mu$ only.
We expect the transfer, since each item has a $\mu^\star$-analogue with
$B_{\mu^\star}$, but it has not been written out or verified.
\end{remark}

\begin{remark}[Later observations, not verified and not used]\label{rem:III-later}
The following observations come from an analysis made after the results of
Section~\ref{sec:III-inf} were refereed; they have not been independently
checked, and nothing in this paper depends on them.
(a) \emph{Sign of the data at weak peaks.}  The sign of the data at an
active weak peak, $\epsilon_k=\mathrm{vs}_k\sgn(e_k-\Delta'_m)$ by
[I,~Lemma~8.27], is now used in Definition~\ref{def:III-H1} and
Lemma~\ref{lem:III-H1suff}; see the discussion before
Definition~\ref{def:III-H1}.
(b) \emph{Heuristic.}  The same analysis suggests that the weak-peak pushes
of Step~$0$ can be dropped if the block scalars are computed with all weak
peaks frozen as non-peaks, which would make (W$^*$) empty and leave the
Step-A residual (A$^*$/U$^*$) as the only obstruction.  This is not proved
here.
\end{remark}

\begin{remark}[Evidence]\label{rem:III-evidence}
Nothing proved in Parts~I--III points to a counterexample to density; none of
the verifications, and none of the finite numerical models used as sanity
checks (the clamp equation, the $\alpha$-mass identity, the bank equations
of the tuning lemma, the lowering transfer, the solvability of the prepared
$d$-consistency system), produced one.  Equally, nothing proved here settles
Problem~\ref{prob:III-LSC} for any $N$.
\end{remark}

\begin{thebibliography}{99}
\bibitem[I]{PartI}
\emph{Recovery of mates on Mart\'in norms and the resonance obstruction},
research note (Part~I of this study), October 2026.
\bibitem[II]{PartII}
\emph{Recovery of mates on Mart\'in norms II: companions, clean sub-windows
and a designed operator}, research note (Part~II of this study), October 2026.
\bibitem{BCR}
J.~Bochnak, M.~Coste and M.-F.~Roy, \emph{Real Algebraic Geometry},
Ergebnisse der Mathematik und ihrer Grenzgebiete (3) \textbf{36},
Springer, Berlin, 1998.
\bibitem{DGS}
G.~Debs, G.~Godefroy and J.~Saint Raymond,
\emph{Topological properties of the set of norm-attaining linear
functionals}, Canad. J. Math. \textbf{47} (1995), 318--329.
\bibitem{Hoffman}
A.~J.~Hoffman, \emph{On approximate solutions of systems of linear
inequalities}, J. Research Nat. Bur. Standards \textbf{49} (1952),
263--265.
\bibitem{JW}
W.~B.~Johnson and J.~Wolfe, \emph{Norm attaining operators},
Studia Math. \textbf{65} (1979), 7--19.
\bibitem{KLMW}
V.~Kadets, G.~L\'opez, M.~Mart\'in and D.~Werner,
\emph{Norm attaining operators of finite rank}, arXiv:1905.08272.
\bibitem{KLMW20}
V.~Kadets, G.~L\'opez, M.~Mart\'in and D.~Werner,
\emph{Equivalent norms with an extremely nonlineable set of norm attaining
functionals}, J. Inst. Math. Jussieu \textbf{19} (2020), 259--279.
\bibitem{Lindenstrauss}
J.~Lindenstrauss, \emph{On operators which attain their norm}, Israel J.
Math. \textbf{1} (1963), 139--148.
\bibitem{Martin}
M.~Mart\'in,
\emph{A Banach space whose set of norm-attaining functionals is
algebraically trivial}, J. Funct. Anal. \textbf{288} (2025);
arXiv:2406.07273.
\bibitem{Miranda}
C.~Miranda, \emph{Un'osservazione su un teorema di Brouwer}, Boll. Un.
Mat. Ital. (2) \textbf{3} (1940), 5--7.
\bibitem{Seidenberg}
A.~Seidenberg, \emph{A new decision method for elementary algebra}, Ann.
of Math. (2) \textbf{60} (1954), 365--374.
\bibitem{Tarski}
A.~Tarski, \emph{A Decision Method for Elementary Algebra and Geometry},
2nd ed., University of California Press, Berkeley and Los Angeles, 1951.
\bibitem{vdDries}
L.~van den Dries, \emph{Tame Topology and O-minimal Structures}, London
Mathematical Society Lecture Note Series \textbf{248}, Cambridge University
Press, Cambridge, 1998.
\end{thebibliography}
\end{document}
